
\documentclass[12pt, a4paper]{report}
\usepackage{graphicx} % For including images
\usepackage{titlesec} % For customizing section titles
\usepackage{tocloft} % For customizing table of contents
\usepackage{acro} % For acronyms
\acsetup{make-links=false, list/display=all}
\usepackage{tikz}
\usetikzlibrary{shapes, arrows, positioning, calc}
\usepackage{amssymb}
\usepackage{algorithm}
% \usepackage{algorithmic}
\usepackage{algpseudocode}
\usepackage{soul}    % in preamble
\usepackage{comment}
\usepackage{xcolor}
\usepackage{enumitem}
\usepackage{amsmath}
\usepackage{amssymb}
\usepackage{pifont}
\usepackage{tabularx}
\usepackage{adjustbox}
\usepackage{booktabs}
\usepackage{enumitem}
\usepackage{graphicx}
\usepackage{caption}
\usepackage{longtable}

\newcommand{\cmark}{\ding{51}} % check mark
\newcommand{\xmark}{\ding{55}} % cross mark


% Remove the first geometry package call and keep only the second one with all options
\usepackage[left=1.5in, right=1in, top=1in, bottom=1in, includehead, includefoot]{geometry}
\usepackage{array}
\usepackage{enumitem}
\usepackage{multirow}
\usepackage{booktabs}
\usepackage{ragged2e}
\usepackage{adjustbox}
% graphicx is already loaded above, remove duplicate
\usepackage{float}
\usepackage{subcaption}




%% ____Bibliography____%%
\usepackage[numbers,sort&compress]{natbib}
\usepackage{chapterbib}
\usepackage[hidelinks]{hyperref}
%\hypersetup{colorlinks=true,citecolor=blue,linkcolor=blue,urlcolor=blue}
% Page margins
\usepackage[left=1.5in, right=1in, top=1in, bottom=1in, includehead, includefoot]{geometry}
\usepackage{float}
\usepackage{subcaption}
% Remove page number from the first page
\thispagestyle{empty}

% Customize table of contents, list of figures, and list of tables
\renewcommand{\cfttoctitlefont}{\hfill\Large\bfseries}
\renewcommand{\cftaftertoctitle}{\hfill}
\renewcommand{\cftloftitlefont}{\hfill\Large\bfseries}
\renewcommand{\cftafterloftitle}{\hfill}
\renewcommand{\cftlottitlefont}{\hfill\Large\bfseries}
\renewcommand{\cftafterlottitle}{\hfill}

% Rein in TOC entry indentation so deeper levels (subsection/subsubsection)
% do not drift far to the right of the page
\cftsetindents{chapter}{0em}{2.3em}
\cftsetindents{section}{2.3em}{2.8em}
\cftsetindents{subsection}{5.1em}{3.4em}
\cftsetindents{subsubsection}{8.5em}{4.2em}
\cftsetindents{paragraph}{12.7em}{4.8em}

% Show subsubsections in TOC / numbering (default report class stops at subsection)
\setcounter{secnumdepth}{4}
\setcounter{tocdepth}{4}


% Define acronyms
\DeclareAcronym{VANET}{
  short = VANET,
  long  = Vehicular Ad Hoc Network
}

\DeclareAcronym{AI}{
  short = AI,
  long  = Artificial Intelligence
}

\DeclareAcronym{SDN}{
  short = SDN,
  long  = Software-Defined Networking
}

\DeclareAcronym{SDVN}{
  short = SDVN,
  long  = Software-Defined Vehicular Network
}

\DeclareAcronym{RSU}{
  short = RSU,
  long  = Roadside Unit
}

\DeclareAcronym{GAT}{
  short = GAT,
  long  = Graph Attention Network
}

\DeclareAcronym{NS3}{
  short = NS-3,
  long  = Network Simulator 3
}

\DeclareAcronym{LSTM}{
  short = LSTM,
  long  = Long Short-Term Memory
}

\DeclareAcronym{FHE}{
  short = FHE,
  long  = Fully Homomorphic Encryption
}

\DeclareAcronym{TRS}{
  short = TRS,
  long  = Threshold Ring Signature
}

\DeclareAcronym{BSM}{
  short = BSM,
  long  = Basic Safety Message
}

\DeclareAcronym{MitM}{
  short = MitM,
  long  = Man-in-the-Middle
}

\DeclareAcronym{ITS}{
  short = ITS,
  long  = Intelligent Transportation System
}

\DeclareAcronym{HMAC}{
  short = HMAC,
  long  = Hash-based Message Authentication Code
}

\DeclareAcronym{BFT}{
  short = BFT,
  long  = Byzantine Fault Tolerance
}

\DeclareAcronym{IPFS}{
  short = IPFS,
  long  = InterPlanetary File System
}

\DeclareAcronym{PQC}{
  short = PQC,
  long  = Post-Quantum Cryptography
}

\DeclareAcronym{PKI}{
  short = PKI,
  long  = Public Key Infrastructure
}

\DeclareAcronym{SUMO}{
  short = SUMO,
  long  = Simulation of Urban Mobility
}

\DeclareAcronym{V2V}{
  short = V2V,
  long  = Vehicle-to-Vehicle
}

\DeclareAcronym{V2I}{
  short = V2I,
  long  = Vehicle-to-Infrastructure
}

\DeclareAcronym{V2X}{
  short = V2X,
  long  = Vehicle-to-Everything
}




\begin{document}
\justifying
\include{cover}

\renewcommand{\thepage}{\roman{page}} % Start page numbering in roman

\chapter*{Abstract}


Mobility pattern and trajectory poisoning attacks corrupt the spatio-temporal
learning models of SDVN controllers by injecting kinematically plausible yet
collectively falsified beacon sequences over time, inducing safety-critical
misrouting, phantom congestion, and suppression of collision warnings without
triggering conventional detection mechanisms.
Existing defenses are fundamentally inadequate against this threat class:
threshold-based mobility validators calibrate detection parameters under
static observation assumptions that collapse at high vehicle speeds;
cryptographic authentication schemes verify identity but provide no guarantee
of data truthfulness, leaving authenticated insiders free to poison learning
models; individual vehicle analysis misses coordinated multi-source attacks
whose spatial inconsistency is only detectable relationally; temporal learning
approaches rely on generic network traffic features rather than mobility-specific
sequential dynamics; and no reviewed framework eliminates the centralized trust
bottleneck at the RSU or controller that a single compromise can silently
subvert. To address this, we propose MPTD-PQS, a dual-mode detection and mitigation
framework that operates in a lightweight mode via rule-based attack signature
detection and HMAC beacon integrity enforcement at the RSU, and in a full mode
via a Graph Attention Network (GAT) for spatial anomaly detection over dynamic
vehicle--RSU graphs, an LSTM temporal autoencoder for sequential trajectory
anomaly scoring, and a post-quantum Threshold Ring Signature (TRS) scheme
requiring $t$-of-$n$ RSU collaborative signing to structurally prevent
minority-compromised RSU aggregates from reaching the controller. each RSU encrypts its local aggregate under post-quantum Fully Homomorphic Encryption (FHE) via OpenFHE before TRS ring coordination, preventing
compromised ring members from observing neighbouring RSU aggregate contents,
with the cloud performing blind homomorphic computations across combined
encrypted aggregates without decryption. A permissioned
blockchain with smart contracts (SC-Trust, SC-Revoke) provides Byzantine
fault-tolerant trust management and real-time revocation across both modes,
with IPFS integration bridging real-time and offline detection. Evaluated
via NS-3 and SUMO under trajectory poisoning, Sybil-based mobility pattern
poisoning, MitM, and control-plane compromise scenarios.


\newpage


% Table of Contents
\tableofcontents
\newpage

% List of Figures
\listoffigures
\newpage

% List of Tables
\listoftables
\newpage

% ========== ACRONYMS PAGE HERE ==========
\addcontentsline{toc}{chapter}{Acronyms}
\chapter*{Acronyms}
\printacronyms[heading=none] % Add [heading=none]

% ========================================
\newpage

\renewcommand{\thepage}{\arabic{page}} % Start page numbering in arabic
\setcounter{page}{1} % start page numbering from 1

% Main Content
\chapter{Introduction}
\section{Overview of SDVN}

Vehicular Ad Hoc Networks (\ac{VANET}) were originally designed to support direct vehicle to vehicle and vehicle to infrastructure communication in a fully distributed manner. However, the absence of centralized control in VANET makes global traffic management, scalability, and security enforcement difficult in highly dynamic environments \cite{hartenstein2008}.

Software Defined Networking (SDN) addresses these challenges by separating the control plane from the data plane and enabling centralized and programmable network control. By applying SDN principles to vehicular environments, Software-Defined Vehicular Networks (SDVN) were introduced to provide global network visibility, flexible control, and intelligent decision making capabilities \cite{wahid2022}.

In the context of SDVN, a centralized controller manages the network by harnessing a holistic and global understanding of vehicles, RSUs, and network communications. The system design makes it easy to manage traffic routing, congestion control, handover, and network optimization. To ensure intelligent decision-making, the SDVN heavily depends on learning approaches to evaluate vehicle mobility and path information. Data like vehicle position, speed, direction, and path history is constantly harvested and used by the controller to understand vehicle mobility and predict the future state of the network. These predictions help in improving the management of the network in highly dynamic environments \cite{cabarkapa2019}.


\section{Significance and Impact of Mobility Pattern and Trajectory Poisoning Attacks}
Mobility pattern poisoning and trajectory poisoning attacks are serious threats to SDVN because they directly target the learning process of the controller. In these attacks, attackers inject false but realistic mobility data over time, causing the learning models to gradually produce incorrect results. Since the data follows normal movement behavior, such attacks are difficult to detect using traditional security methods.

If these attacks succeed, the controller may make wrong traffic predictions, inefficient routing decisions, or unstable handover choices. This can lead to increased congestion, reduced network performance, and unreliable communication. In safety related applications, incorrect decisions may even result in dangerous situations. Therefore, detecting and mitigating mobility and trajectory poisoning attacks is essential for maintaining secure and reliable SDVN operation.

\section{Problem Introduction}

SDVN depend on centralized, learning based models to manage vehicle mobility, traffic conditions, and safety related decisions. These models continuously learn from mobility and trajectory data reported by vehicles and RSU to predict future network states. However, the correctness of these decisions depends on the integrity and reliability of the collected mobility data.

Recent studies shows that attackers can exploit this dependency by injecting falsified but realistic mobility information over time, gradually influencing the learning process without triggering conventional security alarms. Such attacks may originate from different entities within the SDVN architecture, including vehicles, roadside units, or even compromised controllers and can affect both data plane observations and control plane decisions.
\section{Limitations of Existing Approaches}

Although several solutions have been proposed for SDVN security,
existing approaches exhibit ten critical limitations. First, most
focus on traditional network attacks and do not address poisoning of
learning models through falsified vehicle mobility and trajectory
data~\cite{arif2020, ravi2023}. Second, cryptographic authentication
schemes verify identity but not data truthfulness, leaving
authenticated insiders free to inject false trajectories~\cite{arif2020}.
Third, most solutions assume a trustworthy SDVN controller, which is
unrealistic when the controller itself may be compromised~\cite{vehcom}.
Fourth, existing detection thresholds are calibrated under static
observation assumptions that collapse at high vehicle speeds, where
the speed-dependent RSU contact duration $\tau_c$
(Eq.~\eqref{eq:contact_duration}) shrinks the observable beacon
sampling window below the minimum needed to accumulate sufficient
detection evidence~\cite{ghaleb2014}. Fifth, individual and pairwise vehicle
analysis cannot detect coordinated multi-source attacks whose spatial
inconsistency is only visible relationally~\cite{ghaleb2014}. Sixth,
no reviewed work applies deep sequential learning to trajectory
time-series, leaving gradual poisoning across beacon sequences
undetected~\cite{raja2021}. Seventh, all reviewed defenses are
reactive, detecting attacks only after malicious data has already
influenced control decisions~\cite{hong2015, raja2021}. Eighth, every
cryptographic scheme in the vehicular security literature relies on
classical hardness assumptions with no post-quantum resilience
evaluation~\cite{faleh2025csasv, ahmed2022blockchain}. Ninth, no framework analyses detection feasibility under
resource-heterogeneous RSU deployments where legacy units cannot
run deep learning inference --- a gap this work identifies and
motivates the deployment-feasible lightweight mode against, though
a dedicated heterogeneous-RSU deployment evaluation remains outside
the current experimental scope.

\section{Problem Statement}
In SDVN-enabled intelligent transportation systems, the high mobility of
vehicles creates dynamically constrained RSU contact windows that adversaries
exploit to inject mobility-consistent yet collectively falsified beacon
sequences at rates that evade per-beacon plausibility checks, progressively
corrupting the controller's spatio-temporal learning models to induce
safety-critical misrouting, phantom congestion, and collision risk
escalation, a threat class that existing authentication-centric and
generic intrusion detection approaches cannot detect or cryptographically
mitigate, particularly when attacks originate from compromised controllers
or coordinated multi-source adversaries.

\section{Objectives and Scope}

The main objectives of this project are:
\begin{itemize}
	\item \textbf{Attack Analysis:}
	To study and model mobility pattern and trajectory poisoning attacks in SDVN, including attacks at both the data plane and control plane by malicious vehicles, compromised RSUs, and malicious controllers.


	\item \textbf{Mitigation Method Design:}
To design a mitigation approach that combines blockchain and
learning based techniques to protect vehicular mobility data, by
analyzing spatial relationships using Graph Attention Networks and
detecting abnormal mobility behavior over time using an LSTM temporal
autoencoder.

\item \textbf{Performance Evaluation:}
To define an evaluation strategy for the effectiveness of the
proposed mitigation framework under different poisoning attack
scenarios using controlled NS-3 and SUMO simulation experiments
across urban, rural, and highway mobility regimes.
\end{itemize}

The scope of this project is limited to the modeling and mitigation of mobility pattern and trajectory poisoning attacks in learning based SDVN. The evaluation of the proposed framework is planned to be conducted using simulation tools to analyze security performance and system behavior under controlled attack scenarios.

Real world deployment, hardware implementation, large scale field testing, and network attacks unrelated to mobility data poisoning are outside the scope of this project.


\subsection{Research Questions}
\label{sec:rqs}

This work addresses the following research questions, each mapping
directly to a contribution, its evaluation experiment, and its
ablation variants.

\begin{description}

\item[\textbf{RQ1.}] How do vehicle mobility parameters --- specifically
RSU contact duration, handover rate, and traffic density --- quantitatively
alter the detection feasibility window and the adversarial drift budget
available to trajectory and mobility pattern poisoning attackers in SDVN,
and do existing threshold-based defenses remain effective under
high-mobility conditions?
(Maps to: mobility amplification formulation; evaluated by
detection feasibility condition Eq.~\eqref{eq:detection_condition} and
drift budget Eq.~\eqref{eq:cumulative_drift} across speed regimes;
primary empirical validation via Experiment~2
(Section~\ref{sec:exp2}, speed $v \in \{10, 60, 100, 140\}$\,km/h)
and Experiment~4 (Section~\ref{sec:exp4}, streak length component
of TL); reported metrics: MCC (Eq.~\ref{eq:mcc}),
TTD (Eq.~\ref{eq:ttd}), TPE (Eq.~\ref{eq:tpe}).)

\item[\textbf{RQ2.}] Does the proposed multi-source dual-plane threat model
reveal attack vectors --- specifically control-plane compromise by a malicious
controller and coordinated multi-RSU Sybil injection --- that single-source
or data-plane-only threat models used in prior works fail to capture, and
does the proposed framework maintain detection effectiveness across all four
attacker classes simultaneously?
(Maps to: threat model contribution; evaluated by per-attacker-class
MCC (Eq.~\ref{eq:mcc}), TTD (Eq.~\ref{eq:ttd}),
CDER (Eq.~\ref{eq:cder}), TDEE (Eq.~\ref{eq:tdee}),
and TPE (Eq.~\ref{eq:tpe})
across malicious vehicle, compromised RSU, MitM, and malicious
controller scenarios; primary empirical validation via
Experiment~5 (Section~\ref{sec:exp5}, per-attack variant table
under default settings).)

\item[\textbf{RQ3.}] Can rule-based attack signature detection, operating
exclusively at the RSU without deep learning inference, achieve practically
sufficient detection performance against both trajectory and mobility
pattern poisoning attacks, and where does it fail relative to the full mode?
(Maps to: attack signature characterization and lightweight mode;
evaluated by MCC (Eq.~\ref{eq:mcc}) and TTD (Eq.~\ref{eq:ttd})
of ablation variant AB5 (full AI layer removed,
Section~\ref{sec:ablation}) across all attack variants and
penetration rates; primary empirical validation via
Experiment~1 (Section~\ref{sec:exp1}) and
Experiment~2 (Section~\ref{sec:exp2}); PBPO
(Eq.~\ref{eq:pbpo}) confirms the 100\,ms real-time constraint
is satisfied by the lightweight-only path.)

\item[\textbf{RQ4.}] Does the integration of GAT spatial anomaly detection
and LSTM temporal autoencoder scoring in the full operating mode
significantly improve detection performance over the lightweight mode
and over each component individually, across varying attacker types,
penetration rates, and vehicle speed regimes?
(Maps to: dual-mode detection framework; evaluated by MCC
(Eq.~\ref{eq:mcc}) and TTD (Eq.~\ref{eq:ttd}) of ablation variants
AB3 (GAT disabled), AB4 (LSTM-AE disabled), and AB5 (full AI layer
removed) compared against the full mode across all scenarios
(Section~\ref{sec:ablation}); primary system-level empirical
validation via Experiment~1 (Section~\ref{sec:exp1})
across penetration rates and intensity levels,
and Experiment~4 (Section~\ref{sec:exp4}) across threat levels.)

\item[\textbf{RQ5.}] Does the post-quantum TRS scheme structurally reduce
the fraction of poisoned aggregates reaching the SDN controller independent
of AI detection accuracy, and what are the associated computational and
communication overheads at the RSU relative to classical signing?
(Maps to: PQC-TRS and FHE mitigation pipeline; evaluated by
PARR (Eq.~\ref{eq:parr}) of ablation variant AB6 (TRS disabled)
vs.\ full mode; RSU-side cryptographic overhead reported as
COO (Eq.~\ref{eq:coo}) compared against ECDSA baseline;
communication overhead reported as BWO\textsubscript{ratio}
(Eq.~\ref{eq:bwo_ratio}) and BWO\textsubscript{scale}
(Eq.~\ref{eq:bwo_scale}) against the LKH unicast baseline
(ablation variant AB11); primary system-level scalability
validation via Experiment~3 (Section~\ref{sec:exp3},
$N_v \in \{100, 200, 300, 400\}$).)

\item[\textbf{RQ6.}] Does the blockchain smart contract trust management
layer --- specifically the BFT-conditioned SC-Revoke contract --- reduce
the control decision error rate and sustained poisoning impact compared to
frameworks relying on centralized controller-side revocation, and does it
do so without incorrectly revoking or demoting honest entities or
introducing unacceptable transaction latency for real-time ITS operation?
(Maps to: blockchain smart contract contribution; evaluated by
CDER (Eq.~\ref{eq:cder}) and false revocation rate
FRR\textsubscript{revoke} (Eq.~\ref{eq:frr_revoke}) and false
demotion rate FRR\textsubscript{demote} (Eq.~\ref{eq:frr_demote})
of ablation variant AB10 (blockchain removed) vs.\ full mode;
on-chain transaction confirmation latency
TCL\textsubscript{confirm} (Eq.~\ref{eq:tcl_confirm}) and
post-revocation reassignment latency
TCL\textsubscript{reassign} (Eq.~\ref{eq:tcl_reassign})
reported as framework-internal overhead measures; ablation
variants AB8 (RSU lifecycle removed) and AB9 (multi-controller
removed) isolate sub-components of the trust management layer;
primary per-attack-class validation via Experiment~5
(Section~\ref{sec:exp5}).)

\end{description}

\section{Contributions}

\begin{enumerate}
\item \textbf{Multi-Source Dual-Plane Threat Model.}
We propose a comprehensive threat model for mobility pattern and trajectory
poisoning attacks in SDVN that considers four simultaneous attacker classes
--- malicious vehicles, compromised RSUs, MitM adversaries, and malicious
or compromised controllers --- operating across both the data plane (beacon
injection and interception) and the control plane (learning model corruption
and control decision manipulation). Unlike prior works that assume a trusted
controller~\cite{raja2021, ghaleb2014}, confine attackers to
the data plane~\cite{hong2015}, or treat poisoning as a
single-source event, this threat model captures the full adversarial
surface of learning-based SDVN and directly motivates the distributed
trust and dual-plane mitigation design of the proposed framework.

\item \textbf{Mobility-Induced Threat Amplification Formulation.}
We formally quantify how vehicle kinematics amplify poisoning attack
effectiveness in SDVN by deriving closed-form expressions for
speed-dependent RSU contact duration (Eq.~\eqref{eq:contact_duration}),
beacon sampling window (Eq.~\eqref{eq:beacon_count}), handover-induced
learning discontinuity (Eq.~\eqref{eq:handover_rate}), mobility-disguised
drift rate (Eqs.~\eqref{eq:drift_constraint}--\eqref{eq:cumulative_drift}),
and Sybil identity capacity (Eqs.~\eqref{eq:sybil_capacity}--\eqref{eq:sybil_fraction}).
This is the first formal characterization of mobility as an attack
amplification vector in SDVN, showing that high-speed scenarios
simultaneously reduce detection evidence and enlarge the adversarial
action space.

\item \textbf{Attack Signature Characterization for Principled Detection.}
We derive nine formally defined attack signatures
(TP-S1 through MP-S4) that characterize the observable fingerprints of
trajectory poisoning and mobility pattern poisoning in the BSM beacon
stream (Eqs.~\eqref{eq:tp_s1}--\eqref{eq:mp_s4}), forming a composite
rule-based anomaly score (Eq.~\eqref{eq:psi_score}). Unlike existing
threshold-based methods~\cite{ghaleb2014} that apply isolated
checks, our signature set is derived from the full kinematic feasibility
space and includes cross-identity coordination tests absent from all
reviewed works.

\item \textbf{Dual-Mode Detection Framework.}
We design a lightweight operating mode executing entirely at the RSU
using HMAC beacon integrity verification and rule-based signature
detection (Algorithm~\ref{alg:lightweight}), and a full operating mode
executing GAT-based spatial anomaly detection over dynamic vehicle--RSU
graphs (Eqs.~\eqref{eq:edge_def1}--\eqref{eq:multihead_embed1}) and LSTM
temporal autoencoder sequential anomaly scoring
(Eqs.~\eqref{eq:ta_input}--\eqref{eq:ta_threshold}) at the SDVN
controller (Algorithm~\ref{alg:full_detection}). Detection scores from
both components are fused into a single decision variable
(Eq.~\eqref{eq:fusion}). No prior SDVN security framework provides
both a deployment-feasible lightweight mode and a high-assurance full
mode evaluated comparatively.

\item \textbf{Post-Quantum FHE-then-TRS Mitigation Pipeline.}
We design a post-quantum cryptographic mitigation pipeline in which
each RSU first encrypts its local aggregate under post-quantum
FHE (BFV/BGV/CKKS via OpenFHE, Eq.~\eqref{eq:fhe_rsu_enc})
before any ring coordination, preventing compromised ring
members from observing neighbouring RSU aggregate contents during
coordination a privacy threat unique to multi-RSU threshold
signing. Ring members then perform homomorphic addition across the
received ciphertexts (Eq.~\eqref{eq:fhe_ring_combine}), and a
post-quantum Threshold Ring Signature (TRS) scheme requiring
$t_{\mathrm{sign}} = f{+}1$ collaborative RSU signing
(Eqs.~\eqref{eq:trs_partial}--\eqref{eq:trs_verify}) is applied to
the combined encrypted aggregate, so that no plaintext crosses
any RSU boundary at any point and any minority-compromised RSU subset
(of size $< t_{\mathrm{sign}}$) is structurally prevented from
injecting a valid poisoned aggregate into the controller, independent
of AI detection. The cloud performs further blind homomorphic
computations across combined encrypted aggregates without decryption
(Eq.~\eqref{eq:fhe_global}), following feasibility demonstrated
in~\cite{mamun2025postquantum}. This is the first framework to
identify and address cross-RSU aggregate leakage during threshold
ring signing coordination using pre-coordination FHE in SDVN.

\item \textbf{Smart Contract-Based Byzantine Fault-Tolerant Trust Management.}
We design two smart contracts
SC-Trust for continuous trust score management across all three
entity types (vehicles: Eq.~\eqref{eq:trust_update};
RSUs: Eq.~\eqref{eq:rsu_trust};
controllers: Eq.~\eqref{eq:ctrl_trust}),
and SC-Revoke for RSU-direct BFT-conditioned identity revocation
(Eq.~\eqref{eq:bft_revoke_rsu}) that excludes the controller
as an evidence source, preventing a compromised controller from
suppressing legitimate revocations or fabricating false ones
on a permissioned consortium blockchain with IPFS off-chain beacon
storage, bridging real-time per-beacon detection and offline
controller-side deep learning analysis.
SC-Trust enforces a three-state lifecycle on each RSU trusted,
probationary, and client driven by a peer quorum disagreement
signal (Eq.~\eqref{eq:rsu_misbehave}) that prevents demoted RSUs
from inflating the BFT evidence count and maintains the
$|\mathcal{R}_{\mathrm{trusted}}(t)| \geq 3f{+}1$ safety floor
at all times.
When a controller's trust score falls below the revocation threshold,
SC-Revoke updates the on-chain trusted-controller set
$\mathcal{C}_{\mathrm{trusted}}(t)$ (Eq.~\eqref{eq:trusted_ctrl_set})
and RSUs automatically reassign to any available trusted controller
without manual failover.
Unlike prior blockchain applications in vehicular
networks~\cite{sehar2023blockchain, yahiatene2018towards} that
use blockchain only for data storage, our smart contracts encode
the full detection-to-revocation state machine for vehicles, RSUs,
and controllers with unified BFT safety guarantees parameterised
by a single security parameter $f$.

\end{enumerate}

\section{Novelty Statement}
\label{sec:novelty}

This work advances SDVN security on two levels. At the problem level, we are the first to formally characterize mobility pattern and trajectory poisoning as a mobility-amplified, multi-source, dual-plane threat, deriving closed-form bounds that show how vehicle speed, handover rate, and traffic density simultaneously shrink the detection evidence window and enlarge the adversarial drift budget, and proposing a threat model that---unlike all reviewed works---treats malicious vehicles, compromised RSUs, MitM adversaries, and compromised controllers as concurrent poisoning sources across both data and control planes.

At the solution level, we are the first to propose a unified framework combining four independently novel contributions. First, an attack-signature-derived dual-mode detection framework pairing lightweight rule-based signature scoring at the RSU with a full-mode multi-task GAT---trained on $K=7$ parallel attack-type classification heads sharing a common spatial encoder---and an LSTM temporal autoencoder at the controller, where the predicted attack type selects attack-conditioned fusion weights and signature reliability weights at inference, giving each detector component the discriminative emphasis appropriate for the identified attack class. Second, a pre-coordination FHE-blinded post-quantum threshold ring signature pipeline in which RSUs encrypt their local aggregates under FHE before ring coordination begins, eliminating cross-RSU aggregate leakage during share exchange, with TRS signing applied to the homomorphically combined ciphertext so that minority-compromised RSU subsets are structurally blocked from assembling a valid threshold signature and their aggregates never reach the controller, independent of AI detection outcome. Third, a three-state RSU trust lifecycle driven by a quorum-disagreement misbehavior signal evaluated against the currently trusted RSU set only, preventing demoted RSUs from continuing to inflate the BFT evidence count while their cryptographic keys remain valid. Fourth, a permissioned blockchain with smart contracts encoding the complete detection-to-revocation state machine under RSU-direct Byzantine fault-tolerant consensus, where controller compromise is identified by a directional conflict indicator that fires when the controller suppresses anomaly alerts independently raised by RSUs---penalizing alarm suppression, not disagreement---and revocation is triggered by RSU peer votes submitted directly to the blockchain without controller relay, with automatic RSU reassignment to the on-chain trusted-controller set $C_{\text{trusted}}(t)$, collectively eliminating the centralized trust bottleneck including controller-level compromise that every reviewed defense retains.

\section{Implications of the Study}

The results of the research have many implications for research and practice related to software-defined vehicular networks (SDVN). This research clearly articulates the importance of considering the patterns of mobility and trajectory poisoning as multi-source and time-evolving attacks, which brings into focus the issue of security that is currently unattended by the available solutions for the security of SDVN.

In practice, this proposed framework provides design insights for designing resilient SDVNs that do not rely on a trusted controller. The combination of blockchain trust, spatial learning, and temporal anomaly detection concepts can be used for future intelligent transportation systems to improve safety, reliability, and security from insider threats and controller-level attacks.

From a research perspective, this paper opens new research avenues for studying the issue of security threats in vehicular networks based on learning, specifically concerning dynamic attacks. The proposed approach for the problem definition and solution can be applied in general SDN-based networks that make control decisions based on the data obtained by learning.

\section{Organization of the Report}

The remainder of this report is organized as follows. Chapter 2 presents a detailed review of existing literature related to SDVN security, mobility data validation, and poisoning attacks, and identifies the research gaps addressed in this work. Chapter 3 describes the proposed methodology, including the system model, threat model, attack scenarios, and the proposed mitigation framework. Chapter 4 outlines the evaluation plan used to assess the effectiveness of the proposed approach. Finally, Chapter 5 concludes the report and discusses future research directions.


\chapter{Literature Review}
\label{chap:litreview}

Software Defined Vehicular Networks (SDVNs) are viewed as an upcoming solution to the complex problems associated with the conventional vehicular network, especially those pertaining to the dynamic nature of the network and the efficient use of network resources. Although the benefits of the network are many, it poses serious risks, especially concerning the integrity of the data associated with the network's mobility. This chapter reviews the current state of the art to determine the shortcomings of the techniques employed.

\section{Background}
\label{background}
Software-Defined Vehicular Networks is the extension of the traditional SDN concept to mobile networks by splitting the control and data planes, which allows the centralized management of the network by means of SDN Controllers communicating with the vehicle and Road Side Units through the OpenFlow protocol. The proposed defense system combines the strengths of three different techniques, namely the Graph Attention Networks, which learn the spatial relationships among the vehicle network by assigning importance weights to the neighboring nodes through the use of self-attention mechanisms, which facilitate the identification of the coordinated attacks, the Temporal Autoencoders, which encode the trajectory data into the latent space to detect the sequential attacks through the analysis of the reconstruction error, and the Blockchain technology, which provides the distributed consensus and the immutable ledger to overcome the single point of failure of the centralized architecture. This combination of techniques can overcome the inherent limitations of the present state of the art of SDVN security solutions, which are based on the centralized trust management, the analysis of the vehicle data without considering the spatial patterns, and the rule-based validation, which is incapable of identifying the complex sequential poisoning attacks.

\section{Previous Work}
\label{sec:previous-work}

\subsection{SDVN Architecture and Security Foundations}
\label{subsec:sdvn-arch-security}

The three-layered SDVN architecture comprising infrastructure plane (vehicles and RSUs), control plane (SDN controllers), and application plane (network services) was comprehensively documented by \cite{bhatia2019}. Their survey showed how OpenFlow enables treating vehicles and RSUs as programmable switches, providing centralized control that dramatically improves resource management and handles network heterogeneity across DSRC, cellular, and WiFi technologies. Despite these architectural advantages, the authors made a critical observation: ``security is still the major fear that is preventing the SDVN from becoming widely acceptable'' \cite{bhatia2019}.

The security vulnerabilities they identified are particularly relevant to this research. Beyond the well-known DDoS attacks on controllers and flow table overflow issues, they documented the complete absence of trustworthiness evaluation mechanisms for participating vehicles. More critically, their survey revealed three gaps that directly motivate the current work:

	 No \acs{AI} or machine learning-based security mechanisms had been proposed for SDVN security \cite{bhatia2019},
	 Graph-based network analysis for attack detection was entirely missing \cite{bhatia2019},
	 No specific defenses existed against trajectory poisoning attacks \cite{bhatia2019}.

The trajectory-prediction approach introduced by \cite{he2016} demonstrates both the potential and the security risks of mobility-aware SDVN management. Their system extended OpenFlow to include position, velocity, and heading data, allowing the Status Manager to predict future vehicle locations and the Topology Manager to pre-install flow rules proactively. In Network Simulator 3 integrated with SUMO traffic simulation, this approach achieved 10-15\% better packet delivery by switching between OLSR (for dense, slow traffic) and GPSR (for sparse, fast-moving scenarios). However, their architecture diagram included a ``security module'' that was never implemented, evaluated, or even described. The entire prediction mechanism operates on a dangerous assumption: that all vehicles honestly report their positions and velocities. This scenario clearly provides an attack vector: malicious vehicles may provide erroneous trajectory information to impact network dynamics; but in reference \cite{he2016}, there are no methods for identifying or mitigating this tampering. In survey article \cite{ravi2023}, 164 research papers on intelligent SDVN were surveyed, presenting novel ideas such as Digital Twin Networks for real-time observation and prediction. While they have thoroughly described ML-based applications for routing improvement and QoS improvement, all applications considered gave priority to performance improvement over security. They have accepted "security gaps in SDVN due to centralized architecture" in both abstract and discussion sections but failed to provide any solutions for this issue. In this survey article as well, there were no descriptions regarding blockchain-based implementations, graph-based neural networks for pattern observation, or time-based models for trajectory validation.
\subsection{Security Threats in SDVN}
\label{subsec:security-threats}

The literature provides the most detailed taxonomy of threats for SDN-based VANETs, which covers 19 types of attacks that include the application, control, and data layers. This paper is very significant in that it points out the exact problem that this research addresses. The authors claimed that "false position information" and "illusion attacks," which spread fake position information, are challenging and crucial issues in SDN-based VANETs. The authors explained that some vehicles that move slowly send fake information about their positions in order to create virtual congestions that affect the routing path of other vehicles. The authors did not propose any solution for this problem but just analyzed it in their paper \cite{arif2020}.

Their results reveal a critical flaw in current cryptographic protections. Public Key Infrastructure and authentication through certificates are able to ensure a message is sent from a genuine vehicle -- namely, they are able to verify the message's sender is authentic. But they are not able to ensure the message is genuine. An authentic vehicle with genuine credentials can transmit spurious GPS location information, spurious speeds, and entirely fabricated tracks, and these will withstand any and all cryptographic verification. Authentication therefore verifies authenticity, not truthfulness.

The difference pointed out in \cite{arif2020} between network topology and mobility topology is very important.They discussed topology poisoning attacks and referenced solutions like TopoGuard, but these protect network connectivity (which devices connect to which switches), not geographic movement patterns. A vehicle could maintain correct network attachments while systematically falsifying its location, speed, and trajectory. The authors acknowledged this gap but left it completely unaddressed.

\subsection{Network Topology Defense}
\label{subsec:network-topology-defense}

Literature exposed critical vulnerabilities in SDN topology management and demonstrated attacks against all eight major controllers (Floodlight, OpenDaylight, Beacon, POX, NOX, Ryu, Maestro, OpenIRIS) \cite{hong2015}. Their Host Location Hijacking attack spoofed host identifiers to falsify switch port attachments, while their Link Fabrication attack created illusions of non-existent switch connections using forged LLDP packets. Both succeeded because topology update mechanisms lacked authentication and integrity protection \cite{hong2015}.

TopoGuard, their defense system implemented in 700 lines of Java within Floodlight, provides three key protections:

\begin{enumerate}
	\item \textbf{Host migration verification}: Requires Port Down signals before accepting movement (precondition) and uses ICMP probes with 1-second timeout to confirm hosts disappeared from previous locations (post condition)
	\item \textbf{LLDP authentication}: Applies HMAC-SHA256 signatures to protect switch-to-switch discovery
	\item \textbf{Port property management}: Tracks whether each port connects to switches or hosts, preventing hosts from participating in LLDP propagation
\end{enumerate}

Their evaluation showed effective attack prevention with only 0.005-0.431ms overhead per operation \cite{hong2015}.

However, TopoGuard's scope differs fundamentally from mobility attack defense. In their system, host ``location'' means network attachment point: which switch and which port. Their Host Profile data structure contains MAC address, IP address, VLAN, switch DPID, and port number. Notice what's missing: no GPS coordinates, no velocity vectors, no heading information, no trajectory history. TopoGuard cannot detect an adversary who injects false geographic positions showing a vehicle kilometers from its actual location, fabricates velocities claiming impossible speeds, or constructs fake trajectory sequences. The distinction is categorical: TopoGuard addresses topology structure poisoning (network connectivity) but not data content poisoning (falsified mobility information) \cite{hong2015}.
\newline
Moreover, TopoGuard concentrates all trust in the SDN controller. All verification, HMAC signature generation, state machine management, and ICMP probe execution happen within the controller using controller-held secret keys. If the controller is compromised through software vulnerabilities, insider access, or physical attacks, the entire security system fails because there's no distributed verification or redundancy. \cite{hong2015} acknowledged that powerful adversaries who completely mute host-generated traffic could evade detection, and that physical-layer verification lies beyond their scope.

\subsection{Authentication and Intrusion Detection}
\label{subsec:authentication-ids}

Raja et al.\cite{raja2021} developed a two-level security architecture for SDVN combining RSU-based Group Authentication using Elliptic Curve Cryptography with a private Collaborative Intrusion Detection System. Their RGA scheme reduces per-message overhead by providing group identity-key pairs to authenticated vehicles, enabling efficient V2V communication. The collaborative IDS uses Alternating Direction Method of Multipliers for distributed learning across RSUs, with differential privacy adding noise to protect vehicle information and Paillier homomorphic encryption enabling computation on encrypted data. They achieved 96.81\% detection accuracy on the NSL-KDD dataset \cite{raja2021}.

The critical limitation lies in what they're detecting. The paper explicitly acknowledges that ``some rogue vehicles after authentication pose as legitimate VANET users and send messages to interrupt network communication'' \cite{raja2021}---exactly the insider threat scenario where authenticated vehicles inject false data. But their solution uses generic network intrusion detection, not mobility-specific validation. The NSL-KDD dataset contains TCP SYN floods, port scans, buffer overflows---computer network intrusion signatures, not vehicular mobility patterns. An IDS trained to recognize packet-level attacks won't recognize when a vehicle reports a position on a valid road segment, a velocity within normal ranges, and a plausible acceleration, yet these individually valid values form an impossible trajectory sequence.

Furthermore, their privacy mechanisms don't address data integrity. Differential privacy adds noise to protect training data but doesn't verify whether the underlying data is genuine or poisoned---it just makes the poisoning private. Homomorphic encryption enables secure computation but doesn't validate data authenticity: a vehicle can encrypt false trajectories, and the encrypted computation processes that false data without detection. The detection operates reactively, identifying attacks after they manifest as patterns in network traffic, with no real-time evaluation of incoming mobility data before network use \cite{raja2021}.

\section{Related Work}
\label{sec:related-work}
\subsection{Mobility Pattern Analysis}
\label{subsec:mobility-pattern-analysis}
\justifying
This approach pioneered using mobility patterns themselves for misbehavior detection . Their system constructs Location Time Tables from periodic beacons containing position and timestamp information, giving each vehicle distributed local knowledge of surrounding vehicles' claimed locations. The detection framework has four modules working sequentially \cite{ghaleb2014}:

\begin{enumerate}

	\item \textbf{Security verification}: Authenticates message origins using group signatures \cite{ghaleb2014}.
	\item \textbf{Location verification}: Checks four conditions: positions must fall within road map boundaries (from Google Maps or OpenStreetMap), lie within transmission range (1km WAVE limit), not overlap with other vehicles' positions (preventing Sybil attacks), and satisfy speed-of-light constraints\cite{ghaleb2014}.
	\item \textbf{Movement analysis}: Compares current versus previous views to detect abnormal area differences indicating suspicious position changes\cite{ghaleb2014}.
	\item \textbf{Plausibility detection}: Validates speeds against physical constraints (0-240 km/h range, road-specific limits)\cite{ghaleb2014}.
\end{enumerate}

With an 80-car simulation, it was possible to differentiate between real and fake Post-Crash Notifications by observing their behavioral patterns together. If the cars were receiving real alerts of an accident, their trajectory plots showed joint behavior such as changing lanes, slowing down, and reacting together. However, fake alerts showed no change in behavioral patterns because drivers do not act on information that is not visible to them \cite{ghaleb2014}.

The authors themselves admitted the use of "a simple and coarse grain method in classifying good and malicious vehicles" \cite{ghaleb2014}. The admission of this fact is quite surprising as it reveals some limitations of the approach in question. The detection system is completely dependent on thresholds, which could be easily bypassed by a sophisticated malicious user. A malicious user knowing about a threshold of 0-240 km/h in terms of speed could report a speed of 110 km/h, which is still valid but could be a false report. If a malicious user is aware of road-map verification, they could feed false location information on valid roads. By knowing about constraints related to the speed of light, a malicious user could guarantee consistency in distance and time relationships. In these scenarios, the system could verify every value independently without realizing that all values represent a malicious path in combination. The approach examines vehicles individually (location verification and plausibility verification) or compared to a previous view in a pairwise fashion; however, it does not use a graph relational modeling approach. If multiple malicious vehicles coordinate to inject mutually consistent false data, where each vehicle's claims support the others' fabrications, individual checks might validate each separately without detecting the collectively impossible spatial configuration. The paper explicitly states that ``vehicles need to wait until they get enough evidence in its LTT table'' \cite{ghaleb2014}, introducing delays dangerous for split-second safety decisions.

Additionally, the system relies on centralized RSU trust for road map distribution and clock synchronization. A compromised RSU could distribute falsified maps, causing legitimate vehicles reporting correct positions to be flagged as misbehaving when their genuine coordinates fall outside the false map boundaries. Report misbehavior from Vehicle to RSUs at journey's end, delayed and potentially lossy if communication fails or malicious RSUs selectively ignore reports.

Aside from security threats in terms of mobility patterns in SDN-VANETs, other security concerns have been considered by researchers. Arif et al. \cite{arif2020} have presented a comprehensive classification of security threats in SDN-VANETs from DoS to GPS spoofing in different planes of a network. While they have considered false location data as a major security threat, their survey did not include measures to counter such threats. Likewise, Hong et al. \cite{hong2015} have presented topology poisoning attacks in SDN controllers in which attackers create network links or misrepresent host location data. They have developed TopoGuard to counter such attacks through cryptographic validation but considered "location" as a network attachment point instead of geographic location to preserve mobility data integrity.

A different approach is taken by Raja et al. \cite{raja2021} for collaborative intrusion detection in the SDVN, obtaining excellent performance through distributed machine learning while maintaining privacy. Nonetheless, their IDS is based on common network traffic patterns, but not vehicle mobility behavior, and is reactive, but not proactive. Taken as a whole, these papers demonstrate sound bases for secure SDVN but indicate the need for ensuring the integrity of the mobility information itself, as cars may be verified as authentic but broadcast erroneous paths.

Yamamoto et al.~\cite{yamamoto2020modeling} proposed a velocity vector field
combined with a variance tensor field to model and predict manned-vehicle
trajectories in open environments, enabling collision-risk estimation and
anomaly detection when observed trajectories deviate from predicted occupancy
regions. While this work introduces valuable trajectory prediction primitives
applicable to moving vehicles, it operates in an isolated sensing context with
no network adversarial model, no multi-source attacker consideration, and no
connection to SDN-controlled vehicular infrastructure. It cannot detect
coordinated multi-vehicle poisoning, and the variance tensor formulation
assumes honest sensor inputs, rendering it ineffective when the trajectory
data itself is the attack surface.

Ercan et al.~\cite{ercan2022misbehavior} proposed a machine-learning-based
misbehavior detection scheme for position falsification attacks in VANETs,
deriving features such as position change and speed deviation from
exchanged beacons and training supervised classifiers on the VeReMi
dataset. While effective at distinguishing position falsification
variants, the approach classifies beacons on individual per-message
features without spatial-relational modeling across vehicles, does not
address coordinated multi-source poisoning, operates as offline
classification rather than a real-time SDVN control-plane defense, and
assumes a trusted detection authority, leaving the centralized trust
bottleneck of SDVN unaddressed.


\subsection{Detection and Mitigation Techniques in Related Domains}
\label{sec:related_techniques}

\subsubsection{Graph Attention Networks for Intrusion Detection}
Frenken et al.~\cite{frenken2025kdgat} proposed KD-GAT, representing
Controller Area Network (CAN) bus traffic as sliding-window graphs and
training a multi-layer Graph Attention Network with jumping-knowledge
aggregation, then distilling it into a student model only 6.32\% the
teacher's parameter count, achieving 99.97\% detection accuracy on the
Car-Hacking dataset at resource-constrained deployment scales. Their
subsequent work~\cite{frenken2025multistage} further introduces a
multi-stage pipeline using a Variational Graph Autoencoder (VGAE) for
unsupervised anomaly pre-filtering followed by a knowledge-distilled GAT
classifier, improving F1-score by up to 55\% on highly imbalanced CAN
intrusion datasets. These works validate the power of GAT architectures
for detecting coordinated anomalies in graph-structured communication
systems but are limited to in-vehicle CAN networks and do not generalize
to dynamic, geographic SDVN vehicle-RSU graphs where node membership
changes continuously due to mobility.

\subsubsection{Data-Centric Machine Learning for Vehicular Misbehavior Detection}
Sharma and Liu~\cite{sharma2021machine} proposed a data-centric
misbehavior detection model for the Internet of Vehicles that engineers
plausibility and consistency features from basic safety messages and
evaluates multiple supervised classifiers on the VeReMi dataset,
achieving strong position-falsification detection without relying on
vehicle hardware trust. While effective as a data-centric detector, the
model classifies each message largely in isolation, captures neither
temporal trajectory drift across beacon sequences nor spatial
coordination among vehicles, and provides no distributed trust or
mitigation pathway, remaining a reactive centralized detector.

\subsubsection{Blockchain-Based Trust and Distributed Consensus in Vehicular Networks}
Yahiatene and Rachedi~\cite{yahiatene2018towards} combined SDN and
blockchain for vehicular social network security through a hierarchical
controller architecture, using a Distributed Miners Connected Dominating Set
algorithm to dynamically select trusted miners and eliminate the SDN
controller as a single point of failure. Sehar et al.~\cite{sehar2023blockchain}
proposed VBCA, a consortium blockchain consensus replacing Proof-of-Work with
a lightweight protocol reducing transaction latency for V2V/V2I data integrity.
These works demonstrate blockchain's viability in vehicular networks but do
not integrate smart contract-based anomaly scoring or address mobility
data poisoning at the learning model level.

\subsubsection{Threshold Ring Signatures and Post-Quantum Cryptography for VANET}
Ahmed et al.~\cite{ahmed2022blockchain} combined a threshold ring signature
(TRS) scheme requiring $t$-of-$n$ cooperative signing without revealing
individual identities with a blockchain incentive trust model to validate
traffic events in VANETs, demonstrating resistance to collusion attacks.
Zhong et al.~\cite{zhong2022broadcast} proposed identity-based broadcast
encryption for V2I communication enabling a single ciphertext to serve a
vehicle group, significantly reducing per-message overhead. Faleh et
al.~\cite{faleh2025csasv} proposed CSAS-V, a pairing-free certificateless
Schnorr aggregate signature scheme for scalable VANET authentication. While
these cryptographic advances address message authenticity and group privacy,
none provide post-quantum security guarantees. Mamun et al.~\cite{mamun2025postquantum}
provided the first experimental evaluation of lattice-based homomorphic
encryption schemes (BFV, BGV, CKKS via OpenFHE) for privacy-preserving
Infrastructure-to-Infrastructure communication in ITS, demonstrating
feasibility under 128-bit post-quantum security with RSU-Cloud link
benchmarks. Their work directly motivates the use of post-quantum fully
homomorphic encryption for secure aggregate transmission in the proposed
framework, though it addresses only privacy and provides no attack detection
capability.

\newpage
\section{Gaps in Literature}
\label{sec:gaps-literature}

\subsection{Literature Gap Analysis}
\label{subsec:gap-analysis}

\subsubsection{The Missing Link: Mobility Data Content Verification}
\label{subsec:missing-link}

After reviewing all this previous work, one gap stands out above all others: no existing mechanism specifically verifies the correctness of mobility data content. Multiple researchers acknowledge the problem: \cite{arif2020} call it a ``hard and critical issue,'' \cite{hong2015} demonstrate topology poisoning, \cite{raja2021} acknowledge ``rogue vehicles after authentication'', but none provide comprehensive operational solutions for adversaries who systematically inject false GPS coordinates, fabricate velocities, or construct fictitious trajectories.

The problem manifests in several attack scenarios that current defenses cannot handle.

Current cryptographic schemes prove identity but not data truthfulness. A vehicle's certificate confirms authorization to participate but provides no validation that claimed positions correspond to actual locations, that reported velocities match actual speeds, or that trajectory sequences represent actual paths rather than strategic fabrications \cite{arif2020}.

\subsubsection{Centralized Trust Creates Single Points of Failure}
\label{subsec:centralized-trust}

Every system examined relies on centralized trust. SDN controllers centralize all topology management, flow rule installation, and routing decisions \cite{bhatia2019}. If compromised, adversaries can poison topology databases, install malicious flow rules, suppress security alerts, and falsify statistics without remaining validation mechanisms.

RSU-based architectures show parallel vulnerabilities. \cite{ghaleb2014} and \cite{raja2021} both trust RSUs for certificate distribution, road map provisioning, and clock synchronization. Compromised RSUs could distribute falsified maps, selectively ignore misbehavior reports, or inject false timestamps. \cite{raja2021}'s collaborative IDS assumes honest majority, but becomes vulnerable when attackers compromise sufficient RSUs to poison training data or vote for accepting attack patterns.

Blockchain technology offers distributed consensus eliminating single trusted entities through Byzantine fault tolerance. Multiple surveys mention blockchain as ``future work'' \cite{bhatia2019,arif2020,ravi2023}, but none provide architectural designs showing how to integrate blockchain consensus into SDVN mobility validation, protocol specifications for vehicle/RSU/controller participation, or performance evaluations demonstrating whether blockchain latency can meet real-time requirements.

\subsubsection{Individual Analysis Misses Coordinated Attacks}
\label{subsec:individual-analysis}

Current approaches analyze vehicles in isolation or through pairwise comparisons, completely missing the rich relational structure in vehicular networks. Vehicle movements are fundamentally interdependent: following vehicles can't overtake without lane changes, parallel-lane vehicles show correlated speed adjustments, intersection approaches are constrained by signals and right-of-way, highway platoons exhibit characteristic spacing and synchronized patterns \cite{ghaleb2014,he2016}.

Graph Attention Networks (GATs) provide a natural architecture for modeling these dependencies. GATs can learn which spatial configurations represent normal traffic: vehicles in adjacent lanes traveling at similar speeds, safe following distances determined by speed and conditions, platoon coordination patterns.

Attention mechanisms let the network automatically focus on most relevant neighbors without manual specification. Different scenarios require different relationships: on highways, leading/following vehicles matter most; at intersections, crossing paths become critical; in parking areas, nearby stationary vehicles are relevant.

Despite this natural fit and GAT's proven success in social network analysis, molecular property prediction, and traffic forecasting, no SDVN research applies GATs to mobility pattern analysis \cite{bhatia2019,ravi2023}. This is a missed opportunity: these architectural advances suit a fundamentally graph-structured problem. The absence prevents detecting coordinated attacks where multiple adversaries inject mutually consistent but collectively implausible data.

\subsubsection{No Temporal Deep Learning for Sequential Patterns}
\label{subsec:no-temporal-dl}

Vehicle trajectories are inherently sequential time-series data where temporal dependencies carry crucial legitimacy information. A single position provides minimal validation context beyond boundary checks. But a sequence over 10-30 seconds reveals movement patterns distinguishing genuine behavior from fabrications \cite{ghaleb2014,he2016}.

TopoGuard concentrates all verification in the controller using controller-held HMAC keys and controller-managed state machines \cite{hong2015}, such that controller compromise defeats all security simultaneously. RSU-based architectures show parallel vulnerabilities where compromised RSUs can distribute falsified maps or poison collaborative IDS training data \cite{raja2021}, \cite{ghaleb2014}.

These architectures have proven success in financial fraud detection (anomalous transaction patterns), healthcare (unusual patient trajectories), and cyber security (network traffic sequences). Yet their application to vehicular trajectory validation remains unexplored in SDVN security literature \cite{bhatia2019,ravi2023}.

Existing prediction work uses classical models: linear extrapolation, Kalman filtering, traffic flow models \cite{he2016}. These assume honest input and optimize performance (reduce overhead, enable proactive rules), not security. They predict likely future positions to improve efficiency, not detect whether historical positions are genuine or falsified. \cite{raja2021} used logistic regression on packet features, not deep sequential models for trajectory time-series.

\subsubsection{Reactive Detection Allows Damage before Response}
\label{subsec:reactive-detection}

Existing systems identify attacks only after malicious data has been injected, propagated, and potentially influenced decisions \cite{ghaleb2014,raja2021,hong2015}. Mobility pattern analysis requires accumulating ``sufficient evidence in its LTT table'' over multiple beacon intervals \cite{ghaleb2014}. Intrusion detection matches behavior to learned signatures after attacks manifest in traffic patterns. Topology poisoning detection validates updates after receiving them, with Hong et al.'s ICMP probing adding 1 second timeout.

Proactive defense requires different architectures. Rather than waiting for evidence accumulation, proactive systems would continuously score incoming data's anomaly level in real-time, immediately isolating suspicious data before operational impact. Rather than detecting only after damage manifests, they would learn normal patterns comprehensively and predict likely future states, flagging subtle deviations preemptively. Rather than reacting to known signatures from past attacks, they would adaptively learn from evolving threats through online learning, continuously updating models and adjusting thresholds dynamically.

No surveyed SDVN research demonstrates such proactive prevention where mobility data is evaluated and potentially quarantined at reception before any operational use \cite{bhatia2019,arif2020}.

\subsubsection{No Integration - Security Mechanisms Work in Isolation}
\label{subsec:no-integration}

Perhaps the most critical gap is the complete absence of comprehensive frameworks combining multiple complementary defenses \cite{bhatia2019,ravi2023}. Each work addresses one isolated aspect:

\begin{itemize}
	\item Cryptographic authentication of identities \cite{raja2021}
	\item Network intrusion detection through collaborative learning \cite{raja2021}
	\item Mobility pattern threshold checking \cite{ghaleb2014}
	\item Network topology validation via LLDP verification \cite{hong2015}
\end{itemize}

No research synthesizes multiple layers into holistic systems where individual weaknesses are compensated by complementary strengths. \cite{bhatia2019} explicitly stated that existing solutions ``focus on detection without mitigation'' and lack ``comprehensive framework combining multiple defense layers.'' \cite{ravi2023} acknowledged needing to integrate ``blockchain, AI, and advanced analytics'' but provided no architectural designs, protocol specifications, or implementation guidance.

A truly comprehensive defense would require the integration of distributed trust (removing points of failure via blockchain consensus), spatial pattern analysis (graph attention networks identifying attacks via relational anomalies), sequence analysis (deep learning methods identifying poisoning attacks), and real-time proactive scoring (blocking dubious data use before validation is completed). This integrated synthesis has yet to be demonstrated in any current work.

\subsubsection{Detection Thresholds Calibrated Without Mobility Constraints}
\label{subsec:thresholds-no-mobility}

All threshold-based and signature-based detection works reviewed, including ~\cite{ghaleb2014}, calibrate detection parameters
(speed bounds, position deviation limits, beacon evidence accumulation windows)
under implicit assumptions of sufficient observation time per vehicle. None
of these works analyse how vehicle speed constrains RSU contact duration or
how handover frequency fragments trajectory sequences across RSU boundaries.
Consequently, thresholds derived in these works are systematically calibrated
against a static or low-mobility assumption that does not reflect real SDVN
operating conditions, leaving the detection logic untested and likely
inadequate in high-speed vehicular scenarios.

\subsubsection{Cryptographic Protections Restricted to Classical Hardness Assumptions}
\label{subsec:classical-hardness}

Every cryptographic scheme reviewed for vehicular network
security~\cite{faleh2025csasv,zhong2022broadcast, ahmed2022blockchain,
raja2021} relies exclusively on classical hardness assumptions, namely
elliptic curve discrete logarithm or integer factorisation problems. The
vehicular security literature has not evaluated whether these protections
remain viable under quantum adversaries, nor has any SDVN mobility data
protection mechanism been assessed for post-quantum resilience. This is
a recognised open problem in the broader cryptographic community yet has
received no attention in the SDVN mobility security sub-field.

\subsubsection{Absence of Deployment Feasibility Evaluation Across Heterogeneous RSU Infrastructure}
\label{subsec:heterogeneous-rsu}

All reviewed frameworks, including~\cite{raja2021, ghaleb2014,
hong2015, yahiatene2018towards}, are evaluated under a single
uniform deployment assumption where all RSUs are presumed to have equivalent
computational resources. None of the reviewed works analyse how their
proposed detection or mitigation mechanisms perform under resource-heterogeneous
conditions where some RSUs may be computationally constrained legacy units
incapable of running deep learning inference or executing complex cryptographic
operations in real time. This omission means that the practical field
deployability of every reviewed solution remains entirely undemonstrated.

\newpage

\subsection{Literature Comparison}
\label{sec:literature-comparison}

\begin{table*}[!h]
\centering
\caption{Comparative Summary of Related Works Against the Proposed Framework}
\label{tab:comparison}
\renewcommand{\arraystretch}{1.2}
\resizebox{\textwidth}{!}{%
\begin{tabular}{|l|c|c|c|c|c|c|l|l|}
\hline
\textbf{Reference} &
\textbf{\shortstack{Mobility\\
Aware\\
Model}} &
\textbf{\shortstack{Distributed\\
Trust}} &
\textbf{\shortstack{Graph\\
Analysis}} &
\textbf{\shortstack{Temporal\\
Learning}} &
\textbf{\shortstack{Lightweight\\
Mode}} &
\textbf{PQC} &
\textbf{Attack Scope} &
\textbf{Key Limitation} \\
\hline
Ghaleb et al.~\cite{ghaleb2014} & Partial & No & No & No & Yes & No &
Position falsification, Sybil &
Threshold bypass; individual analysis only \\
\hline
Hong et al.~\cite{hong2015} & No & No & No & No & No & No &
SDN topology poisoning &
Network attachment only; not geographic mobility data \\
\hline
Raja et al.~\cite{raja2021} & No & No & No & No & No & No &
Generic network intrusion &
NSL-KDD dataset; not mobility-specific \\
\hline
Yamamoto et al.~\cite{yamamoto2020modeling} & Yes & No & No & Partial & No & No &
Single-vehicle trajectory anomaly &
No network adversary model; no SDVN context \\
\hline
Ercan et al.~\cite{ercan2022misbehavior} & Partial & No & No & No & No & No &
Position falsification &
Per-beacon ML features; individual analysis; no distributed trust \\
\hline
Sharma \& Liu~\cite{sharma2021machine} & Partial & No & No & No & No & No &
Position falsification (IoV) &
Message-level classification; no temporal/spatial or trust model \\
\hline
Frenken et al.~\cite{frenken2025kdgat} & No & No & Yes & No & Yes & No &
In-vehicle CAN bus intrusion &
In-vehicle only; not SDVN vehicle--RSU graph \\
\hline
Ahmed et al.~\cite{ahmed2022blockchain} & No & Yes & No & No & No & No &
Traffic event validation &
Classical TRS; no trajectory poisoning detection \\
\hline
Mamun et al.~\cite{mamun2025postquantum} & No & No & No & No & No & Yes &
I2I communication privacy &
Privacy only; no attack detection capability \\
\hline
\textbf{Proposed} & \textbf{Yes} & \textbf{Yes} & \textbf{Yes} & \textbf{Yes} &
\textbf{Yes} & \textbf{Yes} &
\textbf{\shortstack[l]{Trajectory poisoning,\\
Mobility pattern poisoning\\
(multi-source, dual-plane)}} &
--- \\
\hline
\end{tabular}%
}
\end{table*}

As shown in Table~\ref{tab:comparison}, the literature comparison reveals the following critical gaps addressed by the proposed framework:

\textbf{Lack of Distributed Trust:}  Most of the existing literature requires a centralized entity (for instance, RSU or honest majorities). The proposed system is the only one that tackles the single point of failure problem using blockchain consensus, thus providing a distributed trust system.

\textbf{No Graph-Based Analysis:}
 The existing literature does not represent the spatial relationship between vehicles as a dynamic graph. The proposed Graph Attention Network (GAT) module reduces the fabrication of coordinated mobility pattern data through learning pattern relationships in the network of vehicles.

\textbf{Lack of Temporal Learning:} Current solutions depend on static thresholds \cite{ghaleb2014} or primitive classifiers \cite{raja2021}, which are not capable of identifying complex pattern anomalies. The proposed method uses a temporal autoencoder to learn normal trajectory dynamics to identify gradual poisoning.

\textbf{Lack of Proactive Detection:}The major part of the literature reviewed is based on a reactive approach, detecting attacks only after malicious data affects the network. The framework provides real-time scoring of anomalies to remove malicious mobility data before it affects SDN control decisions.

\textbf{Lack of Comprehensive Integration:} The current state-of-the-art solutions tend to focus on individual parts such as authentication mechanisms, topology safeguarding, or traditional intrusion prevention. The proposed strategy will comprehensively integrate blockchain concepts for trust mechanisms, GAT for graph-based reasoning tasks, and temporal autoencoders for validation.

\section{Summary}
\label{sec:summary}

This chapter reviewed SDVN architecture, threats, and countermeasures to establish how existing research addresses mobility data integrity. The picture that emerges is consistent: although \cite{arif2020} explicitly calls these "hard and critical issues," current defenses against trajectory and mobility pattern poisoning remain inadequate, for six interrelated reasons. First, authentication verifies identity but not data validity, so an authenticated insider can still inject fake trajectories undetected. Second, most solutions assume a trustworthy controller and RSU, with no blockchain-based or otherwise distributed alternative addressing this single point of failure. Third, vehicle-level analysis misses coordinated attacks that only become visible through graph-based, relational methods. Fourth, no reviewed work applies deep sequential learning to trajectory time-series, leaving attacks undetected where each individual sample looks legitimate but the sequence as a whole is not. Fifth, detection across the field is reactive rather than proactive. Sixth, existing methods are developed independently rather than integrated into a single framework. No prior SDVN work combines blockchain-based distributed consensus, Graph Attention Networks for spatial analysis, and temporal autoencoders for sequential analysis to address trajectory poisoning; combining the three is the gap this thesis targets.


\chapter{Methodology}

\section{Research Design}

This research follows a modeling based methodology to design and analyze a learning based mitigation framework for mobility pattern and trajectory poisoning attacks in SDVN. The overall methodology consists of four main stages:
\begin{itemize}
    \item System modelling: SDVN architecture, entity definitions,
          and BSM beacon communication model
    \item Attack modelling: trajectory poisoning, mobility pattern
          poisoning, and control-plane attack formalisation
    \item Mitigation framework design: dual-mode detection pipeline
          with post-quantum cryptographic protection
    \item Evaluation planning: simulation setup, metrics, and
          baseline comparison
\end{itemize}

The proposed system is modeled as consisting of vehicles, Roadside Units (RSUs), and a distributed set of SDVN controllers $\mathcal{C} = \{c_1, c_2, \ldots, c_K\}$, none of which is individually authoritative. Vehicles periodically generate mobility reports containing position, speed, direction, and time information. These reports are forwarded to the controller through RSUs and are used by the controller to learn mobility models for network management and decision making.

\section{Threat Model}
\label{sec:threat_model}
This work considers a realistic threat model for SDVN, with a focus on poisoning attacks that target learning-based mobility and trajectory management. Unlike traditional attacks that disrupt communication, these attacks aim to corrupt the data used by the SDVN controller for learning and decision making, while normal communication behavior stay unchanged.

The SDVN architecture is assumed to consist of vehicles, Roadside
Units (RSUs), and a set of $K$ registered controllers
$\mathcal{C} = \{c_1, c_2, \ldots, c_K\}$, all of which may act
as potential attack sources. RSUs are initially assigned to
controllers based on geographic proximity or load-balancing
through standard SDVN mechanisms. The blockchain continuously
maintains trust scores for all registered controllers. When a
controller is detected as malicious, RSUs previously assigned to
it reassign to any currently trusted controller from $\mathcal{C}$
without disrupting network operation. Since the Cloud/ITS Server
feeds aggregate results into controller decision pipelines, it is
treated as part of the control plane threat boundary and is not
assumed to be unconditionally trusted.
Attacks may target both the data plane, where mobility and
trajectory information is transmitted, and the control plane, where
learning processes and control decisions are generated. It is also
assumed that controllers rely on learned mobility and trajectory
patterns for safety related and traffic management decisions.

The framework assumes the number of simultaneously compromised
RSUs $f$ satisfies $f \leq \lfloor (n-1)/3 \rfloor$, which is
equivalent to the standard BFT condition $n \geq 3f+1$ and
correctly tolerates $f$ Byzantine faults for all valid $n$.
The threshold parameters are set as $t_{\mathrm{sign}} = f+1$ for
TRS signing and $t_{\mathrm{decrypt}} = f+2$ for FHE decryption.
Setting $t_{\mathrm{decrypt}} = f+2$ ensures that $f$ RSUs and the
Cloud, totalling $f+1$ shares, cannot reach the decryption threshold
alone, preventing a coalition of $f$ compromised RSUs plus a
compromised Cloud from performing unilateral decryption. This
maintains consistency across TRS signing, FHE decryption, and BFT
revocation under a unified security parameterisation.

Under this threat model, attackers may be present at different components of the SDVN architecture. The following attacker types are considered:

\begin{itemize}
	\item \textbf{Malicious Vehicles:} Vehicles with valid identities that intentionally send false mobility or trajectory data to influence the learning process.
	\item \textbf{Compromised Roadside Units (RSUs):} RSUs that intercept, alter, or generate false mobility data before forwarding it to the controller.
	\item \textbf{Man in the Middle (MitM) Adversaries:} Attackers that intercept and modify mobility messages exchanged between vehicles and RSUs.
	\item \textbf{Malicious or Compromised Controllers:} SDVN controllers that deliberately poison their internal learning models and issue incorrect control decisions.
\end{itemize}

Attackers are assumed to be capable of launching gradual and time
series poisoning attacks, including identity impersonation and
Sybil type behaviors. These attacks are designed to maintain
realistic mobility patterns, allowing them to attack both data plane
observations and control plane decisions without triggering sudden or
obvious anomalies.

The threat model assumes that standard cryptographic mechanisms for authentication and secure communication are not broken by attackers. Physical layer attacks such as jamming and radio interference are not considered. The scope of this work is limited to data poisoning attacks that affect learning based mobility management, rather than Dos attack or availability attacks.

This threat model defines the assumptions used for the attack scenarios presented in the next section and directly guides the design of the proposed mitigation framework.
\newpage
\section{SDVN System Modeling}

The SDVN is modeled as a system consisting of vehicles, RSUs,
and a set of $K$ registered controllers
$\mathcal{C} = \{c_1, \ldots, c_K\}$ whose trust scores are
continuously maintained on the blockchain via SC-Trust.
Each RSU $r_j$ maintains a current controller assignment
$c_{\mathrm{assigned}}(r_j) \in \mathcal{C}$, initially
determined by geographic proximity or load-balancing.
The trusted controller set maintained by the blockchain at time
$t$ is shown in Equation~\ref{eq:trusted_ctrl_set}:
\begin{equation}
  \mathcal{C}_{\mathrm{trusted}}(t) =
    \bigl\{c_k \in \mathcal{C} :\;
    \tau_{c_k}(t) > \tau_{\min}\bigr\}
  \label{eq:trusted_ctrl_set}
\end{equation}
As shown in Equation~\ref{eq:trusted_ctrl_set}, any controller
whose trust score falls below $\tau_{\min}$ is excluded from the
trusted set and its flow rules are no longer accepted by RSUs.
When $c_{\mathrm{assigned}}(r_j)$ is revoked, RSU $r_j$ queries
the blockchain for $\mathcal{C}_{\mathrm{trusted}}(t)$ and
reassigns to any available controller in that set, maintaining
uninterrupted network operation without manual intervention.
Based on learned mobility information, each assigned controller
generates control actions such as safety alerts, routing
instructions, and traffic management decisions. These control
messages are sent back to vehicles through RSUs. This system
model captures the logical data flow and control interactions in
an SDVN and forms the foundation for attack modeling and
mitigation framework design.


\section{Attack Modeling Based on Defined Scenarios}

Based on the finalized problem sketches, multiple poisoning attack scenarios are systematically modeled in the simulation environment.

\subsection{Trajectory Poisoning Attacks}

Trajectory poisoning attacks aim to corrupt the temporal learning of vehicle movement paths by injecting falsified trajectory information over multiple time steps. These attacks affect how the SDVN controller learns and predicts vehicle trajectories, leading to incorrect safety and traffic management decisions.
\newline
\newline
Figure \ref{fig:trajectory_poisoning_a} shows a data plane trajectory poisoning attack which a compromised RSU change legitimate vehicle trajectory data before forwarding it to the SDVN controller.

\begin{figure}[H]
	\centering
	\includegraphics[width=0.5\textwidth]{images/1(a).eps}
	\caption{Malicious RSU based data-plane trajectory poisoning}
	\label{fig:trajectory_poisoning_a}
\end{figure}

\begin{itemize}
	\item \textbf{Malicious RSU Trajectory Poisoning:}
	As shown in Figure~\ref{fig:trajectory_poisoning_a}, the attack begins when a legitimate vehicle sends correct trajectory information, such as its position and movement path, toward the RSU (Step 1). This trajectory data is first received by RSU as part of normal data plane communication.

	Instead of forwarding the data honestly, the RSU is compromised and acts as the attacker. At Step 2, the malicious RSU modifies the received trajectory data by injecting falsified but realistic trajectory values. These poisoned trajectory packets still follow smooth movement patterns, making them difficult to detect and then the RSU forwards the poisoned trajectory data to the SDVN controller through the control plane communication channel. In the step 3 and 4, same transmission happen like another legitimate vehicle send the trajectory data to controller though the RSU.

	At Step 5, Based on this corrupted input, the controller learns an incorrect trajectory model and is deceived into believing that the traffic situation is safe, even though a real collision risk exists. Using this wrong understanding, the controller generates incorrect control decisions and control packets (Step 5). As a result, safety critical warning messages that should be sent to vehicles are either suppressed or replaced with misleading control instructions (Step 7). Consequently, vehicles do not receive necessary safety alerts, which may lead to unsafe driving decisions and potential collisions, despite the vehicles behaving honestly throughout the process.
\end{itemize}
Figure \ref{fig:trajectory_poisoning_b} shows a data plane trajectory poisoning attack where a malicious vehicle injects falsified but realistic trajectory information into the network.
\begin{figure}[H]
	\centering
	\includegraphics[width=0.5\textwidth]{images/1(b).eps}
	\caption{Malicious vehicle based data-plane trajectory poisoning}
	\label{fig:trajectory_poisoning_b}
\end{figure}

\begin{itemize}
	\item \textbf{Malicious Vehicle Trajectory Poisoning:}
	As shown in Figure~\ref{fig:trajectory_poisoning_b}, the attack starts when a malicious vehicle acts as a legitimate participant in the network but intentionally injects falsified trajectory data. At Step 1, the compromised vehicle generates fake but smooth and realistic trajectory information instead of its true movement data. this poisoned trajectory data is transmitted to the RSU through normal data plane communication. Since the data follows realistic mobility patterns, the RSU does not detect any abnormal behavior and forwards the received information without modification.

	At Step 2, the RSU sends the poisoned trajectory data to the SDVN controller via the control plane communication channel. Over multiple time steps, the controller continuously receives this falsified trajectory information and gradually learns an incorrect trajectory model for the malicious vehicle.

	At the step 3. Based on this corrupted learning, the controller incorrectly assumes that the traffic situation is safe, and send the control signals even though a real collision risk exists . As a result, the controller generates improper control decisions and suppresses or delays safety critical warning messages that should be delivered to nearby vehicles (Step 4).

	Consequently, surrounding vehicles do not receive timely safety alerts, leading to unsafe driving conditions and a higher risk of collision, despite normal network communication behavior being preserved throughout the attack.
\end{itemize}
\newpage
Figure \ref{fig:trajectory_poisoning_c} illustrates a control plane trajectory poisoning attack in which the SDVN controller itself is malicious and poisons the learning process despite receiving correct data.
\begin{figure}[H]
	\centering
	\includegraphics[width=0.5\textwidth]{images/1(c).eps}
	\caption{Control-plane trajectory poisoning}
	\label{fig:trajectory_poisoning_c}
\end{figure}

\begin{itemize}
	\item \textbf{Control-Plane Trajectory Poisoning}
	As shown in Figure~\ref{fig:trajectory_poisoning_c}, this attack occurs entirely at the control plane while the data plane behaves correctly. At Step 1 and Step 2, vehicles honestly generate and transmit correct trajectory data to the RSU through normal data plane communication.

	At Step 3 and 4, the RSU forwards this legitimate trajectory information to the SDVN controller using the control plane channel. Up to this point, no falsified data is injected, and the collected mobility information accurately reflects the real vehicle movements.

	At Step 5 and 6, the attack begins inside the SDVN controller itself. Although it receives correct trajectory data, the controller is malicious or compromised and deliberately poisons its internal learning model. Instead of learning the real vehicle trajectories, the controller constructs a fake trajectory model that incorrectly indicates a safe traffic situation.

	Based on this poisoned learning state, the controller generates incorrect control decisions at Step 5 and Step 6. These control packets are sent back to the RSU and then forwarded to vehicles (Step 7 and Step 8).

	As a result, safety critical warning messages that should alert vehicles about a real collision risk are blocked or suppressed in the data plane. Even though vehicles and RSUs behave honestly and communicate correctly, the malicious controller causes unsafe decisions, leading to potential collisions and traffic hazards.
\end{itemize}

\subsection{Mobility Pattern Poisoning Attacks}
\label{sec:mobility_pattern_poisoning}

Mobility pattern poisoning attacks try to corrupt the SDVN controller's global understanding of traffic flow and vehicle behavior by injecting falsified mobility data over time. Unlike trajectory poisoning, which focuses on individual vehicle paths, mobility pattern poisoning targets aggregated learning outcomes such as congestion estimation, traffic density prediction, and coordinated vehicle behavior.

Figure~\ref{fig:mobility_a} illustrates a Sybil based mobility pattern poisoning attack launched through a compromised RSU.

\begin{figure}[H]
	\centering
	\includegraphics[width=0.5\textwidth]{images/2(a).eps}
	\caption{Sybil based attack via compromised RSU}
	\label{fig:mobility_a}
\end{figure}

\begin{itemize}
	\item \textbf{Sybil Based Mobility Pattern Poisoning via Compromised RSU:}
	As shown in Figure~\ref{fig:mobility_a}, the attack begins at Step 1 and 3 when legitimate vehicles transmit correct mobility data to the RSU through the data plane. These packets contain valid position, speed, and movement information and represent normal network behavior.

	At step 2 and 4, legitimate packets are foreword to the controller.
	Compromised rsu, impersonates multiple fake vehicle identities and then Instead of forwarding only the legitimate packets, the RSU generates additional forged mobility reports using different identities, as indicated at Step 5 and Step 6.

	At Step 5 and 6 , impersonated mobility packets are forwarded to the SDVN controller through the control plane. Since the forged packets follow realistic mobility patterns and carry valid looking identities, the controller treats them as data coming from multiple independent vehicles.

	As a result, the controller aggregates this poisoned mobility information into its global learning model. This leads to false traffic density estimation and incorrect global mobility patterns. Based on this corrupted model, the controller generates incorrect control decisions and sends false control signals back to the RSU at Steps 7 and 8.

	Finally, at Steps~9 and~10, these incorrect control decisions reach the vehicles. This causes distorted traffic management actions and unsafe or inefficient network behavior, even though most vehicles in the network behave honestly.
\end{itemize}

Figure~\ref{fig:mobility_b} illustrates a Sybil based mobility pattern poisoning attack via Malicious Vehicle Impersonation.

\begin{figure}[H]
	\centering
	\includegraphics[width=0.5\textwidth]{images/2(b).eps}
	\caption{Sybil based attack via malicious vehicle impersonation}
	\label{fig:mobility_b}
\end{figure}

\begin{itemize}
	\item \textbf{Sybil Based Mobility Pattern Poisoning via Malicious Vehicle Impersonation:}
	As shown in Figure~\ref{fig:mobility_b}, the attack starts at Step 1 and 3 when a malicious vehicle generates falsified mobility data while pretending to be a legitimate vehicle. The attacker forges multiple mobility packets using stolen or learned identities, even though these vehicles do not actually exist.

	At Step 4, the malicious vehicle transmits  impersonated mobility packets to the RSU through the data plane. These packets contain realistic position, speed, and movement information, allowing them to bypass basic validity checks.

	At Step 2,5 and 6, the RSU forwards both the received mobility packets to the SDVN controller through the control plane without detecting the impersonation. Since the packets carry valid looking identities, the controller assumes that the data is coming from multiple independent vehicles.

	The controller aggregates this poisoned mobility information into its global mobility learning model. As a result, false traffic density and incorrect mobility patterns are learned

	Based on the corrupted learning model, the controller generates incorrect control decisions and sends control packets back to the RSU at Step 7. These control packets are then forwarded to vehicles at Step 8.

	Finally, the vehicles receive incorrect traffic management and control decisions, leading to distorted mobility predictions and unsafe or inefficient network behavior, despite most vehicles in the network behaving honestly.

\end{itemize}

Figure~\ref{fig:mobility_c} illustrates a Data plane MitM based mobility pattern poisoning attack
\begin{figure}[H]
	\centering
	\includegraphics[width=0.5\textwidth]{images/2(c).eps}
	\caption{Data plane MitM based mobility pattern poisoning}
	\label{fig:mobility_c}
\end{figure}

\begin{itemize}
	\item \textbf{Data Plane Mobility Pattern Poisoning via Man in the Middle Attack :}
	As shown in Figure~\ref{fig:mobility_c}, the attack begins at Step 1 when legitimate vehicles transmit correct mobility data, such as position and speed, to the RSU through the data plane. This represents normal communication in the SDVN.

	At Step 2, a Man in the Middle (MitM) attacker positioned in the data plane intercepts these mobility packets before they reach the RSU. The attacker observes the legitimate messages and learns valid vehicle identities and mobility formats.The attacker fabricates new mobility reports by modifying attributes such as location, speed, and acceleration, while keeping the values realistic to avoid detection. These forged mobility packets are generated over multiple time steps to mimic normal driving behavior.
	At Step 4, the poisoned mobility packets are forwarded to the RSU as if they were sent by legitimate vehicles. Since the packets follow valid formats and realistic patterns, the RSU does not detect any anomaly.

	At Step 3 and Step 5, the RSU forwards both legitimate and poisoned mobility data to the SDVN controller through the control plane. The controller treats all received reports as trustworthy inputs from different vehicles.

	At the step 6 The controller aggregates this poisoned mobility information into its global learning model. Over time, the accumulated false data corrupts the learned mobility patterns, and send back to the RSU.

	At the step 7, the vehicle received incorrect control decisions based on the poisoned model, which may lead to unsafe, inefficient, or misleading traffic management actions, even though most vehicles in the network behave honestly.

\end{itemize}


Figure \ref{fig:mobility_d} shows a control plane mobility pattern poisoning attack in which the SDVN controller itself is malicious and corrupts the learned global mobility model despite receiving correct mobility data.

\begin{figure}[H]
	\centering
	\includegraphics[width=0.5 \textwidth]{images/2(d).eps}
	\caption{Control plane mobility pattern poisoning}
	\label{fig:mobility_d}
\end{figure}

\begin{itemize}
	\item \textbf{Control Plane Mobility Pattern Poisoning:}
	As shown in Figure~\ref{fig:mobility_d}, the attack begins at Step 1 and Step 2 when
	vehicles transmit legitimate mobility data, such as position and speed, to the RSU
	through the data plane. This communication represents normal and honest vehicle
	behavior.

	At Step 3 and Step 4, the RSU forwards the collected mobility information to the
	SDVN controller via the control plane. At this stage, both the vehicles and the RSU
	operate correctly and do not inject any false data.

	At Step~5, the SDVN controller, which is malicious or compromised, intentionally
	poisons its internal learning model despite receiving correct mobility inputs. Instead
	of learning the true global mobility patterns, the controller constructs a corrupted
	mobility model that misrepresents traffic conditions and vehicle behavior.

	Based on the poisoned model, the controller generates incorrect control packets, as
	indicated at Step 5, and sends these faulty control decisions back to the RSU through
	the control plane. Finally, the RSU forwards the incorrect control packets to vehicles at Step 6 in the
	data plane. As a result, safety warnings may be suppressed, collision risks may be
	misinterpreted as safe situations, and traffic control actions become incorrect, even
	though all data plane mobility observations were accurate.

\end{itemize}

\subsection{Mobility-Induced Threat Amplification in SDVN}
\label{sec:mobility_amplification}

Unlike static networks, SDVN attack surfaces are fundamentally shaped by
vehicle kinematics. High vehicular mobility constrains RSU contact windows,
accelerates topology changes, and reduces the number of observable beacons per
vehicle--RSU association, all of which directly amplify the effectiveness of
poisoning attacks. We formalize these mobility-induced amplification effects as
follows.

\subsubsection{RSU Contact Duration and Beacon Sampling Window}

Let vehicle $v_i$ travel at speed $s_i$ along a trajectory with perpendicular
distance $d_{\perp}$ from RSU $r_j$ having coverage radius $R_j$. The contact
duration $\tau_c$ is given in Equation~\ref{eq:contact_duration}.
\newline
\newline
This equation quantifies the available time window within which an RSU
can observe and collect beacon evidence from a passing vehicle,
establishing the upper bound on how much trajectory data any single
RSU can accumulate for anomaly detection before the vehicle exits coverage.
\begin{equation}
  \tau_c(v_i, r_j) = \frac{2\sqrt{R_j^2 - d_{\perp}^2}}{s_i},
  \quad d_{\perp} \leq R_j
  \label{eq:contact_duration}
\end{equation}
As shown in Equation~\ref{eq:contact_duration}, $\tau_c$ decreases as vehicle
speed $s_i$ increases or as the vehicle passes farther from the RSU centre,
directly reducing the observation window available for attack detection.
\newline
\newline
Given a beacon broadcast interval $T_b$ (typically 100\,ms per IEEE~802.11p),
$N_b$ is the number of beacon messages the RSU $r_j$ can observe from vehicle
$v_i$ during a single contact (one pass through coverage). The number of
observable beacons $N_b$ per RSU contact is given in Equation~\ref{eq:beacon_count}.
\newline
\newline
This equation translates the contact duration into a discrete count of
observable beacon messages, directly determining whether the temporal
autoencoder receives sufficient sequence length to achieve reliable
detection power during a single RSU contact.
\begin{equation}
  N_b(v_i, r_j) = \left\lfloor \frac{\tau_c(v_i, r_j)}{T_b} \right\rfloor
  \label{eq:beacon_count}
\end{equation}
\newline
\newline
As shown in Equation~\ref{eq:beacon_count}, as $s_i$ increases, $\tau_c$
decreases and $N_b$ shrinks, directly reducing the evidence available for
sequential anomaly detection. Let $N_{\min}$ be the minimum beacon sequence
length required for a temporal detector to achieve target detection power.
\newline
\newline
The detection feasibility condition is given in Equation~\ref{eq:detection_condition}.This condition formally defines when sequential anomaly detection is
feasible at a given RSU: it identifies the speed threshold beyond which
the lightweight and full mode detectors must defer to cross-RSU
offline analysis, motivating the dual-mode architecture design.
\newline
\newline
\begin{equation}
  N_b(v_i, r_j) \geq N_{\min}
  \label{eq:detection_condition}
\end{equation}
As shown in Equation~\ref{eq:detection_condition}, when this condition is
violated due to high speed, the temporal autoencoder receives insufficient
sequence length per RSU contact, creating an evasion window exploitable
by adversaries.

\subsubsection{Handover-Induced Learning Discontinuity}

The handover rate $\lambda_h(s_i)$ of vehicle $v_i$
across RSUs is given in Equation~\ref{eq:handover_rate}. This equation models how frequently a vehicle transitions between RSU
coverage zones, quantifying the degree to which trajectory sequences
are fragmented across RSU records and thus unavailable to any single
RSU's sequential detector which is a fragmentation effect that attackers
deliberately exploit to evade cross-beacon correlation.
\begin{equation}
  \lambda_h(s_i) = \frac{s_i}{\mathbb{E}[2\sqrt{R^2 - d_{\perp}^2}]}
  \label{eq:handover_rate}
\end{equation}
As shown in Equation~\ref{eq:handover_rate}, high $\lambda_h$
means vehicle trajectory sequences are split across multiple RSU records,
disrupting the temporal continuity required for sequential poisoning detection.
A malicious vehicle exploiting this effect injects falsified beacons only
within a single RSU contact window $[t_0,\, t_0 + \tau_c]$ and remains honest
otherwise, evading cross-RSU correlation where $\mathbb{E}[\cdot]$ is the expectation over the perpendicular distance
distribution.

\subsubsection{Mobility-Disguised Poisoning Drift Rate}

In trajectory poisoning, an adversary injects a false position offset
$\boldsymbol{\delta}(t)$ that grows gradually over time. To avoid triggering
threshold-based detectors, the attacker must ensure the per-step drift
constraint given in Equation~\ref{eq:drift_constraint}.
\newline
\newline
This equation characterizes the per-beacon falsification budget
available to a trajectory poisoning attacker who wishes to remain
undetected by threshold-based checks, establishing why gradual
poisoning is structurally harder to detect than abrupt position jumps
and motivating the cumulative drift score in TP-S5.
\begin{equation}
  \|\boldsymbol{\delta}(t) - \boldsymbol{\delta}(t-T_b)\| \leq \epsilon_{\max}
  \label{eq:drift_constraint}
\end{equation}
As shown in
Equation~\ref{eq:drift_constraint}, the attacker is constrained to inject
only small incremental offsets per beacon, forcing a slow gradual drift
that mimics legitimate movement to evade threshold-based detectors where $\epsilon_{\max}$ is the maximum allowable positional deviation per
beacon interval inferred from speed plausibility checks.
\newline
\newline
The cumulative falsified displacement after $k$ beacons is given in
Equation~\ref{eq:cumulative_drift}. This equation quantifies the total falsified displacement an attacker
can accumulate over an entire RSU contact window while staying within
per-beacon kinematic bounds, revealing that the attack budget grows
with vehicle speed and directly motivating the cumulative drift score
detector TP-S5 rather than per-beacon threshold checks alone.
\begin{equation}
  \|\boldsymbol{\delta}(t_0 + k T_b)\| \leq k \cdot \epsilon_{\max}
  \label{eq:cumulative_drift}
\end{equation}
As shown in Equation~\ref{eq:cumulative_drift}, the attack budget
per RSU contact is therefore bounded by $N_b \cdot \epsilon_{\max}$.
At higher vehicle speeds, $N_b$ decreases (Equation~\ref{eq:beacon_count}),
but $\epsilon_{\max}$ increases because speed plausibility thresholds widen
at high velocity, giving the attacker a larger per-beacon drift allowance.
This creates a speed-proportional attack amplification effect unique to
vehicular environments.

\subsubsection{Sybil Identity Capacity Under Mobility Constraints}

A compromised RSU performing Sybil-based mobility pattern poisoning
(Section~\ref{sec:mobility_pattern_poisoning}) can sustain at most
$K_{\text{sybil}}$ ghost identities whose beacons collectively remain
spatially consistent. Given the RSU coverage area $\mathcal{A}_j = \pi R_j^2$
and a minimum inter-vehicle spacing $d_{\min}$ enforced by spatial plausibility
checks, the Sybil identity capacity is bounded as given in
Equation~\ref{eq:sybil_capacity}.
\begin{equation}
  K_{\text{sybil}} \leq \left\lfloor \frac{\mathcal{A}_j}{\pi (d_{\min}/2)^2} \right\rfloor
  \label{eq:sybil_capacity}
\end{equation}
As shown in Equation~\ref{eq:sybil_capacity}, the maximum number of ghost
identities is directly bounded by the RSU coverage area and the minimum
spacing constraint, meaning a larger coverage area or looser spacing
allows more Sybil identities to coexist undetected.
Since $\mathcal{A}_j$ is fixed but the number of legitimate vehicles in
$\mathcal{A}_j$ varies with traffic density $\rho_v$ (vehicles per $m^2$),
the fraction of Sybil identities relative to total reported identities is
given in Equation~\ref{eq:sybil_fraction}.
\begin{equation}
  f_{\text{sybil}} = \frac{K_{\text{sybil}}}{\rho_v \cdot \mathcal{A}_j + K_{\text{sybil}}}
  \label{eq:sybil_fraction}
\end{equation}
As shown in Equation~\ref{eq:sybil_fraction}, at low traffic density
(e.g., highway scenarios with high $s_i$), $\rho_v \cdot \mathcal{A}_j$
is small, so $f_{\text{sybil}}$ approaches 1, meaning a compromised RSU
can dominate the controller's view of local traffic. This quantifies why
mobility conditions that reduce traffic density directly amplify the Sybil
attack's impact on the SDVN controller's learning model.
\newline
\newline
Figure~\ref{fig:dual_mobility_threat} presents the combined mobility-induced threat amplification effects, illustrating how RSU contact duration and the adversarial drift budget vary with vehicle speed across urban, rural, and highway regimes.
\begin{figure}[H]
    \centering
    \includegraphics[width=1.0\linewidth]{images/FigAdiag}
    \caption{Mobility-Induced Threat Amplification.}
    \label{fig:dual_mobility_threat}
\end{figure}
As shown in Figure~\ref{fig:dual_mobility_threat}, as vehicle speed
increases from urban to highway regimes, the RSU contact duration
$\tau_c$ decreases while the adversarial drift budget $N_b \times
\epsilon_{\max}$ increases, with their crossover point marking the
Detection Infeasible Zone where $N_b < N_{\min}$ and sequential
anomaly detection becomes unreliable.


\subsection{Attack Signature Characterization}
\label{sec:attack_signatures}

A Basic Safety Message (BSM) beacon broadcast by vehicle
$v_i$ at time $t$ is defined as the tuple given in
Equation~\ref{eq:bsm}. This tuple formally defines the observable input to all detection
procedures, establishing which kinematic fields are available at the
RSU for signature evaluation and making explicit that the detection
framework operates solely on standard beacon content
without requiring additional sensors or infrastructure.
\begin{equation}
  \mathbf{b}_i(t) = \bigl(p_i(t),\; s_i(t),\; \theta_i(t),\;
                    a_i(t),\; t,\; \mathrm{ID}_i\bigr)
  \label{eq:bsm}
\end{equation}
As shown in
Equation~\ref{eq:bsm}, each BSM encodes the full kinematic state of
vehicle $v_i$, forming the observable beacon stream
$\{\mathbf{b}_i(t)\}$ from which the detectable signatures of each
attack class are characterised where $p_i(t) \in \mathbb{R}^2$ is the GPS position, $s_i(t)$ is speed,
$\theta_i(t)$ is heading angle, $a_i(t)$ is acceleration, and
$\mathrm{ID}_i$ is the pseudonymous vehicle identifier.

Table~\ref{tab:signature-naming} fixes the canonical naming used for
the remainder of this report: each of the nine rule-based signatures
is referred to by its TP-S/MP-S identifier only, and each of the seven
evaluated attack scenarios (as used in Experiment~5,
Section~\ref{sec:res-exp5}, and throughout the Results chapter) is
referred to by its scenario number (\texttt{a1}--\texttt{a7}) only. The two identifier
schemes index different things --- a signature is a detection rule,
an attack scenario is a simulated adversary configuration --- and are
kept deliberately distinct rather than used interchangeably, since a
single attack scenario may trigger several signatures simultaneously
and a single signature may fire under several scenarios.

\begin{table}[H]
\centering
\caption{Canonical mapping between attack scenarios and their primary
rule-based detection signature.}
\label{tab:signature-naming}
\renewcommand{\arraystretch}{1.2}
\begin{tabular}{|l|l|l|l|}
\hline
\textbf{Scenario} & \textbf{Primary signature} & \textbf{Plane} & \textbf{Compromised entity} \\
\hline
a1 & TP-S1 & Data  & Malicious vehicle \\
\hline
a2 & TP-S2 & Data  & Compromised RSU \\
\hline
a3 & MP-S1 & Data  & Compromised RSU (Sybil) \\
\hline
a4 & MP-S2 & Data  & Vehicle (impersonation) \\
\hline
a5 & TP-S3 & Control & Malicious controller \\
\hline
a6 & MP-S3 & Data  & MitM adversary \\
\hline
a7 & MP-S4 & Control & Malicious controller \\
\hline
\end{tabular}
\end{table}
TP-S4 and TP-S5 (cumulative drift residual and its threshold,
Table~\ref{tab_notations}) are supplementary checks within
TP-DETECT that refine detection across all trajectory-poisoning
scenarios rather than identifying a single scenario of their own, so
they do not appear in Table~\ref{tab:signature-naming}.

\subsubsection{Trajectory Poisoning Signatures}

\textbf{TP-S1 (Kinematic Position Feasibility).}
A reported position change must satisfy the speed-of-light constraint
over beacon interval $T_b$. The kinematic position feasibility violation
is defined in Equation~\ref{eq:tp_s1}.
\newline
\newline
This signature detects the most aggressive form of trajectory
poisoning --- a single-beacon position jump that violates physical
speed limits --- serving as the first and fastest gate in TP-DETECT
that eliminates obviously falsified beacons before the more
computationally expensive cumulative checks are applied.
\begin{equation}
  \bigl\|p_i(t) - p_i(t-T_b)\bigr\| > s_{\max} \cdot T_b
  \label{eq:tp_s1}
\end{equation}
As shown in
Equation~\ref{eq:tp_s1}, a position displacement exceeding
$s_{\max} \cdot T_b$ in one beacon interval flags a kinematic
feasibility violation where $s_{\max}$ is the road-class maximum speed.
\newline
\newline
\textbf{TP-S2 (Heading Rate Deviation).}
A realistic vehicle cannot change heading faster than a maximum angular
velocity $\omega_{\max}$ (rad/s). The heading rate deviation signature
is given in Equation~\ref{eq:tp_s2}.
\newline
\newline
This signature detects unrealistic heading changes that cannot be
produced by a road-constrained vehicle, catching trajectory
falsifications where an attacker maintains a plausible position
sequence but fabricates direction information to mislead path
prediction models.
\begin{equation}
  \Delta\theta_i(t) = \bigl|\theta_i(t) - \theta_i(t-T_b)\bigr|
  > \omega_{\max} \cdot T_b
  \label{eq:tp_s2}
\end{equation}
As shown in Equation~\ref{eq:tp_s2}, a heading change exceeding
$\omega_{\max} \cdot T_b$ over one beacon interval indicates an
unrealistic turn rate inconsistent with vehicle dynamics.
\newline
\newline
\textbf{TP-S3 (Kinematic Bound Violation).}
Vehicle acceleration and speed are physically bounded by road friction,
vehicle dynamics, and the road-class speed limit. The kinematic bound
violation signature is given in Equation~\ref{eq:tp_s3}.
\newline
\newline
This signature detects physically impossible acceleration or speed
values, providing a complementary kinematic check to TP-S1. It catches
falsifications where position increments appear plausible but the implied
speed change violates vehicle dynamic limits, and---critically---attacks
that report a constant yet physically impossible speed: a steady
extreme speed produces zero implied acceleration (evading the acceleration
term) on a plausible position (evading TP-S1), so the reported speed
magnitude is the only direct tell.
\begin{equation}
  \bigl|a_i(t)\bigr| > a_{\max}
  \;\;\lor\;\;
  s_i(t) > s_{\max}
  \label{eq:tp_s3}
\end{equation}

As shown in Equation~\ref{eq:tp_s3}, any reported acceleration exceeding
$a_{\max}$ or reported speed exceeding $s_{\max}$ is physically impossible
and flags a falsified beacon, where $a_{\max}$ is the maximum achievable
acceleration and $s_{\max}$ is the road-class maximum speed
(Eq.~\ref{eq:tp_s1}). Because $s_{\max}$ is road-class specific, honest
urban and rural traffic remains well below it, while for high-speed
(highway) road classes $s_{\max}$ is set to that class's speed limit so
that legitimate high-speed vehicles are not flagged.
\newpage
\textbf{TP-S4 (Dead-Reckoning Residual).}
Using the previous beacon state, a predicted position
is computed as given in Equation~\ref{eq:dead_reckoning}.
\newline
\newline
This prediction step reconstructs the expected vehicle position from
its previously reported kinematic state, providing a ground-truth
reference derived entirely from the vehicle's own prior beacons
against which the current reported position is compared to detect
injected displacement.
\begin{equation}
  \hat{p}_i(t) = p_i(t-T_b) + s_i(t-T_b)\cdot T_b \cdot
  \begin{bmatrix}\cos\theta_i(t-T_b)\\ \sin\theta_i(t-T_b)\end{bmatrix}
  \label{eq:dead_reckoning}
\end{equation}
As shown in Equation~\ref{eq:dead_reckoning}, $\hat{p}_i(t)$ estimates
where the vehicle should be based solely on its previously reported
speed and heading.
\newline
\newline
The dead-reckoning residual $r_i(t)$ is given in
Equation~\ref{eq:tp_s4}. This residual quantifies the discrepancy between what the vehicle
reported and where its own prior kinematics predict it should be,
providing a per-beacon anomaly signal that catches falsifications
individually plausible by road-map checks but inconsistent with the
vehicle's own movement history.
\begin{equation}
  r_i(t) = \bigl\|p_i(t) - \hat{p}_i(t)\bigr\|
  \label{eq:tp_s4}
\end{equation}
As shown in Equation~\ref{eq:tp_s4}, $r_i(t)$ measures the deviation
between the reported position and the kinematically predicted position,
with a large residual indicating injected false displacement.
\newline
\newline
\textbf{TP-S5 (Cumulative Drift Score).}
The cumulative drift score $D_i(k)$ for detecting
gradual trajectory poisoning is given in
Equation~\ref{eq:tp_s5}.
\newline
\newline
This score aggregates dead-reckoning residuals over a sliding beacon
window to detect gradual trajectory poisoning that evades per-beacon
checks, directly addressing the slow-drift attack strategy formalized
in Equation~\ref{eq:cumulative_drift} where individual beacons appear
kinematically valid but their sequence reveals systematic falsification.
\begin{equation}
  D_i(k) = \frac{1}{k}\sum_{j=1}^{k} r_i(t - (k-j)\cdot T_b)
  \label{eq:tp_s5}
\end{equation}
As shown in Equation~\ref{eq:tp_s5}, $D_i(k)$ averages the
dead-reckoning residuals over the last $k$ beacons. A drift alert
is issued when $D_i(k) > \delta_{\mathrm{th}}$, where
$\delta_{\mathrm{th}}$ is a road-class calibrated threshold.

\newpage
\subsubsection{Mobility Pattern Poisoning Signatures}

\textbf{MP-S1 (Identity Density Anomaly).}
The number of distinct vehicle IDs reported within RSU $r_j$'s coverage
must not exceed the Sybil identity capacity bound derived in
Equation~\ref{eq:sybil_capacity}. The identity density anomaly check
is given in Equation~\ref{eq:mp_s1}.
\newline
\newline
This check uses the Sybil capacity bound from
Equation~\ref{eq:sybil_capacity} as a data-driven threshold to flag
RSU coverage zones where the number of reported vehicle identities
exceeds what is physically possible, enabling Sybil detection without
requiring knowledge of which specific identities are forged.

\begin{equation}
  \bigl|\{\mathrm{ID}_i : p_i(t) \in \mathcal{A}_j\}\bigr|
  > K_{\mathrm{sybil}} + \rho_v \cdot \mathcal{A}_j
  \label{eq:mp_s1}
\end{equation}
As shown in
Equation~\ref{eq:mp_s1}, when the observed number of distinct IDs
exceeds the sum of the Sybil capacity bound and the expected legitimate
vehicle count, an identity density anomaly is flagged where $\mathcal{A}_j$ denotes the geographic coverage area of RSU $r_j$
and $\rho_v$ is the expected vehicle spatial density.
\newline
\newline
\textbf{MP-S2 (Synchronized Beacon Timing).}
Sybil identities generated by a single source exhibit suspiciously
correlated transmission timestamps. For identity pair
$(\mathrm{ID}_a, \mathrm{ID}_b)$, the synchronized beacon timing
signature is given in Equation~\ref{eq:mp_s2}.
\newline
\newline
This signature exploits the implementation artifact that Sybil
identities generated by a single source share the same hardware clock,
producing statistically improbable beacon timing correlations that
expose co-sourced identity pairs even when their reported positions
and speeds are individually plausible.
\begin{equation}
  \frac{1}{K}\sum_{k=1}^{K}
  \mathbf{1}\!\left[|t_a^{(k)} - t_b^{(k)}| < \tau_{\mathrm{sync}}\right]
  > \rho_{\mathrm{sync}}
  \label{eq:mp_s2}
\end{equation}
As shown in Equation~\ref{eq:mp_s2}, synchronization
is flagged when the fraction of beacon pairs transmitted within
$\tau_{\mathrm{sync}}$ of each other exceeds the co-occurrence
threshold $\rho_{\mathrm{sync}}$ where $\tau_{\mathrm{sync}}$ is a sub-millisecond timing resolution
threshold and $\rho_{\mathrm{sync}} \in (0,1)$ is a co-occurrence
rate threshold.
\newline
\newline
Implementation note. An early implementation of MP-S2 contained
a dead-code path in which the synchronised-timing check was
constructed but never invoked from the combined-attack evaluation
loop, so MP-S2's contribution to $\psi_i(t)$ was silently zero
whenever multiple attack types were active simultaneously. This has
been identified and corrected; all combined-attack results reported
in this document reflect the corrected implementation.
\newline
\newline
\textbf{MP-S3 (Regional Speed Distribution Shift).}
Mobility pattern poisoning shifts the observed regional speed
distribution $\mathcal{P}_t$ away from the historical baseline
$\mathcal{P}_{\mathrm{hist}}$. The regional speed distribution shift
signature is given in Equation~\ref{eq:mp_s3}.
\newline
\newline
This signature detects MitM and mobility pattern poisoning attacks
by measuring how much the current regional speed distribution has
shifted from historical norms, catching large-scale traffic
manipulation that affects the aggregate learned by the controller
even when individual beacon modifications are individually subtle.
\begin{equation}
  D_{\mathrm{KL}}\!\left(\mathcal{P}_t \,\|\, \mathcal{P}_{\mathrm{hist}}\right)
  = \sum_{s} \mathcal{P}_t(s) \log
    \frac{\mathcal{P}_t(s)}{\mathcal{P}_{\mathrm{hist}}(s)}
  > \kappa_{\mathrm{th}}
  \label{eq:mp_s3}
\end{equation}

As shown in Equation~\ref{eq:mp_s3}, when the KL-divergence between
the current and historical speed distributions exceeds
$\kappa_{\mathrm{th}}$, a regional mobility pattern shift is detected where $\kappa_{\mathrm{th}}$ is a calibrated KL-divergence threshold.
\newline
\newline
\textbf{MP-S4 (Ghost Vehicle Transit Impossibility).}
A vehicle identity $\mathrm{ID}_i$ appearing at non-overlapping RSUs
$r_j$ and $r_k$ within a transit time that violates kinematics is
classified as a ghost identity. The ghost vehicle transit impossibility
signature is given in Equation~\ref{eq:mp_s4}.
\newline
\newline
This signature detects ghost identities that appear at
geographically separated RSUs within a transit time that violates
physical speed limits, catching Sybil attacks that a single RSU
cannot detect in isolation because each individual appearance is
locally plausible.
\begin{equation}
  \frac{d(r_j, r_k)}{|t_j - t_k|} > s_{\max}
  \;\Longrightarrow\; \mathrm{ID}_i \text{ is a ghost identity}
  \label{eq:mp_s4}
\end{equation}
As shown in Equation~\ref{eq:mp_s4}, if the implied travel
speed between two RSU appearances exceeds $s_{\max}$, the identity
is physically impossible and is flagged as a ghost vehicle where $d(r_j, r_k)$ is the Euclidean distance between the two RSU
positions.

\subsubsection{Composite Rule-Based Anomaly Score}

Each signature test produces a binary flag weighted by its
per-attack-type reliability $w_k^{(\hat{k}_i)}$, learned in Phase~2
alongside the fusion weights. In the lightweight mode, a single
calibrated weight vector $w_k$ is used since the attack type is
unknown. In the full mode, the predicted attack type $\hat{k}_i$
(Eq.~\ref{eq:gat_atk_type}) selects the weight vector most
appropriate for that attack class, up-weighting the signatures that
are discriminative for class $\hat{k}_i$ and suppressing signatures
that do not fire for it. The attack-type-conditioned composite
rule-based anomaly score for vehicle $v_i$ at time $t$ is given in
Equation~\ref{eq:psi_score}.
\begin{equation}
  \psi_i(t) = \frac{
  \displaystyle\sum_{s=1}^{|\mathcal{S}|}
  w_s^{(\hat{k}_i)} \cdot
  \mathbf{1}[\text{signature } s \text{ violated}]}
 {\displaystyle\sum_{s=1}^{|\mathcal{S}|} w_s^{(\hat{k}_i)}}
  \label{eq:psi_score}
\end{equation}

As shown in Equation~\ref{eq:psi_score}, $\psi_i(t)$ is a
reliability-weighted fraction of violated signatures, where the
weight vector $\{w_s^{(\hat{k}_i)}\}$ is selected by the predicted
attack type $\hat{k}_i$. In the lightweight operating mode,
$\hat{k}_i$ is unavailable and a single calibrated weight vector
$\{w_s\}$ (learned from aggregated Phase~2 data) is used instead.
In the lightweight mode, $\psi_i(t) > \psi_{\mathrm{th}}$ directly
triggers an isolation response without AI inference.
 In the lightweight mode, $\psi_i(t) > \psi_{\mathrm{th}}$
directly triggers an isolation response without AI inference where $\mathcal{S} = \{\text{TP-S1}, \ldots, \text{MP-S4}\}$ is the
full signature set, $\psi_{\mathrm{th}}$ is the composite anomaly score
threshold, and $\delta_{\mathrm{th}}$ is the drift alert threshold.
As shown in Equation~\ref{eq:psi_score}, $\psi_i(t)$ is a
reliability-weighted fraction of violated signatures, ranging from 0
to 1.

The four detection procedures that formalise these signature checks
are presented within the mitigation framework in
Section~\ref{sec:dual_mode}.


\section{Proposed Mitigation Framework}


\subsection{Cryptographic Mitigation Scope Analysis}
\label{sec:crypto_scope}

Before formulating the mitigation framework, we determine whether
cryptographic mechanisms alone can fully mitigate each attack variant,
as this dictates their placement in the detection-mitigation pipeline.

Standard authentication schemes (PKI, certificate-based) verify identity
but not data truthfulness~\cite{arif2020}. A vehicle that possesses
valid credentials can transmit false GPS coordinates, speeds, and trajectories
that pass all cryptographic verifications. Therefore, cryptography alone cannot
mitigate trajectory poisoning or mobility pattern poisoning originating from
authenticated malicious vehicles or compromised RSUs.

However, a post-quantum threshold ring signature (TRS) scheme requiring
collaborative signing by at least $t$-of-$n$ neighbouring RSUs~\cite{ahmed2022blockchain}
structurally prevents any minority-compromised RSU subset (of size $< t$) from
generating a valid aggregate signature. This mitigation acts before
detection: a poisoned aggregate that fails threshold verification is rejected
at the aggregation step without requiring AI-based anomaly scoring.
Accordingly, TRS is placed \textbf{pre-detection} in the pipeline.
TRS verification binds each submission to an authorised ring
identity (Eq.~\eqref{eq:trs_verify}); it attests who signed,
not the numerical validity of the aggregate they contributed.
Insider poisoning by an authorised ring member, in which a field of
$\mathbf{A}_j(t)$ is perturbed while remaining inside the valid
SUMO-derived envelope, is therefore not rejected by TRS and is not
caught by the pre-encryption self-check of
Algorithm~\ref{alg:pq_fhe_trs}, since that check is administered by
the contributing RSU itself and a malicious member may bypass it.
We instead treat this case as a bounded residual. Its impact
is contained by three independent mechanisms. First, static
plausibility envelopes applied both before encryption and after
threshold decryption cap the magnitude of any single contribution,
so an in-bounds perturbation cannot exceed the design-time field
range. Second, because $n \geq 3f+1$ bounds the number of
compromised rings at $\leq f$, the global aggregate
$\bar{X}_{\mathrm{global}}$ is a mean over $M$ rings dominated by the
$\geq 2f+1$ honest contributions, so a single poisoned ring shifts
the mean only within a bounded residual rather than controlling it.
Third, across successive windows the SC-Trust score of a
persistently anomalous RSU decays (Eq.~\eqref{eq:trust_update}) and
triggers $2f+1$ BFT revocation (Eq.~\eqref{eq:bft_revoke_rsu}),
removing the source. The unavoidable limitation is that a single
in-window, in-bounds poisoned contribution is bounded in magnitude
but not per-RSU attributable from the blinded aggregate alone; it is
mitigated temporally rather than eliminated instantaneously.

Post-quantum fully homomorphic encryption (FHE) via BFV, BGV, or CKKS
schemes~\cite{mamun2025postquantum} addresses a precise privacy threat inherent
to the TRS ring coordination design: when $t$-of-$n$ neighboring RSUs
collaborate to collectively sign an aggregate, each RSU in the signing ring
is exposed to the aggregate contents of every other participating RSU, which
it has no legitimate reason to observe. A compromised RSU in the ring therefore
learns traffic patterns from neighboring zones without detection. FHE eliminates
this leakage by requiring each RSU to encrypt its own local aggregate before
sharing it with the signing ring. Ring RSUs perform homomorphic addition across
ciphertexts, producing a combined encrypted aggregate on which TRS is applied.
Since computation proceeds entirely on ciphertext, no RSU in the signing ring
ever observes another RSU's plaintext aggregate. FHE is therefore placed
\textbf{pre-TRS coordination} and is irreplaceable by simpler encryption in
this role, as the ring must compute on ciphertexts without decryption.

For the lightweight pipeline, HMAC-based beacon tagging provides a
computationally cheap integrity anchor at RSU ingestion, preventing
undetected in-transit modification (MitM scenarios) without requiring
AI inference.

The AI components (GAT and Temporal Autoencoder) constitute purely
\textbf{detection logic}, operating on received beacon sequences to
produce anomaly scores. The blockchain layer serves \textbf{both}
detection (through smart contract-based consensus scoring) and mitigation
(through on-chain identity revocation and trust enforcement).
Table~\ref{tab:pipeline_placement} summarises these placements.

\begin{table}[H]
\centering
\caption{Detection vs.\ Mitigation Placement of Framework Components}
\label{tab:pipeline_placement}
\renewcommand{\arraystretch}{1.3}
\begin{tabular}{|p{3.2cm}|p{1.8cm}|p{3.4cm}|p{3.0cm}|}
\hline
\textbf{Component} & \textbf{Role} & \textbf{Placement} & \textbf{Mode} \\
\hline
Rule-based signatures  & Detection   & Pre-model           & Lightweight \& Full \\
\hline
HMAC beacon tagging    & Mitigation  & Pre-detection       & Lightweight \\
\hline
GAT                    & Detection   & Post-ingestion      & Full only \\
\hline
Temporal Autoencoder   & Detection   & Post-ingestion      & Full only \\
\hline
TRS (post-quantum)     & Mitigation  & Pre-detection       & Full only \\
\hline
FHE (post-quantum)     & Mitigation  & Pre-TRS coordination & Full only \\
\hline
Blockchain (ledger)    & Both        & RSU-direct, continuous      & Lightweight \& Full \\
\hline
Smart contracts        & Both        & RSU-triggered, event-driven & Lightweight \& Full \\
\hline
\end{tabular}
\end{table}
As shown in Table~\ref{tab:pipeline_placement}, the framework components
are distributed across pre-detection, post-ingestion, and continuous
placement stages, with cryptographic components (HMAC, TRS, FHE)
handling mitigation and AI components (GAT, Temporal Autoencoder)
handling detection, while the blockchain and smart contracts serve both
roles across both lightweight and full operational modes.

\subsection{Group Key Management via Logical Key Hierarchy}
\label{sec:lkh}

The framework requires two classes of keys: vehicle session keys
$\{K_i\}$ for HMAC beacon integrity (Equation~\ref{eq:hmac}), and RSU
ring key material $\{sk_j\}$ for TRS partial signing
(Equation~\ref{eq:trs_partial}). We adopt a Logical Key Hierarchy (LKH)
to manage both classes efficiently under vehicle join, leave, and RSU
membership changes, following the group key management principle
underlying broadcast encryption for V2I~\cite{zhong2022broadcast}.

\subsubsection{LKH Tree Structure}

The LKH is organised as a balanced binary key tree $\mathcal{T}$ of
depth $d$, where leaf nodes represent individual entities (vehicles or
RSUs) and internal nodes hold group encryption keys. Let
$\mathcal{K}_u$ denote the key assigned to node $u \in \mathcal{T}$.
The set of keys held by each leaf entity $e_i$ along its root path
is given in Equation~\ref{eq:lkh_keyset}.
\begin{equation}
  \mathcal{K}_{e_i} = \{\mathcal{K}_{u} :
    u \in \mathrm{path}(u_i \to \mathrm{root})\}
  \label{eq:lkh_keyset}
\end{equation}
As shown in Equation~\ref{eq:lkh_keyset}, each entity holds exactly
the keys along its path to the root, meaning that rekeying only
requires updating nodes on one root path rather than all members.
The root key $\mathcal{K}_{\mathrm{root}}$ serves as the current group
key shared among all active members. A broadcast message encrypted
under $\mathcal{K}_{\mathrm{root}}$ is decryptable by all legitimate
members and no revoked entity.

The FHE key shares $\{sk_j^{\mathrm{share}}\}_{j=1}^{n}$ produced
during threshold setup (Eq.~\eqref{eq:thresh_keygen}) are distributed
to RSU ring members via the same authenticated LKH ring key
synchronisation channel, eliminating the need for a separate key
distribution protocol for FHE share delivery.

\subsubsection{Vehicle Session Key Establishment}

When vehicle $v_i$ registers with RSU $r_j$, it is assigned a leaf
node $u_i$ in a local LKH subtree maintained per RSU coverage zone.
The RSU derives $v_i$'s session key as a keyed hash of the leaf key
and a session nonce $\eta_i$. The vehicle session key derivation during
RSU registration is given in Equation~\ref{eq:session_key}.
\newline
\newline
This derivation produces a per-session HMAC key unique to each
vehicle registration, ensuring that a session key compromised at one
RSU contact cannot be reused to forge authenticated beacons at
subsequent RSU contacts after the vehicle moves on.
\begin{equation}
  K_i = \mathrm{KDF}(\mathcal{K}_{u_i},\; \eta_i,\; \mathrm{ID}_i)
  \label{eq:session_key}
\end{equation}
As shown in Equation~\ref{eq:session_key}, the session
key $K_i$ binds the entity's leaf key, a fresh nonce, and its identity,
ensuring per-session uniqueness where $\mathrm{KDF}(\cdot)$ is a key derivation function (e.g.,
HKDF-SHA256). This key is used directly in
Equation~\ref{eq:hmac} for per-beacon HMAC tagging.

\subsubsection{Key Update on Vehicle Leave or Revocation}

When vehicle $v_i$ leaves the RSU coverage zone or is revoked via
SC-Revoke (Equation~\ref{eq:bft_revoke_rsu}), the RSU must update all
keys on $v_i$'s root path to preserve forward secrecy. The number
of re-encrypted key messages required upon vehicle leave or revocation
is given in Equation~\ref{eq:lkh_rekey}.
\newline
\newline
This logarithmic cost expression justifies the LKH design choice
over individual unicast rekeying: as vehicle density increases,
the LKH overhead grows only logarithmically while unicast rekeying
grows linearly, making LKH the only scalable option for high-density
urban deployments.
\begin{equation}
  N_{\mathrm{rekey}} = \log_2 |\mathcal{V}_j|
  \label{eq:lkh_rekey}
\end{equation}
where $|\mathcal{V}_j|$ is the current number of vehicles in the RSU
zone. As shown in Equation~\ref{eq:lkh_rekey}, the logarithmic rekey
cost derived from the LKH tree depth provides a significant efficiency
advantage over individual unicast rekeying, which costs
$O(|\mathcal{V}_j|)$ messages~\cite{zhong2022broadcast}.

\subsubsection{RSU Ring Key Generation via Distributed Key Generation (DKG)}

To eliminate the controller as a single point of failure on the
key plane, TRS signing keys and FHE key shares are generated via
a Distributed Key Generation (DKG) protocol executed directly
among the RSU ring members and the Cloud, with no single dealer
and no controller involvement. Each participant contributes a
random secret polynomial, broadcasts a commitment to the
polynomial coefficients, exchanges encrypted share evaluations,
and verifies received shares against commitments. The blockchain
serves as the authenticated public bulletin board for commitment
publication, anchoring key generation in the same trust layer as
identity registration.

The DKG protocol producing individual signing key pairs per RSU
--- one secret key and one public key per participant --- without
any central authority or single dealer is shown in
Equation~\ref{eq:dkg_trs}:
\begin{equation}
  \bigl(\{sk_j,\; pk_j\}_{j=1}^{n}\bigr)
  \leftarrow \mathrm{DKG}(\mathcal{R},\; n,\; t_{\mathrm{sign}})
  \label{eq:dkg_trs}
\end{equation}
As shown in Equation~\ref{eq:dkg_trs}, each RSU receives its own
individual key pair $(sk_j, pk_j)$. The per-RSU public keys
$\{pk_j\}_{j=1}^{n}$ are used directly in TRS verification via
Eq.~\eqref{eq:trs_verify}, maintaining full consistency with the
verification step. No RSU and no controller observes any other
RSU's secret key at any point during or after key generation.

The distributed FHE key generation producing shares for all
participants without any single dealer is shown in
Equation~\ref{eq:dkg_fhe}:
\begin{equation}
  (pk_{\mathrm{FHE}},\; \{sk_j^{\mathrm{share}}\}_{j=1}^{n},\;
   sk_{\mathrm{cloud}}^{\mathrm{share}},\; evk)
  \leftarrow \mathrm{FHE\text{-}DKG}
  (\mathcal{R} \cup \{\mathrm{Cloud}\},\; n+1,\;
   t_{\mathrm{decrypt}})
  \label{eq:dkg_fhe}
\end{equation}
As shown in Equation~\ref{eq:dkg_fhe}, no controller participates
in TRS or FHE key generation at any stage. Vehicle session keys
$K_i$ continue to be derived locally at each RSU via
Eq.~\eqref{eq:session_key}, as the vehicle--RSU session key
relationship is independent of ring-level key generation. The FHE
key shares $\{sk_j^{\mathrm{share}}\}_{j=1}^{n}$ produced during
threshold setup are distributed to RSU ring members via the same
authenticated LKH ring key synchronisation channel, eliminating
the need for a separate key distribution protocol for FHE share
delivery.


\subsection{Dual-Mode Detection and Mitigation Pipeline}
\label{sec:dual_mode}

The SDVN controller contributes full-mode GAT and temporal autoencoder
anomaly scores to the blockchain as one untrusted peer submission, subject
to the same BFT verification as RSU submissions
(Equation~\ref{eq:ctrl_evidence}). Since the threat model includes a
compromised controller, no blockchain decision trust score update,
revocation, or flow rule enforcement relies exclusively on
controller-supplied evidence. RSUs independently submit beacon-level
anomaly evidence directly to the blockchain
(Equation~\ref{eq:rsu_evidence}), and revocation is triggered by RSU
peer consensus alone (Equation~\ref{eq:bft_revoke_rsu}).

The proposed framework operates in two modes. The lightweight mode
executes entirely at the RSU using rule-based signature detection and
HMAC-based integrity protection, requiring no AI inference. The full
mode augments this with GAT-based spatial detection and temporal autoencoder
anomaly scoring at the SDVN controller, together with post-quantum TRS and FHE
cryptographic mitigation. Both modes share the blockchain trust layer and
smart contract arbitration.

Figure~\ref{fig:dual_mode_arch} presents the dual-mode system
architecture of the proposed MPTD-PQS framework.

\begin{figure}[H]
    \centering
    \includegraphics[width=\linewidth]{images/Fig1diag}
    \caption{Dual-mode system architecture of the proposed MPTD-PQS
    framework.}
    \label{fig:dual_mode_arch}
\end{figure}

As shown in Figure~\ref{fig:dual_mode_arch}, the architecture is
organised into three horizontal layers that reflect the actual
deployment boundaries of the framework. At the bottom, the vehicle
layer contains five representative vehicles that broadcast BSM beacons
over the V2I channel and establish per-vehicle session keys $K_i$
through a KDF exchange with their serving RSU. The middle layer,
enclosed by the lightweight mode boundary (orange dashed border),
contains the RSU cluster comprising RSU~1 through RSU~3. Each RSU
executes an identical four-stage pipeline: the HMAC Verifier
(Equation~\ref{eq:hmac}) checks beacon integrity first, the Signature
Checker evaluates all nine rule-based signatures TP-S1 through MP-S4
via the \texttt{LW-DETECT} procedure, the TRS Partial Signer
(Equation~\ref{eq:trs_partial}) accumulates partial threshold signatures
from peer RSUs over the LKH ring key synchronisation channel, and the
LKH Key Store maintains the per-vehicle session keys $K_i$ and the
RSU signing key $sk_j$. When an anomaly is detected, each RSU submits
a direct evidence tuple $\mathcal{E}_j(t)$ to the Blockchain Smart
Contracts (SC-Trust / SC-Revoke) via the red direct-submission
channel; the blockchain issues a CRL update back to the RSU cluster
by direct broadcast without routing through the controller. Beacon
windows are stored off-chain in the IPFS Off-Chain Beacon Store
together with their cryptographic hashes, decoupling real-time RSU
processing from controller-side deep learning analysis. The top layer,
enclosed by the full mode boundary (blue solid border), contains the
SDVN Controller, which retrieves beacon windows $\{\mathbf{X}_i(t)\}$
from IPFS, computes GAT spatial anomaly scores and temporal autoencoder
reconstruction errors, and fuses them via the Fusion Engine
(Equation~\ref{eq:fusion}) before the Decision Engine submits a single
untrusted peer evidence entry to the blockchain indicated by the
grey dashed arrow labelled Untrusted Peer Submission not
required for revocation. On the right, FHE-encrypted, TRS-verified
aggregates are forwarded to the Cloud/ITS Server over the encrypted
channel only after the threshold signature $\sigma_{\mathrm{TRS}}$
has been assembled across the RSU cluster, ensuring privacy and
structural integrity guarantees are enforced before any data leaves
the RSU boundary.

Before detailing the individual detection, cryptographic, and
trust-management mechanisms of the proposed framework, this section
consolidates the mathematical notation used throughout the remainder
of the chapter. Table~\ref{tab_notations} presents a comprehensive
summary of every symbol introduced in the design, organised so that
each notation can be quickly located alongside the equation in which
it is first defined. Grouping the symbols in this manner allows the
subsequent formulations---which span the lightweight rule-based
signatures, the deep-learning detection pipeline, and the
post-quantum cryptographic primitives---to be read without repeatedly
restating the meaning of each variable.

\footnotesize
\begin{longtable}{|m{3.6cm}|m{9.0cm}|}
	\caption{A summary of the notations.}\label{tab_notations}\\
	\hline
	\textbf{Notation} & \textbf{Description} \\
	\hline
	\endfirsthead

	\multicolumn{2}{c}{\tablename\ \thetable\ -- continued from previous page} \\
	\hline
	\textbf{Notation} & \textbf{Description} \\
	\hline
	\endhead

	\hline
	\multicolumn{2}{r}{continued on next page} \\
\hline
	\endfoot

	\hline
	\endlastfoot

	$v_i,\ r_j,\ c_k$; $\mathcal{R},\ \mathcal{C}$; $\mathcal{R}_{\text{trusted}}(t),\ \mathcal{C}_{\text{trusted}}(t)$; $c_{\text{assigned}}(r_j)$; $\mathcal{V}_j(t)$ & Vehicle $i$, RSU $j$, controller $k$; set of RSUs and set of $K$ registered controllers $\mathcal{C}=\{c_1,\dots,c_K\}$; trusted-RSU endorser set and trusted-controller set at time $t$; RSU $r_j$'s currently assigned controller; vehicles within RSU $r_j$'s coverage zone \\
\hline
	$f,\ n$; $t_{\text{sign}},\ t_{\text{decrypt}}$; $\mathcal{S},\ \mathcal{D}$; $M$; $T_w,\ T_{\text{rev}}$ & Tolerated Byzantine faults and number of RSUs ($n\!\ge\!3f{+}1$); TRS signing threshold ($f{+}1$) and FHE decryption threshold ($f{+}2$); TRS signing subset and decryption coalition; number of ring aggregates; revocation temporal window and required low-trust epochs \\
\hline

$\mathcal{R}^{\mathrm{obs}}_i(t)$;\;
$\mathcal{R}^{\mathrm{obs}}_{c_k}(t)$ &
Trusted RSUs that observed vehicle $v_i$ in epoch $t$:
$\{r_j \in \mathcal{R}_{\mathrm{trusted}}(t)
 : v_i \in \mathcal{V}_j(t)\}$;
trusted RSUs assigned to controller $c_k$ in epoch $t$:
$\{r_j \in \mathcal{R}_{\mathrm{trusted}}(t)
 : c_{\mathrm{assigned}}(r_j) = c_k\}$.
Normalising denominators in
Eqs.~\eqref{eq:trust_update}--\eqref{eq:ctrl_trust};
trust is unchanged when the observing set is empty. \\
\hline
	$\mathbf{b}_i(t)$; $p_i(t),\ s_i(t),\ \theta_i(t),\ a_i(t)$; $\text{ID}_i$; $T_b$ & BSM beacon tuple of vehicle $i$ (Eq.~3.10); GPS position, speed, heading angle, and acceleration; pseudonymous vehicle identifier; beacon broadcast interval (100\,ms) \\
\hline
	$\tau_c,\ N_b,\ N_{\min},\ \lambda_h$; $R_j,\ d_\perp,\ \mathcal{A}_j$; $\rho_v,\ d_{\min}$; $K_{\text{sybil}},\ f_{\text{sybil}}$; $\boldsymbol{\delta}(t),\ \epsilon_{\max}$ & RSU contact duration, observable-beacon count, minimum detectable sequence length, handover rate (Eq.~3.2--3.5); RSU coverage radius, perpendicular distance, coverage area; traffic density and minimum inter-vehicle spacing; Sybil identity capacity and Sybil fraction (Eq.~3.8--3.9); injected position-offset drift and per-beacon drift bound \\
\hline
	$s_{\max},\ \omega_{\max},\ a_{\max}$; $\hat{p}_i(t),\ r_i(t),\ D_i(k),\ \delta_{\text{th}}$; $D_{\text{KL}},\ \mathcal{P}_t,\ \mathcal{P}_{\text{hist}},\ \kappa_{\text{th}}$; $\tau_{\text{sync}},\ \rho_{\text{sync}}$ & Road-class maximum speed, heading rate, and acceleration bounds (TP-S1--S3); dead-reckoned position, residual, cumulative drift score and its threshold (TP-S4--S5); KL-divergence of current vs.\ historical speed distribution and its threshold (MP-S3); beacon-timing and co-occurrence thresholds (MP-S2) \\

\hline
$\psi_i(t)$;\;
$\bar{\psi}_i(t)$;\;
$\psi_{\mathrm{th}}$;\;
$w_s^{(k)},\;w_s$;\;
$\theta_{\mathrm{conf}}^{(k)}$;\;
$\bar{\lambda}_j$ &
Composite rule-based anomaly score per beacon
(Eq.~\eqref{eq:psi_score}); its per-detection-window mean
$\bar{\psi}_i(t) = \tfrac{1}{L}\sum_{l=1}^{L}\psi_i(t_l)$,
normalised as $\min(\bar{\psi}_i/\psi_{\mathrm{th}},1)$ in the
full-mode fusion (Eq.~\eqref{eq:fusion}); its rule-based alert
threshold; attack-type-conditioned signature reliability weight
for signature $s$ under class $k$ (full mode) and aggregate
calibrated weight (lightweight mode);
per-attack-class GAT confidence gate threshold;
fallback fusion weight (average over all $K$ sets) \\
\hline
	$K_i,\ \nu_i,\ \Delta_{\text{HMAC}}$; $\text{MAC}_i(t)$ & Vehicle session key, monotonic anti-replay nonce, freshness window; HMAC beacon-integrity tag (Eq.~3.38) \\
\hline
	$\mathcal{T},\ \mathcal{K}_u,\ \mathcal{K}_{\text{root}}$; $\eta_i,\ N_{\text{rekey}}$; $sk_j,\ pk_j$ & LKH key tree, per-node key and root (group) key; session nonce and logarithmic rekey cost (Eq.~3.22--3.24); RSU TRS signing/verification key pair from DKG (Eq.~3.25) \\
\hline
	$\mathcal{G}(t),\mathcal{V}(t),\mathcal{E}(t)$; $R_{\max},\ \phi_{\max},\ \mathbf{v}_i$; $\tilde{x}_i^{(f)},\ \alpha_{ij}^{(h)}$; $\mathbf{W}_h,\mathbf{a}_h,H$; $\mathbf{x}_i',\ \mathcal{S}_i(t)$ & Dynamic vehicle--RSU graph with nodes and mobility-aware edges; max.\ range, heading-divergence bound and velocity vector (Eq.~3.27); z-score--normalised feature and head-$h$ attention coefficient; head weight matrix, attention vector, number of heads; node embedding and GAT spatial anomaly score (Eq.~3.39) \\
\hline
$y_i^{(k)},\;\hat{y}_i^{(k)}$;\;
$K$;\;
$\hat{k}_i$;\;
$\mathbf{w}_{\mathrm{cls}}^{(k)},\,b_{\mathrm{cls}}^{(k)}$;\;
$w_k^{+},\,w_k^{-}$;\;
$\mathcal{V}_{\mathrm{train}},\,\mathcal{V}_k^{+},\,\mathcal{V}_k^{-}$;\;
$\mathcal{L}_{\mathrm{GAT}}$ &
Ground-truth and predicted binary node label for attack class $k$
($y_i^{(k)}{=}1$ attacking with class $k$, $y_i^{(k)}{=}0$ normal);
number of attack classes ($K{=}7$);
predicted attack type (Eq.~\eqref{eq:gat_atk_type});
per-class GAT classification head parameters, discarded after training
(Eq.~\eqref{eq:gat_cls});
per-class inverse-frequency class weights
(Eq.~\eqref{eq:gat_loss});
labeled training node set and its per-class positive and negative
subsets;
multi-task GAT supervised cross-entropy training loss
(Eq.~\eqref{eq:gat_loss}) \\
\hline
	$\mathbf{X}_i(t),\ L,\ d$; $\mathbf{f}_l,\mathbf{u}_l,\mathbf{g}_l,\mathbf{o}_l,\mathbf{c}_l,\mathbf{h}_l$; $\hat{\mathbf{X}}_i(t),\ \mathcal{L}_{\text{total}},\ \beta$; $\varepsilon_i(t),\ \theta_{ae},\ \theta_{\mathcal{S}},\ \kappa,\ \alpha_{\text{freeze}}$ & Beacon-window input, window length and feature dimension; LSTM forget/input/candidate/output gates, cell and hidden states (Eq.~3.31--3.36); reconstruction, mobility-constrained loss and kinematic-penalty weight (Eq.~3.37); reconstruction error, offline MAD-based detection threshold, GAT spatial normalisation threshold (Eq.~\ref{eq:fusion}) scale factor and manifold-freeze threshold (Eq.~3.42--3.43) \\
\hline
   $\Phi_i(t),\;
    \lambda_j^{(k)},\;
    \Phi_{\mathrm{th}}$;\;
    $\theta_{\mathcal{S}}$;\;
    $\mathrm{flag}_i(t)$;\;
    $\mathrm{flag}_i^{\mathrm{GAT}}$ &
    Fused anomaly score in $[0,1]$; attack-conditioned fusion weight
    for component $j$ under attack class $k$
    ($\sum_j\lambda_j^{(k)}{=}1$, $\lambda_j^{(k)}{\geq}0.05$;
    $K$ sets learned in Phase~2, selected at inference by
    $\hat{k}_i$, Eq.~\eqref{eq:fusion}--\eqref{eq:flag});
    GAT spatial score normalisation threshold (Phase~1b, locked);
    binary detection flag (Eq.~\eqref{eq:flag});
    GAT multi-task binary detection flag (Eq.~\eqref{eq:gat_det_flag}) \\
\hline
	$\mathbf{A}_j(t),\ c_j,$\newline$\text{Enc}(\mathbf{A}_{\text{ring}})$;\newline$pk_{\text{FHE}},\ sk_j^{\text{share}},$\newline$sk_{\text{cloud}}^{\text{share}},\ evk$ & RSU local plaintext aggregate, its FHE ciphertext, and homomorphic ring aggregate (Eq.~3.46--3.48); FHE public key, per-RSU and Cloud threshold key shares, and evaluation key (Eq.~3.26, 3.56) \\
\hline
	$m,\ \nu_{\mathcal{S}},\ \text{ID}_{\mathcal{S}},\ h(\mathcal{S})$; $\sigma_j,\ \sigma_{\text{TRS}},\ \Delta_{\text{TRS}}$ & Signed TRS message and its ring nonce, ring identity and member-set hash (Eq.~3.49); per-RSU partial signature, aggregated threshold ring signature, and verification freshness window (Eq.~3.50--3.53) \\
\hline
	$\bar{X}_{\text{global}},\ pd_j,\ pd_{\text{cloud}}$; $L^{\text{SUMO}},\ U^{\text{SUMO}}$ & Decrypted global aggregate, per-RSU and Cloud partial decryptions (Eq.~3.54--3.58); offline SUMO-derived lower/upper plausibility-envelope bounds \\
\hline
$\tau_i(t),\ \tau_{c_k}(t),\ \alpha,$\newline
$\tau_{\min},\ \tau_{\text{warn}}$;\newline
$\mathcal{E}_j(t),\ \mathcal{E}_c(t),$\newline
$\sigma_j^{\text{sub}},\ \sigma_c^{\text{sub}}$;\newline
$\text{conflict}_j(t),\ \text{flag}_c$;\newline
$n_{\text{endorse}},\ \tau_{\text{init}}$ &
Vehicle and controller trust scores, EMA decay factor, revocation and probationary thresholds (Eq.~3.59--3.60); RSU and controller evidence tuples with their signatures (Eq.~3.61--3.62); directional controller--RSU conflict indicator and controller-compromise flag (Eq.~3.66--3.67); registration endorsement count and initial trust (Alg.~8) \\
\hline

\end{longtable}
\normalsize

As summarised in Table~\ref{tab_notations}, the notations are grouped
by their functional role within the framework. The first group defines
the core network entities---the vehicles, RSUs, and the set of
registered controllers together with the trusted-controller set
$\mathcal{C}_{\text{trusted}}(t)$ and the per-RSU controller
assignment. The next group captures the Byzantine fault-tolerance and
threshold parameters that govern the signing and decryption coalitions,
followed by the beacon and kinematic quantities that constitute each
vehicle's reported state. Subsequent groups cover the nine rule-based
trajectory- and mobility-poisoning signatures and their composite score
$\psi_i(t)$, the HMAC and LKH key-management variables, and the
deep-learning components comprising the GAT spatial detector, the LSTM
temporal autoencoder, and their fused anomaly score $\Phi_i(t)$. The
final groups list the post-quantum threshold cryptographic primitives
---spanning the per-RSU FHE encryption, the homomorphic ring aggregate,
the threshold ring signature $\sigma_{\text{TRS}}$, and the threshold
decryption coalition---together with the blockchain trust-management
variables that drive the SC-Trust and SC-Revoke decisions. Each symbol
is annotated with the equation in which it first appears, so that the
table serves as a single reference point for the detailed formulations
presented in the following subsections.

Figure~\ref{fig:seq_lightweight} presents the UML sequence diagram for the lightweight operating mode, tracing the per-beacon message flow across the vehicle, RSU, SC-Trust smart contract, and controller entities during a single beacon epoch.

\begin{figure}[H]
    \centering
    \includegraphics[width=0.8\linewidth]{images/Fig2diag}
    \caption{UML sequence diagram for lightweight real-time detection
    across one beacon epoch.}
    \label{fig:seq_lightweight}
\end{figure}

As shown in Figure~\ref{fig:seq_lightweight}, the sequence begins
with the vehicle broadcasting $\mathbf{b}_i(t) \| \mathrm{MAC}_i(t)$,
where the HMAC tag is computed according to Equation~\ref{eq:hmac}
and binds the beacon payload, timestamp, and vehicle identity under
session key $K_i$. The RSU first verifies this tag; any failure causes
immediate rejection before the rule-based signatures are evaluated,
forming a fail-fast integrity gate. Beacons that pass verification
are then scored against the nine rule-based signatures TP-S1 through
MP-S4 (Equations~\ref{eq:tp_s1}--\ref{eq:mp_s4}), yielding the
composite anomaly score $\psi_i(t)$ from Equation~\ref{eq:psi_score}.
The diagram's alt fragment captures the two outcomes of this
check. On the normal path ($\psi_i(t) \leq \psi_{\mathrm{th}}$), the
beacon is accepted locally at the RSU with no blockchain interaction
and no controller involvement, keeping latency well within the
100\,ms beacon interval $T_b$. On the alert path
($\psi_i(t) > \psi_{\mathrm{th}}$), the RSU submits a signed evidence
tuple $\mathcal{E}_j(t)$ directly to \texttt{SC-Trust}, which updates
the vehicle's cumulative trust score $T_i$ and evaluates the
revocation condition: revocation is issued only if $T_i < T_{\min}$
has persisted for $T_{\mathrm{rev}}$ epochs and at least
$2f+1$ RSUs have submitted corroborating evidence, satisfying the BFT
endorsement condition in Equation~\ref{eq:bft_revoke_rsu}. When both
conditions hold, \texttt{SC-Trust} emits \texttt{Revoke($\mathrm{ID}_i$)}
and broadcasts a CRL update directly to all RSUs, bypassing the
controller to ensure revocation propagates even under controller
compromise.

\subsubsection{GAT Architecture Design for Vehicle--RSU Mobility Graphs}
\label{sec:gat_design}

\textbf{Graph Construction.}
The vehicle--RSU communication graph $\mathcal{G}(t)$ is constructed with
a mobility-aware edge definition. An edge $(i,j) \in \mathcal{E}(t)$ is
instantiated between nodes $v_i$ and $v_j$ (or $v_i$ and $r_j$) only when
two conditions are jointly satisfied. The mobility-aware edge definition
for the vehicle--RSU communication graph is given in
Equation~\ref{eq:edge_def1}.
\newline
\newline
This edge definition restricts the GAT's attention computation to
kinematically relevant vehicle pairs those within communication
range and traveling in compatible directions preventing the spatial
anomaly detector from being confused by oncoming or crossing traffic
whose mobility patterns are legitimately dissimilar.
\begin{equation}
  (i,j) \in \mathcal{E}(t) \iff
  \|p_i(t) - p_j(t)\| \leq R_{\max}
  \;\land\;
  \cos^{-1}\!\left(\frac{\mathbf{v}_i \cdot \mathbf{v}_j}
    {\|\mathbf{v}_i\|\|\mathbf{v}_j\|}\right) \leq \phi_{\max}
  \label{eq:edge_def1}
\end{equation}
As shown in Equation~\ref{eq:edge_def1}, this edge definition
excludes kinematically irrelevant neighbours (e.g., oncoming vehicles on
opposite lanes) from the attention computation, reducing noise in the
spatial anomaly score where $R_{\max}$ is the maximum V2V/V2I communication range,
$\mathbf{v}_i = s_i[\cos\theta_i, \sin\theta_i]^\top$ is the velocity
vector of $v_i$, and $\phi_{\max}$ is a maximum heading divergence
threshold.
\newline
\newline
\textbf{Node Feature Normalization.}
BSM fields have heterogeneous scales (GPS coordinates in degrees, speed
in m/s, acceleration in m/s$^2$). To prevent scale dominance in the
attention computation, the rolling z-score normalization for node features
is given in Equation~\ref{eq:feature_norm1}.
\newline
\newline
This normalization ensures that GPS coordinates, speed, and
acceleration contribute equally to the attention computation
regardless of their raw scale differences, preventing larger-magnitude
features from dominating the spatial anomaly score and obscuring
subtle mobility pattern deviations.
\begin{equation}
  \tilde{x}_i^{(f)}(t) = \frac{x_i^{(f)}(t) - \mu_f(t)}{\sigma_f(t) + \epsilon}
  \label{eq:feature_norm1}
\end{equation}
As shown in
Equation~\ref{eq:feature_norm1}, each feature is normalised relative
to the current population of vehicles in the RSU zone, making the
representation robust to varying traffic densities where $\mu_f(t)$ and $\sigma_f(t)$ are the rolling mean and standard
deviation of feature $f$ across all vehicles in $\mathcal{V}_j(t)$,
and $\epsilon = 10^{-8}$ prevents division by zero.
\newline
\newline
\textbf{Multi-Head Attention.}
To capture multiple relational aspects of vehicle mobility simultaneously
(e.g., speed correlation, heading alignment, proximity), the normalized
attention coefficient computation for the $h$-th head is given in
Equation~\ref{eq:multihead_attn1}.
\newline
\newline
This attention coefficient determines how much each neighbouring
vehicle influences a node's spatial embedding, enabling the GAT to
automatically learn which spatial relationships  such as platoon
spacing or intersection right-of-way  are most diagnostic of normal
versus poisoned mobility patterns without manual feature engineering.
\begin{equation}
  \alpha_{ij}^{(h)} = \frac{\exp\!\left(
    \mathrm{LeakyReLU}(\mathbf{a}_h^\top
    [\mathbf{W}_h\tilde{\mathbf{x}}_i \| \mathbf{W}_h\tilde{\mathbf{x}}_j])
  \right)}
  {\sum_{k\in\mathcal{N}(i)}\exp\!\left(
    \mathrm{LeakyReLU}(\mathbf{a}_h^\top
    [\mathbf{W}_h\tilde{\mathbf{x}}_i \| \mathbf{W}_h\tilde{\mathbf{x}}_k])
  \right)}
  \label{eq:multihead_attn1}
\end{equation}
\newline
\newline
As shown in Equation ~\ref{eq:multihead_attn1}, this formula computes how much attention node i should pay to its neighbour node j during the h-th attention head  essentially a relevance score between two vehicles (or a vehicle and an RSU) in the mobility graph.
\newline
\newline
The multi-head node embedding is given in
Equation~\ref{eq:multihead_embed1} where $\mathbf{W}_h$ is the learnable weight matrix for head $h$,
$\mathbf{a}_h$ is the attention weight vector, and $\mathcal{N}(i)$
is the neighbourhood of node $i$ defined by Equation~\ref{eq:edge_def1}.
As shown in Equation~\ref{eq:multihead_attn1}, the softmax-normalised
attention coefficient $\alpha_{ij}^{(h)}$ measures the relative
importance of neighbour $j$ to node $i$ for head $h$.
\begin{equation}
  \mathbf{x}'_i = \Big\|_{h=1}^{H}
  \sigma\!\left(\sum_{j\in\mathcal{N}(i)} \alpha_{ij}^{(h)}
  \mathbf{W}_h\tilde{\mathbf{x}}_j\right)
  \label{eq:multihead_embed1}
\end{equation}
As shown in
Equation~\ref{eq:multihead_embed1}, the embedding $\mathbf{x}'_i$
aggregates attention-weighted neighbour features from all heads,
capturing diverse spatial relationships simultaneously. The final
layer uses averaging instead of concatenation to produce a fixed-size
output regardless of $H$ where $\|$ denotes concatenation across all $H$ heads.
\newline
\newline
\textbf{GAT Training Objective.}
\newline
\newline
The attention weights $\mathbf{W}_h$ and $\mathbf{a}_h$ are trained
under $K$ parallel supervised binary cross-entropy tasks, one per
attack class, using ground-truth per-attack-type labels obtained
directly from the NS-3 attack configuration. For each vehicle $v_i$
in the labeled training node set $\mathcal{V}_{\mathrm{train}}$, the
ground-truth label for attack class $k$ is $y_i^{(k)} \in \{0,1\}$,
where $y_i^{(k)} = 1$ denotes that vehicle $v_i$ is executing an
attack of type $k$ at the current timestep; $K=7$ matches the seven
attack variants defined in the threat model
(Section~\ref{sec:threat_model}). A separate linear classification
head maps the shared multi-head embedding $\mathbf{x}'_i$ to a
per-attack detection probability. Multi-task training allows the
shared encoder to benefit from all $K$ tasks simultaneously, providing
each attack-specific head with the relational spatial features it
needs while enriching the shared representation through cross-attack
regularisation. The per-attack-type detection probability produced by
classification head $k$ is given in Equation~\ref{eq:gat_cls}.
\newline
\newline
This formulation gives the proposed method the same discriminative
information per attack class as a specialist baseline trained on that
class, while additionally benefiting from the shared encoder trained
across all classes simultaneously --- an advantage unavailable to
per-attack specialist architectures.
\newline
\newline
This prediction converts the learned spatial embedding into a
node-level maliciousness score, enabling the cross-entropy loss to
directly supervise the attention weights to separate normal and
anomalous spatial graph configurations.
\begin{equation}
  \hat{y}_i^{(k)} = \sigma\!\left(
    \mathbf{w}_{\mathrm{cls}}^{(k)\top} \mathbf{x}'_i
    + b_{\mathrm{cls}}^{(k)}\right),
    \quad k = 1,\ldots,K
  \label{eq:gat_cls}
\end{equation}
As shown in Equation~\ref{eq:gat_cls}, each head produces an
independent sigmoid probability $\hat{y}_i^{(k)}$ from a linear
projection of the shared embedding $\mathbf{x}'_i$, where
$\mathbf{w}_{\mathrm{cls}}^{(k)}$ and $b_{\mathrm{cls}}^{(k)}$ are
the weight vector and bias of the $k$-th head.
At inference, the predicted attack type is given in
Equation~\ref{eq:gat_atk_type}:
\begin{equation}
  \hat{k}_i = \arg\max_{k \in \{1,\ldots,K\}} \hat{y}_i^{(k)}
  \label{eq:gat_atk_type}
\end{equation}
As shown in Equation~\ref{eq:gat_atk_type}, $\hat{k}_i$ identifies
the most likely attack class for node $v_i$, which is used to select
the attack-conditioned fusion weights in Equation~\ref{eq:fusion}.
The binary detection flag from the GAT is given in
Equation~\ref{eq:gat_det_flag}:
\begin{equation}
  \mathrm{flag}_i^{\mathrm{GAT}} =
    \bigvee_{k=1}^{K}\,
    \mathbf{1}\!\left[\hat{y}_i^{(k)} > 0.5\right]
  \label{eq:gat_det_flag}
\end{equation}
As shown in Equation~\ref{eq:gat_det_flag}, vehicle $v_i$ is flagged
by the GAT component if any single attack-type head exceeds the
detection threshold; the combined binary detection is strictly at
least as sensitive as any individual head.
\newline
\newline
The multi-task class-weighted binary cross-entropy training loss,
summed across all $K$ attack-type heads, is given in
Equation~\ref{eq:gat_loss}.
\newline
\newline
Summing $K$ independent tasks trains the shared encoder on the full
diversity of attack signatures simultaneously. Per-task class weights
$w_k^+, w_k^-$ are computed independently per attack type, preventing
rare attack classes from being dominated by more frequent ones.
\begin{equation}
  \mathcal{L}_{\mathrm{GAT}} = \sum_{k=1}^{K}
    \frac{-1}{|\mathcal{V}_{\mathrm{train}}|}
    \sum_{v_i \in \mathcal{V}_{\mathrm{train}}}
    \Bigl[
      w_k^{+}\,y_i^{(k)} \log \hat{y}_i^{(k)}
      + w_k^{-}\,(1-y_i^{(k)})\log(1-\hat{y}_i^{(k)})
    \Bigr]
  \label{eq:gat_loss}
\end{equation}
As shown in Equation~\ref{eq:gat_loss},
$\mathcal{L}_{\mathrm{GAT}}$ is the sum of $K$ independent
class-weighted binary cross-entropy losses, where
$w_k^{+} = |\mathcal{V}_{\mathrm{train}}|\,/\,(2|\mathcal{V}_k^{+}|)$
and $w_k^{-} = |\mathcal{V}_{\mathrm{train}}|\,/\,(2|\mathcal{V}_k^{-}|)$
are the inverse-frequency class weights for attack type $k$, with
$\mathcal{V}_k^{+}$ and $\mathcal{V}_k^{-}$ the positive and negative
subsets for class $k$ respectively.
\newline
\newline
Training proceeds over the first temporal segment of each SUMO
trace, run through NS-3 with attack injection active so that
ground-truth labels are available from the simulator configuration.
After training converges, the classification head
$(\mathbf{w}_{\mathrm{cls}}, b_{\mathrm{cls}})$ is discarded and
only the encoder weights
$\{\mathbf{W}_h,\,\mathbf{a}_h\}_{h=1}^{H}$ are retained. The
same temporal segment is then re-run through NS-3 with no attack
injection to produce clean-condition embeddings, from which the
per-node statistics $\bar{\mathbf{x}}'_i$ and
$\boldsymbol{\sigma}'_i$ are computed and locked before evaluation.


\subsubsection{Temporal Autoencoder Design for Mobility Sequences}

\textbf{LSTM Gating Formulation.}
The encoder LSTM processes the beacon sequence
$\mathbf{X}_i(t) \in \mathbb{R}^{L\times d}$ step by step. At each
step $l$, the LSTM gating equations governing information flow are
given in Equations~\ref{eq:lstm_forget}--\ref{eq:lstm_hidden}.
\begin{align}
  \mathbf{f}_l &= \sigma(\mathbf{W}_f[\mathbf{h}_{l-1},
    \mathbf{b}_i(t_l)] + \mathbf{c}_f)
  \label{eq:lstm_forget} \\
  \mathbf{u}_l &= \sigma(\mathbf{W}_u[\mathbf{h}_{l-1},
    \mathbf{b}_i(t_l)] + \mathbf{c}_u)
  \label{eq:lstm_input} \\
  \mathbf{g}_l &= \tanh(\mathbf{W}_g[\mathbf{h}_{l-1},
    \mathbf{b}_i(t_l)] + \mathbf{c}_g)
  \label{eq:lstm_gate} \\
  \mathbf{c}_l &= \mathbf{f}_l \odot \mathbf{c}_{l-1}
    + \mathbf{u}_l \odot \mathbf{g}_l
  \label{eq:lstm_cell} \\
  \mathbf{o}_l &= \sigma(\mathbf{W}_o[\mathbf{h}_{l-1},
    \mathbf{b}_i(t_l)] + \mathbf{c}_o)
  \label{eq:lstm_output} \\
  \mathbf{h}_l &= \mathbf{o}_l \odot \tanh(\mathbf{c}_l)
  \label{eq:lstm_hidden}
\end{align}

As shown in Equation~\ref{eq:lstm_forget}, the forget gate $\mathbf{f}_l$
determines what portion of the previous cell state is discarded based on
the current beacon input. As shown in Equation~\ref{eq:lstm_input}, the
input gate $\mathbf{u}_l$ controls how much new information is written
into the cell. As shown in Equation~\ref{eq:lstm_gate}, the candidate
update $\mathbf{g}_l$ computes the new content to be potentially stored
in the cell state. As shown in Equation~\ref{eq:lstm_cell}, the cell
state $\mathbf{c}_l$ is updated by combining the forget-gated previous
state with the input-gated candidate, enabling the encoder to retain
long-range temporal dependencies across the beacon sequence. As shown in
Equation~\ref{eq:lstm_output}, the output gate $\mathbf{o}_l$ filters
what portion of the cell state is exposed as the hidden state. As shown
in Equation~\ref{eq:lstm_hidden}, the hidden state $\mathbf{h}_l$ is
produced by applying the output gate to the cell state, forming the
temporal representation passed to the next step. The encoder latent
representation is the final hidden state: $\mathbf{h}_i(t) = \mathbf{h}_L$ where $\mathbf{f}_l$, $\mathbf{u}_l$, $\mathbf{o}_l$ are the forget,
input, and output gates; $\mathbf{c}_l$ is the cell state; $\mathbf{g}_l$
is the candidate cell update; and $\odot$ is element-wise multiplication.
\newline
\newline
\textbf{Mobility-Constrained Reconstruction Loss.}
Standard autoencoders minimize only reconstruction error. Since vehicle
mobility follows physical constraints
(Equations~\ref{eq:tp_s1}--\ref{eq:tp_s3}), the mobility-constrained
reconstruction loss used during training is given in
Equation~\ref{eq:ae_loss1}.
\newline
\newline
This loss function trains the autoencoder to reconstruct only
physically plausible mobility sequences by penalizing reconstructions
that violate kinematic constraints, ensuring that the learned normal
manifold is sensitive to the specific class of mobility-consistent
yet collectively falsified sequences that poisoning attacks produce.
\begin{equation}
  \mathcal{L}_{\mathrm{total}} = \underbrace{\frac{1}{L}\sum_{l=1}^{L}
    \|\mathbf{b}_i(t_l) - \hat{\mathbf{b}}_i(t_l)\|^2}_{\text{reconstruction}}
  + \beta \underbrace{\frac{1}{L}\sum_{l=2}^{L}
    \max\!\left(0,\; r_i(t_l) - \delta_{\mathrm{th}}\right)^2}_{\text{kinematic penalty}}
  \label{eq:ae_loss1}
\end{equation}
 As shown in
Equation~\ref{eq:ae_loss1}, the total loss combines standard
reconstruction error with a kinematic penalty term, ensuring the
autoencoder learns a normal mobility manifold that is physically
consistent and sensitive to kinematically implausible reconstructions
produced by poisoned sequences where $\beta$ controls the kinematic
penalty weight and $r_i(t_l)$ is the dead-reckoning residual from
Equation~\ref{eq:tp_s4}.

\subsubsection{Lightweight Mode: HMAC Beacon Integrity and Rule-Based Detection}

Each vehicle $v_i$ appends an HMAC tag to its BSM beacon using
session key $K_i$ established during RSU registration. The tag
includes a per-vehicle monotonically increasing nonce $\nu_i$
alongside the timestamp to prevent replay both outside and within
any freshness window. The HMAC beacon integrity tag binding the
beacon content, timestamp, identity, and nonce is given in
Equation~\ref{eq:hmac}:
\begin{equation}
  \mathrm{MAC}_i(t) = \mathrm{HMAC}_{K_i}
  \!\bigl(\mathbf{b}_i(t) \,\|\, t \,\|\,
  \mathrm{ID}_i \,\|\, \nu_i\bigr)
  \label{eq:hmac}
\end{equation}
As shown in Equation~\ref{eq:hmac}, the tag binds the beacon
content, timestamp, vehicle identity, and nonce under $K_i$, so
any in-transit modification invalidates the tag. The RSU rejects
any beacon whose nonce $\nu_i$ is not strictly greater than the
last accepted nonce $\nu_i^{\mathrm{last}}$ for $\mathrm{ID}_i$.
At RSU handoff, the RSU retrieves $\nu_i^{\mathrm{last}}$ from
a fast shared RSU-cluster cache passed RSU-to-RSU at handoff,
keeping nonce lookup within the 100\,ms beacon budget while
providing consistent nonce state across zone boundaries. The RSU verifies $\mathrm{MAC}_i(t)$
upon beacon receipt, rejecting any beacon that fails verification
before applying the rule-based anomaly score $\psi_i(t)$ from
Equation~\ref{eq:psi_score}. An alert is issued if
$\psi_i(t) > \psi_{\mathrm{th}}$, and the vehicle's trust score is
immediately decremented on-chain via a lightweight smart contract call.

Algorithm~\ref{alg:lightweight} presents the LW-DETECT procedure,
which forms the complete lightweight mode detection pipeline executed
at the RSU for every incoming beacon within a single 100\,ms epoch.

\begin{algorithm}[H]
\caption{Lightweight Rule-Based Beacon Detection (\texttt{LW-DETECT})}
\begin{algorithmic}[1]
  \small
  \Procedure{\texttt{LW-DETECT}}{$\mathbf{b}_i(t), \mathbf{b}_i(t{-}T_b), K_i,
    \psi_{\mathrm{th}}, \mathbf{w}$}
\State Verify $\mathrm{MAC}_i(t)$ via Eq.~\eqref{eq:hmac};
         \textbf{if} fails \textbf{then return}
         $\mathrm{ALERT}$ (authentication failure)
  \State \textbf{if} $|t_{\mathrm{receive}} - t|
         > \Delta_{\mathrm{HMAC}}$
         \textbf{then return} $\mathrm{ALERT}$
         (timestamp stale)
  \State \textbf{if} $\nu_i \leq \nu_i^{\mathrm{last}}$
         \textbf{then return} $\mathrm{ALERT}$
         (nonce replay rejected)
  \State $\nu_i^{\mathrm{last}} \gets \nu_i$ (update local
         cache; push to cluster cache if handoff detected)
  \State Compute signatures TP-S1 to MP-S4
         using Eqs.~\eqref{eq:tp_s1}--\eqref{eq:mp_s4}
  \State $\psi_i(t) \gets \sum_k w_k \cdot \mathbf{1}[\text{sig.}~k~\text{violated}]
         \;/\; \sum_k w_k$ \quad (Eq.~\eqref{eq:psi_score})
  \If{$\psi_i(t) > \psi_{\mathrm{th}}$}
       \State Buffer $\mathcal{E}_j(t) \gets (\mathrm{ID}_i,\,
           \psi_i(t),\, t,\, h(\mathbf{b}_i(t)),\,
           \sigma_j^{\mathrm{sub}})$ in local evidence queue;
           flush to \texttt{SC-Trust} per Tier~3 commit policy
           ($T_{\mathrm{batch}}$, $N_{\mathrm{commit}}$);
           Tier~1 revocation commits remain synchronous
           \quad (Eq.~\eqref{eq:rsu_evidence})
    \State \textbf{return} $\mathrm{ALERT}(\mathrm{ID}_i, \psi_i(t))$
  \EndIf
  \State \textbf{return} $\mathrm{PASS}$
  \EndProcedure
\end{algorithmic}
\label{alg:lightweight}
\end{algorithm}

As shown in Algorithm~\ref{alg:lightweight}, LW-DETECT first verifies
the HMAC tag, then evaluates all nine rule-based signatures to produce
the composite anomaly score $\psi_i(t)$. If $\psi_i(t) >
\psi_{\mathrm{th}}$, the RSU calls \texttt{SC-Trust.Update} on-chain
and raises an alert. Normal beacons are passed with no blockchain
interaction, keeping latency within the 100\,ms beacon interval $T_b$.
The deliberate ordering of HMAC verification before signature scoring
is a fail-fast design decision: any beacon that fails integrity
verification is immediately rejected without expending computational
resources on the nine signature checks, which is critical for meeting
the 100\,ms beacon interval constraint $T_b$ at high beacon arrival
rates. The absence of blockchain interaction on the normal path is
equally deliberate writing to the ledger only on anomaly detection
ensures that the lightweight mode scales to dense urban deployments
where hundreds of beacons per second may arrive at a single RSU without
saturating the consensus layer. The composite threshold $\psi_{\mathrm{th}}$ therefore acts as a
gating condition that separates routine beacon processing, which is
entirely local to the RSU, from the exception path that triggers
distributed trust management.
TP-DETECT, SYB-DETECT, and MITM-DETECT return violation flag
vectors that feed into the composite score $\psi_i(t)$ via
Equation~\ref{eq:psi_score}; the sole ALERT decision point is
LW-DETECT at $\psi_i(t) > \psi_{\mathrm{th}}$, except the HMAC
failure path in MITM-DETECT which triggers immediate rejection
before composite scoring, since a tag failure is a conclusive
indicator of tampering rather than a probabilistic signal
requiring further evidence.

Algorithm~\ref{alg:tp_detect} presents the TP-DETECT procedure, which
evaluates the five trajectory poisoning signatures TP-S1 through TP-S5
for each incoming beacon at the RSU.

\begin{algorithm}[H]
\caption{Trajectory Poisoning Detection (\texttt{TP-DETECT})}
\begin{algorithmic}[1]
  \small
  \Procedure{\texttt{TP-DETECT}}{$\mathbf{b}_i(t), \mathbf{b}_i(t{-}T_b),
    \delta_{\mathrm{th}}, k$}
  \State Check TP-S1: $\mathrm{viol}_1 \gets
    \mathbf{1}\!\left[\|p_i(t)-p_i(t{-}T_b)\| > s_{\max}\cdot T_b\right]$
  \State Check TP-S2: $\mathrm{viol}_2 \gets
    \mathbf{1}\!\left[|\theta_i(t)-\theta_i(t{-}T_b)| > \omega_{\max}\cdot T_b\right]$
  \State Check TP-S3: $\mathrm{viol}_3 \gets
    \mathbf{1}\!\left[|a_i(t)| > a_{\max}
    \;\lor\; s_i(t) > s_{\max}\right]$
  \State Compute $\hat{p}_i(t)$ via dead-reckoning Eq.~\eqref{eq:dead_reckoning}
  \State Compute residual $r_i(t) \gets \|p_i(t) - \hat{p}_i(t)\|$
         \quad (Eq.~\eqref{eq:tp_s4})
  \State Update drift window: $D_i(k) \gets \frac{1}{k}\sum_{j=1}^{k} r_i(t{-}(k{-}j)T_b)$
         \quad (Eq.~\eqref{eq:tp_s5})
  \State $\mathrm{viol}_4 \gets \mathbf{1}\!\left[D_i(k) > \delta_{\mathrm{th}}\right]$
 \State \textbf{return}
    $[\mathrm{viol}_1,\mathrm{viol}_2,\mathrm{viol}_3,\mathrm{viol}_4]$
    \Comment{Flag vector fed into $\psi_i(t)$ in LW-DETECT}
  \EndProcedure
\end{algorithmic}
\label{alg:tp_detect}
\end{algorithm}

As shown in Algorithm~\ref{alg:tp_detect}, the procedure checks
kinematic feasibility (TP-S1), heading rate (TP-S2), acceleration
bound (TP-S3), and dead-reckoning residual (TP-S4) sequentially for
each incoming beacon. Any single violation raises an ALERT-TP with
the vehicle identity and evidence flags; otherwise a PASS is returned
with no blockchain interaction. The sequential rather than parallel
evaluation of the four checks is intentional: TP-S1 through TP-S3
are instantaneous single-beacon tests that can be resolved in
microseconds, while TP-S4 and TP-S5 require maintaining a sliding
window of past residuals, so early termination on TP-S1 through TP-S3
avoids unnecessary window maintenance for obviously malicious beacons.
The cumulative drift score $D_i(k)$ in TP-S5 is the primary detector
for the gradual poisoning strategy described in
Section~\ref{sec:mobility_amplification}, where an adversary stays
within per-beacon kinematic bounds but accumulates displacement over
the full RSU contact window $N_b \cdot \epsilon_{\max}$. Together,
the four violation flags provide complementary coverage: TP-S1 through
TP-S3 catch aggressive single-beacon falsifications, while TP-S5
catches the slow drift attacks that individually valid beacons enable.

Algorithm~\ref{alg:syb_detect} presents the SYB-DETECT procedure,
which detects Sybil-based mobility pattern poisoning by checking
identity density, beacon synchronisation, and ghost transit
impossibility across the RSU coverage zone.

\begin{algorithm}[H]
\caption{Sybil-Based Mobility Pattern Poisoning Detection (\texttt{SYB-DETECT})}
\begin{algorithmic}[1]
  \small
  \Procedure{\texttt{SYB-DETECT}}{$\mathcal{B}_j(t), \mathcal{A}_j,
    K_{\mathrm{sybil}}, \tau_{\mathrm{sync}}, \rho_{\mathrm{sync}}$}
  \State Extract identity set $\mathcal{ID}_j \gets
    \{\mathrm{ID}_i : p_i(t) \in \mathcal{A}_j\}$
  \State Check MP-S1: $\mathrm{viol}_1 \gets
    \mathbf{1}\!\left[|\mathcal{ID}_j| > K_{\mathrm{sybil}}
    + \rho_v \cdot \mathcal{A}_j\right]$
    \quad (Eq.~\eqref{eq:mp_s1})
  \For{each pair $(\mathrm{ID}_a, \mathrm{ID}_b) \in \mathcal{ID}_j^2$}
    \State Compute synchronization score via Eq.~\eqref{eq:mp_s2}
    \If{score $> \rho_{\mathrm{sync}}$}
      \State $\mathrm{viol}_2 \gets 1$;
             flag pair $(\mathrm{ID}_a, \mathrm{ID}_b)$ as co-sourced
    \EndIf
  \EndFor
  \State Check MP-S4 for each $\mathrm{ID}_i$: ghost transit test
         via Eq.~\eqref{eq:mp_s4} across neighbouring RSUs
  \State $\mathrm{viol}_3 \gets \mathbf{1}\!\left[\exists\, \mathrm{ID}_i
    \text{ failing ghost test}\right]$
\State \textbf{return}
    $[\mathrm{viol}_1,\mathrm{viol}_2,\mathrm{viol}_3]$
    \Comment{Flag vector fed into $\psi_i(t)$ in LW-DETECT}
  \EndProcedure
\end{algorithmic}
\label{alg:syb_detect}
\end{algorithm}

As shown in Algorithm~\ref{alg:syb_detect}, the procedure first
extracts the active identity set within the RSU coverage zone, then
evaluates identity density anomaly (MP-S1), pairwise beacon
synchronisation score (MP-S2), and ghost vehicle transit impossibility
(MP-S4). A single positive flag triggers an ALERT-SYB submitted
directly to the blockchain. The pairwise synchronisation loop in
lines~4--9 has quadratic complexity in the number of identities, which
is bounded by $K_{\mathrm{sybil}} + \rho_v \cdot \mathcal{A}_j$ from
Equation~\ref{eq:sybil_capacity}; in practice this remains tractable
because the MP-S1 check in line~3 filters the identity set before the
loop executes, meaning the loop only runs when the identity count is
already suspicious. The ghost transit test in line~10 requires
inter-RSU coordination to obtain appearance timestamps from
neighbouring RSUs, which is handled through the shared blockchain
ledger rather than direct RSU-to-RSU messaging, ensuring that the
test remains consistent even when some neighbouring RSUs are offline
or compromised. The three checks together address distinct Sybil
attack strategies: MP-S1 targets volume-based Sybil injection, MP-S2
targets timing-correlated co-sourced identities, and MP-S4 targets
teleporting ghost identities that no single-RSU check can detect.

Algorithm~\ref{alg:mitm_detect} presents the MITM-DETECT procedure,
which verifies HMAC beacon integrity as its first gate, followed by a
regional speed distribution shift check (MP-S3) to detect
MitM-injected mobility data.

\begin{algorithm}[H]
\caption{MitM Mobility Pattern Poisoning Detection (\texttt{MITM-DETECT})}
\begin{algorithmic}[1]
  \small
  \Procedure{\texttt{MITM-DETECT}}{$\mathbf{b}_i(t), K_i,
    \kappa_{\mathrm{th}}, \mathcal{P}_{\mathrm{hist}}$}
  \State Verify HMAC: $\mathrm{valid} \gets
    \mathrm{HMAC}_{K_i}(\mathbf{b}_i(t) \| t \| \mathrm{ID}_i)$
    \quad (Eq.~\eqref{eq:hmac})
  \If{$\neg\,\mathrm{valid}$}
    \State \textbf{return} $\mathrm{ALERT\_MITM}(\mathrm{ID}_i,\;
      \texttt{MAC\_FAIL})$
  \EndIf
  \State Compute current regional speed distribution $\mathcal{P}_t$
         from $\{s_i(t)\}_{v_i \in \mathcal{V}_j}$
  \State Check MP-S3: $\mathrm{viol}_1 \gets
    \mathbf{1}\!\left[D_{\mathrm{KL}}(\mathcal{P}_t \| \mathcal{P}_{\mathrm{hist}})
    > \kappa_{\mathrm{th}}\right]$
    \quad (Eq.~\eqref{eq:mp_s3})
  \State Run TP-S4 residual check on modified fields
         via Eq.~\eqref{eq:tp_s4}
  \State $\mathrm{viol}_2 \gets \mathbf{1}\!\left[r_i(t) > \delta_{\mathrm{th}}\right]$
\State \textbf{return}
    $[\mathrm{viol}_1,\mathrm{viol}_2]$
    \Comment{MP-S3 and TP-S4 flags fed into $\psi_i(t)$;
    HMAC failure above is the only immediate ALERT here}
  \EndProcedure
\end{algorithmic}
\label{alg:mitm_detect}
\end{algorithm}

As shown in Algorithm~\ref{alg:mitm_detect}, the HMAC verification
step acts as a hard gate --- any tag failure immediately returns
ALERT-MITM without further processing. Only HMAC-valid beacons
proceed to the KL-divergence distribution check and dead-reckoning
residual test, ensuring in-transit modifications are detected before
any rule-based scoring is applied. The hard-gate design is justified
by the MitM threat model: an in-transit attacker who modifies beacon
content necessarily invalidates the HMAC tag computed by the
originating vehicle under its session key $K_i$, so a tag failure is
a conclusive indicator of tampering rather than a probabilistic signal
requiring further evidence. The KL-divergence check MP-S3 in line~7
serves a complementary role for the case where the MitM attacker
replays unmodified beacons from a different region rather than
modifying individual fields, which preserves per-beacon HMAC validity
but shifts the regional speed distribution away from the historical
baseline $\mathcal{P}_{\mathrm{hist}}$. The dead-reckoning residual
check in line~9 then catches field-level modifications that are
individually plausible but inconsistent with the vehicle's previously
reported kinematic state, providing a three-layer detection structure
that covers the full space of MitM attack strategies.


\subsubsection{Full Mode: GAT-Based Spatial Anomaly Detection}
The vehicle--RSU communication structure is modelled as a dynamic graph
$\mathcal{G}(t) = (\mathcal{V}(t), \mathcal{E}(t))$ using the
mobility-aware edge definition given in Equation~\ref{eq:edge_def1},
the rolling z-score normalisation given in
Equation~\ref{eq:feature_norm1}, and the multi-head attention
computation given in Equations~\ref{eq:multihead_attn1} and
\ref{eq:multihead_embed1}, as defined in
Section~\ref{sec:gat_design}. The spatial anomaly score measuring
node $i$'s deviation from expected normal conditions is given in
Equation~\ref{eq:gat_score}.
\begin{equation}
  \mathcal{S}_i(t) = \left\|
    \frac{\mathbf{x}'_i(t) - \bar{\mathbf{x}}'_i}
         {\boldsymbol{\sigma}'_i}
  \right\|_2
  \label{eq:gat_score}
\end{equation}
As shown in Equation~\ref{eq:gat_score},
$\mathcal{S}_i(t)$ measures how far the current spatial embedding
deviates from the normal embedding distribution, with high values
flagging spatially anomalous vehicles, where $\bar{\mathbf{x}}'_i$
and $\boldsymbol{\sigma}'_i$ are the per-node mean and standard
deviation computed from clean attack-free embeddings produced after
GAT weights are locked, and are not updated during evaluation.

\subsubsection{Full Mode: Temporal Autoencoder Sequential Anomaly Detection}

The sliding window input sequence of $L$ consecutive beacons for the
temporal autoencoder is given in Equation~\ref{eq:ta_input}.
\newline
\newline
This input formulation presents the beacon sequence as a time-series
matrix to the LSTM encoder, enabling the temporal autoencoder to learn
normal trajectory dynamics across a configurable observation window
whose length $L$ is set to satisfy the detection feasibility condition
in Equation~\ref{eq:detection_condition}.
\begin{equation}
  \mathbf{X}_i(t) = \bigl[\mathbf{b}_i(t-L\cdot T_b),\,\ldots,\,
                            \mathbf{b}_i(t)\bigr]
                    \in \mathbb{R}^{L\times d}
  \label{eq:ta_input}
\end{equation}
As shown in Equation~\ref{eq:ta_input}, the input captures the recent
trajectory of vehicle $v_i$ as a time-series matrix of $L$ consecutive
beacons. The LSTM encoder processes $\mathbf{X}_i(t)$ step by step
using the gating equations defined in
Equations~\ref{eq:lstm_forget}--\ref{eq:lstm_hidden}, producing the
latent representation $\mathbf{h}_i(t) = \mathbf{h}_L$.
\newline
\newline
The decoder
reconstruction of the full input sequence from this latent
representation is given in Equation~\ref{eq:ta_decoder}. This reconstruction step converts the compressed latent
representation back into a full beacon sequence, enabling anomaly
detection by reconstruction error: sequences that deviate from
learned normal mobility patterns produce higher errors than
legitimate trajectories.
\begin{equation}
  \hat{\mathbf{X}}_i(t) =
    \mathrm{LSTM}_{\mathrm{dec}}\!\bigl(\mathbf{h}_i(t)\bigr)
  \label{eq:ta_decoder}
\end{equation}
As shown in Equation~\ref{eq:ta_decoder}, the decoder reconstructs
the full beacon sequence from the compressed latent representation,
enabling comparison between reported and expected mobility behaviour.
The autoencoder is trained using the mobility-constrained
reconstruction loss defined in Equation~\ref{eq:ae_loss1}.
\newline
\newline
At
inference, the per-window reconstruction error is given in
Equation~\ref{eq:ta_error}. This error metric quantifies how well the autoencoder can reproduce
a given beacon sequence from its latent representation, serving as
the primary anomaly signal for the temporal component: poisoned
sequences that lie off the learned normal manifold produce
systematically higher reconstruction errors than genuine trajectories.
\begin{equation}
  \varepsilon_i(t) = \frac{1}{L}\sum_{l=1}^{L}
    \bigl\|\mathbf{b}_i(t_l) -
    \hat{\mathbf{b}}_i(t_l)\bigr\|^2
  \label{eq:ta_error}
\end{equation}
As shown in Equation~\ref{eq:ta_error}, $\varepsilon_i(t)$ measures
the average per-step reconstruction error over the beacon window,
with higher values indicating sequences that deviate from the learned
normal mobility manifold.
\newline
\newline
To prevent an attacker from drifting the detection baseline by
slowly poisoning the live beacon stream, the detection threshold
is calibrated entirely from offline reconstruction errors computed
on SUMO-generated clean mobility traces, using robust statistics
resistant to outlier contamination. The robust detection threshold
using median and Median Absolute Deviation (MAD) of offline
clean-data reconstruction errors is given in
Equation~\ref{eq:ta_threshold}:
\begin{equation}
  \theta_{\mathrm{ae}} =
    \mathrm{median}(\mathcal{E}_{\mathrm{clean}})
    + \kappa \cdot
    \mathrm{MAD}(\mathcal{E}_{\mathrm{clean}})
  \label{eq:ta_threshold}
\end{equation}

As shown in Equation~\ref{eq:ta_threshold}, the threshold is fixed
from clean data before deployment and is not updated from live
inputs where $\mathcal{E}_{\mathrm{clean}} =
\{\varepsilon_i^{\mathrm{clean}}\}$ is the offline reconstruction
error set and $\mathrm{MAD}(\mathcal{E}) =
\mathrm{median}(|\varepsilon -
\mathrm{median}(\mathcal{E})|)$. The autoencoder's normal mobility manifold is similarly
trained offline on SUMO traces and frozen before deployment.
Online model updates are suspended whenever the fraction of flagged
vehicles in any RSU zone exceeds a configurable freeze threshold
$\alpha_{\mathrm{freeze}}$, preventing adversarial corruption of
the learned manifold during active attacks.

\subsubsection{Combined Detection Score and Fusion}

In the full mode, the combined detection score fusing rule-based,
spatial, and temporal anomaly signals is given in
Equation~\ref{eq:fusion}.
\newline
\newline
This fusion equation combines the complementary strengths of the
three detection components, rule-based kinematic checks, GAT
spatial relational analysis, and LSTM temporal sequential analysis
 into a single decision variable, ensuring that attacks evading
any one component are still detected by the remaining two.
By normalising $\psi_i(t)$ by the rule-based alert threshold
$\psi_{\mathrm{th}}$, $\mathcal{S}_i(t)$ by the clean-data spatial
calibration threshold $\theta_{\mathcal{S}}$, and $\varepsilon_i(t)$
by the clean-data reconstruction threshold $\theta_{\mathrm{ae}}$,
each term reaches exactly $1.0$ when the component alone attains its
independent single-mechanism alerting level.
Each normalised term is then clipped at $1.0$ before summation.
Since $\lambda_1 + \lambda_2 + \lambda_3 = 1$, this clipping makes
$\Phi_i(t)$ a genuine weighted average in $[0,1]$, giving
$\Phi_{\mathrm{th}} = 0.5$ the principled interpretation of
at least half the weighted evidence signals attack.
Without this clip, a single component with a large raw value dominates
the sum, rendering $\Phi_{\mathrm{th}}$ an arbitrary threshold on an
unbounded quantity. Without normalising $\psi_i(t)$, the raw rule
score is bounded well below $\Phi_{\mathrm{th}}$ regardless of the
learned weight $\lambda_1$, preventing the rule term from contributing
to the fusion decision even when lightweight detection succeeds.
\begin{equation}
  \Phi_i(t) =
    \lambda_1^{(\hat{k}_i)}\,\min\!\left(\frac{\bar{\psi}_i(t)}{\psi_{\mathrm{th}}},\,1\right)
  + \lambda_2^{(\hat{k}_i)}\,\min\!\left(\frac{\mathcal{S}_i(t)}{\theta_{\mathcal{S}}},\,1\right)
  + \lambda_3^{(\hat{k}_i)}\,\min\!\left(\frac{\varepsilon_i(t)}{\theta_{\mathrm{ae}}},\,1\right)
  \label{eq:fusion}
\end{equation}
As shown in Equation~\ref{eq:fusion}, $\Phi_i(t) \in [0,1]$ is a
genuine weighted average of three clipped normalised components,
where $\theta_{\mathcal{S}}$ is locked from a high percentile of
$\mathcal{S}_i(t)$ computed over clean Phase~1b data.
Here, $\bar{\psi}_i(t) = \tfrac{1}{L}\sum_{l=1}^{L}\psi_i(t_l)$ is
the mean composite rule score across the $L$-beacon detection window,
providing temporal consistency with the window-level LSTM-AE score
$\varepsilon_i(t)$.
When the per-window mean rule score satisfies
$\bar{\psi}_i(t) \geq \psi_{\mathrm{th}}$, the normalised rule term
reaches its cap of $1.0$, contributing $\lambda_1$ to $\Phi_i(t)$;
if $\lambda_1 > \Phi_{\mathrm{th}}$, the rule signal alone is
sufficient to trigger a full-mode alert, ensuring the full mode
subsumes the lightweight mode for sustained attacks.
An additional feasibility constraint $\lambda_3^{(k)} < \lambda_1^{(k)}
+ \lambda_2^{(k)}$ is enforced during the Phase~2 SLSQP optimisation
for every attack class $k$, preventing the LSTM-AE weight from
dominating the fusion outright --- a pathology in which a silent,
untriggered autoencoder component could still cap $\Phi_i(t)$ below
$\Phi_{\mathrm{th}}$ for classes where the reconstruction error is not
discriminative. This constraint was found to be necessary to resolve
an impossible detection ceiling observed on attack classes a1 and a5.
\newline
\newline
The binary detection decision is given in
Equation~\ref{eq:flag} where $\lambda_1^{(k)} + \lambda_2^{(k)} + \lambda_3^{(k)} = 1$ and
each $\lambda_j^{(k)} \geq 0.05$ for all $j, k$; $K$ independent
weight sets are learned in Phase~2, one per attack class, by
partitioning the Phase~2 labelled validation segment by ground-truth
attack type and running $K$ constrained SLSQP optimisations in
parallel. At inference, $\hat{k}_i$ from Equation~\ref{eq:gat_atk_type}
selects the weight set that best matches the predicted attack class,
so each detector component is weighted according to its expected
discriminative power for that specific attack type rather than an
average across all attacks.
\newline
\newline
This binary decision converts the continuous fused anomaly score
into an actionable flag that triggers the on-chain trust update and
revocation pipeline, translating the detection output into a concrete
network management action without requiring manual operator
intervention.
\begin{equation}
  \mathrm{flag}_i(t) = \mathbf{1}\!\left[\Phi_i(t) > \Phi_{\mathrm{th}}\right]
    \;\lor\; \mathrm{flag}_i^{\mathrm{GAT}}
  \label{eq:flag}
\end{equation}
As shown in Equation~\ref{eq:flag}, a vehicle is flagged as malicious
when either the fused anomaly score exceeds $\Phi_{\mathrm{th}}$ or any
individual GAT attack-type head independently fires
($\mathrm{flag}_i^{\mathrm{GAT}}$, Eq.~\ref{eq:gat_det_flag}), triggering
the on-chain trust update and revocation pipeline. The OR condition
prevents a case where a single head is confidently correct but its
signal is diluted below $\Phi_{\mathrm{th}}$ by the weighted-average
fusion.

When the GAT attack-type prediction confidence is low, committing
to the wrong weight set actively degrades detection. A per-head
confidence gate is therefore applied: if
$\hat{y}_i^{(\hat{k}_i)} < \theta_{\mathrm{conf}}^{(\hat{k}_i)}$,
the fusion falls back to the unweighted average weight set
$\bar{\lambda}_j = \frac{1}{K}\sum_{k=1}^{K}\lambda_j^{(k)}$
rather than selecting the argmax weight set. Each
$\theta_{\mathrm{conf}}^{(k)}$ is calibrated independently per
attack-type head on the Phase~2 validation set to maximise MCC,
since different heads exhibit different confidence calibration.

Algorithm~\ref{alg:full_detection} presents the FULL-DETECT procedure,
which integrates GAT spatial anomaly detection, LSTM temporal
autoencoder scoring, and rule-based signatures into the complete
full-mode detection pipeline executed at the SDVN controller per
sliding window.

\begin{algorithm}[H]
\caption{Full-Mode Spatio-Temporal Detection (\texttt{FULL-DETECT})}
\begin{algorithmic}[1]
  \small
  \Procedure{\texttt{FULL-DETECT}}{$\mathcal{G}(t), \{\mathbf{X}_i(t)\},
    \Phi_{\mathrm{th}}, \{\boldsymbol{\lambda}^{(k)}\}_{k=1}^{K}$}
  \State Retrieve beacon windows $\{\mathbf{X}_i(t)\}$ from IPFS
         using on-chain hashes
  \For{each $v_i \in \mathcal{V}(t)$}
      \State Compute $\psi_i(t_l)$ via \texttt{LW-DETECT} for each
         $l \in \{1,\ldots,L\}$;
         \;$\bar{\psi}_i(t) \gets \tfrac{1}{L}
         \sum_{l=1}^{L}\psi_i(t_l)$
         \Comment{per-window mean rule score}
    \State Compute GAT embedding $\mathbf{x}'_i(t)$ via
           Eqs.~\eqref{eq:multihead_attn1}--\eqref{eq:multihead_embed1}
    \State Compute $\hat{y}_i^{(k)}$ for $k=1,\ldots,K$ via
           Eq.~\eqref{eq:gat_cls};
           $\hat{k}_i \gets \arg\max_k \hat{y}_i^{(k)}$
           \quad (Eq.~\eqref{eq:gat_atk_type})
    \State Compute spatial score $\mathcal{S}_i(t)$ via
           Eq.~\eqref{eq:gat_score}
    \State Encode $\mathbf{h}_i(t) \gets \mathrm{LSTM}_{\mathrm{enc}}(\mathbf{X}_i(t))$
           \quad (Eqs.~\eqref{eq:lstm_forget}--\eqref{eq:lstm_hidden})
    \State Reconstruct $\hat{\mathbf{X}}_i(t) \gets
           \mathrm{LSTM}_{\mathrm{dec}}(\mathbf{h}_i(t))$
           \quad (Eq.~\eqref{eq:ta_decoder})
    \State Compute $\varepsilon_i(t)$ via Eq.~\eqref{eq:ta_error}
    \State $\Phi_i(t) \gets
               \lambda_1^{(\hat{k}_i)}\min(\bar{\psi}_i/\psi_{\mathrm{th}},1)
             + \lambda_2^{(\hat{k}_i)}\min(\mathcal{S}_i/\theta_{\mathcal{S}},1)
             + \lambda_3^{(\hat{k}_i)}\min(\varepsilon_i/\theta_{\mathrm{ae}},1)$
             \State $\mathrm{flag}_i(t) \gets \mathbf{1}[\Phi_i(t) > \Phi_{\mathrm{th}}] \lor \mathrm{flag}_i^{\mathrm{GAT}}$
           \Comment{OR-fused decision; fixes GAT-flag wiring gap} \quad (Eq.~\eqref{eq:flag})
           \Comment{attack-conditioned weights selected by $\hat{k}_i$;
           clipped weighted average in $[0,1]$;
           $\psi_{\mathrm{th}},\theta_{\mathcal{S}},\theta_{\mathrm{ae}}$
           fixed offline; $\Phi_{\mathrm{th}}=0.5$ is the evidence midpoint}
           \quad (Eq.~\eqref{eq:fusion})

    \State Submit $\mathcal{E}_c(t) \gets (\mathrm{ID}_i, \Phi_i(t), t,
       h(\mathbf{X}_i(t)), \sigma_c^{\mathrm{sub}})$
       to blockchain as untrusted peer \quad (Eq.~\eqref{eq:ctrl_evidence})
    \Comment{SC-Trust and SC-Revoke decisions driven by RSU evidence only}
  \EndFor
  \State \textbf{return} $\{\mathrm{flag}_i(t)\}$
  \EndProcedure
\end{algorithmic}
\label{alg:full_detection}
\end{algorithm}

As shown in Algorithm~\ref{alg:full_detection}, for each vehicle the
procedure computes the rule-based score via a LW-DETECT sub-call, the
GAT spatial embedding, and the LSTM reconstruction error, then fuses
them into $\Phi_i(t)$. The controller submits its result to the
blockchain as one untrusted peer vote; all SC-Trust and SC-Revoke
decisions remain driven by RSU evidence only. The IPFS retrieval step
in line~2 is architecturally significant: by storing beacon windows
off-chain with only their content hashes recorded on the ledger, the
framework bridges the real-time RSU-side lightweight detection and the
offline controller-side deep learning analysis without requiring the
controller to maintain a direct connection to each RSU, and enables
reanalysis of historical windows on demand if detection thresholds are
recalibrated. The controller's submission being treated as one
untrusted peer vote rather than a privileged input reflects the threat
model assumption in Section~\ref{sec:threat_model} that the controller
itself may be compromised: a malicious controller that submits a
falsified $\Phi_i(t)$ score is outvoted by the $2f+1$ honest RSU
submissions required by Equation~\ref{eq:bft_revoke_rsu}. The
per-vehicle loop scales linearly with the number of vehicles in
$\mathcal{V}(t)$, and in practice the GAT forward pass over the full
vehicle--RSU graph is the dominant computational cost, motivating the
sliding window update period $W = L \cdot T_b$ rather than per-beacon
execution at the controller.


\subsection{Full Mode Mitigation: Post-Quantum TRS and FHE Aggregate
Protection}
\label{sec:pqc_mitigation}

Figure~\ref{fig:seq_fullmode} presents the corresponding UML sequence
diagram for one full-mode detection epoch. .

\begin{figure}[H]
    \centering
    \includegraphics[width=\linewidth,
                     height=0.80\textheight,
                     keepaspectratio]{images/Sequence Diagram Full-Mode Detection and PQC Mitigation}
    \caption{UML sequence diagram for the full operating mode across
    one detection epoch. Five lifelines }
    \label{fig:seq_fullmode}
\end{figure}

While Figure~\ref{fig:seq_lightweight} captures the per-beacon timing of
the lightweight path at a single RSU, the full operating mode spans
five distinct entities and combines real-time RSU-side processing with
offline controller-side deep learning inference. As shown in
Figure~\ref{fig:seq_fullmode}, five lifelines---Vehicle, RSU,
Controller~$c_{\text{assigned}}(r_j)$, IPFS off-chain store, Cloud/ITS
Server, and the Blockchain (SC-Trust\,/\,SC-Revoke)---participate across
a single detection epoch. As a precondition, the threshold signing keys
and FHE key shares are established once by distributed key generation
among the RSU ring and the Cloud (Equations~\ref{eq:dkg_trs} and
\ref{eq:dkg_fhe}), with no controller involvement; each entity then
registers via SC-Register and is confirmed on-chain only after BFT
endorsement by $2f{+}1$ RSU peers.

The epoch begins on the data plane with the vehicle broadcasting its
beacon, which the RSU first authenticates through the HMAC tag, freshness
window, and monotonic anti-replay nonce (Equation~\ref{eq:hmac}); stale
or replayed beacons are dropped at this fail-fast gate. The RSU evaluates
the nine rule-based signatures and the composite score
$\psi_i(t)$ (\texttt{LW-DETECT}); a direct evidence tuple
$\mathcal{E}_j(t)$ is submitted to the blockchain only on the alert
path (Equation~\ref{eq:rsu_evidence}), preserving the lightweight
skip-on-pass behaviour, while the beacon window itself is stored
off-chain in IPFS with only its hash retained on-chain.

The privacy-preserving aggregation path then executes in the order
mandated by the threat model. Each RSU first FHE-encrypts its own
local aggregate, $c_j = \mathrm{Enc}(\mathbf{A}_j)$
(Equations~\ref{eq:rsu_aggregate}--\ref{eq:fhe_rsu_enc}), before
any ring sharing (step~5a). The ring members then exchange ciphertexts
and homomorphically combine them into
$\mathrm{Enc}(\mathbf{A}_{\text{ring}}) = \bigoplus_j c_j$
(Equation~\ref{eq:fhe_ring_combine}), and the post-quantum threshold ring
signature is produced over the combined ciphertext only
(Equations~\ref{eq:trs_message}--\ref{eq:trs_aggregate}, step~5b)---so no
plaintext mobility data ever crosses the ring. The signed aggregate
$(\mathrm{Enc}(\mathbf{A}_{\text{ring}}), \sigma_{\text{TRS}})$ is sent to
the Cloud, which first checks message freshness
(Equation~\ref{eq:trs_fresh}) and then verifies the threshold ring
signature (Equation~\ref{eq:trs_verify}). This verification is a hard
gate: any aggregate failing it is rejected before homomorphic
computation begins, and only a passing aggregate proceeds to global
homomorphic aggregation (Equation~\ref{eq:fhe_global}).

Recovering the plaintext result is itself a threshold operation. The
Cloud contributes one partial decryption (Equation~\ref{eq:partial_dec_cloud}),
but its participation is optional; each member of the decryption coalition
$\mathcal{D}$ computes its own partial using its key share
(Equation~\ref{eq:partial_dec}), and the global aggregate is reconstructed
by threshold decryption only when $|\mathcal{D}| \ge t_{\text{decrypt}} = f{+}2$
(Equation~\ref{eq:thresh_dec}). Consequently neither the Cloud alone nor
any sub-threshold subset of RSUs can decrypt unilaterally; if the Cloud is
offline, $f{+}2$ RSU shares suffice on their own. The decrypted aggregate
$\bar{X}_{\text{global}}$ is delivered to local ITS services and returned
to the Cloud for analytics.

The AI detection path runs at the assigned controller. It retrieves the
beacon windows $\{\mathbf{X}_i(t)\}$ from IPFS, computes the GAT spatial
anomaly score $\mathcal{S}_i(t)$ (Equation~\ref{eq:gat_score}) and the
LSTM-autoencoder temporal reconstruction error $\varepsilon_i(t)$
(Equation~\ref{eq:ta_error}), and fuses them into $\Phi_i(t)$; an alert is
raised when $\Phi_i(t) > \Phi_{\text{th}}$
(Equations~\ref{eq:fusion}--\ref{eq:flag}). Crucially, the controller
submits its result to the blockchain as a single untrusted peer
vote (Equation~\ref{eq:ctrl_evidence}): revocation is driven by RSU
evidence, requiring $2f{+}1$ distinct RSU submissions over the temporal
window (Equation~\ref{eq:bft_revoke_rsu}), not by the controller alone.

Finally, SC-Trust updates the controller trust score $\tau_{c_k}(t)$ using
the directional conflict indicator $\text{conflict}_j(t)$
(Equations~\ref{eq:ctrl_trust} and \ref{eq:conflict_pred}); if
$\tau_{c_k}$ remains below $\tau_{\min}$ for $T_{\text{rev}}$ epochs,
SC-Revoke issues \texttt{ControllerRevoked} and updates the trusted-controller
set $\mathcal{C}_{\text{trusted}}(t)$ on-chain (Equation~\ref{eq:trusted_ctrl_set}),
prompting affected RSUs to reassign $c_{\text{assigned}}(r_j)$ to any trusted
controller. In the optional CP-DETECT branch, when the controller-compromise
flag is set (Equation~\ref{eq:ctrl_flag}) the controller is excluded from
consensus and the RSUs fall back to rule-based control (\texttt{LW-DETECT});
the resulting CRL update is broadcast directly to all RSUs, bypassing the
controller, and the on-chain trust scores are updated
(Equation~\ref{eq:trust_update}).

\subsubsection{RSU Aggregate Construction and Pre-Coordination FHE Encryption}

Each RSU $r_j$ first computes a local plaintext aggregate from the beacon
messages received within its own coverage zone $\mathcal{V}_j(t)$. The local
plaintext aggregate computed by RSU $r_j$ from its own coverage zone is shown
in Equation~\eqref{eq:rsu_aggregate}:
\begin{equation}
  \mathbf{A}_j(t) = \frac{1}{|\mathcal{V}_j(t)|}
    \sum_{v_i \in \mathcal{V}_j(t)} \mathbf{b}_i(t)
  \label{eq:rsu_aggregate}
\end{equation}
As shown in Equation~\eqref{eq:rsu_aggregate}, the aggregate represents the
field-wise mean of beacon values across all vehicles in the RSU's own coverage
zone. Since RSU $r_j$ legitimately serves $\mathcal{V}_j(t)$, computing this
aggregate in plaintext introduces no privacy concern at this stage.

The privacy threat arises during TRS ring coordination, where $t$-of-$n$
neighboring RSUs must collaborate to sign a combined aggregate, exposing each
RSU's aggregate to all ring participants. To prevent cross-RSU aggregate
leakage, each RSU encrypts its own local aggregate using FHE before sharing
it with the ring. The FHE encryption of the local aggregate by RSU $r_j$ prior
to ring sharing is shown in Equation~\eqref{eq:fhe_rsu_enc}:
\begin{equation}
  c_j = \mathrm{Enc}(pk_{\mathrm{FHE}},\; \mathbf{A}_j(t))
  \label{eq:fhe_rsu_enc}
\end{equation}
As shown in Equation~\eqref{eq:fhe_rsu_enc}, each RSU independently encrypts
its local aggregate before any inter-RSU sharing occurs, ensuring that no ring
participant receives another participant's plaintext aggregate at any point
during coordination.

Upon receiving encrypted aggregates from all $t$ signing ring members, each
RSU performs homomorphic addition across the received ciphertexts to produce
a combined encrypted aggregate without decryption. The combined encrypted
aggregate produced by homomorphic addition across the signing ring is shown
in Equation~\eqref{eq:fhe_ring_combine}:
\begin{equation}
  \mathrm{Enc}(\mathbf{A}_{\mathrm{ring}}) =
    \bigoplus_{j \in \mathcal{S}} c_j
  \label{eq:fhe_ring_combine}
\end{equation}
As shown in Equation~\eqref{eq:fhe_ring_combine}, the combined aggregate is
produced entirely on ciphertext no RSU decrypts any other RSU's contribution
at any point during this computation.

\subsubsection{Post-Quantum Threshold Ring Signature on Combined Encrypted Aggregate}

The TRS is applied to the combined encrypted aggregate rather than any
individual RSU's plaintext data. The message committed to the TRS signing process binds the combined
encrypted aggregate to the signing ring identity, timestamp, and a
per-ring monotonically increasing nonce $\nu_{\mathcal{S}}$ to
prevent replay of valid signatures within any freshness window.
The TRS signed message incorporating timestamp and nonce for full
replay protection is shown in Equation~\eqref{eq:trs_message}:
\begin{equation}
  m = \bigl(\mathrm{Enc}(\mathbf{A}_{\mathrm{ring}}),\;
      t,\; \nu_{\mathcal{S}},\;
      \mathrm{ID}_{\mathcal{S}},\;
      h(\mathcal{S})\bigr)
  \label{eq:trs_message}
\end{equation}
As shown in Equation~\eqref{eq:trs_message}, the signed message
contains the combined ciphertext, timestamp, nonce, signing ring
identity, and a hash of the signing member set --- no plaintext
aggregate content is exposed during signing. The cloud verifier
checks that $\nu_{\mathcal{S}}$ strictly exceeds the last accepted
nonce for ring $\mathcal{S}$, in addition to the timestamp
freshness check in Eq.~\eqref{eq:trs_fresh}.
Each RSU $r_j$ in the signing subset $\mathcal{S}$ generates a partial
signature over $m$ using its lattice-based Dilithium secret key. The partial
signature generated by each ring member is shown in
Equation~\eqref{eq:trs_partial}:
\begin{equation}
  \sigma_j = \mathrm{Sign}(sk_j,\; m),
  \quad r_j \in \mathcal{S}
  \label{eq:trs_partial}
\end{equation}
As shown in Equation~\eqref{eq:trs_partial}, each partial signature is computed
independently without requiring access to any other RSU's plaintext data. The
partial signatures are then aggregated into a single compact threshold ring
signature. The aggregated threshold ring signature combining all partial
signatures is shown in Equation~\eqref{eq:trs_aggregate}:
\begin{equation}
  \sigma_{\mathrm{TRS}} =
    \mathrm{Aggregate}\!\left(\{\sigma_j\}_{j\in\mathcal{S}}\right)
  \label{eq:trs_aggregate}
\end{equation}
As shown in Equation~\eqref{eq:trs_aggregate}, the aggregate signature is
compact and conceals which specific $t$ members signed. Verification at the
cloud or SDN controller confirms that at least $t$ legitimate RSUs endorsed
the combined encrypted aggregate. The TRS verification condition at the cloud
is shown in Equation~\eqref{eq:trs_verify}:
\begin{equation}
  \mathrm{Verify}\!\left(\{pk_j\}_{j=1}^{n},\; m,\;
  \sigma_{\mathrm{TRS}}\right) \in \{0,1\}
  \label{eq:trs_verify}
\end{equation}
As shown in Equation~\eqref{eq:trs_verify}, a result of 0 indicates
that fewer than $t_{\mathrm{sign}}$ honest RSUs endorsed the combined
aggregate, causing the cloud to reject the input as potentially
poisoned.

To prevent replay of valid but stale aggregate submissions, the
verifier enforces a freshness window $\Delta_{\mathrm{TRS}}$ on the
timestamp $t_m$ embedded in $m$. The TRS freshness acceptance
condition rejecting stale submissions regardless of signature
validity is shown in Equation~\eqref{eq:trs_fresh}:
\begin{equation}
  \bigl|t_{\mathrm{verify}} - t_m\bigr|
  \leq \Delta_{\mathrm{TRS}}
  \label{eq:trs_fresh}
\end{equation}
As shown in Equation~\eqref{eq:trs_fresh}, any
$(\mathrm{Enc}(\mathbf{A}_{\mathrm{ring}}),\,\sigma_{\mathrm{TRS}})$
pair whose timestamp falls outside the acceptance window is rejected
at the cloud, preventing an attacker from replaying a previously
valid aggregate to re-inject stale poisoned data.
\subsubsection{Cloud-Side Blind Homomorphic Computation and Result Decryption}

Following TRS verification, the cloud receives combined encrypted aggregates
from multiple RSU signing rings and performs further homomorphic computations
across them without decryption. The global encrypted result computed
homomorphically across $M$ ring aggregates at the cloud is shown in
Equation~\eqref{eq:fhe_global}:
\begin{equation}
  \mathrm{Enc}(\bar{X}_{\mathrm{global}}) =
    \frac{1}{M}\bigoplus_{k=1}^{M}
    \mathrm{Enc}(\mathbf{A}_{\mathrm{ring},k})
  \label{eq:fhe_global}
\end{equation}
As shown in Equation~\eqref{eq:fhe_global}, the cloud computes the global
aggregate entirely on ciphertext -- no individual RSU aggregate is ever
decrypted at the cloud, preserving each RSU's aggregate privacy throughout
cloud-side computation. Since the Cloud holds only one threshold share and cannot decrypt
independently, the Cloud returns
$\mathrm{Enc}(\bar{X}_{\mathrm{global}})$ to the originating RSU
cluster for cooperative threshold decryption. The Cloud computes its
own partial decryption using its key share as one of the required $t$
cooperating participants. The partial decryption computed by the Cloud
using its share is shown in Equation~\eqref{eq:partial_dec_cloud}:
\begin{equation}
  pd_{\mathrm{cloud}} = \mathrm{PDec}\!\left(
    sk_{\mathrm{cloud}}^{\mathrm{share}},\;
    \mathrm{Enc}(\bar{X}_{\mathrm{global}})\right)
  \label{eq:partial_dec_cloud}
\end{equation}
As shown in Equation~\eqref{eq:partial_dec_cloud}, the Cloud's
partial decryption contribution is computed using only its own share
and transmitted to the RSU cluster along with
$\mathrm{Enc}(\bar{X}_{\mathrm{global}})$. No entity decrypts
the result independently at any stage.


\subsubsection{Threshold FHE Decryption at RSU Cluster}
\label{sec:thresh_fhe}

To enable RSUs to locally recover the global result without any single
RSU holding the complete FHE secret key which would allow that RSU
to unilaterally decrypt other RSUs' encrypted aggregates received during
ring coordination the FHE secret key is split into $n$ shares across
all RSU ring members under a $(t,n)$ threshold secret sharing scheme.
Since the Cloud belongs to the control plane threat boundary
(Section~\ref{sec:threat_model}), it must not hold unilateral
decryption capability. The FHE secret key is therefore split into
$n+1$ shares across all $n$ RSU ring members and the Cloud, under
a $(t, n+1)$ threshold secret sharing scheme. The threshold FHE key generation producing a shared public key,
$n$ RSU key shares, one Cloud key share, and an evaluation key
is shown in Equation~\eqref{eq:thresh_keygen}:
\begin{equation}
  \left(
    pk_{\mathrm{FHE}},\;
    \{sk_j^{\mathrm{share}}\}_{j=1}^{n},\;
    sk_{\mathrm{cloud}}^{\mathrm{share}},\;
    evk
  \right)
  \leftarrow \mathrm{ThGen}(1^{\lambda},\; n+1,\;
  t_{\mathrm{decrypt}})
  \label{eq:thresh_keygen}
\end{equation}
As shown in Equation~\eqref{eq:thresh_keygen}, neither any individual RSU nor the Cloud holds the full secret key $sk_{\mathrm{FHE}}$. Each entity holds exactly one share, which is computationally insufficient for decryption alone. Homomorphic aggregation operations at the Cloud use only the evaluation key $evk$ and require no secret key, so removing the Cloud's full key does not affect its ability to perform homomorphic computations across ring aggregates.

Upon receiving $\mathrm{Enc}(\bar{X}_{\mathrm{global}})$ and
$pd_{\mathrm{cloud}}$ from the Cloud (when the Cloud participates),
each RSU $r_j$ in the decryption coalition $\mathcal{D}$ computes
a partial decryption using its own key share. When the Cloud
participates, it counts as one of the required
$t_{\mathrm{decrypt}}$ shares, so only $t_{\mathrm{decrypt}}-1 =
f+1$ additional RSU shares are needed. If the Cloud is
unavailable, $t_{\mathrm{decrypt}} = f+2$ RSU shares alone
suffice. The partial decryption
computed independently by each cooperating RSU is shown in
Equation~\eqref{eq:partial_dec}:
\begin{equation}
  pd_j = \mathrm{PDec}\!\left(sk_j^{\mathrm{share}},\;
    \mathrm{Enc}(\bar{X}_{\mathrm{global}})\right),
  \quad r_j \in \mathcal{D},\quad |\mathcal{D}| \geq t-1
  \label{eq:partial_dec}
\end{equation}
As shown in Equation~\eqref{eq:partial_dec}, each partial decryption
uses only the RSU's own share without any RSU learning another's share
or being able to proceed to full decryption independently. The threshold combination of any $t_{\mathrm{decrypt}}$ partial
decryptions from the combined pool
$\mathcal{R} \cup \{\mathrm{Cloud}\}$ to recover the global
aggregate is shown in Equation~\eqref{eq:thresh_dec}:
\begin{equation}
  \bar{X}_{\mathrm{global}} =
    \mathrm{ThDec}\!\left(\{pd_e\}_{e \in \mathcal{D}}\right),
    \quad \mathcal{D} \subseteq
    \mathcal{R} \cup \{\mathrm{Cloud}\},\;
    |\mathcal{D}| \geq t_{\mathrm{decrypt}}
  \label{eq:thresh_dec}
\end{equation}
As shown in Equation~\eqref{eq:thresh_dec}, any
$t_{\mathrm{decrypt}} = f+2$ participants from the combined pool
may cooperate for decryption. Cloud participation is sufficient
but not mandatory if the Cloud is unavailable, $t_{\mathrm{decrypt}}$
RSU shares alone suffice, eliminating the availability single
point of failure that mandatory Cloud participation would create.
When the Cloud participates, its one share counts toward the
required $t_{\mathrm{decrypt}}$, reducing the RSU share
requirement to $t_{\mathrm{decrypt}} - 1 = f+1$.
Since $f$ RSUs plus the Cloud total only $f+1$ shares, which is
below $t_{\mathrm{decrypt}} = f+2$, no attacker controlling $f$
RSUs and the Cloud can decrypt unilaterally.
TRS signing uses $t_{\mathrm{sign}} = f+1$
(Eq.~\eqref{eq:trs_partial}) and BFT revocation uses $2f+1$
(Eq.~\eqref{eq:bft_revoke_rsu}), providing consistent fault
tolerance across all three mechanisms.

Following threshold decryption, $\bar{X}_{\mathrm{global}}$ is
available locally at the RSU cluster for ITS services such as
congestion monitoring, eco-driving guidance, and dynamic routing,
and is also returned to the Cloud for cloud-side analytics use.
Since $\bar{X}_{\mathrm{global}}$ is a mean of RSU-level means,
no individual RSU's zone-level aggregate is recoverable from it,
preserving zone-level privacy even after plaintext delivery.

Algorithm~\ref{alg:pq_fhe_trs} presents the PQ-FHE-TRS procedure, which
formalises the pre-coordination FHE encryption, homomorphic aggregation,
and post-quantum threshold ring signing steps executed by the RSU cluster
to securely process mobility data without exposing local plaintexts.

\begin{algorithm}[htb]
\caption{Pre-Coordination FHE Encryption and PQ-TRS Signing
(\texttt{PQ-FHE-TRS})}
\begin{algorithmic}[1]
  \small
  \Procedure{\texttt{PQ-FHE-TRS}}{$\mathcal{V}_j(t),\;
    \mathcal{R},\; \{sk_j\},\; sk_j^{\mathrm{share}},\;
    pk_{\mathrm{FHE}},\; evk,\; t$}
  \Comment{$sk_j^{\mathrm{share}}$ distributed via LKH at setup;
           no RSU holds full $sk_{\mathrm{FHE}}$}
  \State Compute local plaintext aggregate $\mathbf{A}_j(t)$
         via Eq.~\eqref{eq:rsu_aggregate}
  \State \textbf{if} $\exists\, f \in \mathcal{F} :
         A_j^{(f)}(t) \notin
         [L_f^{\mathrm{SUMO}},\; U_f^{\mathrm{SUMO}}]$
         \textbf{then} log field anomaly to SC-Trust and abort
         \Comment{Pre-enc.\ range check against SUMO-derived bounds}
  \State Encrypt local aggregate before ring sharing:
         $c_j \gets \mathrm{Enc}(pk_{\mathrm{FHE}},\;
         \mathbf{A}_j(t))$
         \quad (Eq.~\eqref{eq:fhe_rsu_enc})
  \State Broadcast $c_j$ to all $r_k \in \mathcal{S}$,
         $r_k \neq r_j$
  \State Receive $\{c_k\}_{k \in \mathcal{S},\; k \neq j}$
         from ring peers
  \State Compute combined encrypted aggregate:
         $\mathrm{Enc}(\mathbf{A}_{\mathrm{ring}}) \gets
         \bigoplus_{j \in \mathcal{S}} c_j$
         \quad (Eq.~\eqref{eq:fhe_ring_combine})
  \State Form TRS message: $m \gets
         (\mathrm{Enc}(\mathbf{A}_{\mathrm{ring}}),\;
         t,\; \nu_{\mathcal{S}},\; \mathrm{ID}_{\mathcal{S}},\;
         h(\mathcal{S}))$
         \quad (Eq.~\eqref{eq:trs_message})
  \For{each $r_j \in \mathcal{S}$}
    \State $\sigma_j \gets \mathrm{Sign}(sk_j,\; m)$
           \quad (Eq.~\eqref{eq:trs_partial})
  \EndFor
  \State $\sigma_{\mathrm{TRS}} \gets
         \mathrm{Aggregate}(\{\sigma_j\}_{j \in \mathcal{S}})$
         \quad (Eq.~\eqref{eq:trs_aggregate})
  \State Transmit $(\mathrm{Enc}(\mathbf{A}_{\mathrm{ring}}),\;
         \sigma_{\mathrm{TRS}})$ to cloud
  \State \textbf{return} $\sigma_{\mathrm{TRS}},\;
         \mathrm{Enc}(\mathbf{A}_{\mathrm{ring}})$
  \EndProcedure
\end{algorithmic}
\label{alg:pq_fhe_trs}
\end{algorithm}

As shown in Algorithm~\ref{alg:pq_fhe_trs}, the procedure first computes the
local plaintext aggregate and immediately encrypts it under the FHE public key
(line~3) before any inter-RSU sharing occurs. This structural reordering ensures
that when RSUs exchange data to coordinate the threshold ring signature
(lines~4--5), they operate entirely on ciphertexts. The ring members then
homomorphically combine these ciphertexts into a single encrypted aggregate
without decryption (line~6). Finally, the procedure binds this combined
ciphertext to the signing ring identity, coordinates the partial signing across
the $t$-of-$n$ subset (lines~8--11), and aggregates them into the final compact
threshold signature $\sigma_{\mathrm{TRS}}$ (line~12) for transmission to the
cloud. This ensures that the cloud receives only privacy-protected,
threshold-verified data.


\begin{algorithm}[htb]
\caption{Threshold FHE Decryption at RSU Cluster
(\texttt{THRESH-DEC})}
\begin{algorithmic}[1]
  \small
  \Procedure{\texttt{THRESH-DEC}}{$\mathrm{Enc}
    (\bar{X}_{\mathrm{global}}),\;
    \{sk_e^{\mathrm{share}}\}_{e \in \mathcal{D}},\;
    t_{\mathrm{decrypt}}$}
  \Comment{$\mathcal{D} \subseteq \mathcal{R} \cup
    \{\mathrm{Cloud}\}$; Cloud not mandatory}
  \If{$|\mathcal{D}| < t_{\mathrm{decrypt}}$}
    \State \textbf{abort} --- insufficient shares for
           decryption
  \EndIf
  \For{each $e \in \mathcal{D}$}
    \State $pd_e \gets \mathrm{PDec}(sk_e^{\mathrm{share}},\;
           \mathrm{Enc}(\bar{X}_{\mathrm{global}}))$
           \hfill (Eq.~\eqref{eq:partial_dec})
    \State Broadcast $pd_e$ to all $e' \in \mathcal{D}$,
           $e' \neq e$
  \EndFor
  \State $\bar{X}_{\mathrm{global}} \gets
         \mathrm{ThDec}(\{pd_e\}_{e \in \mathcal{D}})$
         \hfill (Eq.~\eqref{eq:thresh_dec})
  \If{$\bar{X}_{\mathrm{global}} \in
      [\mathbf{L}^{\mathrm{SUMO}},\; \mathbf{U}^{\mathrm{SUMO}}]$}
    \State $\bar{X}_{\mathrm{global}}^{\mathrm{valid}} \gets
           \bar{X}_{\mathrm{global}}$
           \Comment{accept; update last-valid aggregate}
  \Else
    \State log aggregate anomaly to SC-Trust (best-effort)
    \State $\bar{X}_{\mathrm{global}} \gets
           \bar{X}_{\mathrm{global}}^{\mathrm{valid}}$
           \Comment{reject; reuse last valid window. Static
           envelope, no live ground truth, no per-RSU attribution}
  \EndIf
  \State Use $\bar{X}_{\mathrm{global}}$ for local ITS
         services and return to Cloud for analytics
  \State \textbf{return} $\bar{X}_{\mathrm{global}}$
  \EndProcedure
\end{algorithmic}
\label{alg:thresh_dec}
\end{algorithm}


\subsection{Blockchain Trust Layer with Smart Contract Design}
\label{sec:blockchain}

The blockchain layer is implemented on a permissioned consortium
chain (e.g., Hyperledger Fabric) shared among RSUs and the SDN
controller, with RSUs acting as endorsing peers and the controller
as the ordering service.Although every registered RSU holds a peer certificate and is therefore
eligible to endorse, the endorsing committee is not the full RSU
set. At runtime the system selects the $2f+1$ endorsers as the
highest-trust members of the trusted-RSU set $R_{\text{trusted}}(t)$,
re-evaluating this ranking periodically (\textasciitilde$1$\,Hz) from the
live on-chain trust scores maintained by SC-Trust. This keeps endorsement
on the most reliable peers and avoids the cost of treating all RSUs as
active endorsers. When trust scores are uninitialised or tied (e.g.\ at
bootstrap), the selection falls back to a uniform random draw over
$R_{\text{trusted}}(t)$. IPFS is integrated for off-chain storage
of raw beacon sequences: each RSU stores beacon windows on IPFS and
records only the content hash on-chain, enabling real-time retrieval
by the GAT and temporal autoencoder components without incurring
on-chain storage costs.

\subsubsection{Registration Smart Contract (SC-Register)}
Before participating in the framework, each vehicle $v_i$ or RSU $r_j$
submits its identity, public key, and LKH leaf key hash directly to the
blockchain, as shown in Algorithm \ref{alg:sc_register}. SC-Register commits the registration only if at least 2f+1 endorsements are collected from the top-trust 2f+1 members of $R_{\text{trusted}}(t)$, using the same BFT threshold defined in Eq.~\eqref{eq:bft_revoke_rsu}. This eliminates any reliance on a centralized registration authority -- a single compromised RSU or controller cannot unilaterally onboard a fabricated identity. SC-Trust processes anomaly evidence exclusively for identities with a valid
on-chain registration record, and SC-Revoke can only revoke registered
identities, closing the identity lifecycle on-chain from registration
through revocation.

\begin{algorithm}[htb]
\caption{Registration Smart Contract (\texttt{SC-Register})}
\begin{algorithmic}[1]
  \small
  \Procedure{\texttt{SC-Register}}{$\mathrm{ID}_i,\; pk_i,\;
    h(\mathcal{K}_{u_i}),\; \{\sigma_j^{\mathrm{endorse}}\}_{j \in \mathcal{S}}$}
  \If{$\mathrm{ID}_i$ already exists on-chain}
    \State \textbf{reject} -- duplicate identity
  \EndIf
  \State $n_{\mathrm{endorse}} \gets \sum_{r_j \in \mathcal{R}}
    \mathbf{1}\!\left[\mathrm{Verify}(pk_j,\;
    \mathrm{ID}_i \| pk_i \| h(\mathcal{K}_{u_i}),\;
    \sigma_j^{\mathrm{endorse}})\right]$
  \If{$n_{\mathrm{endorse}} < 2f+1$}
    \State \textbf{reject} -- insufficient RSU endorsements
  \EndIf
  \State Commit $(\mathrm{ID}_i,\; pk_i,\;
    h(\mathcal{K}_{u_i}),\; t_{\mathrm{reg}})$ to ledger
  \State Set $\tau_i \gets \tau_{\mathrm{init}}$,\;
    $\mathrm{status}_i \gets \texttt{ACTIVE}$
  \State \textbf{emit} $\texttt{Registered}(\mathrm{ID}_i,\;
    t_{\mathrm{reg}})$
  \EndProcedure
\end{algorithmic}
\label{alg:sc_register}
\end{algorithm}

\subsubsection{Trust Score Smart Contract (SC-Trust)}

The on-chain trust score $\tau_i(t)$ update for vehicle $v_i$ after
each detection epoch is given in Equation~\ref{eq:trust_update}. This exponential moving average update accumulates anomaly evidence
from multiple RSUs across time, preventing a single false positive
from triggering revocation while ensuring that sustained malicious
behaviour gradually reduces the trust score below the revocation
threshold regardless of how slowly the attack progresses.
\begin{equation}
  \tau_i(t) =
  \begin{cases}
    \alpha\,\tau_i(t-1)
    + (1-\alpha)\!\left(1 -
      \dfrac{1}{\bigl|\mathcal{R}^{\mathrm{obs}}_i(t)\bigr|}
      \displaystyle\sum_{r_j \in \mathcal{R}^{\mathrm{obs}}_i(t)}
      \psi_j^{(i)}(t)\right)
    & \mathcal{R}^{\mathrm{obs}}_i(t) \neq \emptyset \\[6pt]
    \tau_i(t-1) & \mathcal{R}^{\mathrm{obs}}_i(t) = \emptyset
  \end{cases}
  \label{eq:trust_update}
\end{equation}
where
$\mathcal{R}^{\mathrm{obs}}_i(t)
  \triangleq \{r_j \in \mathcal{R}_{\mathrm{trusted}}(t)
  : v_i \in \mathcal{V}_j(t)\}$
is the set of currently trusted RSUs that observed vehicle $v_i$
in epoch~$t$, $\alpha \in (0,1)$ is the EMA decay factor, and
trust is held constant when no trusted RSU observed the vehicle.
In the lightweight mode, $\psi_j^{(i)}(t)$ from
Eq.~\eqref{eq:psi_score} is used directly; in the full mode,
$\Phi_i(t)$ from Eq.~\eqref{eq:fusion} substitutes for
$\psi_j^{(i)}(t)$.

SC-Trust maintains identical trust score records for all
registered controllers in $\mathcal{C}$ alongside vehicles.
RSU evidence about controller behaviour is the binary conflict
indicator $\mathrm{conflict}_j(t)$ defined in
Eq.~\eqref{eq:conflict_pred}. The trust score update for
controller $c_k$ driven by RSU conflict evidence submitted
directly to the blockchain is shown in
Equation~\eqref{eq:ctrl_trust}:
\begin{equation}
  \tau_{c_k}(t) = \alpha\,\tau_{c_k}(t-1) + (1-\alpha)\,
    \left(1 -
      \frac{1}{\bigl|\mathcal{R}^{\mathrm{obs}}_{c_k}(t)\bigr|}
      \sum_{r_j \in \mathcal{R}^{\mathrm{obs}}_{c_k}(t)}
      \mathrm{conflict}_j(t)\right)
  \label{eq:ctrl_trust}
\end{equation}
where
$\mathcal{R}^{\mathrm{obs}}_{c_k}(t)
  \triangleq \{r_j \in \mathcal{R}_{\mathrm{trusted}}(t)
  : c_{\mathrm{assigned}}(r_j) = c_k\}$
is the set of currently trusted RSUs assigned to controller
$c_k$ in epoch~$t$.
As shown in Equation~\eqref{eq:ctrl_trust}, a controller whose
decisions systematically conflict with the RSUs assigned to it
accumulates decreasing trust. When $\tau_{c_k}(t)$ remains below
$\tau_{\min}$ for $T_{\mathrm{rev}}$ consecutive epochs,
SC-Trust emits a revocation request to SC-Revoke, triggering
RSU reassignment to $\mathcal{C}_{\mathrm{trusted}}(t)$ as
defined in Eq.~\eqref{eq:trusted_ctrl_set}. As shown in
Equation~\ref{eq:trust_update}, the trust score is updated as an
exponential moving average, where higher anomaly scores from RSUs
reduce $\tau_i(t)$, and sustained low scores eventually trigger
revocation. The SC-Trust contract enforces the following state
machine: $\tau_i \in [\tau_{\min}, 1]$; if
$\tau_i(t) < \tau_{\mathrm{warn}}$ the vehicle is placed in a
probationary state receiving reduced routing priority;
if $\tau_i(t)$ remains below $\tau_{\min}$ for $T_{\mathrm{rev}}$
consecutive epochs, SC-Trust emits a \texttt{RevocationRequest}
event to SC-Revoke.

\textbf{Tiered Commit Strategy.}
Per-beacon synchronous blockchain commits are impractical at the
100\,ms beacon rate: each Hyperledger Fabric consensus round-trip
would block the RSU processing pipeline and inflate end-to-end
latency. A four-tier commit strategy is therefore adopted to balance
security guarantees against operational feasibility.

Tier~1 --- Immediate synchronous commit.
Operations whose outcome must be reflected in
$\mathcal{R}_{\mathrm{trusted}}(t)$ or
$\mathcal{C}_{\mathrm{trusted}}(t)$ before the next epoch is
processed are committed synchronously: \texttt{Revoke}($\mathrm{ID}_i$)
and \texttt{ControllerRevoked} events from SC-Revoke, and RSU
controller-reassignment updates to $c_{\mathrm{assigned}}(r_j)$.
These events change which RSU evidence counts toward the BFT
threshold and must therefore achieve ledger finality immediately.

Tier~2 --- Per-epoch commit (every $L \cdot T_b$ seconds).
Trust score updates that cross a threshold boundary --- the
trusted-to-probationary transition ($\tau < \tau_{\mathrm{warn}}$)
and the probationary-to-client transition ($\tau < \tau_{\min}$)
--- are committed once per detection epoch, as are TRS aggregate
signatures $\sigma_{\mathrm{TRS}}$ and RSU lifecycle state changes.
One Fabric transaction per epoch per event type replaces per-beacon
writes for these events.

Tier~3 --- Batched periodic commit (every $T_{\mathrm{batch}}$
seconds or $N_{\mathrm{commit}}$ items).
Anomaly evidence submissions $\mathcal{E}_j(t)$
(Eq.~\ref{eq:rsu_evidence}) that have not triggered a threshold
crossing are accumulated in a local per-RSU evidence buffer and
flushed to the ledger as a single batched Fabric transaction every
$T_{\mathrm{batch}}$ seconds or when the buffer reaches
$N_{\mathrm{commit}}$ entries, whichever occurs first.
The flush executes asynchronously in a background thread without
blocking the simulation or detection pipeline.
The BFT security guarantee of Equation~\ref{eq:bft_revoke_rsu}
is preserved: the revocation condition counts distinct RSU
submissions within temporal window $T_w$, and since
$T_{\mathrm{batch}} \ll T_w$, all batched submissions arrive
within the evidence accumulation window regardless of batching.

Tier~4 --- Off-chain only (IPFS, no chain commit).
Raw beacon windows are stored on IPFS with only their content
hash committed on-chain as already specified; individual per-beacon
$\psi_i(t)$ scores, GAT and LSTM-AE detection scores $\Phi_i(t)$,
and controller evidence submissions $\mathcal{E}_c(t)$
(Eq.~\ref{eq:ctrl_evidence}) are stored off-chain and their hashes
included in the next Tier~3 batch, since these are untrusted peer
contributions that do not require immediate ledger finality.

The design trade-off introduced by Tiers~2 and~3 is that trust
score transitions are reflected on-chain at epoch or batch
granularity rather than per-beacon. The revocation decision
itself (Tier~1) remains immediate, so the consequential
network action --- exclusion from $\mathcal{R}_{\mathrm{trusted}}(t)$
--- is unaffected by the batching delay.


SC-Trust additionally maintains a trust score $\tau_{r_j}(t)$ for
every registered RSU $r_j \in \mathcal{R}$, instantiating the RSU
decay-and-revocation mechanism anticipated in the cryptographic
mitigation scope analysis. Because a compromised RSU is an endorsing
peer rather than a monitored vehicle, its evidence cannot be scored
against an external anomaly detector; instead it is scored against the
collective decision of its honest peers. For each vehicle $v_i$ that
$r_j$ reports on within window $T_w$, let
$\mathrm{flag}_i^{\mathrm{rsu},j}$ be $r_j$'s own binary verdict
(Eq.~\eqref{eq:rsu_bin_flag}), and let $q_i(t)\in\{0,1\}$ be the
$2f+1$ quorum outcome of Eq.~\eqref{eq:bft_revoke_rsu} computed
only over the current trusted endorser set
$\mathcal{R}_{\mathrm{trusted}}(t)$. The per-window RSU misbehaviour
signal is the mean disagreement between the two, shown in
Equation~\eqref{eq:rsu_misbehave}:
\begin{equation}
  m_j(t) = \frac{1}{|\mathcal{V}_j(t)|}
    \sum_{v_i \in \mathcal{V}_j(t)}
    \bigl|\,\mathrm{flag}_i^{\mathrm{rsu},j}(t) - q_i(t)\,\bigr|
  \label{eq:rsu_misbehave}
\end{equation}
where $\mathcal{V}_j(t)$ is the set of vehicles $r_j$ reported on in
the window. Equation~\eqref{eq:rsu_misbehave} penalises both
under-reporting (a malicious vehicle the trusted quorum caught
but $r_j$ cleared, $\mathrm{flag}=0,\,q_i=1$) and
over-reporting (a benign vehicle the trusted quorum cleared but
$r_j$ flagged, $\mathrm{flag}=1,\,q_i=0$). Evaluating $q_i(t)$ against
$\mathcal{R}_{\mathrm{trusted}}(t)$ alone prevents a coalition of
demoted RSUs from mutually validating one another's fabricated
verdicts. The RSU trust score then decays by the same exponential
moving average used for vehicles (Eq.~\eqref{eq:trust_update}), driven
by $m_j(t)$, as shown in Equation~\eqref{eq:rsu_trust}:
\begin{equation}
  \tau_{r_j}(t) = \alpha\,\tau_{r_j}(t-1)
    + (1-\alpha)\,\bigl(1 - m_j(t)\bigr)
  \label{eq:rsu_trust}
\end{equation}
SC-Trust enforces a three-state lifecycle on each RSU, reusing the
vehicle thresholds. While $\tau_{r_j}(t)\in[\tau_{\min},1]$ the RSU is
trusted: its evidence counts toward the $2f+1$ quorum of
Eq.~\eqref{eq:bft_revoke_rsu} and it is eligible to endorse. If
$\tau_{r_j}(t)<\tau_{\mathrm{warn}}$ it enters a probationary
trusted state. If $\tau_{r_j}(t)<\tau_{\min}$ it is demoted to the
client state: it may still submit evidence---so its score can
recover---but those submissions carry zero quorum weight and it may
not endorse. A client-state RSU whose subsequent $m_j(t)$ realigns
with $\mathcal{R}_{\mathrm{trusted}}(t)$ and whose score climbs back
to $\tau_{\min}$ is promoted to the trusted state. If $\tau_{r_j}(t)$
remains below $\tau_{\min}$ for $T_{\mathrm{rev}}$ consecutive windows,
SC-Trust emits a \texttt{RevocationRequest} for $r_j$ to SC-Revoke,
permanently removing it from $\mathcal{R}_{\mathrm{trusted}}$.The BFT safety guarantee requires $|\mathcal{R}_{\mathrm{trusted}}(t)| \geq 3f+1$ at all times; if demotion events would violate this floor, the lowest-trust probationary RSU is retained as trusted until a replacement is registered and endorsed.


\subsubsection{Revocation Smart Contract (SC-Revoke)}

Since the threat model includes a compromised controller, SC-Revoke
is triggered exclusively by RSU peer evidence submitted directly to
the blockchain, bypassing the controller entirely. The signed evidence
tuple submitted by each RSU $r_j$ directly to the blockchain is given
in Equation~\ref{eq:rsu_evidence}.
\newline
\newline
This signed evidence tuple enables each RSU to submit tamper-evident
anomaly reports directly to the blockchain without routing through
the controller, ensuring that a compromised controller cannot
suppress or modify RSU observations before they reach the
revocation smart contract.
\begin{equation}
  \mathcal{E}_j(t) = \bigl(\mathrm{ID}_i,\;
    \psi_j^{(i)}(t),\; t,\; h(\mathbf{b}_i(t)),\;
    \sigma_j^{\mathrm{sub}}\bigr)
  \label{eq:rsu_evidence}
\end{equation}
As shown in
Equation~\ref{eq:rsu_evidence}, each RSU independently submits a
signed evidence tuple containing the vehicle identity, anomaly score,
timestamp, beacon hash, and its own verifiable signature, ensuring
tamper-evident RSU-side evidence on-chain where $\sigma_j^{\mathrm{sub}}$ is RSU $r_j$'s signature over the
submission, verifiable on-chain using $pk_j$. The supplementary evidence
tuple that the controller may submit is given in
Equation~\ref{eq:ctrl_evidence}.
\newline
\newline
This submission structure treats the controller's anomaly score as
one additional peer vote subject to the same BFT verification as
RSU submissions, preventing a compromised controller from gaining
disproportionate influence over revocation decisions while still
incorporating full-mode AI detection results when the controller
is honest.
\begin{equation}
  \mathcal{E}_c(t) = \bigl(\mathrm{ID}_i,\;
    \Phi_i(t),\; t,\; h(\mathbf{X}_i(t)),\;
    \sigma_c^{\mathrm{sub}}\bigr)
  \label{eq:ctrl_evidence}
\end{equation}
As shown in Equation~\ref{eq:ctrl_evidence}, the controller's
submission is treated as one additional peer vote and never a
prerequisite for revocation.
\newline
\newline
Since a vehicle typically interacts with only one RSU at a time,
requiring $2f+1$ simultaneous RSU observations of the same vehicle
is geometrically infeasible under real mobility. Revocation is
therefore defined over the number of distinct RSUs that
have independently flagged vehicle $v_i$ within a temporal window
$T_w$, accumulating evidence across all RSUs a vehicle visits
over time. Counting each RSU at most once removes the risk of a
single malicious RSU inflating the evidence count through
repeated submissions. The revocation condition requiring at least
$2f+1$ distinct RSUs to have independently flagged vehicle $v_i$
within window $T_w$ is given in Equation~\ref{eq:bft_revoke_rsu}:
\begin{equation}
    \left|\left\{r_j \in \mathcal{R}_{\mathrm{trusted}}(t) :\;
    \exists\, t' \in [t - T_w,\, t],\;
    \psi_j^{(i)}(t') > \psi_{\mathrm{th}}\right\}\right|
  \geq 2f + 1
  \;\Longrightarrow\;
  \text{emit } \texttt{Revoke}(\mathrm{ID}_i)
  \label{eq:bft_revoke_rsu}
\end{equation}
As shown in Equation~\ref{eq:bft_revoke_rsu}, each RSU is counted
at most once regardless of how many beacon intervals it reports
anomalies for vehicle $v_i$, so no single RSU can inflate the
evidence count. Revocation requires $2f+1$ independent RSU
witnesses accumulated over time rather than simultaneously,
making it reachable under normal single-RSU-at-a-time mobility
while remaining robust against single-source false positives. The controller compromise detection condition is defined at the
binary decision level to avoid undefined comparisons between
heterogeneous continuous scores. The binary detection flag
produced by the controller for vehicle $v_i$ is shown in
Equation~\eqref{eq:ctrl_bin_flag}:
\begin{equation}
  \mathrm{flag}_i^{\mathrm{ctrl}} =
    \mathbf{1}\!\left[\Phi_i(t) > \Phi_{\mathrm{th}}\right]
  \label{eq:ctrl_bin_flag}
\end{equation}
As shown in Equation~\eqref{eq:ctrl_bin_flag}, the controller
produces a binary alert or no-alert decision from its fusion
score. The binary detection flag produced independently by RSU
$r_j$ from its local rule-based score is shown in
Equation~\eqref{eq:rsu_bin_flag}:
\begin{equation}
  \mathrm{flag}_i^{\mathrm{rsu},j} =
    \mathbf{1}\!\left[\psi_j^{(i)}(t) > \psi_{\mathrm{th}}\right]
  \label{eq:rsu_bin_flag}
\end{equation}
As shown in Equation~\eqref{eq:rsu_bin_flag}, each RSU
independently produces its own binary decision without consulting
the controller. A conflict is defined directionally: only
the case where the controller classifies a vehicle as benign
while the RSU flags it as anomalous constitutes a conflict. The
reverse case the controller flagging a vehicle that the
lightweight RSU rules miss is the controller performing
superior detection and must not be penalised. The directional
conflict indicator that fires only when the controller fails to
flag what RSUs flag is shown in
Equation~\eqref{eq:conflict_pred}:
\begin{equation}
  \mathrm{conflict}_j(t) =
    \bigl(1 - \mathrm{flag}_i^{\mathrm{ctrl}}\bigr)
    \cdot
    \mathrm{flag}_i^{\mathrm{rsu},j}
  \label{eq:conflict_pred}
\end{equation}
As shown in Equation~\eqref{eq:conflict_pred}, a conflict fires
only when $\mathrm{flag}_i^{\mathrm{ctrl}} = 0$ and
$\mathrm{flag}_i^{\mathrm{rsu},j} = 1$, identifying controller
suppression without penalising superior full-mode detection.
The controller compromise flag aggregating binary conflicts
across RSU peers is shown in Equation~\eqref{eq:ctrl_flag}:
\begin{equation}
  \mathrm{flag}_c = \mathbf{1}\!\left[
    \sum_{r_j \in \mathcal{R}}
    \mathrm{conflict}_j(t) \geq f + 1\right]
  \label{eq:ctrl_flag}
\end{equation}
As shown in Equation~\eqref{eq:ctrl_flag}, when at least $f+1$
independent RSU--controller disagreements are detected,
SC-Revoke sets $\mathrm{flag}_c = 1$, excludes
$\mathcal{E}_c(t)$ from all subsequent voting rounds, and
delegates routing decisions to RSU-level rule-based control.
RSUs then automatically reassign to any available trusted
controller from $\mathcal{C}_{\mathrm{trusted}}(t)$ per
Eq.~\eqref{eq:trusted_ctrl_set}, with no manual failover
required. Upon revocation, SC-Revoke broadcasts a signed
\texttt{CRL-Update} directly to all RSUs, without controller
relay.

SC-Revoke applies identically to controllers in $\mathcal{C}$.
When SC-Trust triggers revocation of controller $c_k$, SC-Revoke
marks $c_k$ as untrusted on-chain, updates
$\mathcal{C}_{\mathrm{trusted}}(t)$ per
Eq.~\eqref{eq:trusted_ctrl_set}, and broadcasts a
\texttt{ControllerRevoked}($\mathrm{ID}_{c_k}$) event to all
RSUs. RSUs assigned to $c_k$ query the blockchain for
$\mathcal{C}_{\mathrm{trusted}}(t)$ and independently reassign
to any available trusted controller, maintaining uninterrupted
network operation without a designated replacement or manual
failover.
\newline
\newline
Algorithm~\ref{alg:cp_detect} presents the CP-DETECT procedure,
which verifies the post-quantum TRS aggregate signature at the RSU
level and cross-checks the controller's blockchain submission against
RSU peer evidence to detect control-plane compromise.

\begin{algorithm}[H]
\caption{Control-Plane Poisoning Detection (\texttt{CP-DETECT})}
\begin{algorithmic}[1]
  \small
  \Procedure{\texttt{CP-DETECT}}{$\mathbf{A}_j(t), \sigma_{\mathrm{TRS}},
    \{pk_j\}_{j=1}^{n}, \{\mathcal{E}_j(t)\}_{r_j \in \mathcal{R}}$}
  \State Verify TRS at RSU level:
         $\mathrm{valid} \gets
         \mathrm{Verify}(\{pk_j\}, m_j, \sigma_{\mathrm{TRS}})$
         \quad (Eq.~\eqref{eq:trs_verify})
  \If{$\neg\,\mathrm{valid}$}
    \State RSUs emit $\mathrm{ALERT\_CP}(\mathrm{ID}_{r_j},
           \texttt{TRS\_FAIL})$ directly to blockchain
    \State \textbf{return}
  \EndIf
  \State Retrieve controller submission $\mathcal{E}_c(t)$
         from blockchain
  \State For each $r_j \in \mathcal{R}$: compute
         $\mathrm{conflict}_j(t) =
         (1 - \mathrm{flag}_i^{\mathrm{ctrl}}) \cdot
         \mathrm{flag}_i^{\mathrm{rsu},j}$
         \hfill (Eq.~\eqref{eq:conflict_pred})
  \State Compute $\mathrm{flag}_c$ via
         Eq.~\eqref{eq:ctrl_flag}
  \If{$\mathrm{flag}_c = 1$}
    \State SC-Revoke excludes $\mathcal{E}_c(t)$ from consensus
    \State RSUs reassign to $\mathcal{C}_{\mathrm{trusted}}(t)$
           (Eq.~\eqref{eq:trusted_ctrl_set}); no manual failover
    \State \textbf{return} $\mathrm{ALERT\_CP}
           (\mathrm{ID}_c, \texttt{CTRL\_COMPROMISED})$
  \EndIf
  \State \textbf{return} $\mathrm{PASS}$
  \EndProcedure
\end{algorithmic}
\label{alg:cp_detect}
\end{algorithm}

As shown in Algorithm~\ref{alg:cp_detect}, a TRS verification failure
immediately causes RSUs to emit an alert directly to the blockchain,
bypassing the controller entirely. If verification passes but more
than $f+1$ RSUs report a conflict with the controller submission,
SC-Revoke excludes the controller from consensus and activates
RSU-level rule-based fallback control. The fallback control mechanism
delegates routing and safety decisions to the distributed rule-based
signature detection running at each RSU independently, using the
lightweight mode pipeline from Algorithm~\ref{alg:lightweight}, which
ensures that the network continues to operate at a reduced but
functional security level even when the controller is excluded from
consensus. The conflict detection in lines~8--10 distinguishes between
a malicious controller and a merely slow or partitioned one by
requiring $f+1$ conflicting RSU submissions rather than a single
discrepancy: a controller that is temporarily delayed will produce
consistent evidence once its submission arrives, whereas a malicious
controller will consistently submit scores that contradict the
RSU-observed anomaly levels across multiple epochs. The signed
evidence structure in Equation~\ref{eq:rsu_evidence} is what makes
this cross-checking possible on-chain without requiring RSUs to
communicate directly with each other, as each RSU's submission is
independently verifiable using its public key $pk_j$ recorded on the
permissioned ledger.

CP-DETECT's conflict computation combines two distinct vantage points
on the same epoch, and the framework's logging deliberately keeps them
separate rather than collapsing them into a single verdict. The
RSU-view verdict $\mathrm{flag}_i^{\mathrm{rsu},j}$ is produced locally
at each RSU $r_j$ from its own lightweight or full-mode detection
pipeline, before anything is forwarded to the controller; the
controller-view verdict $\mathrm{flag}_i^{\mathrm{ctrl}}$ is the
decision the controller reports back on-chain after receiving the
RSU's submission. This two-vantage-point separation reflects the
threat model directly: a compromised controller can misreport its own
verdict without being able to alter what an honest RSU already
observed and signed. Every logged beacon-evaluation record is tagged
with which vantage point produced it, so that the conflict indicator
$\mathrm{conflict}_j(t)$ (Eq.~\ref{eq:conflict_pred}) can be
recomputed and audited independently of the live CP-DETECT run.

\subsubsection{Real-Time vs.\ Offline Detection Mode}

The proposed framework explicitly supports two temporal modes.
Real-time detection executes the rule-based signature check
(Algorithm~\ref{alg:lightweight}) at the RSU per beacon arrival
(latency $\leq T_b = 100$\,ms), the TRS signing per aggregate epoch,
and the SC-Trust update after each epoch, all without waiting for
controller-side AI inference. Offline (semi-real-time) detection
executes the GAT (Algorithm~\ref{alg:full_detection}) and temporal
autoencoder at the controller over a sliding window updated every
$W = L \cdot T_b$ seconds, with results fed back to SC-Trust as
a correction signal. IPFS-stored beacon sequences enable the
controller to retrieve historical windows on demand for reanalysis,
bridging the real-time and offline modes without data loss.

Figure~\ref{fig:detection_flow_1_2} presents the detection logic flow diagrams for trajectory poisoning detection (TP-DETECT) and Sybil-based mobility pattern poisoning detection (SYB-DETECT), showing the sequential signature checks and decision branching for each attack class.

\begin{figure}[H]
    \centering
    \includegraphics[width=0.9\linewidth]{images/Fig5A}
    \caption{Detection flow diagrams for TP-DETECT and SYB-DETECT.}
    \label{fig:detection_flow_1_2}
\end{figure}
As shown in Figure~\ref{fig:detection_flow_1_2}, TP-DETECT (left)
evaluates four signature checks sequentially TP-S1 (position jump),
TP-S2 (heading rate), TP-S3 (acceleration bound), and TP-S5
(cumulative drift score) where any violation feeds into the
composite score $\psi_i(t)$; if $\psi_i(t) > \psi_{\mathrm{th}}$
an ALERT is raised and signed evidence is submitted directly to the
blockchain via $\mathcal{E}_j(t)$, otherwise a PASS is returned with
no blockchain interaction. SYB-DETECT (right) checks MP-S1 (identity
density anomaly), MP-S2 (synchronized beacon timing), and MP-S4
(ghost vehicle transit impossibility) before applying the same
composite score threshold decision.
\newline
\newline

\begin{figure}[H]
    \centering
    \includegraphics[width=0.9\linewidth]{images/Fig5B}
    \caption{Detection flow diagrams for MITM-DETECT and CP-DETECT.}
    \label{fig:detection_flow_3_4}
\end{figure}
As shown in Figure~\ref{fig:detection_flow_3_4}, MITM-DETECT (left)
first verifies the HMAC tag a failure immediately triggers an ALERT
without further processing then checks MP-S3 (KL-divergence speed
distribution shift) and the dead-reckoning residual before computing
the composite score. CP-DETECT (right) operates at the aggregate
level: TRS threshold verification failure triggers a TRS\_FAIL alert
emitted directly by RSUs to the blockchain; if verification passes,
controller submission conflicts with more than $f+1$ RSU peers trigger
CTRL\_COMPROMISED and activate RSU fallback control; finally, the
temporal autoencoder reconstruction error is checked against the
adaptive threshold, raising AE\_ANOMALY if exceeded.


\chapter{Experimental Setup}
\label{chap:evaluation}

\section{Baseline Selection and Benchmarking Rationale}
\label{sec:baselines}

\subsection{External Baselines}

Three external baselines are selected from post 2020 literature that
address mobility data anomaly detection or vehicular network poisoning
in related settings.

\textbf{B1 - Ghaleb et al.\ (2014)~\cite{ghaleb2014}:}
Selected as the most directly comparable work addressing mobility
pattern misbehavior detection in vehicular networks using rule based
threshold checks. Despite its age, no post 2020 SDVN specific mobility
poisoning detector exists in the literature, making this the canonical
baseline for the detection component. It represents the class of
threshold based approaches against which our signature characterized
detection must demonstrate superiority.

\textbf{B2 - Ercan et al.\ (2022)~\cite{ercan2022misbehavior}:}
Selected as a post 2020 machine learning baseline that specifically
targets position falsification misbehavior detection in VANETs,
representing feature engineered supervised classification over beacon
data. Comparison against B2 evaluates whether our mobility specific
spatial (GAT) and temporal (LSTM-AE) formulation provides detection
gains over engineered per beacon features classified in isolation.

\textbf{B3 - Sharma and Liu (2021)~\cite{sharma2021machine}:}
Selected as a post 2020 data centric machine learning misbehavior
detection baseline for the Internet of Vehicles using plausibility based
feature engineering on basic safety messages. Comparison against B3
isolates the contribution of our graph relational spatial fusion and
sequential temporal modeling over a strong but message level
data centric classifier.


\subsection{Ablation Study Design}
\label{sec:ablation}

Eleven ablation variants are defined to isolate every independently
ablatable design component of the proposed framework. Each variant
removes or replaces exactly one component while holding all others
fixed, enabling clean attribution of observed performance changes
to the removed component. The five cryptographic and blockchain
variants (AB6--AB11) use ablation-only metrics (PARR, COO, BWO,
TCL, FRR) that cannot be evaluated against external baselines
B1--B3 because those baselines share no equivalent mechanism.
The six detection variants (AB1--AB5, AB8--AB9) additionally
report metrics that are directly comparable across baselines.
Table~\ref{tab:ablation-design} summarises the full design.

\begin{table}[H]
\centering
\caption{Ablation study design: eleven independent component
isolation experiments. Metric abbreviations follow
Section~\ref{sec:metrics}.}
\label{tab:ablation-design}
\renewcommand{\arraystretch}{1.2}
\resizebox{\textwidth}{!}{%
\begin{tabular}{|l|l|p{3.8cm}|p{3.8cm}|p{4.5cm}|}
\hline
\textbf{ID} &
\textbf{Subsystem} &
\textbf{Component removed / modified} &
\textbf{X variable} &
\textbf{Y metrics} \\
\hline
AB1 & Detection & Rule-based signatures (TP-S1--MP-S4) &
  Attack penetration rate $\rho_a$ &
  MCC, TTD, PBPO \\
\hline
AB2 & Detection & HMAC beacon integrity and nonce anti-replay &
  MitM penetration rate $\rho_{\mathrm{MitM}}$ &
  MCC, TTD, CDER \\
\hline
AB3 & Detection & GAT spatial detector &
  Coordinated attacker count $n_{\mathrm{coord}}$ &
  MCC, TTD, CDER \\
\hline
AB4 & Detection & LSTM-AE temporal detector &
  Stealth drift cap $\epsilon_{\max}$ &
  MCC, TTD, CDER, TPE \\
\hline
AB5 & Detection & Full AI layer (GAT + LSTM-AE) &
  Vehicle speed regime &
  MCC, TTD, CDER, TDEE, TPE, PBPO \\
\hline
AB6 & Cryptography & TRS threshold signing gate &
  Compromised RSU fraction $f/n$ &
  PARR, CDER, COO \\
\hline
AB7 & Cryptography & FHE pre-coordination encryption &
  RSU ring size $n$ &
  COO, TRS signing latency (PQ vs.\ ECDSA), BWO\textsubscript{ratio}, DKG setup latency \\
\hline
AB8 & Trust & RSU three-state trust lifecycle &
  Malicious RSU fraction $f/n$ &
  MCC, CDER, FRR\textsubscript{revoke}, FRR\textsubscript{demote} \\
\hline
AB9 & Trust & Multi-controller architecture and controller revocation &
  Controller compromise onset $t_{\mathrm{comp}}/T_{\mathrm{sim}}$ &
  MCC, CDER \\
\hline
AB10 & Trust & Blockchain trust layer (SC-Trust / SC-Revoke) &
  Revocation-triggering vehicle count &
  CDER, FRR\textsubscript{revoke}, FRR\textsubscript{demote},
  TCL\textsubscript{confirm}, TCL\textsubscript{reassign} \\
\hline
AB11 & Key Mgmt. & LKH replaced with unicast key distribution &
  Vehicle zone density $|\mathcal{V}_j|$ &
  BWO\textsubscript{scale}, PBPO \\
\hline
\end{tabular}}
\end{table}

\begin{description}

\item[\textbf{AB1 --- Rule-Based Signatures Removed.}]
The nine rule-based attack signatures (TP-S1 through MP-S4) and the
composite anomaly score $\psi_i(t)$ (Eq.~\ref{eq:psi_score}) are
disabled. HMAC beacon integrity verification, GAT, LSTM-AE, and the
blockchain trust layer remain fully active. The system relies
entirely on the HMAC integrity gate and AI detection with no
rule-based pre-filter. Attack penetration rate $\rho_a$ is varied
across $\{0.05, 0.10, 0.20, 0.30, 0.40\}$. MCC quantifies detection
quality loss without the kinematic pre-filter; TTD quantifies the
detection delay increase when early kinematic violations can no longer
trigger immediate lightweight alerts; PBPO captures the per-beacon
overhead change when the signature scoring stage is bypassed.
Isolates the rule-based signatures' contribution to detection speed
and lightweight beacon pre-filtering before AI inference.

\item[\textbf{AB2 --- HMAC Removed.}]
HMAC tag generation, freshness window verification, and nonce
anti-replay enforcement (Eq.~\ref{eq:hmac}) are disabled. Rule-based
signatures, GAT, LSTM-AE, and the blockchain trust layer remain
fully active. MitM adversaries can modify beacon content in transit
without triggering the integrity gate before rule-based scoring.
MitM penetration rate $\rho_{\mathrm{MitM}}$ is varied across
$\{0.05, 0.10, 0.20, 0.30, 0.40\}$. MCC captures the detection
quality loss under unrestricted in-transit modification; TTD captures
how much slower detection becomes when modified beacons bypass the
fail-fast HMAC gate and must be caught by downstream components;
CDER captures the end-to-end control decision degradation from
undetected in-transit poisoning. Isolates HMAC's specific
contribution against data-plane interception, which is the only
mechanism that catches modification before any rule-based or AI
scoring is applied.

\item[\textbf{AB3 --- GAT Disabled.}]
The GAT spatial detector is removed; $\lambda_2^{(k)} = 0$ for all
$k$, and the remaining fusion weights $\lambda_1^{(k)}, \lambda_3^{(k)}$
are renormalized to sum to one within each attack-class set. HMAC, rule-based signatures, LSTM-AE, and blockchain
remain fully active. The number of simultaneously coordinating
attack vehicles $n_{\mathrm{coord}}$ is varied across $\{1, 2, 5,
10, 20\}$ while total penetration rate is held fixed at
$\rho_a = 0.20$. Higher $n_{\mathrm{coord}}$ increases the spatial
collective inconsistency that the GAT is specifically designed to
detect. MCC captures detection quality loss on coordinated attacks;
TTD captures the delay increase when neighborhood-level spatial
inconsistency can no longer be detected at the embedding level;
CDER captures the end-to-end control decision impact of undetected
coordinated poisoning. Isolates the GAT's unique contribution to
detecting multi-vehicle spatial configurations that individually
pass all per-beacon and per-vehicle checks.

\item[\textbf{AB4 --- LSTM-AE Disabled.}]
The LSTM temporal autoencoder is removed; $\lambda_3^{(k)} = 0$ for
all $k$, and the remaining fusion weights $\lambda_1^{(k)}, \lambda_2^{(k)}$
are renormalized within each attack-class set.
HMAC, rule-based signatures, GAT, and blockchain remain fully active.
The stealth perturbation magnitude $\epsilon_{\max}$ is varied across
$\{0.1, 0.2, 0.3, 0.4, 0.5\}$\,m per beacon while penetration rate
is fixed at $\rho_a = 0.20$. Lower $\epsilon_{\max}$ produces
trajectories that individually pass all kinematic checks and produce
spatially plausible embeddings, isolating the regime where only
cumulative temporal drift is detectable. MCC captures detection
quality loss on near-invisible stealth attacks; TTD captures the
detection delay increase when sequential temporal drift can no longer
be accumulated into a reconstruction error signal; CDER captures
the control decision degradation; TPE captures the trajectory
prediction accuracy lost when gradual drift goes undetected.
Isolates the LSTM-AE's unique contribution to catching gradual
temporal trajectory poisoning that stays within per-beacon kinematic
bounds throughout the attack.

\item[\textbf{AB5 --- Full AI Layer Removed (Lightweight Mode Only).}]
Both GAT and LSTM-AE are removed; $\lambda_2^{(k)} = \lambda_3^{(k)} = 0$
for all $k$.
Detection relies solely on HMAC and rule-based signatures. The
vehicle speed regime is the independent variable, evaluated across
the three scenarios: urban ($\leq 30$\,km/h), rural
($30$--$80$\,km/h), and highway ($> 80$\,km/h). This directly
exposes the Detection Infeasible Zone (Eq.~\ref{eq:detection_condition}):
at high speed, the beacon sampling window $N_b$ falls below
$N_{\min}$ and rule-based cumulative drift detection degrades
while AI-based sequential anomaly scoring would recover detection
power through multi-RSU offline analysis. MCC across speed regimes
shows the detection quality degradation pattern; TTD shows how
detection delay increases in high-speed regimes; CDER and both
attack-impact metrics (TDEE, TPE) quantify the control-plane
consequences of this degradation; PBPO confirms that the
lightweight mode satisfies $\leq T_b = 100$\,ms throughout.
Isolates the entire AI layer's added value over rule-based
detection and validates the mobility amplification formulation.

\item[\textbf{AB6 --- TRS Threshold Signing Gate Removed.}]
Post-quantum threshold ring signature generation and verification
(Eqs.~\ref{eq:trs_partial}--\ref{eq:trs_verify}) are disabled.
Aggregates are forwarded to the cloud without threshold endorsement.
FHE, detection pipeline, and blockchain remain active. The
compromised RSU fraction $f/n$ is varied across $\{0/4, 1/4, 2/4,
3/4\}$ (where $n = 4$ per the simulation configuration).
PARR measures the fraction of poisoned aggregates that would have
been structurally rejected at the TRS gate, showing it drops to
zero by definition without TRS; the comparison between AB6 and
full-mode PARR isolates TRS's structural blocking contribution
independent of AI detection accuracy. CDER captures the
downstream control decision degradation from unblocked poisoned
aggregates. COO quantifies the cryptographic computation saved
by removing TRS, establishing its per-epoch cost. Ablation
only.

\item[\textbf{AB7 --- FHE Pre-Coordination Encryption Removed.}]
FHE encryption of local RSU aggregates before ring sharing
(Eq.~\ref{eq:fhe_rsu_enc}) and homomorphic ring combination
(Eq.~\ref{eq:fhe_ring_combine}) are disabled. RSUs exchange
plaintext aggregates during TRS ring coordination. TRS threshold
signing, detection pipeline, and blockchain remain active. RSU
ring size $n$ is varied across $\{4, 6, 8, 10, 12\}$ to expose how
the post-quantum cryptographic pipeline scales with ring
membership, benchmarked against a classical ECDSA-class
(Shamir-Schnorr-P256) baseline at each ring size. COO quantifies the
per-epoch computation cost of the full pipeline directly; the
per-RSU TRS signing latency is additionally compared against the
classical baseline to isolate the fixed post-quantum signing
premium; BWO\textsubscript{ratio} quantifies the ciphertext
expansion factor of the ring coordination message traffic; and the
one-time DKG plus key-generation setup latency
(Eqs.~\ref{eq:dkg_trs}--\ref{eq:dkg_fhe}) is reported as an
amortised deployment cost. Together, these metrics justify FHE and
post-quantum TRS deployment by making their cost explicit against
the privacy and structural-integrity guarantees they provide.
Ablation only.

\item[\textbf{AB8 --- RSU Three-State Trust Lifecycle Removed.}]
The RSU misbehavior signal (Eq.~\ref{eq:rsu_misbehave}), RSU
trust EMA update (Eq.~\ref{eq:rsu_trust}), and three-state
demotion/revocation mechanism for RSUs are disabled. RSUs remain
permanently trusted once registered, regardless of observed
behavior. All vehicle and controller trust mechanisms, detection
pipeline, and cryptographic components remain active. The malicious
RSU fraction $f/n$ is varied across $\{0/4, 1/4, 2/4, 3/4\}$.
MCC captures detection quality loss when compromised RSUs are
never demoted despite systematic misbehavior, since their evidence
continues to count toward the BFT quorum with equal weight.
CDER captures the end-to-end control decision degradation from
sustained insider RSU poisoning that the TRS alone cannot prevent
(compromised ring members hold valid keys and produce valid
signatures). FRR\textsubscript{revoke} and
FRR\textsubscript{demote} validate that the lifecycle is safe:
honest RSUs must not be incorrectly demoted or revoked by the
misbehavior signal under normal operating conditions. Isolates
the RSU three-state trust lifecycle's contribution to containing
insider RSU threats that exist within an otherwise valid key ring.

\item[\textbf{AB9 --- Multi-Controller Architecture Removed.}]
The trusted controller set $\mathcal{C}_{\mathrm{trusted}}(t)$
(Eq.~\ref{eq:trusted_ctrl_set}), controller trust score updates
(Eq.~\ref{eq:ctrl_trust}), and controller revocation pathway are
disabled. A single fixed controller replaces the distributed set;
controller compromise has no recovery path. Detection pipeline,
cryptographic components, and RSU/vehicle trust mechanisms remain
active. Controller compromise onset time
$t_{\mathrm{comp}}/T_{\mathrm{sim}}$ is varied across $\{0.1,
0.25, 0.5, 0.75\}$ (fraction of total simulation time elapsed
before compromise). Earlier compromise onset gives the malicious
controller more time to corrupt control decisions before the
simulation ends. MCC captures detection quality degradation when
the SDVN controller is a permanent unrecoverable attack source
that cannot be excluded from evidence handling. CDER captures the
cumulative control decision corruption that accumulates without a
reassignment mechanism. Isolates the multi-controller architecture
and controller revocation pathway's contribution to resilience
under sustained control-plane attacks.

\item[\textbf{AB10 --- Blockchain Trust Layer Removed.}]
SC-Trust and SC-Revoke are disabled. Trust decisions revert to
centralized controller-side revocation with no BFT guarantee and
no controller-bypass path. Detection pipeline, cryptographic
components, and entity trust score computations remain active but
cannot trigger on-chain revocation events. The number of distinct
honest vehicles, RSUs, and controllers that accumulate enough
evidence to trigger revocation is varied from $\{5, 10, 20, 50,
100\}$ over the simulation duration by adjusting the attack
penetration rate. CDER quantifies the end-to-end control decision
degradation from non-BFT, controller-dependent revocation.
FRR\textsubscript{revoke} and FRR\textsubscript{demote} confirm
that safety is also preserved: without the on-chain BFT threshold,
a single malicious controller can suppress legitimate revocations
or fabricate false ones, which should surface as elevated CDER and
abnormal FRR patterns. TCL\textsubscript{confirm} and
TCL\textsubscript{reassign} drop to zero by definition (no
blockchain transactions), establishing the latency baseline for
the full-mode TCL comparison. Isolates the blockchain trust
layer's BFT-conditioned revocation contribution independent of
detection quality. Ablation only.

\item[\textbf{AB11 --- LKH Replaced with Unicast Key Distribution.}]
The Logical Key Hierarchy (Eq.~\ref{eq:lkh_keyset}) is replaced
with direct per-vehicle unicast rekeying, incurring $O(|\mathcal{V}_j|)$
rekey messages per vehicle leave or revocation event instead of
the $O(\log_2|\mathcal{V}_j|)$ LKH cost (Eq.~\ref{eq:lkh_rekey}).
All other components remain unchanged. Vehicle zone density
$|\mathcal{V}_j|$ is varied across $\{10, 25, 50, 100, 200\}$ to
produce the empirical growth curve for both LKH and unicast.
BWO\textsubscript{scale} (Eq.~\ref{eq:bwo_scale}) records the
measured rekey message count for each zone size under both
approaches and compares the empirical growth curves against the
theoretical $\log_2|\mathcal{V}_j|$ and $|\mathcal{V}_j|$
predictions. PBPO captures whether the additional unicast rekey
traffic affects the per-beacon detection latency budget at high
zone densities. Isolates LKH's logarithmic scalability advantage
and empirically validates the closed-form complexity claim in
Equation~\ref{eq:lkh_rekey}. Ablation only.

\end{description}


\section{Primary Evaluation Metrics}
\label{sec:metrics}

Eleven metrics are defined across four evaluation dimensions.
Six metrics are evaluated against external baselines B1--B3 and
internal ablation variants; five metrics are ablation-study
only because the external baselines share no equivalent mechanism
(e.g.\ no TRS gate, no blockchain, no on-chain revocation), making
a fair head-to-head comparison undefined.
Table~\ref{tab:metric-rq} maps each metric to the research question
it primarily answers.

\begin{table}[H]
\centering
\caption{Metric-to-RQ mapping and evaluation scope.}
\label{tab:metric-rq}
\renewcommand{\arraystretch}{1.15}
\begin{tabular}{|l|l|l|l|}
\hline
\textbf{Metric} & \textbf{Eq.} & \textbf{Primary RQ} & \textbf{Scope} \\
\hline
MCC   & \ref{eq:mcc}          & RQ2, RQ3, RQ4        & All baselines \\
\hline
TTD   & \ref{eq:ttd}          & RQ1, RQ3, RQ4        & All baselines \\
\hline
CDER  & \ref{eq:cder}         & RQ2, RQ6             & All baselines \\
\hline
TDEE  & \ref{eq:tdee}         & RQ2, RQ4, RQ5        & All baselines \\
\hline
TPE   & \ref{eq:tpe}          & RQ1, RQ4             & All baselines \\
\hline
PBPO  & \ref{eq:pbpo}         & RQ3, RQ5             & All baselines \\
\hline
PARR  & \ref{eq:parr}         & RQ5                  & Ablation only \\
\hline
FRR   & \ref{eq:frr_revoke}, \ref{eq:frr_demote} & RQ6 & Ablation only \\
\hline
COO   & \ref{eq:coo}          & RQ5                  & Ablation only \\
\hline
BWO   & \ref{eq:bwo_ratio}, \ref{eq:bwo_scale}   & RQ5, RQ6 & Ablation only \\
\hline
TCL   & \ref{eq:tcl_confirm}, \ref{eq:tcl_reassign} & RQ6 & Ablation only \\
\hline
\end{tabular}
\end{table}

\subsection{Detection Performance Metrics}

\textbf{Matthews Correlation Coefficient (MCC)} is used as the
primary detection quality metric, robust to class imbalance between
attack and normal beacon samples. The MCC is given in
Equation~\ref{eq:mcc}.
\newline
\newline
MCC is selected as the sole detection quality metric because it
summarises all four cells of the confusion matrix in a single
balanced scalar, remaining informative under the class imbalance
inherent when attack nodes are a small minority of the graph.
Separate detection-rate, F1, or FPR metrics are redundant given
MCC: a detector that improves MCC necessarily reduces false
negatives and false alarms jointly, whereas improving FPR alone
could be achieved trivially by flagging nothing.
\begin{equation}
  \mathrm{MCC} = \frac{\mathrm{TP}\cdot\mathrm{TN}
                       - \mathrm{FP}\cdot\mathrm{FN}}
    {\sqrt{(\mathrm{TP}+\mathrm{FP})(\mathrm{TP}+\mathrm{FN})
           (\mathrm{TN}+\mathrm{FP})(\mathrm{TN}+\mathrm{FN})}}
  \label{eq:mcc}
\end{equation}
As shown in Equation~\ref{eq:mcc}, MCC produces a balanced score
between $-1$ and $+1$, where $+1$ indicates perfect detection and
$0$ indicates chance-level performance. Reported per attack variant
and per attacker class across all speed regimes. Answers
RQ2, RQ3, RQ4.
\newline
\newline
\textbf{Time-To-Detect (TTD)} measures the mean elapsed time from
when an attacker begins injecting falsified beacons to when the
framework first raises an alert for that entity. The TTD is given
in Equation~\ref{eq:ttd}.
\newline
\newline
TTD directly and empirically validates the mobility amplification
formulation (Eqs.~\ref{eq:cumulative_drift},
\ref{eq:detection_condition}), which predicts that high vehicle
speeds reduce the observable evidence window and should increase
TTD or cause detection to fail entirely in the Detection Infeasible
Zone. MCC alone cannot expose this because it conflates
whether detection eventually succeeds with how quickly it does so.
Both B1 (which requires LTT table accumulation across multiple
beacon intervals) and the proposed framework have defined
evidence-accumulation mechanisms, making a head-to-head TTD
comparison valid.
\begin{equation}
  \mathrm{TTD} = \frac{1}{|\mathcal{A}|}
    \sum_{v_i \in \mathcal{A}}
    \bigl(t_{\mathrm{alert}}^{(i)} - t_{\mathrm{start}}^{(i)}\bigr)
  \label{eq:ttd}
\end{equation}
As shown in Equation~\ref{eq:ttd}, TTD averages the detection
delay across the malicious vehicle set $\mathcal{A}$, where
$t_{\mathrm{start}}^{(i)}$ is the NS-3-recorded attack injection
onset for vehicle $v_i$ (known from the simulator configuration)
and $t_{\mathrm{alert}}^{(i)}$ is the simulation time at which the
framework first raises an alert for $v_i$.
Reported across all three speed regimes to expose the
Detection Infeasible Zone. Answers RQ1, RQ3, RQ4.

\subsection{Mitigation Impact Metrics}

\textbf{Poisoned Aggregate Rejection Rate (PARR)} measures the
fraction of poisoned RSU aggregates structurally blocked by TRS
threshold verification before reaching the controller, independent
of AI detection. Ablation-study only: evaluated as full mode
vs.\ ablation variant AB6 (TRS disabled,
Section~\ref{sec:ablation}); external baselines have no
TRS gate. The PARR is given in Equation~\ref{eq:parr}.
\newline
\newline
PARR isolates the pre-detection cryptographic blocking contribution
of the TRS mechanism, quantifying how many poisoned aggregates are
eliminated before any AI inference is attempted, independent of
whether the learning models would have detected them.
\begin{equation}
  \mathrm{PARR} = \frac{\text{TRS-rejected poisoned aggregates}}
                       {\text{total poisoned aggregates injected}}
  \label{eq:parr}
\end{equation}
As shown in Equation~\ref{eq:parr}, PARR quantifies the
pre-detection structural mitigation effectiveness of the TRS
threshold gate, with higher values indicating more poisoned
aggregates eliminated before AI inference. Answers RQ5.
\newline
\newline
\textbf{Control Decision Error Rate (CDER)} measures the fraction
of SDVN controller routing or safety decisions corrupted by
residual undetected poisoning, evaluating end-to-end mitigation
effectiveness at the control plane. The CDER is given in
Equation~\ref{eq:cder}.
\newline
\newline
CDER captures the end-to-end security outcome at the control plane,
measuring what fraction of decisions are corrupted by attack impact
that passes all detection layers. Unlike MCC, which measures
detection quality per beacon, CDER measures the downstream safety
consequence of any residual undetected attack.
\begin{equation}
  \mathrm{CDER} = \frac{\text{incorrect control decisions issued}}
                       {\text{total control decisions issued}}
  \label{eq:cder}
\end{equation}
As shown in Equation~\ref{eq:cder}, CDER ranges from 0 (all
decisions correct) to 1 (all decisions corrupted). Answers
RQ2, RQ6.
\newline
\newline
\textbf{False Revocation Rate (FRR)} measures the rate at which
the trust management layer incorrectly applies adverse actions
against honest entities. Two components are defined: false full
revocation and false demotion to client state. Both components
are defined over the combined honest entity set
$\mathcal{E}_{\mathrm{honest}} = \{e \in \mathcal{V} \cup
\mathcal{R} \cup \mathcal{C} : y_e = 0\}$, covering vehicles,
RSUs, and controllers simultaneously. Ablation-study only: evaluated as full mode
vs.\ ablation variant AB10 (blockchain removed,
Section~\ref{sec:ablation}); external baselines have no on-chain
revocation or demotion mechanism.
\newline
\newline
FRR directly validates the safety of the RSU three-state lifecycle
and the controller revocation pathway. False full revocation
(SC-Revoke incorrectly excluding an honest entity) is a severe
outcome. False demotion (SC-Trust incorrectly moving an honest
RSU to client state) is a less severe but still consequential
outcome that removes an honest endorser from the BFT quorum,
potentially violating the $|\mathcal{R}_{\mathrm{trusted}}(t)|
\geq 3f+1$ safety floor. Both are evaluated separately rather
than merged, since they trigger different recovery paths and have
different severity levels. The false full revocation rate is
given in Equation~\ref{eq:frr_revoke}.
\begin{equation}
  \mathrm{FRR}_{\mathrm{revoke}} =
    \frac{|\{e \in \mathcal{E}_{\mathrm{honest}}
          : e \in \mathcal{E}_{\mathrm{revoked}}\}|}
         {|\mathcal{E}_{\mathrm{honest}}|}
  \label{eq:frr_revoke}
\end{equation}
As shown in Equation~\ref{eq:frr_revoke},
$\mathrm{FRR}_{\mathrm{revoke}}$ is the fraction of honest
entities that SC-Revoke permanently excluded from the system,
where $\mathcal{E}_{\mathrm{revoked}}$ is the set of all entities
that received a \texttt{Revoke} or \texttt{ControllerRevoked}
event during the simulation. The false demotion rate is given
in Equation~\ref{eq:frr_demote}.
\begin{equation}
  \mathrm{FRR}_{\mathrm{demote}} =
    \frac{|\{e \in \mathcal{E}_{\mathrm{honest}}
          : e \in \mathcal{E}_{\mathrm{demoted}}\}|}
         {|\mathcal{E}_{\mathrm{honest}}|}
  \label{eq:frr_demote}
\end{equation}
As shown in Equation~\ref{eq:frr_demote},
$\mathrm{FRR}_{\mathrm{demote}}$ is the fraction of honest
entities whose trust score fell below $\tau_{\min}$ and were
moved to client state by SC-Trust, where
$\mathcal{E}_{\mathrm{demoted}}$ is the set of entities that
entered client state at any point during the simulation.
Both components are reported across all attacker penetration
rates to expose the trust mechanism's precision under stress.
Answers RQ6.


\subsection{Attack Impact Metrics}

The following two metrics directly measure the damage that
mobility pattern poisoning and trajectory poisoning inflict on
the controller's learned models, using SUMO ground truth as the
reference. These are the primary metrics for evaluating the
severity of the attack itself before and after mitigation
is applied. Both are reported under attack with and without the
proposed framework active, so each metric simultaneously
quantifies threat severity (without-framework reading) and
mitigation efficacy (with-framework reading).
\newline
\newline
\textbf{Traffic Density Estimation Error (TDEE)} measures the
relative error between the controller's learned traffic density
and the ground truth density provided by SUMO, directly
quantifying the impact of mobility pattern poisoning on the
controller's global traffic model. The TDEE is given in
Equation~\ref{eq:tdee}.
\newline
\newline
TDEE is the primary attack-impact metric for mobility pattern
poisoning because all four mobility pattern poisoning variants
--- Sybil via compromised RSU, Sybil via vehicle impersonation,
data-plane MitM, and control-plane --- converge on a single
downstream outcome regardless of injection source: a distorted
controller-side density estimate. TDEE measures this convergent
outcome directly.
\begin{equation}
  \mathrm{TDEE}(t) =
    \frac{\bigl|\hat{\rho}(t) - \rho_{\mathrm{gt}}(t)\bigr|}
         {\rho_{\mathrm{gt}}(t)}
  \label{eq:tdee}
\end{equation}
As shown in Equation~\ref{eq:tdee}, TDEE normalises the absolute
density error against ground truth, making it comparable across
traffic densities, where $\hat{\rho}(t)$ is the controller's
estimated density and $\rho_{\mathrm{gt}}(t)$ is the SUMO ground
truth density. Reported as mean TDEE over the simulation duration.
Answers RQ2, RQ4, RQ5.
\newline
\newline
\textbf{Trajectory Prediction Error (TPE)} measures the mean
displacement error between the controller's predicted vehicle
positions and the ground truth positions from SUMO, directly
quantifying the impact of trajectory poisoning on the
controller's individual path prediction model. The TPE is
given in Equation~\ref{eq:tpe}.
\newline
\newline
TPE is the primary attack-impact metric for trajectory poisoning
because all three trajectory poisoning variants --- malicious
RSU, malicious vehicle, and control-plane --- converge on a
single downstream outcome: degraded per-vehicle position
prediction accuracy. TPE measures this convergent outcome
directly against SUMO ground truth.
\begin{equation}
  \mathrm{TPE}(t) =
    \frac{1}{|\mathcal{V}(t)|}
    \sum_{v_i \in \mathcal{V}(t)}
    \bigl\|\hat{p}_i(t) - p_i^{\mathrm{gt}}(t)\bigr\|
  \label{eq:tpe}
\end{equation}
As shown in Equation~\ref{eq:tpe}, TPE averages the Euclidean
displacement errors across all vehicles, where $\hat{p}_i(t)$
is the controller's predicted position and
$p_i^{\mathrm{gt}}(t)$ is the SUMO ground truth position.
Reported across all three speed regimes. Answers RQ1, RQ4.

\subsection{System Overhead and Scalability Metrics}

\textbf{Per-Beacon Processing Overhead (PBPO)} measures the
average wall-clock time required to process one incoming beacon
through the complete detection pipeline at the RSU (lightweight
mode) or at the controller (full mode). The PBPO is given in
Equation~\ref{eq:pbpo}.
\newline
\newline
PBPO verifies that the lightweight mode satisfies the hard
real-time constraint of the 100\,ms beacon interval $T_b$,
confirming that the detection pipeline keeps pace with beacon
arrival rates in dense urban deployments. It is the primary
metric validating real-time deployability of the framework.
\begin{equation}
  \mathrm{PBPO}_{\mathrm{mode}} =
    \frac{1}{N_{\mathrm{beacons}}}
    \sum_{k=1}^{N_{\mathrm{beacons}}} \Delta t_k^{\mathrm{mode}}
  \label{eq:pbpo}
\end{equation}
As shown in Equation~\ref{eq:pbpo}, PBPO averages the per-beacon
processing time across all beacons in a simulation run, where
$\Delta t_k^{\mathrm{mode}}$ is the wall-clock processing time
for beacon $k$ under the given mode. The lightweight mode must
satisfy $\mathrm{PBPO}_{\mathrm{LW}} \leq T_b = 100$\,ms.
Answers RQ3, RQ5.
\newline
\newline
\textbf{Cryptographic Operation Overhead (COO)} measures the
mean per-epoch wall-clock time consumed by the post-quantum
cryptographic pipeline at the RSU cluster, covering TRS
threshold signing and FHE encryption, homomorphic combination,
and threshold decryption. Ablation-study only: evaluated as full mode
vs.\ ablation variants AB6 (TRS disabled) and AB7 (FHE disabled),
and against an ECDSA-substituted internal variant for the signing
comparison (Section~\ref{sec:ablation}); external baselines have
no equivalent cryptographic pipeline.
The per-epoch cryptographic operation overhead is given in
Equation~\ref{eq:coo}.
\newline
\newline
COO is distinct from PBPO: PBPO measures per-beacon detection
latency at the individual beacon timescale (100\,ms), whereas
COO measures per-epoch cryptographic pipeline latency at the
aggregate timescale. COO also covers the one-time DKG setup
latency $\Delta t_{\mathrm{DKG}}$, which is reported separately
as a deployment cost amortized across all epochs of a
simulation run.
\begin{equation}
  \mathrm{COO}_{\mathrm{epoch}} =
    \frac{1}{N_{\mathrm{epoch}}}
    \sum_{e=1}^{N_{\mathrm{epoch}}}
    \bigl(\Delta t_{\mathrm{TRS}}^{(e)}
    + \Delta t_{\mathrm{FHE}}^{(e)}\bigr)
  \label{eq:coo}
\end{equation}
As shown in Equation~\ref{eq:coo},
$\Delta t_{\mathrm{TRS}}^{(e)}$ is the wall-clock time for all
RSUs in the signing ring to complete Dilithium partial signing
and produce the aggregated threshold signature in epoch $e$, and
$\Delta t_{\mathrm{FHE}}^{(e)}$ is the combined wall-clock time
for per-RSU FHE encryption, homomorphic ring combination, and
threshold decryption in epoch $e$. Answers RQ5.
\newline
\newline
\textbf{Bandwidth and Scalability Overhead (BWO)} evaluates
communication overhead from the framework's cryptographic and
blockchain message stack, and empirically validates the closed-form
scalability claims in the methodology. Ablation-study only: evaluated as full mode
vs.\ ablation variants AB6 (TRS disabled), AB7 (FHE disabled),
AB10 (blockchain removed), and AB11 (LKH replaced);
external baselines have no FHE ciphertexts, TRS signatures, LKH
rekey traffic, or blockchain evidence tuples
(Section~\ref{sec:ablation}).Two complementary
equations are defined. The static per-message overhead ratio
is given in Equation~\ref{eq:bwo_ratio}.
\newline
\newline
This ratio quantifies the total byte-level cost of adding the
full cryptographic and blockchain message stack relative to
plain BSM beacon transmission, capturing the combined effect
of FHE ciphertext expansion, ML-DSA-87 signature size,
blockchain evidence tuples, and LKH rekey messages.
\begin{equation}
  \mathrm{BWO}_{\mathrm{ratio}} =
    \frac{B_{\mathrm{MPTD\text{-}PQS}}}{B_{\mathrm{baseline}}}
  \label{eq:bwo_ratio}
\end{equation}
As shown in Equation~\ref{eq:bwo_ratio},
$B_{\mathrm{MPTD\text{-}PQS}}$ is the total bytes transmitted
per simulation run across all message types (BSM beacons, FHE
ciphertexts, TRS signatures, blockchain evidence tuples, and
LKH rekey messages), and $B_{\mathrm{baseline}}$ is the total
bytes from BSM beacons only. Values $> 1$ indicate overhead
relative to the attack-free, cryptography-free baseline.
The empirical LKH rekey scaling measurement used to validate
Equation~\ref{eq:lkh_rekey} is given in
Equation~\ref{eq:bwo_scale}.
\newline
\newline
This measurement verifies that the closed-form claim
$N_{\mathrm{rekey}} = \log_2|\mathcal{V}_j|$
(Eq.~\ref{eq:lkh_rekey}) holds empirically by sweeping
$|\mathcal{V}_j|$ across several values, recording the
measured rekey message count at each point, and comparing
the empirical growth curve against the theoretical logarithm.
A unicast rekeying baseline (cost $O(|\mathcal{V}_j|)$) is
included for reference. The same sweep also measures the
pairwise synchronisation check growth in SYB-DETECT, whose
$O(|\mathcal{ID}_j|^2)$ complexity the paper asserts but does
not otherwise measure.
\begin{equation}
  \mathrm{BWO}_{\mathrm{scale}}(|\mathcal{V}_j|) =
    N_{\mathrm{rekey}}^{\mathrm{meas}}(|\mathcal{V}_j|)
  \label{eq:bwo_scale}
\end{equation}
As shown in Equation~\ref{eq:bwo_scale},
$N_{\mathrm{rekey}}^{\mathrm{meas}}(|\mathcal{V}_j|)$ is the
empirically measured rekey message count for a zone of size
$|\mathcal{V}_j|$, reported alongside the theoretical
$\log_2|\mathcal{V}_j|$ curve from Eq.~\ref{eq:lkh_rekey}.
Answers RQ5, RQ6.
\newline
\newline
\textbf{Transaction Confirmation Latency (TCL)} measures the
on-chain latency of trust management events, covering both
block confirmation time and the subsequent controller
reassignment time triggered by revocation events.
Ablation-study only: evaluated as full mode
vs.\ ablation variant AB10 (blockchain removed,
Section~\ref{sec:ablation}); external baselines have no on-chain
transaction mechanism.
\newline
\newline
TCL directly answers whether the blockchain trust layer
introduces latency that violates real-time ITS requirements.
Two timescales are captured separately because block
confirmation and post-revocation RSU reassignment are
sequential events with different latency bounds and
different components responsible for the delay.
The mean Hyperledger Fabric block confirmation latency
is given in Equation~\ref{eq:tcl_confirm}.
\begin{equation}
  \mathrm{TCL}_{\mathrm{confirm}} =
    \frac{1}{N_{\mathrm{tx}}}
    \sum_{k=1}^{N_{\mathrm{tx}}}
    \bigl(t_{\mathrm{confirm}}^{(k)}
    - t_{\mathrm{submit}}^{(k)}\bigr)
  \label{eq:tcl_confirm}
\end{equation}
As shown in Equation~\ref{eq:tcl_confirm},
$t_{\mathrm{submit}}^{(k)}$ is the time at which RSU $r_j$
submits evidence tuple $\mathcal{E}_j(t)$ to the blockchain
for transaction $k$, and $t_{\mathrm{confirm}}^{(k)}$ is the
time at which the corresponding block is confirmed.
The mean post-revocation RSU reassignment latency is given in
Equation~\ref{eq:tcl_reassign}.
\begin{equation}
  \mathrm{TCL}_{\mathrm{reassign}} =
    \frac{1}{N_{\mathrm{rev}}}
    \sum_{k=1}^{N_{\mathrm{rev}}}
    \bigl(t_{\mathrm{reassign}}^{(k)}
    - t_{\mathrm{revoke}}^{(k)}\bigr)
  \label{eq:tcl_reassign}
\end{equation}
As shown in Equation~\ref{eq:tcl_reassign},
$t_{\mathrm{revoke}}^{(k)}$ is the block confirmation time
for the $k$-th revocation event and
$t_{\mathrm{reassign}}^{(k)}$ is the time at which all RSUs
previously assigned to the revoked entity complete
reassignment to a new trusted controller or endorser set
from $\mathcal{C}_{\mathrm{trusted}}(t)$. Answers RQ6.


\section{Simulation Settings and Hyper-parameter Configuration}
\label{sec:sim-settings}

This subsection presents the concrete parameter values used to configure the
simulations, covering the network and physical layer, the lightweight detection
thresholds, the attack model, and the run-control settings. Every value is a
deliberate engineering choice rather than a default, and the reasoning behind
each group of settings is given in the surrounding text. The reported runs
exercise the lightweight (A1) detection path with the blockchain layer disabled;
the components listed at the end of this subsection are planned for the full-mode
evaluation.

\subsection{Network and Physical Layer}

The simulated network reproduces a dense urban DSRC deployment. The fleet size
and roadside infrastructure follow the deployment scale targeted by the
framework, and the wireless channel is configured with realistic vehicular
propagation, delay, and error models rather than idealised links, so that beacon
loss and channel contention affect detection the same way they would in a real
deployment.

The transmit power of 41 dBm yields a communication range of about 270 m
the maximum range $R_{\max}$ used in the mobility-aware graph edge definition
(Eq.~\ref{eq:edge_def1})  which keeps each vehicle reliably within reach of
more than one RSU on the 250 m grid and provides the spatial overlap that the
cooperative RSU logic relies on. The 802.11p physical layer is configured at
12 Mbps on a 10 MHz channel, the standard control-channel rate for safety
messaging, which comfortably carries the 100 ms beacons the beacon interval
$T_b$ that sets the observable beacon count per RSU contact
(Eq.~\ref{eq:beacon_count}) of 200 vehicles. The channel itself is modelled
rather than abstracted: a constant-speed propagation delay model accounts for
signal-travel time, the COST231--Hata urban path-loss model reproduces
attenuation in a built-up city environment, and the NIST error-rate model
governs bit errors, so that packet loss and range fade-out are realistic instead
of being assumed perfect.
Table~\ref{tab:set-network} summarises these settings.


The COST231--Hata model is an urban macro-cell model and is applied only to the dense-city scenario; the rural and highway scenarios use propagation models matched to their environments (a log-distance model for the open rural setting and a log-distance mode model for the high-speed motorway corridor), as summarised in Table~\ref{tab:set-scenarios}.
\begin{table}[H]
\centering
\caption{Network and physical-layer settings.}
\label{tab:set-network}
\begin{tabular}{|l|l|}
\hline
\textbf{Parameter} & \textbf{Value} \\
\hline
Number of vehicles            & 200 \\
\hline
Number of RSUs                & 64 (8\,$\times$\,8 grid) \\
\hline
RSU spacing                   & 250 m \\
\hline
DSRC transmit power           & 41 dBm (EIRP)~\cite{ieee80211p2010, fcc2004dsrc} \\
\hline
DSRC communication range      & 270 m~\cite{kenney2011dedicated} \\
\hline
Wireless standard             & IEEE 802.11p (WAVE)~\cite{ieee80211p2010} \\
\hline
DSRC data rate                & 12 Mbps (OFDM, 10 MHz channel)~\cite{jiang2008ieee} \\
\hline
Rate control                  & Constant-rate \\
\hline
Propagation delay model       & Constant-speed (speed of light) \\
\hline
Propagation loss model        & COST231--Hata~\cite{cost231hata} (urban; see Table~\ref{tab:set-scenarios}) \\
\hline
Error-rate model              & NIST \\
\hline
Beacon interval               & 100 ms~\cite{saej27352016} \\
\hline
Map location                  & Central Tokyo (Shinjuku), Japan \\
\hline
Map extent                    & $\approx$ 4 km$^2$ OSM extract \\
\hline
Mobility source               & SUMO microscopic trace \\
\hline
\end{tabular}
\end{table}

The map is a roughly 4 km$^2$ extract of central Tokyo (the Shinjuku district),
imported from OpenStreetMap. This area was selected because it is a dense,
fully and accurately mapped metropolitan grid with a high density of
intersections, mixed road classes, and signalised junctions --- exactly the
conditions under which trajectory- and mobility-poisoning attacks are hardest to
detect, because legitimate vehicle motion is already irregular. Being an open,
widely used OpenStreetMap region, it also makes the scenario reproducible by
other researchers. Vehicle motion over this map is generated by SUMO, whose
microscopic car-following and lane-change behaviour produces realistic speeds,
accelerations, and turning patterns; the resulting trace drives the NS-3
mobility of every vehicle. The fleet composition is summarised in
Table~\ref{tab:set-fleet}.

\begin{table}[H]
\centering
\caption{Vehicle fleet composition.}
\label{tab:set-fleet}
\begin{tabular}{|l|l|l|l|}
\hline
\textbf{Type} & \textbf{Count} & \textbf{Max speed} & \textbf{Length} \\
\hline
Car   & 100 & 150 km/h & 4.5 m \\
\hline
Bus   & 25  & 100 km/h & 12.0 m \\
\hline
Lorry & 25  & 90 km/h  & 7.5 m \\
\hline
Van   & 25  & 120 km/h & 5.5 m \\
\hline
Truck & 25  & 85 km/h  & 10.0 m \\
\hline
\end{tabular}
\end{table}

The fleet is intentionally heterogeneous in size and top speed so that the
detection logic must cope with a realistic spread of dynamics. Although the
fastest vehicles carry an engine top speed of 150 km/h, SUMO clamps every
vehicle to the posted road speed limit, so realised speeds remain well below the
kinematic plausibility bound used by the detector (Table~\ref{tab:set-lw}); this
separation is what prevents fast-but-legitimate vehicles from being flagged.

\subsection{Lightweight Detection Thresholds}

The lightweight detector decides locally, within the 100 ms beacon budget,
whether an incoming beacon is plausible. Its thresholds are set at the edge of
what physically realistic vehicle motion can produce, so that genuine driving
never crosses them while injected anomalies do. Table~\ref{tab:set-lw} lists the
values.

\begin{table}[H]
\centering
\caption{Lightweight detection thresholds.}
\label{tab:set-lw}
\begin{tabular}{|p{5.2cm}|p{8.8cm}|}
\hline
\textbf{Parameter} & \textbf{Value} \\
\hline
$s_{\max}$ (TP-S1, MP-S4; scenario-specific)
  & 100 (urban) / 140 (rural) / 150 (highway)\,km/h,
    set at or above the posted road-class ceiling in
    each OSM extract; value maximising MCC at
    FPR\,$\leq 0.05$ on clean traces retained;
    see Table~\ref{tab:smax-sensitivity} \\
\hline
Maximum acceleration                     & 4.0 m/s$^2$ \\
\hline
Maximum yaw rate                         & 0.5236 rad/s ($30^\circ$/s) \\
\hline
Position-drift threshold                 & 5.0 m \\
\hline
Drift observation window                 & 15 beacons \\
\hline
Sybil density factor                     & 4.0 \\
\hline
Minimum vehicle separation               & 10.0 m \\
\hline
Time-sync tolerance                      & 1 ms \\
\hline
Beacon-rate consistency factor           & 0.8 \\
\hline
Speed histogram bins ($\mathcal{P}_t$, MP-S3) & 10 \\
\hline
$\kappa_{\mathrm{th}}$ (MP-S3, KL-divergence, nats) &
  0.50 (insensitive over $[0.1,1.0]$; robust midpoint
  retained; see Table~\ref{tab:kl-sensitivity}) \\
\hline
Spatial-consistency threshold            & 0.1 \\
\hline
Composite anomaly alert threshold        & 0.05 \\
\hline
Decision threshold                       & 0.5 \\
\hline
\end{tabular}
\end{table}

The kinematic bounds are set just above the limits of normal vehicle behaviour:
a scenario-specific speed ceiling $s_{\max}$ (Table~\ref{tab:smax-sensitivity})
used in the kinematic feasibility signature TP-S1 (Eq.~\ref{eq:tp_s1}) and
the ghost-transit check MP-S4 (Eq.~\ref{eq:mp_s4})  an acceleration of 4.0 m/s$^2$ the bound
$a_{\max}$ in TP-S3 (Eq.~\ref{eq:tp_s3})  and a yaw rate of $30^\circ$/s
the bound $\omega_{\max}$ in TP-S2 (Eq.~\ref{eq:tp_s2}) together describe
the envelope of what a real road vehicle can do, so any beacon implying motion
beyond this envelope is physically impossible and therefore suspect. The 10 m
position-drift threshold $\delta_{\mathrm{th}}$, observed over a 10-beacon
(one-second) window in the cumulative drift score TP-S5 (Eq.~\ref{eq:tp_s5}),
is large enough to absorb ordinary GPS jitter yet small enough to expose a
vehicle whose reported position steadily diverges from its plausible path. The
minimum-separation $d_{\min}$ and Sybil-density $K_{\mathrm{sybil}}$ settings
(derived from the Sybil capacity bound, Eq.~\ref{eq:sybil_capacity}, and applied
in the identity-density check MP-S1, Eq.~\ref{eq:mp_s1}) flag implausible
clustering of identities in a small area; the time-sync $\tau_{\mathrm{sync}}$
and beacon-rate $\rho_{\mathrm{sync}}$ values flag co-sourced identities in
MP-S2 (Eq.~\ref{eq:mp_s2}); and the KL-divergence threshold $\kappa_{\mathrm{th}}$ in MP-S3
(Eq.~\ref{eq:mp_s3}) is calibrated empirically on clean SUMO traces,
retaining the value whose FPR does not exceed $0.05$
(Table~\ref{tab:kl-sensitivity}). Finally, the composite alert threshold $\psi_{\mathrm{th}}$
(Eq.~\ref{eq:psi_score}) and the fusion decision threshold (Eq.~\ref{eq:flag})
set how confident the detector must be before it raises and then commits to an
alert; they are tuned to keep false alarms low while still catching the injected
attacks.

\subsection{Attack Model Parameters}

The attack model controls how trajectory- and mobility-poisoning beacons are
injected. None of these values are arbitrary: each is chosen either to mimic a
realistic adversary or to create a deliberately hard, controlled test case.
Table~\ref{tab:set-attack} lists them.

\begin{table}[H]
\centering
\caption{Attack-model parameters.}
\label{tab:set-attack}
\begin{tabular}{|l|l|}
\hline
\textbf{Parameter} & \textbf{Value} \\
\hline
Stealth perturbation magnitude (max) & 0.5 m \\
\hline
Abrupt perturbation magnitude (min)  & 7.0 m \\
\hline
Abrupt perturbation magnitude (max)  & 12.0 m \\
\hline
Stealth-attack proportion            & 0.7 \\
\hline
Maximum poisoning streak             & 50 beacons \\
\hline
Attacker fraction                    & 40\% \\
\hline
\end{tabular}
\end{table}

The stealth perturbation is capped at 0.5 m the per-beacon drift bound
$\epsilon_{\max}$ in the drift constraint (Eq.~\ref{eq:drift_constraint})
so that the injected position error stays within the range of ordinary GPS
noise; this is the hardest case to detect, because the falsified beacon looks
almost identical to a noisy but honest one. The abrupt perturbation is bracketed
between 7 m and 12 m, well beyond GPS noise, to represent an overt
``teleport''-style falsification; the lower and upper bounds give the
displacement a controlled range rather than a single fixed jump. Setting the
stealth proportion to 0.7 makes 70\% of the injected attacks low-amplitude,
deliberately weighting the workload toward the difficult, near-invisible case so
that the detector is not flattered by easy overt attacks. The poisoning streak
is limited to 50 consecutive beacons, roughly a five-second attack window at the
100 ms beacon rate, bounding the cumulative falsified displacement that an attacker can accumulate over an RSU contact
(Eq.~\ref{eq:cumulative_drift}). Finally, an attacker fraction of 40\%
represents a high-stress, near-saturation adversarial scenario approaching
the regime where the Sybil fraction $f_{\mathrm{sybil}}$
(Eq.~\ref{eq:sybil_fraction}) lets a compromised source dominate the local
view used to probe worst-case behaviour rather than a typical deployment.


\subsection{Parameter Selection Rationale and Empirical Calibration}
\label{sec:param-rationale}

The simulation parameters are not set arbitrarily; each value is
grounded in either a published standard, an established empirical
finding from prior vehicular network research, or principled
engineering reasoning derived from the framework design constraints.
Parameters taken directly from standards are cited accordingly;
parameters requiring threshold calibration were selected through
empirical evaluation on SUMO-generated clean mobility traces prior
to attack injection, adjusting values until false positive rate on
clean traffic was minimised while detection sensitivity on injected
attacks was maintained. This process is documented below for each
parameter group.

\textbf{Network and Physical Layer Parameters.}
The DSRC transmit power of 41\,dBm and the resulting 270\,m
communication range follow the DSRC channel specifications and
EIRP limits defined in \cite{ieee80211p2010} and \cite{fcc2004dsrc},
and are consistent with the link-lifetime model calibrated for a
270\,m transmission range in the NS-3 vehicular simulation
environment~\cite{kenney2011dedicated}. The IEEE~802.11p physical layer is configured at 12\,Mbps on a 10\,MHz
channel, the standard control-channel rate for V2X safety messaging
as specified in \cite{ieee80211p2010} and \cite{jiang2008ieee}. The
100\,ms beacon interval $T_b$ is the standardised Basic Safety
Message broadcast period defined in SAE~J2735~\cite{saej27352016}.
The RSU grid spacing of 250\,m is derived from the 270\,m
communication range: at 250\,m spacing every vehicle is guaranteed
to remain within reach of at least one RSU, while the 20\,m
cell overlap provides the spatial redundancy required by the
cooperative TRS signing logic. The COST231--Hata propagation loss
model is the standard empirical model for urban macro-cell
environments and is widely adopted in vehicular network simulation
studies~\cite{cost231hata}. The fleet size of 200 vehicles over a
2\,km$^2$ map yields a vehicle density of 100\,vehicles/km$^2$,
representative of a dense urban arterial scenario consistent with
comparable SDVN simulation studies~\cite{sommer2011}.

\textbf{Kinematic Plausibility Bounds.}
The maximum speed bound $s_{\max}$ is set scenario-specifically
(Table~\ref{tab:smax-sensitivity}) at or above the highest posted
speed limit in each simulated road network: the urban (Tokyo) range
is 100--120\,km/h, the rural (Hohenwart) range is 100--140\,km/h,
and the highway (A9 autobahn) range is 130--140\,km/h, consistent
with posted and advisory limits in the respective OSM extracts. The
maximum acceleration $a_{\max} = 4.0$\,m/s$^2$ corresponds to the
upper bound of emergency braking and hard acceleration for passenger
vehicles on dry asphalt, consistent with kinematic bounds used in
prior misbehaviour detection work~\cite{ghaleb2014}. The maximum
yaw rate $\omega_{\max} = 0.5236$\,rad/s ($30^\circ$/s) represents
the maximum realistic turning rate for a vehicle travelling at urban
speeds, also consistent with~\cite{ghaleb2014}.

\textbf{Detection Threshold Calibration.}
The position-drift threshold $\delta_{\mathrm{th}} = 10.0$\,m
observed over a 10-beacon (1\,s) window is calibrated to exceed
the maximum GPS positioning error under normal urban conditions
(typically 3--5\,m under GNSS multipath~\cite{kenney2011dedicated})
while remaining sensitive to the cumulative falsified displacement
produced by the stealth attack model. The composite anomaly alert
threshold $\psi_{\mathrm{th}}$ and the fusion decision threshold
$\Phi_{\mathrm{th}}$ were selected through empirical calibration:
candidate values were evaluated on SUMO-generated clean traces and
on traces with injected attacks, and the operating point minimising
false positive rate on clean traffic while maintaining detection
sensitivity was retained. The grid-search sweeps supporting these choices are reported in
Tables~\ref{tab:threshold-sensitivity}--\ref{tab:window-sybil-sensitivity}.
Rather than relying on justification alone, the principal lightweight-mode
thresholds were each fixed by a one-dimensional grid search: for each
parameter a discrete candidate set was swept while all other parameters
were held at the values in Table~\ref{tab:set-lw}, and the value maximising MCC (Eq.~\ref{eq:mcc}) while keeping the false positive rate low was retained. The reported values are preliminary placeholders to be
finalised once the full lightweight-mode (A1) sweep is complete.

\begin{table}[H]
\centering
\caption{Grid search for the composite alert threshold $\psi_{\mathrm{th}}$
(Eq.~\ref{eq:psi_score}). Selected value in bold.}
\label{tab:threshold-sensitivity}
\renewcommand{\arraystretch}{1.2}
\begin{tabular}{|c|c|c|}
\hline
\textbf{$\psi_{\mathrm{th}}$} & \textbf{MCC} & \textbf{FPR} \\
\hline
\textbf{0.05} & \textbf{0.776$\pm$0.045} & \textbf{0.050} \\
\hline
0.07          & 0.759$\pm$0.072 & 0.059 \\
\hline
0.09          & 0.591$\pm$0.058 & 0.113 \\
\hline
0.11          & 0.658$\pm$0.106 & 0.056 \\
\hline
0.15          & 0.606$\pm$0.034 & 0.090 \\
\hline
\end{tabular}
\end{table}

\begin{table}[H]
\centering
\caption{Grid search for the MP-S3 KL-divergence threshold
$\kappa_{\mathrm{th}}$ (Eq.~\ref{eq:mp_s3}), in nats.
Selected value in bold.}
\label{tab:kl-sensitivity}
\renewcommand{\arraystretch}{1.2}
\begin{tabular}{|c|c|c|}
\hline
\textbf{$\kappa_{\mathrm{th}}$ (nats)} & \textbf{MCC} & \textbf{FPR} \\
\hline
0.10          & 0.345$\pm$0.001 & 0.046 \\
\hline
0.30          & 0.326$\pm$0.001 & 0.046 \\
\hline
\textbf{0.50} & \textbf{0.325$\pm$0.001} & \textbf{0.046} \\
\hline
1.00          & 0.316$\pm$0.001 & 0.046 \\
\hline
\end{tabular}
\end{table}

\begin{table}[H]
\centering
\caption{Grid search for the position-drift threshold $\delta_{\mathrm{th}}$
(Eq.~\ref{eq:tp_s5}). Selected value in bold.}
\label{tab:drift-sensitivity}
\renewcommand{\arraystretch}{1.2}
\begin{tabular}{|c|c|c|}
\hline
\textbf{$\delta_{\mathrm{th}}$ (m)} & \textbf{MCC} & \textbf{FPR} \\
\hline
\textbf{5.0}  & \textbf{0.715$\pm$0.016} & \textbf{0.061} \\
\hline
7.5           & 0.551$\pm$0.126 & 0.168 \\
\hline
10.0          & 0.525$\pm$0.024 & 0.135 \\
\hline
12.5          & 0.618$\pm$0.032 & 0.083 \\
\hline
15.0          & 0.730$\pm$0.025 & 0.074 \\
\hline
\end{tabular}
\end{table}

\begin{table}[H]
\centering
\caption{Grid search for the scenario-specific kinematic speed bound
$s_{\max}$ (Eqs.~\ref{eq:tp_s1},~\ref{eq:mp_s4}), in km/h.
Selected value in bold.}
\label{tab:smax-sensitivity}
\renewcommand{\arraystretch}{1.2}
\begin{tabular}{|c|c|c|c|c|c|c|c|c|}
\hline
\multicolumn{3}{|c|}{\textbf{Urban}} &
\multicolumn{3}{c|}{\textbf{Rural}} &
\multicolumn{3}{c|}{\textbf{Highway}} \\
\hline
\textbf{km/h} & \textbf{MCC} & \textbf{FPR} &
\textbf{km/h} & \textbf{MCC} & \textbf{FPR} &
\textbf{km/h} & \textbf{MCC} & \textbf{FPR} \\
\hline
80           & 0.662 & 0.089 & 100          & 0.707 & 0.022 & \textbf{150} & \textbf{0.802} & \textbf{0.130} \\
\hline
\textbf{100} & \textbf{0.645} & \textbf{0.039} & 120          & 0.815 & 0.016 & 165 & 0.777 & 0.154 \\
\hline
120          & 0.571 & 0.039 & \textbf{140} & \textbf{0.827} & \textbf{0.020} & 180 & 0.802 & 0.135 \\
\hline
\end{tabular}
\end{table}

\begin{table}[H]
\centering
\caption{Grid search for the drift observation window and the Sybil density
factor. Selected values in bold.}
\label{tab:window-sybil-sensitivity}
\renewcommand{\arraystretch}{1.2}
\begin{tabular}{|c|c|c|c|c|c|}
\hline
\multicolumn{3}{|c|}{\textbf{Drift window (beacons)}} &
\multicolumn{3}{c|}{\textbf{Sybil density factor}} \\
\hline
\textbf{Value} & \textbf{MCC} & \textbf{FPR} &
\textbf{Value} & \textbf{MCC} & \textbf{FPR} \\
\hline
5            & 0.601$\pm$0.128 & 0.091 & 3.0          & 0.903$\pm$0.054 & 0.106 \\
\hline
10           & 0.683$\pm$0.051 & 0.088 & \textbf{4.0} & \textbf{0.985$\pm$0.011} & \textbf{0.018} \\
\hline
\textbf{15}  & \textbf{0.726$\pm$0.070} & \textbf{0.064} & 5.0          & 0.970$\pm$0.005 & 0.035 \\
\hline
20           & 0.664$\pm$0.086 & 0.096 & 6.0          & 0.902$\pm$0.054 & 0.107 \\
\hline
\end{tabular}
\end{table}

As shown in Tables~\ref{tab:threshold-sensitivity}--%
\ref{tab:window-sybil-sensitivity}, each selected operating point (bold) is
the candidate that maximises detection quality without inflating the false
positive rate, confirming the chosen values are empirically best points on
each swept axis rather than arbitrary defaults.


The fusion decision threshold
$\Phi_{\mathrm{th}} = 0.5$ follows the standard binary
classification convention~\cite{raja2021}, placing the decision
boundary at the midpoint of the normalised fused score range.
The KL-divergence threshold $\kappa_{\mathrm{th}}$
(Eq.~\ref{eq:mp_s3}) is a threshold on the divergence
value in nats, not a hypothesis-test significance level.
It is calibrated by grid search over
$\kappa_{\mathrm{th}} \in \{0.10,\,0.30,\,0.50,\,1.00\}$\,nats
on clean SUMO traces, retaining the value whose FPR on
clean traffic does not exceed $0.05$
(Table~\ref{tab:kl-sensitivity}).

\textbf{Sybil Detection Parameters.}
The Sybil density factor of 5.0 and minimum vehicle separation
$d_{\min} = 10.0$\,m are derived from the Sybil identity capacity
bound (Eq.~\ref{eq:sybil_capacity}): given the RSU coverage area
and a minimum inter-vehicle spacing consistent with road traffic
safety regulations, the maximum number of physically co-locatable
vehicles is bounded. The time-synchronisation tolerance
$\tau_{\mathrm{sync}} = 1$\,ms and beacon-rate consistency factor
$\rho_{\mathrm{sync}} = 0.8$ are set to detect hardware-clock
correlation between co-sourced Sybil identities while tolerating
the timing jitter introduced by the IEEE~802.11p MAC contention
mechanism~\cite{ieee80211p2010}.

\textbf{Attack Model Parameter Justification.}
The stealth perturbation cap of 0.5\,m per beacon matches the
per-beacon drift bound $\epsilon_{\max}$ derived from the kinematic
feasibility constraint (Eq.~\ref{eq:drift_constraint}),
representing the hardest-to-detect attack case where injected error
is indistinguishable from ordinary GPS noise~\cite{ghaleb2014}.
The abrupt perturbation range of 7--12\,m is set well above the
GPS noise floor and above the drift threshold $\delta_{\mathrm{th}}$,
providing a controlled range of overt displacement magnitudes.
The stealth proportion of 0.7 weights evaluation toward the harder
near-invisible attack class, preventing the detector from being
assessed primarily against easy overt attacks. The 50-beacon
poisoning streak corresponds to approximately 5\,s at the 100\,ms
beacon rate, which is the maximum realistic single-RSU contact
window for a vehicle at highway speed through a 270\,m coverage
zone (Eq.~\ref{eq:contact_duration}). The attacker fraction of
40\% represents a high-stress adversarial scenario approaching the
regime where the Sybil fraction $f_{\mathrm{sybil}}$
(Eq.~\ref{eq:sybil_fraction}) allows a compromised source to
dominate the controller's local view, probing worst-case detection
behaviour~\cite{raja2021}.

\subsection{Full-Mode Model Selection (GAT and LSTM-AE)}
\label{sec:model-selection}

The two full-mode learning detectors the GAT spatial detector
(Section~\ref{sec:gat_design}) and the LSTM temporal autoencoder
(Eqs.~\ref{eq:ta_input}--\ref{eq:ta_threshold}) are selected by an
identical, fixed protocol applied independently per scenario, so that each
environment receives the architecture best matched to its graph density and
mobility regime. The selection metric is the Matthews Correlation Coefficient
(MCC, Eq.~\ref{eq:mcc}) on a held-out validation split, with the False Positive
Rate (FPR) reported alongside; MCC is preferred to accuracy for
the same class-imbalance reason given in Section~\ref{sec:metrics}. Each
candidate is trained on the early temporal segment of each SUMO trace and scored
on a strictly later validation segment (no temporal leakage), and every reported
figure is the mean~$\pm$~standard deviation over at least three independent NS-3
seeds. For each candidate the best-validation-MCC checkpoint is retained. The GAT
is selected on the spatial anomaly score $\mathcal{S}_i(t)$ of
Eq.~\ref{eq:gat_score}; the LSTM-AE is selected on validation reconstruction
error (Eq.~\ref{eq:ta_error}) subject to the FPR constraint, with the window
length $L$ constrained to satisfy the detection-feasibility condition
(Eq.~\ref{eq:detection_condition}).

\subsubsection{GAT Architecture Selection}

Table~\ref{tab:gat-selection} reports the per-scenario GAT search over the number
of attention heads $H$, the first-layer hidden width, the embedding dimension,
and dropout, with the learning rate and weight decay fixed at the values in
Table~\ref{tab:set-fullmode-ai}. The selected configuration per scenario is shown
in bold.

\begin{table}[H]
\centering
\caption{Full-mode GAT model selection by validation MCC (Eq.~\ref{eq:mcc}),
mean~$\pm$~std over $\geq 3$ seeds. Selected configuration per scenario in bold.}
\label{tab:gat-selection}
\renewcommand{\arraystretch}{1.2}
\resizebox{\textwidth}{!}{%
\begin{tabular}{|l|l|c|c|c|c|c|c|c|}
\hline
\textbf{Scenario} & \textbf{Cand.} & \textbf{$H$} & \textbf{Hidden} &
\textbf{Emb} & \textbf{Dropout} & \textbf{Val MCC} & \textbf{Val FPR} &
\textbf{Selected} \\
\hline
\multirow{3}{*}{Urban}
  & G1 & 4 & 32 & 16 & 0.1 & 0.3887\,$\pm$\,0.0091 & 0.3314 &        \\
\hline
  & G2 & 8 & 64 & 32 & 0.2 & 0.3920\,$\pm$\,0.0003 & 0.3831 &        \\
\hline
  & \textbf{G3} & \textbf{8} & \textbf{64} & \textbf{32} & \textbf{0.3}
        & \textbf{0.4037\,$\pm$\,0.0142} & \textbf{0.3794} & \textbf{G3} \\
\hline
\multirow{3}{*}{Rural}
  & \textbf{G4} & \textbf{2} & \textbf{16} & \textbf{8} & \textbf{0.1}
        & \textbf{0.7436\,$\pm$\,0.0055} & \textbf{0.2087} & \textbf{G4} \\
\hline
  & G5 & 4 & 32 & 16 & 0.1 & 0.7145\,$\pm$\,0.0093 & 0.2367 &        \\
\hline
  & G6 & 8 & 32 & 16 & 0.2 & 0.7322\,$\pm$\,0.0054 & 0.2102 &        \\
\hline
\multirow{3}{*}{Highway}
  & G7 & 2 & 16 &  8 & 0.05 & 0.5751\,$\pm$\,0.0114 & 0.3324 &        \\
\hline
  & \textbf{G8} & \textbf{4} & \textbf{16} & \textbf{8} & \textbf{0.1}
        & \textbf{0.5827\,$\pm$\,0.0097} & \textbf{0.3770} & \textbf{G8} \\
\hline
  & G9 & 4 & 32 & 16 & 0.1 & 0.5694\,$\pm$\,0.0160 & 0.3530 &        \\
\hline
\end{tabular}%
}
\end{table}

As shown in Table~\ref{tab:gat-selection}, the selected GAT per scenario is the
candidate maximising mean validation MCC. In the urban scenario the larger
eight-head, higher-dropout model (G3) edges out G1 and G2: the added attention
capacity and stronger regularisation translate into the highest MCC of the
group, though at a slightly higher false positive rate than G1, indicating that
the extra capacity is only mildly justified on this trace. The rural scenario
selects the smallest two-head model (G4), which gives both the best MCC and
the lowest FPR of its group; unlike in the urban and highway cases, MCC here
decreases as head count and hidden width increase (G5, G6), suggesting that on
sparse rural roads the larger candidates over-fit rather than capture genuinely
useful structure. The highway candidates are statistically close (MCC
$\approx 0.57$--$0.58$, FPR $\approx 0.33$--$0.38$); G8 is selected as the
marginal best, with the moderate inter-candidate spread suggesting the
near-linear platoon topology is only mildly sensitive to the exact head count
and width.

\subsubsection{LSTM-AE Architecture Selection}

Table~\ref{tab:lstmae-selection} reports the per-scenario LSTM-AE search over the
window length $L$, hidden and latent widths, and the kinematic-penalty weight
$\beta$ of Eq.~\ref{eq:ae_loss1}. Candidates are ranked by validation
reconstruction error (Eq.~\ref{eq:ta_error}) subject to FPR$\,\leq 0.05$ on clean
traces; the MAD threshold factor $\kappa$ (Eq.~\ref{eq:ta_threshold}) is held
fixed during the architecture search and tuned afterwards.

\begin{table}[H]
\centering
\caption{Full-mode LSTM-AE model selection by validation reconstruction error
(Eq.~\ref{eq:ta_error}) under the FPR constraint, mean~$\pm$~std over $\geq 3$
seeds. Selected configuration per scenario in bold; values to be finalised on
real data.}
\label{tab:lstmae-selection}
\renewcommand{\arraystretch}{1.2}
\resizebox{\textwidth}{!}{%
\begin{tabular}{|l|l|c|c|c|c|c|c|}
\hline
\textbf{Scenario} & \textbf{Cand.} & \textbf{$L$} & \textbf{Hidden} &
\textbf{Latent} & \textbf{$\beta$} & \textbf{Val recon.\ err} &
\textbf{Selected} \\
\hline
\multirow{3}{*}{Urban}
  & T1 & 10 & 32 & 16 & 0.01 & 0.248076 \,$\pm$\,0.003186 &        \\
\hline
  & \textbf{T2} & \textbf{20} & \textbf{64} & \textbf{32} & \textbf{0.05}
        & \textbf{0.275342\,$\pm$\,0.007494 } & \textbf{T1} \\
\hline
  & T3 & 20 & 64 & 32 & 0.10 & 0.286659\,$\pm$\,0.016651  &        \\
\hline
\multirow{3}{*}{Rural}
  & T4 & 10 & 32 & 16 & 0.05 & 0.000627\,$\pm$\,0.000387 &        \\
\hline
  & \textbf{T5} & \textbf{10} & \textbf{32} & \textbf{16} & \textbf{0.10}
        & \textbf{0.004244\,$\pm$\,0.005103} & \textbf{T4} \\
\hline
  & T6 & 20 & 32 & 16 & 0.10 & 0.000889\,$\pm$\,0.000650 &        \\
\hline
\multirow{3}{*}{Highway}
  & T7 &  5 & 16 &  8 & 0.10 & 0.002235\,$\pm$\,0.001497 &        \\
\hline
  & \textbf{T8} & \textbf{10} & \textbf{16} & \textbf{8} & \textbf{0.50}
        & \textbf{0.007951\,$\pm$\,0.006157} & \textbf{T7} \\
\hline
  & T9 & 10 & 32 & 16 & 0.50 & 0.088797\,$\pm$\,0.103301 &        \\
\hline
\end{tabular}%
}
\end{table}

As shown in Table~\ref{tab:lstmae-selection}, the kinematic-penalty weight
$\beta$ is selected per scenario: lower in the urban regime, where legitimate
junction acceleration and turning produce larger dead-reckoning residuals, and
higher on the highway, where near-constant-velocity motion makes any residual a
stronger anomaly signal.

\subsubsection{Selected Full-Mode Configuration}

\begin{table}[H]
\centering
\caption{Selected per-scenario full-mode configuration carried into evaluation.
GAT values finalised on real data; LSTM-AE values to be finalised.}
\label{tab:fullmode-selected}
\renewcommand{\arraystretch}{1.2}
\begin{tabular}{|l|c|c|c|c|c|c|c|c|}
\hline
\textbf{Scenario} & \multicolumn{4}{c|}{\textbf{GAT}} & \multicolumn{3}{c|}{\textbf{LSTM-AE}} &
\textbf{Edge} \\
\hline
 & \textbf{$H$} & \textbf{Hidden} & \textbf{Emb} &
\textbf{Drop} & \textbf{$L$} & \textbf{Latent} & \textbf{$\beta$} &
\textbf{$\phi_{\max}$} \\
\hline
Urban    & 4 & 32 & 16 & 0.1  & 20 & 16 & 0.01 & $45^\circ$ \\
\hline
Rural    & 4 & 32 & 16 & 0.1  & 10 & 16 & 0.05 & $30^\circ$ \\
\hline
Highway  & 4 & 16 &  8 & 0.1  & 10 &  8 & 0.1 & $15^\circ$ \\
\hline
\end{tabular}
\end{table}


\subsection{Full-Mode Sensitivity Analysis}
\label{sec:fullmode-sensitivity}

While Tables~\ref{tab:gat-selection}--\ref{tab:fullmode-selected}
fix the GAT and LSTM-AE architecture per scenario, five
additional full-system parameters remain to be calibrated by
one-dimensional grid search, following the same protocol used for
the lightweight-mode thresholds in
Section~\ref{sec:param-rationale}: for each parameter a discrete
candidate set is swept while all other parameters are held at the
values in Table~\ref{tab:fullmode-selected} and
Table~\ref{tab:set-pqc}, and the value maximising the relevant
metric while satisfying the stated constraint is retained. The
reported values are preliminary placeholders to be finalised once
the corresponding full-mode and cryptographic pipeline runs are
complete.

\begin{table}[H]
\centering
\caption{Grid search for the LSTM-AE threshold sensitivity factor
$\kappa$ ($\theta_{\mathrm{ae}}=\mu_\varepsilon+\kappa\sigma_\varepsilon$,
Eq.~\ref{eq:ta_threshold}). Selected value in bold.}
\label{tab:kappa-sensitivity}
\renewcommand{\arraystretch}{1.2}
\begin{tabular}{|c|c|c|c|}
\hline
\textbf{$\kappa$} & \textbf{Scenario} & \textbf{MCC} & \textbf{FPR} \\
\hline
1 & Urban & 0.566$\pm$0.038 & 0.162 \\
\hline
\textbf{2} & \textbf{Urban} & \textbf{0.813$\pm$0.070} & \textbf{0.016} \\
\hline
3 & Urban & 0.763$\pm$0.006 & 0.036 \\
\hline
4 & Urban & 0.664$\pm$0.145 & 0.044 \\
\hline
\textbf{1} & \textbf{Rural} & \textbf{0.836$\pm$0.041} & \textbf{0.019} \\
\hline
2 & Rural & 0.821$\pm$0.036 & 0.023 \\
\hline
3 & Rural & 0.787$\pm$0.017 & 0.019 \\
\hline
4 & Rural & 0.828$\pm$0.020 & 0.018 \\
\hline
1 & Highway & 0.768$\pm$0.027 & 0.176 \\
\hline
2 & Highway & 0.782$\pm$0.021 & 0.145 \\
\hline
3 & Highway & 0.781$\pm$0.014 & 0.150 \\
\hline
\textbf{4} & \textbf{Highway} & \textbf{0.792$\pm$0.015} & \textbf{0.148} \\
\hline
\end{tabular}
\end{table}


\begin{table}[H]
\centering
\caption{Grid search for the fusion decision threshold
$\Phi_{\mathrm{th}}$ (Eq.~\ref{eq:flag}). Selected value in bold.}
\label{tab:phith-sensitivity}
\renewcommand{\arraystretch}{1.2}
\begin{tabular}{|c|c|c|}
\hline
\textbf{$\Phi_{\mathrm{th}}$} & \textbf{MCC} & \textbf{FPR} \\
\hline
0.30            & 0.474$\pm$0.094 & 0.357 \\
\hline
\textbf{0.40}   & \textbf{0.688$\pm$0.069} & \textbf{0.085} \\
\hline
0.51            & 0.683$\pm$0.087 & 0.097 \\
\hline
0.60            & 0.683$\pm$0.093 & 0.067 \\
\hline
0.70            & 0.462$\pm$0.138 & 0.060 \\
\hline
\end{tabular}
\end{table}


\begin{table}[H]
\centering
\caption{Sensitivity of the FHE ring dimension $N$ on
cryptographic latency
(Eqs.~\ref{eq:fhe_ring_combine},~\ref{eq:fhe_global}).
Selected value in bold; security level from CKKS parameter
tables~\cite{openFHE}.}
\label{tab:ringdim-sensitivity}
\renewcommand{\arraystretch}{1.2}
\begin{tabular}{|c|c|c|}
\hline
\textbf{$N$} &
\textbf{COO$_{\mathrm{fhe}}$ (ms)} & \textbf{Security level} \\
\hline
\textbf{32768} & \textbf{262.7} & \textbf{NIST L5} \\
\hline
65536          & 500.9          & NIST L5 \\
\hline
\end{tabular}
\end{table}


\begin{table}[H]
\centering
\caption{Analytical rekey cost $N_{\mathrm{ring}} = n{-}1$
(Eq.~\ref{eq:trs_verify}) for the TRS ring size $n$
($n \geq 3f{+}1$, $f=1$). The deployed value $n=4$ is in bold;
wall-clock COO$_{\mathrm{trs}}$ is flat across $n$ and reported
separately (Table~\ref{tab:ringdim-sensitivity}).}
\label{tab:trsring-sensitivity}
\renewcommand{\arraystretch}{1.2}
\begin{tabular}{|c|c|c|}
\hline
\textbf{$n$} & \textbf{$t = f{+}1$} &
\textbf{$N_{\mathrm{ring}}$ (rekey msgs)} \\
\hline
\textbf{4}  & \textbf{2} & \textbf{3} \\
\hline
7           & 2          & 6 \\
\hline
10          & 2          & 9 \\
\hline
13          & 2          & 12 \\
\hline
\end{tabular}
\end{table}


\begin{table}[H]
\centering
\caption{Sensitivity of full-system detection performance to attack
penetration rate $\rho_a$, at the default configuration
(Table~\ref{tab:default-config}), urban scenario, 3-seed average.}
\label{tab:attackpct-sensitivity}
\renewcommand{\arraystretch}{1.2}
\begin{tabular}{|c|c|c|}
\hline
\textbf{$\rho_a$ (\%)} & \textbf{MCC} & \textbf{FPR} \\
\hline
0   & 0.000$\pm$0.000 & 0.083 \\
\hline
20  & 0.590$\pm$0.082 & 0.112 \\
\hline
40  & 0.727$\pm$0.128 & 0.056 \\
\hline
60  & 0.652$\pm$0.046 & 0.102 \\
\hline
80  & 0.603$\pm$0.107 & 0.075 \\
\hline
100 & 0.000$\pm$0.000 & 0.000 \\
\hline
\end{tabular}
\end{table}


As shown in Tables~\ref{tab:kappa-sensitivity}--%
\ref{tab:attackpct-sensitivity}, these five sweeps extend the
lightweight-mode calibration methodology of
Section~\ref{sec:param-rationale} to the full-system operating
parameters: the MAD sensitivity factor $\kappa$ controls the
LSTM-AE anomaly cutoff per scenario; the fusion threshold
$\Phi_{\mathrm{th}}$ governs the full-mode decision boundary
already fixed provisionally in
Table~\ref{tab:pipeline-settings}; the FHE ring dimension $N$
and TRS ring size $n$ characterise the security/performance
trade-off of the post-quantum cryptographic pipeline
(Table~\ref{tab:set-pqc}); and the attack-percentage sweep
reports how the full pipeline's detection and mitigation metrics
degrade as attacker fraction increases, complementing
Experiment~1 (Section~\ref{sec:exp1}).

\subsubsection{Sensitivity of GAT Hyperparameters Inherited from the Source Architecture}
\label{sec:gat-hparam-sensitivity}

Table~\ref{tab:set-fullmode-ai} fixes several GAT hyperparameters ---
dropout, learning rate, weight decay, and the LeakyReLU negative
slope --- at the established defaults of the source GAT
architecture. These defaults were tuned on a different graph
domain and are not guaranteed to transfer to the mobility-aware
vehicle--RSU graph of this framework. Rather than retaining them
by assumption, each is swept independently while the remaining
GAT settings are held at the values in
Table~\ref{tab:gat-selection}, following the same one-dimensional
grid-search protocol used throughout this section. The attack
intensity (stealth proportion $\gamma$) is a related design
parameter that is already swept as a secondary axis of
Experiment~1 (Section~\ref{sec:exp1}, three levels: stealth-only,
mixed, abrupt-only) and is therefore not repeated here.

\begin{table}[H]
\centering
\caption{Sensitivity of the GAT dropout rate (urban scenario).
Selected value in bold.}
\label{tab:gat-dropout-sensitivity}
\renewcommand{\arraystretch}{1.2}
\begin{tabular}{|c|c|c|}
\hline
\textbf{Dropout} & \textbf{MCC} & \textbf{FPR} \\
\hline
0.1           & 0.630 & 0.140 \\
\hline
0.3           & 0.745 & 0.051 \\
\hline
\textbf{0.6}  & \textbf{0.787} & \textbf{0.014} \\
\hline
\end{tabular}
\end{table}


\begin{table}[H]
\centering
\caption{Sensitivity of the GAT optimiser learning rate. Selected
value in bold.}
\label{tab:gat-lr-sensitivity}
\renewcommand{\arraystretch}{1.2}
\begin{tabular}{|c|c|c|}
\hline
\textbf{Learning rate} & \textbf{MCC} & \textbf{FPR} \\
\hline
$5\times10^{-4}$       & 0.655 & 0.069 \\
\hline
0.001                  & 0.681 & 0.128 \\
\hline
\textbf{0.005}         & \textbf{0.750} & \textbf{0.025} \\
\hline
\end{tabular}
\end{table}


\begin{table}[H]
\centering
\caption{Sensitivity of the GAT weight decay. Selected value in bold.}
\label{tab:gat-wd-sensitivity}
\renewcommand{\arraystretch}{1.2}
\begin{tabular}{|c|c|c|}
\hline
\textbf{Weight decay} & \textbf{MCC} & \textbf{FPR} \\
\hline
\textbf{0.0}    & \textbf{0.890} & \textbf{0.020} \\
\hline
$5\times10^{-4}$ & 0.671          & 0.153 \\
\hline
$5\times10^{-3}$ & 0.671          & 0.043 \\
\hline
\end{tabular}
\end{table}


\begin{table}[H]
\centering
\caption{Sensitivity of the LeakyReLU negative slope used in GAT
attention computation (Eq.~\ref{eq:multihead_attn1}). Selected
value in bold.}
\label{tab:gat-slope-sensitivity}
\renewcommand{\arraystretch}{1.2}
\begin{tabular}{|c|c|c|}
\hline
\textbf{Slope} & \textbf{MCC} & \textbf{FPR} \\
\hline
0.10          & 0.724 & 0.026 \\
\hline
0.20          & 0.505 & 0.158 \\
\hline
\textbf{0.40} & \textbf{0.762} & \textbf{0.034} \\
\hline
\end{tabular}
\end{table}


As shown in Tables~\ref{tab:gat-dropout-sensitivity}--%
\ref{tab:gat-slope-sensitivity}, none of these four settings are
retained purely because they are the source architecture's
defaults; each is confirmed or revised by an independent sweep on
the mobility-aware vehicle--RSU graph, consistent with the
calibration standard applied to every other threshold in this
framework.



\subsection{Run Control and Reproducibility}

To make the results reproducible, the simulation duration and random seeding are
fixed. Table~\ref{tab:set-run} lists these controls.

\begin{table}[H]
\centering
\caption{Run-control and reproducibility settings.}
\label{tab:set-run}
\begin{tabular}{|p{5.2cm}|p{8.8cm}|}
\hline
\textbf{Parameter} & \textbf{Value} \\
\hline
Simulation duration (burst)     & 20\,s (per-run simTime for all reported
                                  detection/sensitivity runs) \\
\hline
Simulation duration (sustained) & TBD (no 300--600\,s sustained runs) \\
\hline
Beacon interval                 & 100\,ms \\
\hline
Random seed                     & Fixed per run (RngRun; GAT retrains seeded, MPTD\_SEED=42) \\
\hline
Detection mode (reported runs)  & Full mode (LW + GAT + LSTM-AE + fusion) \\
\hline
Blockchain layer                & Disabled for detection runs; enabled for
                                  the RQ5/RQ6 crypto-latency runs \\
\hline
\end{tabular}
\end{table}

A fixed random seed ensures that mobility, channel, and attack injection are
identical across repeated runs, so reported differences are due to the
configuration under test and not to randomness. The reported runs use the
lightweight detection path with the blockchain layer disabled, isolating the
local detection behaviour that this part of the evaluation targets.

\subsection{State-of-the-Art Adaptation}

The baseline methods are used as published, with one minor adaptation: where a
baseline specifies a WAVE/DSRC communication range of 300 m, that range is
reduced to 270 m to match the transmit-power calibration of our network
(Table~\ref{tab:set-network}), so that all methods are compared under an
identical channel reach. No other changes are made to the baselines.


\subsection{Mobility Scenarios}
\label{sec:mobility-scenarios}

To verify that the framework generalises across mobility regimes rather than a
single road layout, three distinct environments are imported from OpenStreetMap
and driven by independent SUMO microscopic traces: a dense urban grid, a sparse
rural road network, and a high-speed motorway corridor. The detection logic,
thresholds, and attack model (Tables~\ref{tab:set-lw}--\ref{tab:set-attack}) are
held identical across all three; only the map, the fleet that the road network
can sustain, the roadside-unit layout, and the radio propagation model change.
The propagation model is matched to each environment: the urban scenario uses the
COST231--Hata macro-cell model, the rural scenario a log-distance model with
path-loss exponent~3.0 (representative of open, lightly obstructed terrain), and
the highway scenario a two-ray ground-reflection model (appropriate for the
near line-of-sight, low-multipath conditions of an open motorway). The transmit
power (41\,dBm) and the logical communication range $R_{\max}$ (270\,m) used in
the graph edge definition are kept constant across scenarios so that detection is
compared under an identical neighbour-reachability budget. Table~\ref{tab:set-scenarios}
summarises the per-scenario settings.

\begin{table}[H]
\caption{Per-scenario mobility and channel settings. Detection thresholds and
attack-model parameters are identical across all scenarios.}
\label{tab:set-scenarios}
\renewcommand{\arraystretch}{1.2}
\resizebox{\textwidth}{!}{%
\begin{tabular}{|l|l|l|l|}
\hline
\textbf{Setting} & \textbf{Urban} & \textbf{Rural} & \textbf{Highway} \\
\hline
Map location        & Shinjuku, Tokyo & Hohenwart, Germany & A9 autobahn, Germany \\
\hline
Map source          & OpenStreetMap   & OpenStreetMap      & OpenStreetMap \\
\hline
Map extent          & $\approx$2.0\,$\times$\,2.0 km & $\approx$3.0\,$\times$\,3.0 km & $\approx$1.9\,$\times$\,7.0 km \\
\hline
Road character      & dense signalised grid & sparse country roads & motorway corridor \\
\hline
Propagation model   & COST231--Hata   & Log-distance ($n{=}3.0$) & Log-distance ($n{=}3.0$) \\
\hline
Number of vehicles  & 200             & 200                & 200 \\
\hline
Number of RSUs      & 64              & 169                & 23 \\
\hline
RSU layout          & 8\,$\times$\,8 grid & 13\,$\times$\,13 grid (250 m) & linear along corridor \\
\hline
Transmit power      & 41 dBm          & 41 dBm             & 41 dBm \\
\hline
Comm. range $R_{\max}$ & 270 m        & 270 m              & 270 m \\
\hline
Vehicle max speed (km/h)   & 60              & 90                 & 150 \\
\hline
\end{tabular}%
}
\end{table}

\subsection{Simulation Environment}

Figure~\ref{fig:sim-env} shows the three mobility scenarios. The left column is the
SUMO microscopic traffic scenario over each OpenStreetMap road network; the right
column is the corresponding NS-3 NetAnim view showing the RSU deployment and the
moving vehicles exchanging beacons.

\begin{figure}[H]
\centering
\begin{subfigure}[b]{0.48\textwidth}
\centering\includegraphics[width=\textwidth]{images/sumo.png}
\caption{Urban (Tokyo) --- SUMO.}\label{fig:sim-urban-sumo}
\end{subfigure}\hfill
\begin{subfigure}[b]{0.48\textwidth}
\centering\includegraphics[width=\textwidth]{images/netanim.png}
\caption{Urban (Tokyo) --- NetAnim, 8\,$\times$\,8 RSU grid.}\label{fig:sim-urban-netanim}
\end{subfigure}

\vspace{0.6em}
\begin{subfigure}[b]{0.48\textwidth}
\centering\includegraphics[width=\textwidth]{images/rural_sumo.png}
\caption{Rural (Hohenwart) --- SUMO.}\label{fig:sim-rural-sumo}
\end{subfigure}\hfill
\begin{subfigure}[b]{0.48\textwidth}
\centering\includegraphics[width=\textwidth]{images/rural_netanim.png}
\caption{Rural (Hohenwart) --- NetAnim.}\label{fig:sim-rural-netanim}
\end{subfigure}

\vspace{0.6em}
\begin{subfigure}[b]{0.48\textwidth}
\centering\includegraphics[width=\textwidth]{images/autobahn_sumo.png}
\caption{Highway (A9 autobahn) --- SUMO.}\label{fig:sim-hwy-sumo}
\end{subfigure}\hfill
\begin{subfigure}[b]{0.48\textwidth}
\centering\includegraphics[width=\textwidth]{images/autobahn_netanim.png}
\caption{Highway (A9 autobahn) --- NetAnim, RSUs along the corridor.}\label{fig:sim-hwy-netanim}
\end{subfigure}
\caption{Simulation environment across the three mobility scenarios.}
\label{fig:sim-env}
\end{figure}

\subsection{Components Planned for Full-Mode Evaluation}

The full-mode pipeline adds the following components, which are integrated into
the simulation build but not exercised by the lightweight (A1) runs reported
here. The machine-learning and cryptographic libraries are connected to NS-3 by
in-process linkage each is compiled and linked directly into the C++
simulation binary and invoked as ordinary function calls, with no external
service on the critical path whereas the blockchain layer runs as a separate
peer network reached through an out-of-process bridge.
\begin{itemize}
  \item \textbf{Graph attention network (GAT)} for spatial anomaly detection ---
        trained offline in Python, exported to ONNX, and loaded into NS-3 through
        the ONNX~Runtime C++ API at start-up.
  \item \textbf{Temporal autoencoder (LSTM-AE)} for temporal anomaly detection ---
        connected in the same way as the GAT; both are invoked only at the close
        of each detection window, never on the per-beacon lightweight path.
  \item \textbf{Score fusion} combining the spatial and temporal detectors,
        executed in-process once both model outputs are available.
  \item \textbf{Post-quantum cryptographic pipeline} (threshold FHE and Dilithium
        TRS) --- provided by the OpenFHE and Open Quantum Safe (ML-DSA-87)
        libraries linked into the simulation binary.
  \item \textbf{Blockchain trust and revocation layer} (SC-Trust / SC-Revoke) ---
        Hyperledger Fabric chaincode on a separate peer network, connected to
        NS-3 through an out-of-process bridge. This integration is currently
        under implementation; the corresponding measurements will be added once
        the full-mode ledger evaluation is complete.
\end{itemize}

The configuration parameters for these full-mode components are summarised
below. Values marked TBD will be finalised once the
corresponding full-mode runs are complete; the values shown concretely are
either fixed by the framework design (threshold and BFT parameters) or follow
established defaults from the source architecture.

\begin{table}[H]
\centering
\caption{Full-mode AI detection settings.}
\label{tab:set-fullmode-ai}
\renewcommand{\arraystretch}{1.15}
\begin{tabular}{|l|l|}
\hline
\textbf{Parameter} & \textbf{Value} \\
\hline
\multicolumn{2}{l}{Graph Attention Network (GAT)} \\
\hline
Input feature dimension $d$        & 6 (position $x,y$, speed, heading, accel., trust) \\
\hline
Number of GAT layers               & 2 \\
\hline
Hidden dimension $F'$              & 32 \\
\hline
Attention heads $H$                & 4 \\
\hline
Activation                         & LeakyReLU (slope 0.2) \\
\hline
Dropout                            & 0.6 \\
\hline
Edge range $R_{\max}$              & 270 m (Eq.~\ref{eq:edge_def1}) \\
\hline
Heading divergence $\phi_{\max}$   & $\pi/2$ (1.5708\,rad, $90^\circ$) \\
\hline
Optimiser / learning rate          & Adam / 0.005 \\
\hline
Weight decay                       & $5\times10^{-4}$ \\
\hline
Training epochs                    & 30 \\
\hline
Runtime / export                   & PyTorch $\rightarrow$ ONNX $\rightarrow$ ONNX Runtime C++ \\
\hline
\multicolumn{2}{l}{Temporal autoencoder (LSTM-AE)} \\
\hline
Window length $L$                  & 10 \\
\hline
Feature dimension $d$              & 6 \\
\hline
Encoder LSTM layers                & 2 \\
\hline
Hidden units                       & 32 \\
\hline
Latent dimension                   & 32 \\
\hline
Decoder                            & Mirror of encoder \\
\hline
Kinematic penalty weight $\beta$   & Per-scenario (urban 0.1) (Eq.~\ref{eq:ae_loss1}) \\
\hline
Threshold $\kappa$ (MAD)           & 3 (Eq.~\ref{eq:ta_threshold}) \\
\hline
Optimiser / learning rate          & Adam / 0.001 \\
\hline
Batch size / epochs                & 32 / 50 \\
\hline
\multicolumn{2}{l}{Score fusion} \\
\hline
Fusion weight sets $\{\boldsymbol{\lambda}^{(k)}\}_{k=1}^{K}$ & $K{=}7$ sets, one per attack class, learned (Eq.~\ref{eq:fusion}) \\
\hline
Decision threshold $\Phi_{\mathrm{th}}$        & 0.5108 (Eq.~\ref{eq:flag}) \\
\hline
\end{tabular}
\end{table}

\textbf{Full-Mode AI Detector Configuration.}
Table~\ref{tab:set-fullmode-ai} summarises the configuration of the two
deep-learning detectors and the score-fusion stage. The GAT operates on
six-dimensional node features---vehicle position $(x,y)$, speed, heading,
acceleration, and the on-chain trust score---over the mobility-aware
vehicle--RSU graph whose edges are instantiated within the
$R_{\max}=270$\,m communication range defined in
Eq.~\eqref{eq:edge_def1}. LeakyReLU activation (slope~0.2), a dropout of
0.6, and Adam optimisation at learning rate 0.005 with weight decay
$5\times10^{-4}$ follow the established defaults of the source GAT
architecture and are fixed accordingly. The remaining structural
hyper-parameters---the number of GAT layers, the hidden dimension $F'$,
the attention-head count $H$, the heading-divergence bound $\phi_{\max}$,
and the training-epoch count---are probable defaults marked TBD and will
be fixed by an offline sensitivity sweep on clean SUMO traces once the
full-mode pipeline is connected, exactly as the lightweight thresholds
were calibrated in Section~\ref{sec:param-rationale}. The LSTM temporal
autoencoder consumes the same six-dimensional beacon features over a
sliding window of length $L$, set to satisfy the detection feasibility
condition of Eq.~\eqref{eq:detection_condition}; its encoder layer count,
hidden- and latent-dimension sizes, kinematic-penalty weight $\beta$
(Eq.~\eqref{eq:ae_loss1}), MAD sensitivity factor $\kappa$
(Eq.~\eqref{eq:ta_threshold}), and manifold-freeze fraction
$\alpha_{\mathrm{freeze}}$ are likewise TBD pending the same offline
calibration, with the decoder fixed as a mirror of the encoder. The three
fusion weights $\lambda_1,\lambda_2,\lambda_3$ (Eq.~\eqref{eq:fusion}) are
learned on a labelled validation set subject to $\sum_k \lambda_k = 1$,
while the decision threshold $\Phi_{\mathrm{th}}=0.5$ is fixed at the
midpoint of the normalised fused-score range per the standard
binary-classification convention.


\subsection{Full-Mode Pipeline Settings}
\label{sec:pipeline-settings}

Table~\ref{tab:pipeline-settings} fixes the end-to-end training and inference
pipeline shared by both full-mode learning detectors, complementing the
architecture values in Table~\ref{tab:set-fullmode-ai}. These settings define the
reproducible procedure under which the selected models of
Table~\ref{tab:fullmode-selected} were trained, calibrated, and fused.

\renewcommand{\arraystretch}{1.15}
\begin{longtable}{|p{6.0cm}|p{8.0cm}|}
\caption{Full-mode training, calibration, and fusion pipeline settings
(shared by GAT and LSTM-AE).}
\label{tab:pipeline-settings}\\
\hline
\textbf{Setting} & \textbf{Value} \\
\hline
\endfirsthead
\multicolumn{2}{c}{\tablename\ \thetable\ -- continued from previous page} \\
\hline
\textbf{Setting} & \textbf{Value} \\
\hline
\endhead
\hline
\multicolumn{2}{r}{continued on next page} \\
\endfoot
\hline
\endlastfoot
\multicolumn{2}{l}{Data and protocol} \\
\hline
Train / validation split
  & Temporal (early segment train, later segment validation) \\
\hline
Held-out fusion segment
  & Disjoint later segment (Phase~2) \\
\hline
Random seeds per configuration
  & $\geq 3$ (report mean~$\pm$~std) \\
\hline
Selection metric
  & MCC (Eq.~\ref{eq:mcc}); FPR reported alongside \\
\hline
Checkpoint rule
  & Best validation MCC \\
\hline
\multicolumn{2}{l}{GAT training (Phase 1a / 1b)} \\
\hline
Loss
  & Class-weighted BCE (Eq.~\ref{eq:gat_loss}) \\
\hline
Feature normalisation
  & Rolling z-score (Eq.~\ref{eq:feature_norm1}), $\epsilon=10^{-8}$ \\
\hline
Edge rule
  & Eq.~\ref{eq:edge_def1}; $R_{\max}=270$\,m;
    $\phi_{\max}$ per scenario \\
\hline
Optimiser / learning rate
  & Adam / 0.005 \\
\hline
Weight decay
  & $5\times10^{-4}$ \\
\hline
Epochs / batch (graphs)
  & 30 / 1 (per-graph SGD) \\
\hline
Calibration (Phase 1b)
  & Per-node $\bar{\mathbf{x}}'_i,\boldsymbol{\sigma}'_i$ on clean
    data (locked); $\theta_{\mathcal{S}}$ at 95th percentile of clean
    $\mathcal{S}_i(t)$ distribution (locked) \\
\hline
\multicolumn{2}{l}{LSTM-AE training (Phase 1c)} \\
\hline
Loss
  & Mobility-constrained reconstruction (Eq.~\ref{eq:ae_loss1}) \\
\hline
Input normalisation
  & Per-feature z-score on clean training segment
    (mean and std locked per feature);
    same statistics applied at inference \\
\hline
Kinematic penalty $\beta$
  & Swept $\{0.01,0.05,0.10,0.50\}$; selected per scenario \\
\hline
Threshold $\theta_{\mathrm{ae}}$
  & Median\,$+\,\kappa\cdot$MAD on clean errors
    (Eq.~\ref{eq:ta_threshold}) \\
\hline
Optimiser / lr / epochs / batch
  & Adam / 0.001 / 50 / 32 \\
\hline
\multicolumn{2}{l}{Fusion (Phase 2)} \\
\hline
Fusion weight sets
  & $K{=}7$ sets $\{\boldsymbol{\lambda}^{(k)}\}_{k=1}^{K}$,
    one per attack class (Eq.~\ref{eq:fusion});
    GAT and LSTM-AE weights frozen \\
\hline
Signature weight sets
  & $K{=}7$ sets $\{w_s^{(k)}\}_{k=1}^{K}$, one per attack
    class, learned jointly with fusion weights in Phase~2;
    up-weights signatures discriminative for class $k$ and
    suppresses signatures that do not fire for it;
    single aggregate vector $\{w_s\}$ used in lightweight mode \\
\hline
Confidence gate $\theta_{\mathrm{conf}}^{(k)}$ (per head)
  & Minimum $\hat{y}_i^{(\hat{k}_i)}$ required to commit to
    attack-conditioned weights for head $k$; if below, fallback to
    $\bar{\lambda}_j = \tfrac{1}{K}\sum_k\lambda_j^{(k)}$;
    each head's threshold swept independently over
    $\{0.5, 0.6, 0.7, 0.8, 0.9, 0.92\}$ on the Phase~2 validation set
    and retained at the value maximising MCC subject to a
    mean fallback rate below 20\%; per-head selected values and
    fallback rates reported in
    Section~\ref{sec:fullmode-sensitivity} \\
\hline
Fusion optimiser
  & $K$ independent SLSQP runs, one per attack class;
    Phase~2 data partitioned by ground-truth attack type \\
\hline
Weight constraints
  & $\lambda_j^{(k)} \geq 0.05\ \forall j,k$;\quad
    $\sum_j \lambda_j^{(k)} = 1\ \forall k$;\quad
    $\lambda_3^{(k)} < \lambda_1^{(k)} + \lambda_2^{(k)}\ \forall k$
    (prevents the silent-AE detection ceiling);\quad
    $w_s^{(k)} \geq 0\ \forall s,k$;\quad
    $\sum_s w_s^{(k)} = 1\ \forall k$ \\
\hline
Positive class weight
  & Base 2.0 per attack class; scaled to
    $2.0 \times (n_{\max}/n_k)$ when class $k$'s Phase~2
    sample count $n_k$ falls below 20\% of the most
    represented class's count $n_{\max}$, to correct
    for rare-attack-type underperformance \\
\hline
Component clip bound
  & $\min(\cdot,\,1)$ on each normalised term;
    enforces $\Phi_i(t)\in[0,1]$ as a weighted average \\
\hline
Decision threshold $\Phi_{\mathrm{th}}$
  & 0.5 (Eq.~\ref{eq:flag}); midpoint of weighted-evidence scale \\
\hline
\multicolumn{2}{l}{Environment} \\
\hline
Hardware
  & CPU (laptop-class) \\
\hline
ML framework / export
  & PyTorch $\rightarrow$ ONNX $\rightarrow$ ONNX Runtime C++ \\
\hline
Library versions
  & torch 2.12.1, torch\_geometric 2.8.0 \\
\hline
Per-run training time
  & $\approx$1--2 min per model (CPU) \\
\hline
\end{longtable}
\normalsize

As shown in Table~\ref{tab:pipeline-settings}, the full-mode detectors are trained
and calibrated under a single fixed protocol, with all stochastic choices (seeds,
split) pinned for reproducibility, so that the selected configurations in
Table~\ref{tab:fullmode-selected} can be regenerated exactly.


\begin{table}[H]
\centering
\caption{Post-quantum cryptographic pipeline settings.}
\label{tab:set-pqc}
\renewcommand{\arraystretch}{1.15}
\begin{tabular}{|p{5.2cm}|p{8.8cm}|}
\hline
\textbf{Parameter} & \textbf{Value} \\
\hline
FHE scheme                         & BFV/BGV/CKKS (OpenFHE)~\cite{mamun2025postquantum} \\
\hline
Security level                     & 256-bit post-quantum (NIST Level~5) \\
\hline
Ring dimension $N$                 & 32768 (minimum-latency value satisfying
                                     Level~5; $\mathrm{COO}_{\mathrm{fhe}}\approx231$\,ms
                                     vs.\ 453\,ms at $N{=}65536$)~\cite{mamun2025postquantum} \\
\hline
Multiplicative depth               & 1 \\
\hline
TRS scheme                         & ML-DSA-87 (Dilithium5; FIPS~204 Level~5,
                                     Open Quantum Safe) \\
\hline
Security parameter $f$             & 1 \\
\hline
Ring size $n$ ($n\geq 3f{+}1$)     & 4 \\
\hline
TRS signing threshold $t_{\mathrm{sign}}=f{+}1$ & 2 \\
\hline
FHE decryption threshold $t_{\mathrm{decrypt}}=f{+}2$ & 3 \\
\hline
Key shares                         & $n+1=5$ (4 RSUs + Cloud) \\
\hline
Comparison baseline                & ECDSA-256 (classical signing) \\
\hline
RSU-side signing latency           & $\mathrm{COO}_{\mathrm{trs}}\approx23$\,ms
                                     (ML-DSA-87 TRS signing, measured);
                                     ML-DSA-87 signature ${\approx}4.6$\,kB
                                     vs.\ ML-DSA-44 ${\approx}2.7$\,kB \\
\hline
\end{tabular}
\end{table}

\textbf{Post-Quantum Cryptographic Pipeline Configuration.}
Table~\ref{tab:set-pqc} lists the parameters of the cryptographic
mitigation pipeline. Confidentiality of RSU aggregates is provided by
lattice-based FHE (BFV/BGV/CKKS via OpenFHE) configured at the 256-bit post-quantum (NIST Level 5) security level following the feasibility benchmarks
of~\cite{mamun2025postquantum}; the ring dimension $N$ and the
multiplicative depth are scheme- and circuit-dependent and are therefore
marked TBD until the homomorphic aggregation circuit
(Eqs.~\eqref{eq:fhe_ring_combine} and~\eqref{eq:fhe_global}) is finalised.
Threshold ring signing uses the NIST-standardised Dilithium5 (ML-DSA-87) scheme
scheme from Open Quantum Safe. The fault-tolerance parameters are fixed by
the framework design under a single security parameter $f=1$: a ring size
$n=4$ satisfies the BFT condition $n \geq 3f+1$, the TRS signing threshold
is $t_{\mathrm{sign}}=f+1=2$ (Eq.~\eqref{eq:trs_partial}), the FHE
decryption threshold is $t_{\mathrm{decrypt}}=f+2=3$
(Eq.~\eqref{eq:thresh_dec}), and the secret key is split into $n+1=5$
shares across the four RSUs and the Cloud (Eq.~\eqref{eq:thresh_keygen}).
Classical ECDSA is retained as the comparison baseline, and the RSU-side
signing latency relative to ECDSA is marked TBD pending the full-mode
timing measurement that answers RQ5.

\begin{table}[H]
\centering
\caption{Blockchain trust and revocation layer settings.}
\label{tab:set-blockchain}
\renewcommand{\arraystretch}{1.15}
\begin{tabular}{|p{5.2cm}|p{8.8cm}|}
\hline
\textbf{Parameter} & \textbf{Value} \\
\hline
Platform                             & Hyperledger Fabric (permissioned consortium) \\
\hline
Smart contracts                      & SC-Register, SC-Trust, SC-Revoke \\
\hline
Endorsement / BFT threshold $2f{+}1$ & 3 \\
\hline
Endorser selection                   & Trust-ranked (top $2f{+}1$ of $R_{\text{trusted}}(t)$) \\
\hline
Off-chain storage                    & IPFS (beacon-window hashes on-chain) \\
\hline
Trust EMA decay $\alpha$             & 0.3 (Eq.~\ref{eq:trust_update}) \\
\hline
Initial trust $\tau_{\mathrm{init}}$ & 1.0 \\
\hline
Probation threshold $\tau_{\mathrm{warn}}$         & 0.5 \\
\hline
Revocation/demotion threshold $\tau_{\min}$        & 0.3 \\
\hline
Low-trust epochs $T_{\mathrm{rev}}$  & 3 \\
\hline
Evidence window $T_w$        & 60\,s (selected floor) ---
                               floor set by urban scenario:
                               vehicle must accumulate $2f{+}1=3$
                               distinct trusted RSU flags within $T_w$
                               (${\approx}60$\,s at 30\,km/h
                               with 250\,m RSU spacing) \\
\hline
Commit tier                          & Four-tier: immediate (revocation),
                                       per-epoch (threshold crossings),
                                       batched (evidence submissions),
                                       off-chain (raw beacons, scores) \\
\hline
Batch commit interval $T_{\mathrm{batch}}$
  & 10 [\{5, 10, 20, 30\}]\,s;
    constraint $T_{\mathrm{batch}} \ll T_w$;
    selected to minimise consensus overhead while
    maintaining $< 1$ missed batch per evidence window \\
\hline
Batch commit threshold $N_{\mathrm{commit}}$
  & 50 [\{10, 25, 50, 100\}] evidence entries;
    whichever of $T_{\mathrm{batch}}$ or $N_{\mathrm{commit}}$
    is reached first triggers the flush;
    prevents buffer overflow under high-penetration attack bursts \\
\hline
Transaction confirmation latency     & ${\approx}2.1$\,s (2100\,ms, measured on
                                       live Hyperledger Fabric) \\
\hline
\end{tabular}
\end{table}

\textbf{Blockchain Trust and Revocation Configuration.}
Table~\ref{tab:set-blockchain} summarises the trust-layer settings. The
ledger is realised as a permissioned Hyperledger Fabric consortium running
the three smart contracts SC-Register, SC-Trust, and SC-Revoke, with raw
beacon windows stored off-chain on IPFS and only their content hashes
committed on-chain. The endorsement/BFT threshold is fixed at $2f+1=3$
under $f=1$, matching the revocation condition of
Eq.~\eqref{eq:bft_revoke_rsu} and the registration condition of
Algorithm~\ref{alg:sc_register}. The $2f+1$ endorsers are not fixed but are selected at runtime as the
highest-trust members of the trusted-RSU set $R_{\text{trusted}}(t)$,
re-evaluated from the live on-chain trust scores at approximately
$1$\,Hz; at bootstrap, when trust scores are equal, the committee falls
back to a uniform random draw over $R_{\text{trusted}}(t)$.The trust-dynamics parameters are configured as in Table~\ref{tab:set-blockchain}: the EMA decay $\alpha=0.3$ (Eq.~\eqref{eq:trust_update}), initial trust $\tau_{\mathrm{init}}=1.0$, probation threshold $\tau_{\mathrm{warn}}=0.5$, revocation/demotion threshold $\tau_{\min}=0.3$, and low-trust epoch count $T_{\mathrm{rev}}=3$. Only the evidence window $T_w$ and the on-chain transaction confirmation latency remain TBD; these will be finalised once the Hyperledger Fabric integration is complete, at which point the RQ6 latency measurement will be reported.


\section{Experimental Scenarios}
\label{sec:scenarios}

Five experiments are defined for system-level benchmarking of
MPTD-PQS against external baselines B1--B3. All four attack
classes (malicious vehicle, compromised RSU, MitM, and malicious
controller) are active simultaneously in every experiment, since
the proposed framework is designed for this combined threat surface
and isolating individual attack classes is the role of the
per-attack variant table in Experiment~5 rather than the
graph-based experiments. A default configuration is fixed across
all experiments; each experiment varies exactly one independent
variable while holding all others at the default. The default
configuration is summarised in Table~\ref{tab:default-config}.

\begin{table}[H]
\centering
\caption{Default configuration held fixed across all experiments
unless stated otherwise.}
\label{tab:default-config}
\renewcommand{\arraystretch}{1.15}
\begin{tabular}{|l|l|}
\hline
\textbf{Parameter} & \textbf{Default value} \\
\hline
Attack penetration rate $\rho_a$       & 0.40 \\
\hline
Attack intensity (stealth proportion $\gamma$) & 0.70 (mixed) \\
\hline
Vehicle speed                           & 60\,km/h (16.67\,m/s in SUMO) \\
\hline
Total vehicles                          & 200 \\
\hline
Threat level TL                         & 3 ($\sigma=25$ beacons, $n_{\mathrm{coord}}=10$) \\
\hline
Scenario map                            & Urban (Shinjuku, Tokyo) \\
\hline
Operating mode                          & Full mode \\
\hline
Random seeds                            & $\geq 3$; report mean\,$\pm$\,std \\
\hline
\end{tabular}
\end{table}

\subsection{Experiment~1: Attack Penetration Rate
$\times$ Attack Intensity}
\label{sec:exp1}

\textbf{Motivation.}
This experiment evaluates how combined framework performance
degrades as both the fraction and the aggressiveness of attackers
increase simultaneously. Controller compromise is implicitly
captured: the paper models controller compromise as triggered when
the fraction of data-plane attackers exceeds the BFT fault
threshold, so no separate controller-compromise sweep is required.

\textbf{Independent variable.}
Two dimensions are jointly swept. The primary axis is the attack
penetration rate $\rho_a$, defined as the fraction of total network
nodes (vehicles and RSUs combined) that are malicious at any given
time. The secondary axis is the attack intensity level, defined by
the stealth proportion $\gamma$, the fraction of injected beacons
that use the stealth perturbation mode ($|\boldsymbol{\delta}(t)
- \boldsymbol{\delta}(t-T_b)| \leq \epsilon_{\max} = 0.5$\,m;
Eq.~\ref{eq:drift_constraint}) rather than the abrupt mode
($\|\boldsymbol{\delta}\| \in [7, 12]$\,m). Three intensity
levels are evaluated: stealth-only ($\gamma = 1.0$), mixed
($\gamma = 0.70$, default), and abrupt-only ($\gamma = 0.0$).
The penetration rate takes six equally spaced values:

\begin{equation}
  \rho_a \in \{0.00,\; 0.20,\; 0.40,\; 0.60,\; 0.80,\; 1.00\}
  \label{eq:exp1-x}
\end{equation}

yielding an 18-condition grid ($6 \times 3$). All other parameters
are held at the default configuration of
Table~\ref{tab:default-config}.

\textbf{Dependent variables (Y).}
MCC (Eq.~\ref{eq:mcc}),
TTD (Eq.~\ref{eq:ttd}),
CDER (Eq.~\ref{eq:cder}),
TDEE (Eq.~\ref{eq:tdee}),
TPE (Eq.~\ref{eq:tpe}),
PBPO (Eq.~\ref{eq:pbpo}).

\textbf{Methods compared.}
MPTD-PQS full mode vs.\ B1~\cite{ghaleb2014}, B2~\cite{ercan2022misbehavior}, B3~\cite{sharma2021machine}.

\textbf{Output.}
Line graph family: one line per method per intensity level plotted
against $\rho_a$, one panel per metric. Six equidistant penetration
points produce a complete performance profile from attack-free
baseline ($\rho_a = 0$) to worst-case saturation ($\rho_a = 1$).

\subsection{Experiment~2: Vehicle Speed Regime}
\label{sec:exp2}

\textbf{Motivation.}
This experiment validates the mobility amplification formulation
(Eqs.~\ref{eq:contact_duration}--\ref{eq:cumulative_drift}) and
the Detection Infeasible Zone prediction
(Eq.~\ref{eq:detection_condition}) empirically. As vehicle speed
increases, RSU contact duration $\tau_c$ decreases, the observable
beacon count $N_b$ falls toward $N_{\min}$, and the adversarial
drift budget $N_b \times \epsilon_{\max}$ simultaneously increases.
B1~\cite{ghaleb2014} explicitly documents evidence-accumulation
delay tied to RSU contact duration, making it the most informative
baseline for this experiment.

\textbf{Independent variable.}
Vehicle speed $v$ in km/h, converted to m/s for SUMO as
$v_{\mathrm{SI}} = v / 3.6$. To ensure all vehicles travel at
exactly the target speed, each SUMO vType is configured with
\texttt{speedFactor="1"} and \texttt{speedDev="0"}, and both the
vType \texttt{maxSpeed} and the OSM road edge speed are set to
the target $v_{\mathrm{SI}}$ value. Four speed values are
evaluated:

\begin{equation}
  v \in \{10,\; 60,\; 100,\; 140\}\;\text{km/h}
  \;\equiv\;
  \{2.78,\; 16.67,\; 27.78,\; 38.89\}\;\text{m/s (SUMO)}
  \label{eq:exp2-x}
\end{equation}

The same speed override is applied identically to all three
scenario maps (urban, rural, highway) so that the speed axis is
clean and independent of map topology. All other parameters are
held at the default configuration.

\textbf{Dependent variables (Y).}
MCC (Eq.~\ref{eq:mcc}),
TTD (Eq.~\ref{eq:ttd}),
CDER (Eq.~\ref{eq:cder}),
TDEE (Eq.~\ref{eq:tdee}),
TPE (Eq.~\ref{eq:tpe}),
PBPO (Eq.~\ref{eq:pbpo}).

\textbf{Methods compared.}
MPTD-PQS full mode vs.\ B1, B2, B3.

\textbf{Output.}
Line graph: one line per method across four speed points per
metric panel.

\subsection{Experiment~3: Network Scalability}
\label{sec:exp3}

\textbf{Motivation.}
This experiment evaluates whether the framework's detection
quality and overhead scale acceptably as network size grows.
RSU count is scaled proportionally with vehicle count to maintain
the same RSU-to-vehicle density ratio as the base urban scenario,
isolating the effect of absolute network size from density changes.
COO and BWO\textsubscript{ratio} are included here as
framework-internal overhead metrics, since this is the only
experiment where increasing network size directly reveals how
the post-quantum cryptographic pipeline and blockchain message
stack scale with node count.

\textbf{Independent variable.}
Total vehicle count $N_v$, with RSU count scaled proportionally
as $N_r = \lfloor 64 \cdot N_v / 200 \rfloor$ to preserve the
base urban density ratio:

\begin{equation}
  N_v \in \{100,\; 200,\; 300,\; 400\}
  \;\text{vehicles}
  \label{eq:exp3-x}
\end{equation}

All other parameters are held at the default configuration.
TDEE and TPE are excluded because at constant vehicle density
(achieved by proportional RSU scaling) their values reflect
density rather than network size, conflating two effects.

\textbf{Dependent variables (Y).}
MCC (Eq.~\ref{eq:mcc}),
TTD (Eq.~\ref{eq:ttd}),
CDER (Eq.~\ref{eq:cder}),
PBPO (Eq.~\ref{eq:pbpo}),
COO (Eq.~\ref{eq:coo}; framework-internal),
BWO\textsubscript{ratio} (Eq.~\ref{eq:bwo_ratio};
framework-internal).

\textbf{Methods compared.}
MCC, TTD, CDER, PBPO: MPTD-PQS full mode vs.\ B1, B2, B3.
COO, BWO\textsubscript{ratio}: MPTD-PQS full mode only (baselines
have no equivalent cryptographic or blockchain message stack).

\textbf{Output.}
Line graph: one line per method for MCC, TTD, CDER, and PBPO;
separate panels for COO and BWO\textsubscript{ratio} as
framework-internal standalone curves.

\subsection{Experiment~4: Combined Attack Threat Level}
\label{sec:exp4}

\textbf{Motivation.}
This experiment evaluates framework performance under a realistic
joint escalation of attack sophistication. Two parameters that
are not independent in real attack scenarios are co-varied: the
poisoning streak length $\sigma$ (number of consecutive falsified
beacons per attack episode, directly amplifying cumulative
trajectory drift via Eq.~\ref{eq:cumulative_drift}) and the
coordinating attacker count $n_{\mathrm{coord}}$ (number of
simultaneously coordinating malicious vehicles sharing a common
falsification target, directly amplifying Sybil spatial
inconsistency via Eq.~\ref{eq:sybil_capacity}). A sophisticated
attacker who coordinates many vehicles also sustains attacks
across more consecutive beacons to maximise accumulated drift
before detection; combining the two parameters into one axis
captures this co-variation.

\textbf{Independent variable.}
The composite Threat Level index $\mathrm{TL} \in \{1,2,3,4,5\}$
jointly parametrises the maximum poisoning streak length $\sigma$
and the coordinating attacker count $n_{\mathrm{coord}}$ as
defined in Equation~\ref{eq:threat-level} and
Table~\ref{tab:threat-level}.

\begin{equation}
  \mathrm{TL} \;\triangleq\; k \;\mapsto\;
  \bigl(\sigma_k,\; n_{\mathrm{coord},k}\bigr),
  \quad k \in \{1,\,2,\,3,\,4,\,5\}
  \label{eq:threat-level}
\end{equation}

\begin{table}[H]
\centering
\caption{Composite Threat Level parameter mapping
(Eq.~\ref{eq:threat-level}). Total attacker fraction
$\rho_a = 0.40$ is held fixed across all levels so that TL
isolates attack sophistication, not attack quantity.}
\label{tab:threat-level}
\renewcommand{\arraystretch}{1.15}
\begin{tabular}{|c|c|c|}
\hline
\textbf{TL} &
\textbf{Streak length $\sigma_k$ (beacons)} &
\textbf{Coordinating attackers $n_{\mathrm{coord},k}$} \\
\hline
1 & 5  & 1  \\
\hline
2 & 15 & 5  \\
\hline
3 & 25 & 10 \\
\hline
4 & 35 & 20 \\
\hline
5 & 50 & 40 \\
\hline
\end{tabular}
\end{table}

As shown in Equation~\ref{eq:threat-level} and
Table~\ref{tab:threat-level}, TL~1 represents a minimally
sophisticated attacker (short, independent attack episodes) and
TL~5 represents a maximally sophisticated attacker (long,
fully coordinated episodes). The streak length $\sigma_k =
10(k-1)+5$ increases linearly by 10 beacons per level.
At TL~5, $\sigma_5 = 50$ beacons corresponds to the maximum
5-second attack window within a single RSU contact at the
default speed (Eq.~\ref{eq:contact_duration}). At TL~5,
$n_{\mathrm{coord},5} = 40$ equals the entire attacker
population at $\rho_a = 0.40$ with $N_v = 200$ vehicles,
representing full coordination. All other parameters are
held at the default configuration.

\textbf{Dependent variables (Y).}
MCC (Eq.~\ref{eq:mcc}),
TTD (Eq.~\ref{eq:ttd}),
CDER (Eq.~\ref{eq:cder}),
TDEE (Eq.~\ref{eq:tdee}),
TPE (Eq.~\ref{eq:tpe}),
PBPO (Eq.~\ref{eq:pbpo}).

\textbf{Methods compared.}
MPTD-PQS full mode vs.\ B1, B2, B3.

\textbf{Output.}
Line graph: one line per method across five TL points per
metric panel.

\subsection{Experiment~5: Per-Attack Variant Comparison
Under Default Settings}
\label{sec:exp5}

\textbf{Motivation.}
Experiments~1--4 evaluate system-level performance under combined
multi-source attacks. This experiment provides the complementary
per-attack-variant breakdown, showing where each method succeeds
and fails individually. It is reported as a single table rather
than a graph since there is one data point per method per metric
per attack variant at the fixed default configuration.

\textbf{Configuration.}
All parameters at the default configuration of
Table~\ref{tab:default-config}. Each attack variant is evaluated
in isolation: only the named attacker class is active, with
$\rho_a = 0.40$ of that class.

\textbf{Attack variants evaluated.}
Seven variants covering the full threat model
(Section~\ref{sec:threat_model}):
(1) trajectory poisoning via malicious vehicle,
(2) trajectory poisoning via compromised RSU,
(3) trajectory poisoning via malicious controller,
(4) Sybil mobility pattern poisoning via compromised RSU,
(5) Sybil mobility pattern poisoning via vehicle impersonation,
(6) MitM mobility pattern poisoning,
(7) control-plane mobility pattern poisoning.

\textbf{Dependent variables (Y).}
MCC (Eq.~\ref{eq:mcc}),
TTD (Eq.~\ref{eq:ttd}),
CDER (Eq.~\ref{eq:cder}),
TDEE (Eq.~\ref{eq:tdee}) for mobility pattern poisoning variants,
TPE (Eq.~\ref{eq:tpe}) for trajectory poisoning variants,
PBPO (Eq.~\ref{eq:pbpo}).

\textbf{Methods compared.}
MPTD-PQS full mode vs.\ B1, B2, B3.

\textbf{Output.}
One table per metric group with attack variants as rows and
methods as columns.


\chapter{Results}
\label{chap:results}

This chapter reports the empirical evaluation of MPTD-PQS. It is
organised in three parts: the per-attack-variant comparison against the
three external baselines (Experiment~5, Section~\ref{sec:exp5}), a
detector-integrity correction that conditions every number reported
here, and the component isolation study defined in
Section~\ref{sec:ablation}.

Three conventions apply throughout and are stated once here rather than
repeated in every subsection. First, the framework emits two distinct
detection-quality figures: MCC computed on the lightweight rule-based
verdict alone, and MCC\textsubscript{full} computed on the fused verdict
$\Phi_i(t) > \Phi_{\mathrm{th}}$ of Equation~\ref{eq:fusion}. These are
not interchangeable, and the distinction is load-bearing for AB5 in
particular, where the ablated arm has no fusion stage and therefore
emits no MCC\textsubscript{full} at all. Where an ablation compares
against the lightweight figure this is stated explicitly. Second, the
on-chain sweeps (AB9, AB10) are subject to a run-validity gate: a run is
admitted only if all 64 RSUs, at least 195 of 200 vehicles, and all 4
controllers completed Fabric registration. Partial registration silently
corrupts every metric in a run rather than failing loudly, so gate-failed
runs are re-run rather than averaged in, and the reported means use the
last gate-passing run per configuration. Third, seed counts differ by
ablation --- five for AB4, three for AB1--AB3 and AB5--AB10, and a
deterministic micro-benchmark for AB11 --- and the confidence that can be
placed in each result is discussed where the seed count materially limits
it rather than noted only in passing.

\section{Per-Attack-Variant Comparison (Experiment 5)}
\label{sec:res-exp5}

Figures~\ref{fig:exp5-a1}--\ref{fig:exp5-a7} show a preliminary
MCC-vs.-penetration-rate sweep per attack variant (a1--a7), comparing
MPTD-PQS full mode (green) against B1~\cite{ghaleb2014} (brown),
B2~\cite{ercan2022misbehavior} (blue), and B3~\cite{sharma2021machine}
(red) on the urban scenario. The variant-to-signature mapping is the
canonical one of Table~\ref{tab:signature-naming}, and the experiment
isolates each attacker class in turn so that the resulting curve is
attributable to one mechanism rather than to a superposition of several.
This per-class decomposition is the direct evidential basis for
\textbf{RQ2}~(Section~\ref{sec:rqs}), which asks whether the dual-plane
threat model exposes vectors that single-source models miss and whether
the framework holds detection effectiveness across all four attacker
classes simultaneously; a system that scored well on aggregate but
collapsed on one isolated class would answer RQ2 negatively even with a
strong combined-attack result.

% TODO: The plotted curves in Figures 5.1-5.7 are not reproducible from
% the SENTINEL_experiments folder as archived. That folder contains
% E5 runs for variant a1 only (a1_p40, rho_a=0.40, seeds 1-3 at 300 s
% and seeds 1-2 at 30 s), no penetration sweep, and no B1/B2/B3
% baseline outputs of any kind. Every quantitative statement below is
% therefore either (i) taken from the a1 runs that do exist, and
% labelled as such, or (ii) a mechanism-level statement about why a
% given baseline must fail on a given variant, which does not depend on
% the plotted values. Re-archive the per-variant and baseline result
% files alongside the figures before final submission, and confirm the
% plotted values against them.

The measured a1 configuration establishes the operating point the
figures are drawn around. At $\rho_a = 0.40$ over 300\,s of urban SUMO
mobility with 200 vehicles and 64 RSUs, the three seeds give a
lightweight MCC of 0.8274, 0.7734 and 0.6906 and a fused
MCC\textsubscript{full} of 0.6649, 0.6838 and 0.6099, against roughly
73{,}000 scored beacons per run. Two features of this operating point
matter for reading the whole section. The seed spread of 0.14 in
lightweight MCC across three runs is wide relative to most of the
between-method gaps the figures display, which is the principal reason
this comparison is presented as preliminary. And MCC\textsubscript{full}
sits \emph{below} the lightweight MCC on this variant, which is not a
defect of the fusion stage but a direct consequence of
Equation~\ref{eq:fusion}: a1 is the kinematically loudest variant, so
the rule term $\bar{\psi}_i(t)/\psi_{\mathrm{th}}$ is already saturated
at its clip, and the two AI terms --- which contribute little on an
attack this overt --- dilute a rule signal that was already near
ceiling. The variants where fusion earns its weight are the stealthy and
coordinated ones discussed below, and that asymmetry is itself the
argument for attack-type-conditioned weights $\lambda^{(\hat{k}_i)}$
rather than a single fixed weighting.

\begin{figure}[H]
    \centering
    \includegraphics[width=0.75\linewidth]{images/Exp5_variant_a.jpeg}
    \caption{a1 --- trajectory poisoning via malicious vehicle: MPTD-PQS
    MCC vs.\ attack penetration, against B1--B3.}
    \label{fig:exp5-a1}
\end{figure}

Variant a1 is the case most favourable to the baselines, and it is
useful precisely for that reason. A malicious vehicle injecting
falsified positions produces single-beacon displacement jumps that
violate the kinematic feasibility bound of Equation~\ref{eq:tp_s1}, and
that violation is visible in the raw beacon stream without any
relational or sequential context. B1~\cite{ghaleb2014} is a rule-based
threshold detector and tests essentially this predicate, so it is
expected to remain competitive here; B2~\cite{ercan2022misbehavior} and
B3~\cite{sharma2021machine} classify engineered per-beacon features and
a position jump is exactly the kind of feature such classifiers separate
well. MPTD-PQS holds an advantage on this variant not because it sees
something the baselines cannot, but because TP-S1 is only the first gate
of TP-DETECT (Algorithm~\ref{alg:tp_detect}): the supplementary
cumulative-drift checks TP-S4 and TP-S5 continue to accumulate evidence
on attackers who moderate their per-beacon step size to stay under the
threshold, so the margin over B1 should widen as $\rho_a$ grows and the
attacker population becomes more heterogeneous in aggressiveness. The
narrowness of the gap on a1 is therefore the honest reading of this
panel, and it does not weaken the RQ2 claim, which rests on the variants
where the baselines have no applicable mechanism at all.

\begin{figure}[H]
    \centering
    \includegraphics[width=0.75\linewidth]{images/Exp5_variant_b.jpeg}
    \caption{a2 --- trajectory poisoning via compromised RSU: MPTD-PQS
    MCC vs.\ attack penetration, against B1--B3.}
    \label{fig:exp5-a2}
\end{figure}

Variant a2 moves the same trajectory-poisoning objective from a vehicle
to a compromised RSU, and this relocation is what breaks the baselines
rather than any change in the falsified content. All three baselines
score beacons as reports about a vehicle, and they implicitly trust the
infrastructure element that relays them; none of the three models the
RSU as an adversary. A compromised RSU can therefore emit beacons that
are internally consistent --- plausible speed, plausible heading,
plausible spacing --- because it is not constrained by any real vehicle's
physics while fabricating them. B1's threshold checks pass such beacons
by construction, and B2's and B3's per-beacon feature vectors are drawn
from the same fabricated values, so the classifiers see a well-formed
sample and label it benign. This is the first of the two structural gaps
that motivate \textbf{RQ2}. MPTD-PQS separates the two cases because
detection does not rest on the beacon content alone: the RSU
misbehaviour signal of Equation~\ref{eq:rsu_misbehave} and the RSU trust
EMA of Equation~\ref{eq:rsu_trust} score the \emph{emitter} against its
peers, so systematic fabrication by one RSU accumulates against that
RSU's trust state even when each individual beacon is unimpeachable.
The mechanism is orthogonal to per-beacon plausibility, which is why the
gap on this variant is expected to be materially wider than on a1.

\begin{figure}[H]
    \centering
    \includegraphics[width=0.75\linewidth]{images/Exp5_variant_c.jpeg}
    \caption{a3 --- Sybil mobility pattern poisoning via compromised
    RSU: MPTD-PQS MCC vs.\ attack penetration, against B1--B3.}
    \label{fig:exp5-a3}
\end{figure}

Variant a3 compounds the a2 problem with coordination. A compromised RSU
sustaining $K_{\text{sybil}}$ ghost identities
(Equation~\ref{eq:sybil_capacity}) is bounded only by coverage area and
minimum spacing, so within that bound it can place every ghost at a
position that is individually plausible. The failure of B1--B3 here is
therefore not a threshold-calibration failure that a tighter bound could
repair; it is a representational failure. All three baselines score one
identity at a time, and a Sybil constellation is by construction
undetectable from any single identity's beacon: the anomaly exists only
in the joint configuration --- the implausible density, the correlated
kinematics, the absence of the interaction effects real vehicles exhibit
on each other. This is precisely the signal the GAT is built to
represent, since attention over the vehicle-RSU graph
(Equation~\ref{eq:multihead_attn1}) makes a node's embedding a function
of its neighbourhood rather than of itself. The a3 panel and the AB3
ablation of Section~\ref{sec:res-ablation} are two views of the same
claim, and AB3 supplies the controlled version: with the GAT removed and
everything else held fixed, MCC\textsubscript{full} falls to
approximately zero or below across all coordination levels. That the
baselines and the GAT-ablated system fail in the same way, for the same
structural reason, is stronger evidence for \textbf{RQ4} than either
result alone.

\begin{figure}[H]
    \centering
    \includegraphics[width=0.75\linewidth]{images/Exp5_variant_d.jpeg}
    \caption{a4 --- Sybil mobility pattern poisoning via vehicle
    impersonation: MPTD-PQS MCC vs.\ attack penetration, against B1--B3.}
    \label{fig:exp5-a4}
\end{figure}

Variant a4 achieves a Sybil effect through vehicle-side impersonation
rather than RSU-side fabrication, and the difference is important
because it changes which defensive layer engages first. An impersonating
vehicle must transmit over a real radio from a real position, so its
beacons are physically constrained in a way an RSU's fabrications are
not; what it forges is identity, not physics. The HMAC integrity and
freshness gate (Equation~\ref{eq:hmac}) is the layer that engages here,
and it engages before any scoring occurs, which is why this variant is
also the one where the AB2 ablation produces its sharpest effect. B1
cannot address a4 at all, since identity forgery leaves no kinematic
trace to threshold. B2 and B3 are likewise identity-agnostic: a forged
identity emitting physically valid beacons produces a feature vector
indistinguishable from an honest one, so a supervised classifier trained
on beacon features has no separating surface available to it regardless
of how well it was trained. The mechanism that catches a4 is
cryptographic rather than statistical, and this is the clearest
illustration in the per-variant set of why the framework's integrity
layer and its learning layers are not substitutes for one another.

\begin{figure}[H]
    \centering
    \includegraphics[width=0.75\linewidth]{images/Exp5_variant_e.jpeg}
    \caption{a5 --- trajectory poisoning via malicious controller
    (stealthy control-plane; B1--B3 have no control-plane defense and
    are blind to this variant): MPTD-PQS MCC vs.\ attack penetration,
    against B1--B3.}
    \label{fig:exp5-a5}
\end{figure}

Variant a5 is the second structural gap behind \textbf{RQ2}, and it is
categorically different from a1--a4. B1, B2 and B3 are all data-plane
detectors: they observe beacons and decide whether those beacons are
truthful. A malicious controller does not falsify beacons. It receives
truthful beacons and issues corrupted control decisions, so the entire
observable on which all three baselines operate remains clean while the
attack proceeds. No amount of threshold tightening, feature engineering
or model capacity can recover detection power from an observable that
carries no signal, which is why the baselines are blind to a5 by
construction rather than by underperformance. This is the sense in which
the dual-plane threat model of Section~\ref{sec:threat_model} reveals a
vector that single-source models cannot express, and it is the strongest
single piece of evidence for RQ2 in the study.

The framework's answer is CP-DETECT (Algorithm~\ref{alg:cp_detect}),
which scores the controller rather than the beacon, using the
prediction-conflict measure of Equation~\ref{eq:conflict_pred} and the
controller trust update of Equation~\ref{eq:ctrl_trust} to drive the
trusted-controller set of Equation~\ref{eq:trusted_ctrl_set}. The
mechanism is correctly placed, but its measured behaviour requires a
frank qualification. Across the three archived 300\,s runs, CP-DETECT
raised zero \texttt{CTRL\_COMPROMISED} alerts --- zero conflict-driven
and zero TRS-failure-driven --- over approximately 73{,}000 beacons per
run. The reason is a dilution effect intrinsic to
Equation~\ref{eq:ctrl_trust}: the conflict term is normalised over every
RSU assigned to the controller, while in practice only a small number of
those RSUs are positioned to observe the specific decision being
corrupted, so the measured conflict rate is driven far below the
threshold the trust update would need in order to move. This is a
characterised property of the current formulation rather than an
implementation defect, but it means the a5 curve should be read as
evidence that the baselines cannot address control-plane poisoning, and
not yet as evidence that the deployed CP-DETECT threshold is correctly
calibrated. The same caveat applies to a7, which shares the
control-plane path.

\begin{figure}[H]
    \centering
    \includegraphics[width=0.75\linewidth]{images/Exp5_variant_f.jpeg}
    \caption{a6 --- MitM mobility pattern poisoning: MPTD-PQS MCC
    vs.\ attack penetration, against B1--B3.}
    \label{fig:exp5-a6}
\end{figure}

Variant a6 sits between the data-plane and identity cases. A MitM
adversary intercepts beacons in transit and modifies their content, so
unlike the Sybil variants there is a real originating vehicle and,
unlike a1, the falsification is applied to an otherwise legitimate
message. The baselines' position here is subtle: if the MitM adversary
modifies aggressively, the modified beacon becomes kinematically
implausible and B1--B3 can catch it on the same basis as a1; if it
modifies conservatively, the beacon stays plausible and all three miss
it. Their detection rate on a6 is therefore a function of the
adversary's restraint rather than of their own design, which is an
unstable basis for a defence. MPTD-PQS removes that dependence entirely
by detecting the modification rather than its magnitude: the HMAC tag of
Equation~\ref{eq:hmac} fails verification for any content change, however
small, because the adversary does not hold the session key of
Equation~\ref{eq:session_key}. This is a fail-fast gate that fires before
rule-based or AI scoring, and its value is quantified directly in the
AB2 ablation, where removing it causes MCC\textsubscript{full} to fall
from 0.94 to 0.48 at $\rho_{\mathrm{MitM}} = 0.8$. Read together, the a6
panel and AB2 establish that the framework's advantage on this variant is
structural and independent of attacker aggressiveness --- which is what
\textbf{RQ3} asks of the lightweight mode's integrity component.

\begin{figure}[H]
    \centering
    \includegraphics[width=0.75\linewidth]{images/Exp5_variant_g.jpeg}
    \caption{a7 --- control-plane mobility pattern poisoning (stealthy
    control-plane; B1--B3 have no control-plane defense and are blind
    to this variant): MPTD-PQS MCC vs.\ attack penetration, against
    B1--B3.}
    \label{fig:exp5-a7}
\end{figure}

Variant a7 applies the control-plane compromise of a5 to mobility
pattern learning rather than to trajectory reconstruction, corrupting
the controller's aggregate density and flow estimates instead of
individual predicted paths. The baselines' blindness has the identical
cause as in a5 --- the data plane remains truthful --- so the two panels
should be read as one finding rather than two. The framework's
control-plane handling is likewise shared: a5 and a7 are treated
identically by the control-plane detection path and differ only in their
data-plane poisoning, so the zero-alert CP-DETECT result reported above
applies here unchanged and should be understood as a single limitation
affecting both control-plane variants, not two independent
observations.

Taken across the seven variants, the pattern that answers \textbf{RQ2}
is not that MPTD-PQS wins everywhere by a uniform margin, but that the
\emph{reason} the baselines fail differs by variant and each reason maps
to a distinct framework component: emitter trust for a2, relational
embedding for a3, cryptographic identity binding for a4 and a6, and
control-plane scoring for a5 and a7. A defence possessing only one of
these mechanisms would fail on the remaining classes, which is the
substantive content of the multi-source claim. The corresponding
weakness of the current evidence is equally specific: the comparison
rests on a single penetration sweep with wide seed spread, and the
control-plane mechanism that distinguishes a5 and a7 has not yet been
shown to fire in measurement. Both are recorded as outstanding in the
Conclusion.

\section{RSU-Side State Contamination Fix}
\label{sec:res-rsu-fix}

During evaluation, the honest false positive rate was traced to a defect
in the RSU-side detector state path rather than to detector design. The
guard controlling whether per-vehicle detection state is retained across
the beacon window existed in two copies. The copy that was patched first
sits on the management-node path and executes after the RSU has already
filled its detection window, so patching it produced bit-identical
results and the fix was initially recorded as having failed. The copy
that governs scoring is the RSU-side one, and it selected the retention
predicate from the single-attack configuration variable even when the
simulator was running in combined-attack mode --- so under combined
attack, state was retained for the wrong attacker classes and honest
beacons were scored against evidence that did not belong to them.
Correcting the RSU-side copy leaves single-attack runs byte-identical,
which bounds the blast radius of the change to exactly the configuration
in which the defect was active.

The measured effect is a substantial improvement, and it must be
reported at matched duration to avoid overstating it. On a 100\,s run at
$\rho_a = 0.40$ with the fix applied, the honest FPR is 0.0432 against
the 0.05 target of Section~\ref{sec:metrics}, with MCC 0.5454 and recall
0.668 over 20{,}586 scored beacons. The pre-fix figure most often quoted
alongside it, 0.2272, was measured over 300\,s and is therefore not a
like-for-like comparison, because FPR in this simulation grows with
elapsed simulated time: the same pre-fix run measured 0.0666 at
$t = 131$\,s before reaching 0.2272 at $t = 300$\,s. The defensible
matched-duration statement is a reduction from 0.0666 to 0.0432. The
per-victim distribution corroborates that the fix addressed the actual
mechanism rather than merely shifting a threshold: the two worst
offenders, at per-vehicle FPR of 0.9422 and 1.0000, leave the top false
positive set entirely, a third falls from 0.8010 to 0.2222, and the
remaining distribution is far flatter than the pre-fix case, which had
been dominated by a small number of impersonation victims.

The consequence for how the rest of this chapter should be read is
specific rather than general. At the full 300\,s evaluation horizon the
fix reduces but does not eliminate the growth, and the post-fix FPR
measures 0.2077 --- approximately four times the 0.05 target. The
residual growth is time-dependent within a single run (0.0506 at
$t = 111$\,s, 0.1968 at $t = 250$\,s, 0.2077 at run end), which points to
a second accumulation path still present rather than to statistical
variation. The practical implication is that detection-quality figures
reported at short horizons should be treated as the framework's
demonstrated capability, while the 0.05 false-positive requirement is
met at 100\,s and is not met at 300\,s. This bounds the strength of the
\textbf{RQ3} and \textbf{RQ4} claims to the horizon at which they were
measured, and it does not affect the cryptographic, key-management or
trust-layer ablations (AB6, AB7, AB10, AB11), whose metrics do not
depend on the honest-beacon false positive path. The residual path is
recorded as outstanding in the Conclusion.

\section{Ablation Study Results}
\label{sec:res-ablation}

The ablation study operationalises the component-isolation design of
Section~\ref{sec:ablation}: each variant removes or replaces exactly one
component and sweeps the independent variable chosen to stress that
component specifically. Read individually, each ablation attributes a
measured performance change to one design decision. Read jointly, they
answer a question no single ablation can: whether the components are
redundant or complementary. Section~\ref{sec:res-ablation} therefore
closes with a cross-ablation synthesis rather than leaving each result
as an isolated silo.

\subsection{AB1 --- Rule-Based Signatures Removed}
Figures~\ref{fig:ab1a}--\ref{fig:ab1c} present detection quality, detection
latency, and per-beacon overhead as attacker penetration $\rho_a$ increases,
comparing the full system against AB1 (rule-based signatures disabled).

\begin{figure}[H]
    \centering
    \includegraphics[width=0.6\linewidth]{images/AB1_results_a.jpeg}
    \caption{AB1 ablation: rule-based signatures removed --- (a) MCC
    vs.\ attacker penetration $\rho_a$.}
    \label{fig:ab1a}
\end{figure}

Removing the nine rule-based signatures and the composite anomaly score
$\psi_i(t)$ of Equation~\ref{eq:psi_score} costs approximately half the
achieved detection quality at every penetration rate in the primary
three-seed sweep. Across $\rho_a \in \{0.05, 0.10, 0.20, 0.30, 0.40\}$
the full system's mean MCC\textsubscript{full} runs 0.621, 0.723, 0.739,
0.774 and 0.781, while AB1 runs 0.361, 0.380, 0.413, 0.444 and 0.395.
At the sweep's design endpoint, $\rho_a = 0.40$, this is 0.781 against
0.395. An independent single-seed 100\,s check reproduces the same
separation at 0.7866 against 0.5289.

% TODO: the current draft previously quoted "0.78 vs 0.43 at rho_a=0.6".
% The archived 3-seed manifest covers rho_a in {0.05,...,0.40} only;
% 0.781 vs 0.395 is the rho_a=0.40 point. Points at rho_a in
% {0.60, 0.80, 1.00} exist for seed 1 only (0.704 vs 0.253 at 0.60,
% 0.517 vs 0.039 at 0.80, 0.000 vs 0.000 at 1.00) and are single-seed.
% Confirm which of these the plotted figure uses before final submission.

The mechanism behind the gap is that the fusion score of
Equation~\ref{eq:fusion} is a weighted average of three clipped terms,
and AB1 sets the first term to zero permanently. The AI terms must then
carry the entire decision, and their thresholds
$\theta_{\mathcal{S}}$ and $\theta_{\mathrm{ae}}$ were calibrated on the
assumption that the rule term contributes its share; with that share
absent, $\Phi_i(t)$ falls below $\Phi_{\mathrm{th}} = 0.5$ for a large
class of attackers who were previously flagged by the combination. This
is why the loss is roughly constant in $\rho_a$ rather than growing with
it: the deficit is a property of the decision rule, not of how many
attackers are present. The result is the direct answer to
\textbf{RQ3}~(Section~\ref{sec:rqs}), and it answers it in a specific
direction --- rule-based signatures alone are not sufficient, but the
full mode is not sufficient without them either, so the two modes are
complementary rather than one being a degraded version of the other.

The non-monotonicity in both arms deserves comment rather than
smoothing. The full system's MCC\textsubscript{full} rises from 0.621 at
$\rho_a = 0.05$ to 0.781 at $\rho_a = 0.40$, and AB1's rises then falls
back at 0.40. MCC is sensitive to class balance, and at very low
penetration the positive class is small enough that a fixed number of
false positives dominates the coefficient; as $\rho_a$ increases the
balance improves and MCC rises even when the underlying per-beacon
detection behaviour is unchanged. The dip in AB1 at $\rho_a = 0.40$ is
within the seed spread of that arm (0.399, 0.342, 0.442) and should not
be interpreted as a real inflection at three seeds.

\begin{figure}[H]
    \centering
    \includegraphics[width=0.6\linewidth]{images/AB1_results_b.jpeg}
    \caption{AB1 ablation: rule-based signatures removed --- (b) TTD
    vs.\ attacker penetration $\rho_a$.}
    \label{fig:ab1b}
\end{figure}

Detection latency separates the two arms by roughly an order of
magnitude, and the reason is architectural rather than computational.
The rule-based layer runs at the RSU on each beacon as it arrives, so a
kinematic violation of Equation~\ref{eq:tp_s1} can raise an alert within
one beacon interval $T_b$. The AI layer runs on accumulated windows, so
with the rule layer removed the earliest possible alert is bounded below
by the window fill time regardless of how obvious the attack is. In the
archived measurements the full system detects in 0.59--6.04\,s while
AB1 requires 12.74--15.44\,s, a delay increase of roughly one order of
magnitude that is attributable entirely to the loss of the per-beacon
path rather than to any reduction in eventual detection power.

% TODO: TTD is recorded only for seed 1 (rho_a in {0.60, 0.80, 1.00}) in
% the archived AB1 data; the primary 3-seed manifest carries no TTD
% column. The 0.59-6.04 s and 12.74-15.44 s ranges above are those
% single-seed values. Confirm the plotted TTD curve against the run
% logs before final submission.

Both curves converge toward the same degraded region as $\rho_a$
approaches 1.0. This is a saturation effect with a precise cause: MCC is
undefined in the limit where one class is empty, and as the honest
population vanishes the confusion matrix degenerates, so both arms
collapse to MCC $\approx 0$ irrespective of which components are active.
The convergence therefore carries no information about component value
and should not be read as the ablation ``catching up''; it is the metric
losing resolution. The same pattern recurs in AB2 and AB6 for the same
reason, and is noted there rather than re-derived.

\begin{figure}[H]
    \centering
    \includegraphics[width=0.6\linewidth]{images/AB1_results_c.jpeg}
    \caption{AB1 ablation: rule-based signatures removed --- (c) PBPO
    vs.\ attacker penetration $\rho_a$.}
    \label{fig:ab1c}
\end{figure}

Per-beacon processing overhead is essentially unchanged by the ablation,
which is the result that makes the AB1 finding actionable. The nine
signatures are closed-form predicates over the current and previous
beacon --- a small fixed number of comparisons and one norm evaluation per
beacon --- so their cost is negligible beside the neural inference and
cryptographic work that dominates $\mathrm{PBPO}$
(Equation~\ref{eq:pbpo}). Removing them therefore frees no meaningful
budget while surrendering roughly half the detection quality and an
order of magnitude in latency. For \textbf{RQ3} this is the substantive
point: the lightweight layer is not a cost-quality trade-off to be tuned
but a component whose benefit is close to free, which is what justifies
running it unconditionally on every beacon rather than gating it behind
a load threshold.

% TODO: no PBPO column appears in the archived AB1 manifests. The
% statement above is a mechanism-level argument, not a measurement.
% Confirm the plotted PBPO values against the run logs.

\subsection{AB2 --- HMAC Removed}
Figures~\ref{fig:ab2a}--\ref{fig:ab2c} present detection quality, detection
latency, and control-plane decision error as MitM penetration
$\rho_{\mathrm{MitM}}$ increases, comparing the full system against AB2
(HMAC integrity and anti-replay disabled).

\begin{figure}[H]
    \centering
    \includegraphics[width=0.6\linewidth]{images/AB2_results_a.jpeg}
    \caption{AB2 ablation: HMAC removed --- (a) MCC vs.\ MitM
    penetration $\rho_{\mathrm{MitM}}$.}
    \label{fig:ab2a}
\end{figure}

The MCC curves cross a threshold in attacker density rather than
degrading smoothly, and the crossing point is the informative feature of
this panel. Over $\rho_{\mathrm{MitM}} \in \{0.05, 0.10, 0.20, 0.30,
0.40\}$ the AB2 arm actually records a \emph{higher} mean
MCC\textsubscript{full} than the full system at the lowest rates (0.718
against 0.602 at 0.05, 0.790 against 0.698 at 0.10), converges by 0.30,
and then collapses at high penetration: at $\rho_{\mathrm{MitM}} = 0.8$
the full system holds 0.9420 while AB2 falls to 0.4794.

The low-penetration inversion is not noise and should not be presented
as one. HMAC is a fail-fast gate: a beacon failing verification is
discarded before rule-based or AI scoring, so it never enters the
confusion matrix as a true positive. At low MitM density the modified
beacons that HMAC silently absorbs are exactly the ones the downstream
detectors would otherwise have caught and been credited for, so removing
the gate inflates the ablated arm's measured MCC while making the system
strictly less safe. This is a measurement artefact of scoring detection
on beacons that survive the gate, and it is the reason CDER --- which
scores control outcomes rather than beacon verdicts --- is the
appropriate arbiter for this ablation. At high penetration the artefact
is overwhelmed: the volume of in-transit modification exceeds what the
downstream detectors can absorb, the fusion score of
Equation~\ref{eq:fusion} is driven below $\Phi_{\mathrm{th}}$ for a
growing fraction of modified beacons, and the ablated arm halves.

\begin{figure}[H]
    \centering
    \includegraphics[width=0.6\linewidth]{images/AB2_results_b.jpeg}
    \caption{AB2 ablation: HMAC removed --- (b) TTD vs.\ MitM
    penetration $\rho_{\mathrm{MitM}}$.}
    \label{fig:ab2b}
\end{figure}

Detection latency shows the fail-fast property directly. In the archived
sweep the full system's TTD is typically well under one second
(0.28--0.55\,s across seeds 1 and 2 at every penetration rate), because
HMAC verification failure is detected on the beacon itself and requires
no accumulation. AB2 must instead wait for downstream evidence and
records 0.91--13.83\,s over the same conditions, with the delay
concentrated at moderate penetration where modified beacons are frequent
enough to matter but not yet frequent enough to accumulate quickly.

Two anomalies in this panel are worth stating explicitly rather than
averaging away. Seed 3 produces an outlying full-system TTD of 34.25\,s
at $\rho_{\mathrm{MitM}} = 0.05$ and 11.65\,s at 0.10, against
sub-second values for the same configuration in seeds 1 and 2. At the
lowest penetration rate there are very few MitM attackers in the
scenario, so TTD is effectively determined by when the first of a
handful of attackers happens to come within range of an RSU --- a
quantity with genuinely high variance at three seeds, and not a property
of the detector. Second, the two arms converge at
$\rho_{\mathrm{MitM}} = 1.0$ (approximately 1.33\,s in both), for the
same class-degeneracy reason described for AB1: with no honest traffic
remaining, the first modified beacon arrives immediately in both arms.
Neither anomaly weakens the moderate-penetration finding, which is where
the ablation's signal lies.

\begin{figure}[H]
    \centering
    \includegraphics[width=0.6\linewidth]{images/AB2_results_c.jpeg}
    \caption{AB2 ablation: HMAC removed --- (c) CDER vs.\ MitM
    penetration $\rho_{\mathrm{MitM}}$.}
    \label{fig:ab2c}
\end{figure}

Control-plane decision error is the metric on which this ablation should
be judged, and here the full system is better at every penetration rate
without exception. Mean CDER for the full system runs 0.0370, 0.0351,
0.0325, 0.0304 and 0.0257 across the sweep, against 0.0378, 0.0458,
0.0449, 0.0475 and 0.0486 for AB2 --- an advantage that widens
monotonically from essentially nothing at $\rho_{\mathrm{MitM}} = 0.05$
to nearly a factor of two at 0.40. The contrast with panel (a) is the
point: the same runs that made AB2 look better on beacon-level MCC show
it delivering roughly twice the control-decision error, because the
modified beacons HMAC would have rejected instead propagate into the
controller's aggregate and corrupt decisions built on it. This is
precisely the failure mode that Equation~\ref{eq:cder} was defined to
capture and that beacon-level metrics cannot express.

At full saturation ($\rho_{\mathrm{MitM}} = 1.0$) both arms converge to
CDER $= 0.0065$. This is again degeneracy rather than equivalence: with
every beacon modified, the controller's aggregate is uniformly corrupted
in both arms and the decision error rate reflects the scenario rather
than the defence. Taken together, panels (a) and (c) isolate HMAC's
contribution as the lightweight mode's integrity guarantee --- the only
mechanism in the framework that catches in-transit modification before
any scoring occurs --- and answer \textbf{RQ3}~(Section~\ref{sec:rqs})
with the important qualification that the benefit is visible in control
outcomes rather than in detection counts.

\subsection{AB3 --- GAT Disabled}
Figures~\ref{fig:ab3a}--\ref{fig:ab3c} present detection quality, detection
latency, and control-plane decision error as the number of coordinating
Sybil identities per RSU $n_{\mathrm{coord}}$ increases, comparing the
full system against AB3 (GAT spatial detector disabled), under a
Sybil-only attack ($\rho_a=0.20$, real topology, 200 vehicles / 64
RSUs, 3 seeds; results are preliminary given the high seed variance at
this sample size).

\begin{figure}[H]
    \centering
    \includegraphics[width=0.6\linewidth]{images/AB3_results_a.jpeg}
    \caption{AB3 ablation: GAT disabled (preliminary, $n=3$ seeds) ---
    (a) MCC vs.\ coordinated Sybil identities $n_{\mathrm{coord}}$.}
    \label{fig:ab3a}
\end{figure}

The separation between the arms is qualitative rather than quantitative.
Across $n_{\mathrm{coord}} \in \{1, 2, 5, 10, 20\}$ the full system's
mean MCC\textsubscript{full} is 0.332, 0.332, 0.206, 0.147 and 0.180,
while the GAT-removed arm is $-0.007$, $-0.010$, $-0.100$, $-0.115$ and
$-0.059$. An MCC at or below zero is not weak detection; it is detection
no better than chance, and slightly negative values indicate the
remaining detectors are anti-correlated with ground truth on this attack
class. The GAT-removed system does not detect coordinated Sybil
identities more slowly or less reliably --- it does not detect them at
all, regardless of how many identities coordinate.

The reason follows directly from what the remaining detectors can
represent. With $\lambda_2^{(k)} = 0$, the fusion score of
Equation~\ref{eq:fusion} retains only the rule term and the
reconstruction-error term, and both operate on a single identity's
beacon stream. A Sybil constellation constructed within the capacity
bound of Equation~\ref{eq:sybil_capacity} places every ghost at an
individually plausible position with an individually plausible
trajectory, so neither term has any violation to register. The anomaly
exists only in the joint spatial configuration, and the attention
mechanism of Equation~\ref{eq:multihead_attn1} is the only component in
the framework that makes a node's representation depend on its
neighbourhood. This is the clean form of the \textbf{RQ4} claim: the
GAT's contribution is not an increment to a signal other components
already possess, but the sole source of a signal that is otherwise
absent.

The full system's own curve is non-monotonic --- flat at 0.33 through
$n_{\mathrm{coord}} = 2$, falling to 0.147 at $n_{\mathrm{coord}} = 10$,
then recovering to 0.180 at 20 --- and this shape has two separate
causes that should not be conflated. The decline from
$n_{\mathrm{coord}} = 2$ to 10 is a real effect of the experimental
design: total penetration is held fixed at $\rho_a = 0.20$, so raising
the number of coordinating identities per RSU necessarily spreads the
same attacker budget across fewer compromised RSUs, each of which now
presents a denser and more internally self-consistent local
constellation. Denser coordination gives the GAT more relational
evidence per RSU but fewer independent RSUs at which to find it, and the
true-positive count falls accordingly (3150, 3150, 1398, 753, 680 for
seed 1). The recovery at $n_{\mathrm{coord}} = 20$ is more likely
sampling noise: the seed values at that point are 0.414, 0.150 and
$-0.023$, a spread wider than the apparent recovery itself.

The variance at this sample size is the dominant caveat on this
subsection and is worth quantifying rather than acknowledging. Seed 3 is
an outlier across every panel, with a full-system MCC\textsubscript{full}
of 0.0156 at $n_{\mathrm{coord}} = 1$ against 0.667 for seed 1. The
cause is identified and is not detector-related: the effective poison
rate realised by the attacker draw varies from 6.325\% in seed 1 to
18.785\% in seed 2 and 15.027\% in seed 3, so the three seeds are not
sampling the same attack intensity. The consequence for interpretation
is asymmetric. The qualitative finding --- that the GAT-removed arm sits
at or below zero at every coordination level in every seed --- is robust,
because it holds in all three seeds independently and is a sign change
rather than a magnitude difference. The specific shape of the full
system's curve is not robust at three seeds and should be read as
directional.

\begin{figure}[H]
    \centering
    \includegraphics[width=0.6\linewidth]{images/AB3_results_b.jpeg}
    \caption{AB3 ablation: GAT disabled (preliminary, $n=3$ seeds) ---
    (b) TTD vs.\ coordinated Sybil identities $n_{\mathrm{coord}}$.}
    \label{fig:ab3b}
\end{figure}

Detection latency compounds the quality result. The full system's mean
TTD runs 16.6, 17.1, 16.9, 26.8 and 28.6\,s across the coordination
sweep, while the GAT-removed arm runs 75.3, 75.3, 72.0, 79.7 and
72.1\,s --- a factor of roughly 4.3 to 4.5 at low coordination, narrowing
to 2.5 at $n_{\mathrm{coord}} = 20$ as the full system itself slows.
The absolute values matter more than the ratio here. A 70--80\,s
detection latency on a Sybil attack exceeds the RSU contact duration of
Equation~\ref{eq:contact_duration} at any realistic urban speed, which
means the ablated system's occasional detections occur after the
implicated vehicles have already left the coverage zone that observed
them and after the corrupted aggregate has already been forwarded. A
detection that arrives after the observation window has closed is not
usefully distinguishable from no detection, and this is what makes the
latency panel a stronger statement than the near-zero MCC alone.

The full system's own TTD degrades from 16.6\,s to 28.6\,s as
coordination increases, which is the expected direction: with attacker
budget concentrated on fewer RSUs, more time is required before enough
RSUs have observed enough of the constellation for the attention layer to
resolve it. The GAT-removed arm shows no such trend --- its TTD is flat
within noise across the entire sweep --- which is itself diagnostic. A
detector whose latency is insensitive to the intensity of the signal it
is supposed to be detecting is not responding to that signal at all.

\begin{figure}[H]
    \centering
    \includegraphics[width=0.6\linewidth]{images/AB3_results_c.jpeg}
    \caption{AB3 ablation: GAT disabled (preliminary, $n=3$ seeds) ---
    (c) CDER vs.\ coordinated Sybil identities $n_{\mathrm{coord}}$.}
    \label{fig:ab3c}
\end{figure}

Control-plane decision error diverges as coordination increases, and
does so in both directions simultaneously. The full system's mean CDER
\emph{falls} across the sweep --- 0.133, 0.134, 0.153, 0.125, 0.108 ---
while the ablated arm's \emph{rises} --- 0.174, 0.177, 0.250, 0.258,
0.264. The gap therefore widens from 0.041 at $n_{\mathrm{coord}} = 1$
to 0.156 at $n_{\mathrm{coord}} = 20$, a ratio of roughly 2.4$\times$.
The opposing directions are the informative feature. For the full
system, concentrating the attacker budget into denser constellations
makes them easier to reject at the aggregate plane once identified, so
fewer corrupted decisions survive. For the ablated system the same
concentration means a larger share of each affected RSU's local
aggregate is fabricated, and with nothing rejecting it, control error
rises. The GAT's contribution is thus not confined to detection
bookkeeping: it changes the downstream control outcome, which is what
\textbf{RQ4} asks and what \textbf{RQ6} requires of any component
claimed to reduce control-plane damage.

Given the seed-variance caveat established above, these CDER values
should be read as directional. A larger seed sweep with the attacker
draw held fixed across seeds would be required to establish the
magnitudes, and this is recorded as outstanding in the Conclusion.

\subsection{AB4 --- LSTM-AE Disabled}
Figures~\ref{fig:ab4a}--\ref{fig:ab4d} present detection quality, detection
latency, control-plane decision error, and trajectory prediction error
as the stealth drift bound $\epsilon_{\max}$ increases, comparing the
full system against AB4 (LSTM temporal autoencoder disabled), under
stealth trajectory poisoning ($\rho_a=0.20$, attack TP-S1, 5 seeds).

\begin{figure}[H]
    \centering
    \includegraphics[width=0.6\linewidth]{images/AB4_results_a.jpeg}
    \caption{AB4 ablation: LSTM-AE removed --- (a) MCC vs.\ stealth
    drift bound $\epsilon_{\max}$.}
    \label{fig:ab4a}
\end{figure}

This is the cleanest ablation in the study, because the two arms diverge
along the independent variable rather than merely differing in level.
Across $\epsilon_{\max} \in \{0.1, 0.2, 0.3, 0.4, 0.5\}$\,m per beacon,
the full system's mean MCC\textsubscript{full} climbs from $-0.021$
through 0.019, 0.133, 0.273 to 0.367, while AB4 is flat at 0.024,
0.024, 0.024, 0.025 and 0.025. The AB4 values are not merely similar
across the sweep --- they are numerically identical for
$\epsilon_{\max} \in \{0.1, 0.2, 0.3\}$ and identical again for
$\{0.4, 0.5\}$, seed for seed. The ablated system's response to stealth
drift magnitude is exactly zero.

The mechanism is the cumulative drift bound of
Equation~\ref{eq:cumulative_drift}. By construction, a stealth attacker
injecting at most $\epsilon_{\max}$ per beacon never violates the
per-beacon kinematic feasibility test of Equation~\ref{eq:tp_s1}, so the
rule term of Equation~\ref{eq:fusion} cannot fire no matter how large
the cumulative displacement $\|\boldsymbol{\delta}(t_0 + kT_b)\|$
becomes. The GAT operates on spatial relations at an instant and sees a
vehicle that is exactly where a vehicle could plausibly be. Only the
reconstruction error $\varepsilon_i(t)$ of the temporal autoencoder,
compared against the MAD-calibrated threshold $\theta_{\mathrm{ae}}$ of
Equation~\ref{eq:ta_threshold}, integrates the deviation over the
sequence and therefore grows with $\epsilon_{\max}$. Remove that term
and the growth has nowhere to register. This is the strongest available
evidence for the temporal half of \textbf{RQ4}~(Section~\ref{sec:rqs}):
the LSTM-AE is not improving a detection the other components perform
poorly, it is the only component for which this attack class is visible
at all.

The full system's near-zero MCC\textsubscript{full} at $\epsilon_{\max}
= 0.1$ is a genuine limit rather than a measurement problem, and reading
it correctly matters. At 0.1\,m per beacon the accumulated drift over a
single RSU contact remains within the reconstruction error that clean
SUMO mobility itself produces, so $\varepsilon_i(t)$ does not clear
$\theta_{\mathrm{ae}}$ and the attack is genuinely undetectable by this
mechanism. The threshold is deliberately calibrated from offline clean
traces using median and MAD (Equation~\ref{eq:ta_threshold}) rather than
from live data, specifically so that an attacker cannot walk the
baseline upward by slow poisoning; the cost of that design choice is a
fixed floor of undetectability, and $\epsilon_{\max} = 0.1$ sits beneath
it. That the curve then rises steeply once drift exceeds clean-data
variance is the expected signature of a correctly calibrated robust
threshold.

Seed~4 is a systematic outlier and should be identified rather than
absorbed into the mean. Its TPE is roughly one sixth of the other seeds'
at every $\epsilon_{\max}$ (0.044 to 0.106 against 0.156 to 0.833) and
its CDER is 0.069 against 0.21--0.31 elsewhere, indicating a
substantially weaker attacker draw; its MCC\textsubscript{full} is
consequently pinned flat at 0.1245 for $\epsilon_{\max} \geq 0.3$ while
the other seeds rise. The full system's mean at high $\epsilon_{\max}$
is therefore an underestimate of the achievable value, since one of five
seeds contributes a near-absent attack. This strengthens rather than
weakens the ablation conclusion, because it means the measured
separation understates the true one.

\begin{figure}[H]
    \centering
    \includegraphics[width=0.6\linewidth]{images/AB4_results_b.jpeg}
    \caption{AB4 ablation: LSTM-AE removed --- (b) TTD vs.\ stealth
    drift bound $\epsilon_{\max}$.}
    \label{fig:ab4b}
\end{figure}

The latency panel shows that the ablated system fails by not detecting
rather than by detecting late, and the archived per-seed values make the
distinction explicit. In three of five seeds AB4 records
TTD $= -1$ at every $\epsilon_{\max}$ --- the sentinel for ``no genuine
detection event was ever localised'' --- and in the remaining two it
records 12.5--16.0\,s. The full system, by contrast, detects in every
seed at every drift bound, with mean TTD falling from 6.44\,s at
$\epsilon_{\max} = 0.1$ to 5.26\,s at 0.5 as larger drift clears
$\theta_{\mathrm{ae}}$ sooner within each sequence.

The correct summary of this panel is therefore a censored one: for the
majority of seeds the ablated system's detection time is not a large
number but an undefined one, and averaging the sentinel with the
finite values would produce a figure that describes neither. The
practical reading is that removing the LSTM-AE converts stealth
trajectory poisoning from an attack that is detected slowly into one
that is usually not detected within the run at all. The two seeds that
do produce finite detections are those with the strongest attacker draw,
where cumulative drift eventually grew large enough to trip the
per-beacon kinematic check --- that is, where the attack stopped being
stealthy.

% TODO: the current draft previously described this panel as "fixed at
% the 8 s window ceiling with zero variance across seeds". The archived
% AB4 logs show TTD = -1 in seeds 2, 4 and 5 at every eps_max, and
% 12.49-16.00 s in seeds 1 and 3. Confirm how the plotted figure
% aggregates the -1 sentinel before final submission.

\begin{figure}[H]
    \centering
    \includegraphics[width=0.6\linewidth]{images/AB4_results_c.jpeg}
    \caption{AB4 ablation: LSTM-AE removed --- (c) CDER vs.\ stealth
    drift bound $\epsilon_{\max}$.}
    \label{fig:ab4c}
\end{figure}

Control-plane decision error is identical between the two arms at every
drift bound and in every seed --- 0.2526, 0.2641, 0.2121, 0.0690 and
0.3058 for seeds 1 through 5, reproduced exactly in both --- and this
apparent null result is in fact a validity check on the ablation itself.
CDER as measured here scores the lightweight control-decision plane,
which AB4 does not modify: the ablation sets $\lambda_3^{(k)} = 0$ in
the fusion stage only. Bit-identical CDER across arms therefore confirms
that the ablation changed exactly one thing and left the rest of the
pipeline untouched, which is the precondition for attributing the MCC
and TTD differences to the LSTM-AE rather than to an incidental side
effect. Where a component's removal ought to move a metric and does not,
that is a finding; where it ought not to and does not, that is
evidence the isolation is clean.

A small monotone decline in CDER with increasing $\epsilon_{\max}$ is
visible in both arms (0.2526 to 0.2442 in seed 1). Larger per-beacon
drift produces trajectories that diverge faster from the true path, and
a fraction of the resulting control decisions become wrong in ways that
the lightweight plane's plausibility gating rejects outright rather than
acting upon --- so the denominator of Equation~\ref{eq:cder} shifts
slightly. The effect is under one percentage point across the full
sweep and is not material to any conclusion.

\begin{figure}[H]
    \centering
    \includegraphics[width=0.6\linewidth]{images/AB4_results_d.jpeg}
    \caption{AB4 ablation: LSTM-AE removed --- (d) TPE vs.\ stealth
    drift bound $\epsilon_{\max}$.}
    \label{fig:ab4d}
\end{figure}

Trajectory prediction error is likewise identical between the two
systems and rises close to linearly with the drift bound: mean TPE runs
0.145, 0.260, 0.375, 0.491 and 0.606 as $\epsilon_{\max}$ goes from 0.1
to 0.5. Both properties are expected and both are useful. TPE is a pure
attack-severity measure computed from the injected trajectory against
SUMO ground truth before any detection filter is applied, so it must be
detector-independent; that it is confirms the metric is implemented as
Equation~\ref{eq:tpe} specifies rather than being contaminated by
detection outcomes. The near-linear growth confirms that the stealth
attack model is delivering the drift budget the bound of
Equation~\ref{eq:cumulative_drift} predicts, which validates the
independent variable of this ablation.

The panel's real force emerges when it is read against panel~(a). At
$\epsilon_{\max} = 0.5$ the attack is inflicting four times the
trajectory error it inflicts at 0.1, and the ablated system's detection
quality is unchanged at 0.025 throughout. A defence whose response is
invariant to a fourfold increase in the damage being done is not
mitigating that attack class in any degree, and this pairing --- rising
severity against flat detection --- is the most direct statement of the
LSTM-AE's necessity available in the study.

\subsection{AB5 --- Full AI Layer Removed (Speed Regime)}
Figures~\ref{fig:ab5a}--\ref{fig:ab5f} present all six primary metrics as
vehicle speed increases, comparing the full system (full AI layer
active) against AB5 (lightweight mode only, GAT and LSTM-AE removed),
under combined attack conditions (3 seeds).

One measurement convention must be stated before the panels are read.
The AB5 arm has no fusion stage, so it emits no MCC\textsubscript{full}
and no CDER\textsubscript{full}; the comparison is necessarily between
the full system's fused verdict and the ablated system's lightweight
verdict. This is not an unfair comparison but the only one available,
and the archived data supplies a strong internal control on it: the full
system's \emph{own} lightweight MCC is bit-identical to the AB5 arm's at
every speed and in every seed, as are TDEE and TPE. The ablation
therefore provably changed only the fusion stage, and the differences
reported below are attributable to the AI layer alone.

\begin{figure}[H]
    \centering
    \includegraphics[width=0.6\linewidth]{images/AB5_results_a.jpeg}
    \caption{AB5 ablation: full AI layer removed --- (a) MCC vs.\
    vehicle speed.}
    \label{fig:ab5a}
\end{figure}

Across $v \in \{10, 60, 100, 140\}$\,km/h the full system's mean
MCC\textsubscript{full} is 0.867, 0.833, 0.867 and 0.867, while the
lightweight-only arm's MCC is 0.333, 0.367, 0.367 and 0.267. The full
system is flat in speed; the ablated system is both far lower and
degrades at the highest speed regime. Both features follow from the
mobility amplification formulation of
Section~\ref{sec:mobility_amplification}. As $v$ rises, contact duration
$\tau_c$ falls by Equation~\ref{eq:contact_duration} and the beacon
sampling window $N_b$ shrinks with it, so a detector confined to a
single RSU contact has strictly less evidence to work with at 140\,km/h
than at 10. The lightweight arm is so confined and degrades accordingly.
The full mode is not, because its window-level scoring aggregates across
RSU records, which is why its curve is flat rather than merely higher.
For \textbf{RQ1}, this is the load-bearing observation: mobility does
not degrade detection uniformly, it degrades detectors that lack
cross-RSU aggregation, and that distinction is exactly what
Equation~\ref{eq:detection_condition} predicts.

\begin{figure}[H]
    \centering
    \includegraphics[width=0.6\linewidth]{images/AB5_results_b.jpeg}
    \caption{AB5 ablation: full AI layer removed --- (b) CDER vs.\
    vehicle speed.}
    \label{fig:ab5b}
\end{figure}

The control-decision panel separates two quantities that are easy to
conflate. On the lightweight control plane, CDER is approximately
0.6--0.7 in \emph{both} arms at every speed --- unchanged by the
ablation, as it must be, since the lightweight plane is what remains in
both. The difference appears in the fusion-verdict control error: the
full system records CDER\textsubscript{full} of 0.0--0.1 across all
speeds, while the ablated arm records the $-1$ sentinel, the quantity
being undefined without a fusion stage. The AI layer therefore reduces
control-decision error by roughly an order of magnitude relative to the
plane it operates on, and the lightweight plane's own error rate is
untouched. Read with panel~(a), this establishes that the AI layer's
benefit is not confined to detection bookkeeping but reaches the control
decisions that \textbf{RQ6} concerns itself with.

\begin{figure}[H]
    \centering
    \includegraphics[width=0.6\linewidth]{images/AB5_results_c.jpeg}
    \caption{AB5 ablation: full AI layer removed --- (c) TTD vs.\
    vehicle speed.}
    \label{fig:ab5c}
\end{figure}

This panel is the most direct empirical evidence in the thesis for the
Detection Infeasible Zone that motivates \textbf{RQ1}. The ablated
system records TTD $= -1$ at every speed and in every seed without
exception, meaning it never localises a genuine detection event within
the observation window --- not that it detects slowly. The full system
detects in 0.6--2.1\,s throughout. The distinction between ``late'' and
``never'' is the entire content of
Equation~\ref{eq:detection_condition}: when $N_b(v_i, r_j) < N_{\min}$,
the temporal detector does not receive a sequence long enough to reach
its target detection power, so no amount of additional waiting recovers
the detection, and the correct measurement outcome is a censored one
rather than a large finite latency.

That the sentinel appears at 10\,km/h as well as at 140 is worth
addressing rather than passing over, because a naive reading of the
speed-dependence argument would predict the lightweight arm to succeed at
low speed. The explanation is that the infeasibility here is driven by
the combined-attack condition rather than by speed alone: under combined
attack the poisoning is distributed across attacker classes, so a
lightweight detector confined to per-beacon predicates has insufficient
evidence at every speed, and speed determines how much worse the
situation becomes rather than whether it is bad. The speed-dependent
component of the prediction is visible in panel~(a), where the ablated
arm's MCC drops from 0.367 to 0.267 between 100 and 140\,km/h; the
latency panel shows the harder fact that the lightweight-only
configuration never reaches a localisable detection under these
conditions at all. Both are consistent with
Equation~\ref{eq:detection_condition}, and the latency panel is the
stronger of the two statements.

\begin{figure}[H]
    \centering
    \includegraphics[width=0.6\linewidth]{images/AB5_results_d.jpeg}
    \caption{AB5 ablation: full AI layer removed --- (d) TDEE vs.\
    vehicle speed.}
    \label{fig:ab5d}
\end{figure}

\begin{figure}[H]
    \centering
    \includegraphics[width=0.6\linewidth]{images/AB5_results_e.jpeg}
    \caption{AB5 ablation: full AI layer removed --- (e) TPE vs.\
    vehicle speed.}
    \label{fig:ab5e}
\end{figure}

Traffic density estimation error and trajectory prediction error are
bit-identical between the two arms at every speed and in every seed ---
TDEE runs 0.2--0.8 and TPE 38.9--225.8 across the sweep, with each value
reproduced exactly in both arms. As in AB4 panel~(d), this is a validity
check rather than a null finding: both are attack-severity quantities
measured against SUMO ground truth before any detection filter applies,
so identity across arms confirms they are implemented as
Equations~\ref{eq:tdee} and~\ref{eq:tpe} specify and are not
contaminated by detector state.

The seed-to-seed spread in TPE is large --- 148.1, 38.9 and 114.7 at
60\,km/h --- and is attributable to the attacker draw rather than to
speed, since the same spread appears at every speed regime. This is the
same three-seed limitation noted for AB3, and it means the TPE panel
supports the claim that attack severity is detector-independent but
does not support fine-grained claims about how severity varies with
speed.

\begin{figure}[H]
    \centering
    \includegraphics[width=0.6\linewidth]{images/AB5_results_f.jpeg}
    \caption{AB5 ablation: full AI layer removed --- (f) PBPO vs.\
    vehicle speed.}
    \label{fig:ab5f}
\end{figure}

Per-beacon overhead quantifies what the AI layer costs, and the answer
bounds the deployment argument. The full system's mean
PBPO\textsubscript{Full} is 1.07, 1.73, 1.60 and 1.37\,ms across the
speed sweep, against 0.0--0.1\,ms for the lightweight-only arm; the
lightweight component PBPO\textsubscript{LW} is below 0.1\,ms in both
arms. Enabling GAT and LSTM-AE inference therefore costs roughly
1--2\,ms per beacon, which is one to two percent of the $T_b = 100$\,ms
beacon budget. Both configurations satisfy the real-time constraint of
Equation~\ref{eq:pbpo} with two orders of magnitude of headroom, so the
choice between them is not a latency trade-off. Set against panel~(a),
where the AI layer more than doubles MCC, and panel~(c), where it is the
difference between detection and none, the cost-benefit conclusion for
\textbf{RQ3} and \textbf{RQ4} is unambiguous.

The overhead curve is non-monotonic in speed, peaking at 60\,km/h and
declining thereafter, which merits an explanation rather than being left
as scatter. PBPO is a per-beacon average over the beacons an RSU
actually processes; at higher speeds contact duration $\tau_c$ falls
(Equation~\ref{eq:contact_duration}) and each RSU processes shorter
sequences per vehicle, so fixed per-window setup costs are amortised
over fewer beacons while the batches themselves shrink. The low value at
10\,km/h has the opposite cause: long dwell times produce large batches
over which fixed costs amortise well. The peak at 60\,km/h is where
these two effects cross. This reading is consistent with the measured
values but is an inference from the mechanism rather than a directly
instrumented decomposition, and it is offered as such.

\subsection{AB6 --- TRS Threshold Signing Gate Removed}
Figures~\ref{fig:ab6a}--\ref{fig:ab6c} present poisoned aggregate rejection
rate, control decision error, and cryptographic operation overhead as
the compromised ring-member fraction $f/n$ increases ($n=4$, combined
attack, $\rho_a=0.40$, 60\,km/h, 3 seeds), comparing the full TRS+FHE
pipeline against AB6 (single-signer, no quorum).

\begin{figure}[H]
    \centering
    \includegraphics[width=0.6\linewidth]{images/AB6_results_a.jpeg}
    \caption{AB6 ablation: TRS threshold gate weakened to single-signer
    --- (a) PARR vs.\ compromised ring-member fraction $f/n$.}
    \label{fig:ab6a}
\end{figure}

The two arms trace different degradation curves against ring
compromise, and the difference is the structural content of the
threshold-signing design. The full $t$-of-$n$ gate records mean PARR of
0.967, 0.967, 0.800, 0.500 and 0.000 at $f/n = 0, 0.25, 0.5, 0.75$ and
1.0, while the single-signer alternative records 1.000, 0.800, 0.500,
0.233 and 0.000. At $f/n = 0.5$ this is 0.80 against 0.50 --- the full
gate rejects 60\% more poisoned aggregates than the alternative at the
same level of ring compromise.

The shapes follow from the two schemes' forgery conditions. With a
single signer, an adversary controlling a fraction $f/n$ of the ring
succeeds whenever the designated signer happens to be one of the
compromised members, giving a rejection rate that falls linearly as
$1 - f/n$; the measured 1.000, 0.800, 0.500, 0.233 track that prediction
closely. The threshold scheme of
Equations~\ref{eq:trs_partial}--\ref{eq:trs_verify} requires $t$
independent partial signatures to combine, so a compromised minority
cannot produce a valid aggregate at all and the rejection rate holds
near ceiling until $f$ approaches $t$. The measured plateau at 0.967
through $f/n = 0.25$ and the subsequent knee at 0.5 and 0.75 are the
signature of that threshold being crossed. This is the structural claim
\textbf{RQ5}~(Section~\ref{sec:rqs}) makes: TRS reduces the fraction of
poisoned aggregates reaching the controller by a mechanism that does not
depend on AI detection accuracy at all, and the curve's shape --- a
plateau then a knee, rather than a smooth decline --- is what
distinguishes a quorum from a probabilistic defence.

Both arms reach PARR $= 0.000$ at $f/n = 1.0$, and this convergence is a
correctness property rather than a failure. When every ring member is
compromised there is no honest signer left to withhold a signature, so a
valid threshold signature over a poisoned aggregate can be produced by
the adversary alone. No signing scheme can reject an aggregate that has
been legitimately signed by a fully compromised quorum, and a scheme
that appeared to would be miscalibrated. The convergence therefore
marks the boundary of the BFT assumption rather than a limitation of
this particular construction. Note also that this point lies beyond the
$\{0/4, 1/4, 2/4, 3/4\}$ range specified in
Section~\ref{sec:ablation}; the sweep was extended to $f/n = 1.0$
specifically to locate this boundary empirically.

\begin{figure}[H]
    \centering
    \includegraphics[width=0.6\linewidth]{images/AB6_results_b.jpeg}
    \caption{AB6 ablation: TRS threshold gate weakened to single-signer
    --- (b) CDER vs.\ compromised ring-member fraction $f/n$.}
    \label{fig:ab6b}
\end{figure}

Control decision error mirrors the rejection curve with the sign
reversed, as it should if PARR is doing the work attributed to it. Mean
CDER for the full gate runs 0.500, 0.500, 0.533, 0.600 and 0.700 across
the sweep, against 0.500, 0.567, 0.600, 0.633 and 0.700 for the
single-signer alternative. The gap opens immediately at $f/n = 0.25$,
reaches its maximum of 0.067 at $f/n = 0.5$ --- precisely where the PARR
separation is widest --- narrows to 0.033 at 0.75, and closes entirely at
1.0.

The shape of this gap is the point, and it differs from a simple
``advantage grows with compromise'' reading. The threshold gate helps
most in the regime it was designed for: a compromised minority below the
threshold, where a quorum requirement is exactly the thing that blocks
them. Below $f/n = 0.25$ there is nothing to block, so the arms agree.
Above $f/n = 0.75$ the compromised set is large enough that the quorum
requirement is increasingly satisfiable by the adversary, so the
advantage erodes; and at $f/n = 1.0$ both arms record 0.700, the same
BFT-boundary convergence seen in panel~(a). A defence whose benefit
peaked outside its design regime would indicate a calibration error;
that this one peaks squarely inside it is evidence the threshold $t$ is
correctly chosen for $n = 4$.

The residual CDER of 0.500 at $f/n = 0$ in both arms should not be
mistaken for a TRS failure. With no ring members compromised, the
remaining control decision error originates in the data-plane attack
($\rho_a = 0.40$ combined attack) rather than in the aggregate plane,
and TRS has no jurisdiction over it. This panel measures the increment
TRS contributes above that floor, and the floor is set by the scenario.

\begin{figure}[H]
    \centering
    \includegraphics[width=0.6\linewidth]{images/AB6_results_c.jpeg}
    \caption{AB6 ablation: TRS threshold gate weakened to single-signer
    --- (c) COO vs.\ compromised ring-member fraction $f/n$.}
    \label{fig:ab6c}
\end{figure}

Cryptographic operation overhead for the full TRS+FHE pipeline rises
with ring compromise --- 138.1, 136.5, 165.4, 223.5 and 291.6\,ms per
epoch --- while the TRS-removed reference is flat at approximately
219\,ms, since without threshold signing the per-epoch cost is
FHE-dominated and independent of $f$. The rise in the full arm is driven
by signature verification and re-coordination work that scales with the
number of compromised members whose partial signatures must be
identified and excluded before a valid aggregate can be formed.

The crossover is the feature of this panel that most needs stating,
because it inverts the naive expectation. Below approximately
$f/n = 0.6$ the full TRS+FHE pipeline is \emph{cheaper} than the
FHE-only alternative --- 138\,ms against 219\,ms at $f/n = 0$ --- and only
becomes more expensive above that point, reaching a 33\% premium at
$f/n = 1.0$. Threshold signing does not simply add cost to the FHE
pipeline; the two interact, because a valid threshold endorsement lets
the pipeline short-circuit work that the unendorsed path must perform in
full. The deployment implication for \textbf{RQ5} is favourable and
concrete: across the entire compromise range the framework is designed
to operate in --- a minority of compromised ring members --- threshold
signing is not a cost to be justified against its security benefit, it
is cost-neutral or better, and the premium appears only in the regime
where the BFT assumption has already been violated.

\subsection{AB7 --- FHE Pre-Coordination Encryption Removed}
Figures~\ref{fig:ab7a}--\ref{fig:ab7d} present per-epoch cryptographic
cost, TRS signing latency, communication overhead ratio, and one-time
distributed key generation setup latency as the RSU signing-ring size
$n$ increases ($n \in \{4, 6, 8, 10, 12\}$), comparing the full
post-quantum PQ-FHE-TRS pipeline against a classical ECDSA-class
baseline (Shamir-Schnorr-P256) at each ring size.

\begin{figure}[H]
    \centering
    \includegraphics[width=0.6\linewidth]{images/AB7_results_a.jpeg}
    \caption{AB7 ablation: post-quantum vs.\ classical cryptographic
    overhead --- (a) per-epoch crypto cost (COO$_{\mathrm{epoch}}$)
    vs.\ signing-ring size $n$.}
    \label{fig:ab7a}
\end{figure}

Per-epoch cryptographic cost for the post-quantum pipeline is
non-monotonic in ring size, measuring 238.6, 184.7, 225.7, 265.8 and
320.0\,ms at $n = 4, 6, 8, 10$ and 12 --- a dip to a minimum at $n = 6$
followed by a steady rise. The rise beyond $n = 6$ is straightforward:
each additional ring member contributes one more FHE ciphertext to be
homomorphically combined per epoch under
Equation~\ref{eq:fhe_ring_combine}, so the dominant cost term grows with
membership.

The decomposition is essential to attributing that rise correctly, and
it contradicts the intuitive explanation. The per-epoch cost separates
into a TRS component and an FHE component, and the TRS component is
\emph{flat} at 21.9, 21.9, 22.3, 21.5 and 22.5\,ms across the entire
ring-size sweep. The whole of the variation lives in the FHE component,
which runs 216.7, 162.9, 203.5, 244.3 and 297.5\,ms. Aggregating more
Dilithium partial signatures is therefore not what drives this curve;
homomorphic ciphertext combination is. This matters for
\textbf{RQ5}~(Section~\ref{sec:rqs}) because it localises the scaling
cost of the cryptographic pipeline to FHE specifically, which is the
component this ablation was designed to isolate, and it means ring-size
scaling decisions should be made against FHE ciphertext arithmetic
rather than against signature aggregation.

The dip at $n = 6$ requires an explanation that the FHE-scaling account
alone cannot supply, since ciphertext count increases monotonically with
$n$. The decisive evidence is that the classical arm shows the identical
shape --- 230.3, 167.0, 207.3, 260.5 and 306.7\,ms, with its own minimum
at $n = 6$ --- despite using an entirely different signature scheme. A
feature common to both arms cannot be caused by the property that
distinguishes them. The most plausible remaining explanation is
measurement environment rather than algorithmic behaviour: $n = 4$ was
the first configuration executed within each arm and therefore absorbs
one-time process warm-up costs --- cryptographic context construction,
allocator and page-cache population --- that later configurations inherit
already amortised. Under that account the true monotone FHE scaling
begins at $n = 6$ and the $n = 4$ point is inflated by a fixed startup
term, which is consistent with the $n = 4$ standard deviations being
among the largest in the sweep. This reading is an inference from run
ordering and cross-arm agreement rather than a directly instrumented
measurement, and the $n = 4$ point should accordingly be treated as an
upper bound rather than a precise estimate.

\begin{figure}[H]
    \centering
    \includegraphics[width=0.6\linewidth]{images/AB7_results_b.jpeg}
    \caption{AB7 ablation: post-quantum vs.\ classical cryptographic
    overhead --- (b) TRS signing latency ($\Delta t_{\mathrm{TRS}}$),
    PQ vs.\ ECDSA-class, vs.\ signing-ring size $n$.}
    \label{fig:ab7b}
\end{figure}

The per-RSU signing latency is flat in ring size for both schemes ---
approximately 22.0\,ms for the post-quantum construction against
3.48\,ms for the classical baseline, a ratio of 6.34$\times$ that varies
by less than 5\% across the entire sweep. Flatness is the expected
consequence of Equation~\ref{eq:trs_partial}: each RSU computes its
partial signature over the message independently, with no interaction
with other ring members, so the cost is a function of the signature
scheme and the message, not of how many other signers exist.

The practical significance for \textbf{RQ5} is that the post-quantum
premium is a fixed per-signer constant rather than a scaling penalty.
An 18.5\,ms absolute increase per RSU per epoch does not compound as the
ring grows, does not interact with the BFT security parameter $f$, and
is incurred once per epoch rather than per beacon --- so it does not
consume the $T_b = 100$\,ms per-beacon budget that
Equation~\ref{eq:pbpo} governs. A 6.34$\times$ multiplier sounds severe
when quoted as a ratio and is modest when quoted as a bounded constant
against an epoch-scale budget; the latter is the correct framing for a
deployment decision, and it is what makes post-quantum signing viable
here rather than merely desirable.

\begin{figure}[H]
    \centering
    \includegraphics[width=0.6\linewidth]{images/AB7_results_c.jpeg}
    \caption{AB7 ablation: post-quantum vs.\ classical cryptographic
    overhead --- (c) byte-overhead ratio (BWO$_{\mathrm{ratio}}$)
    vs.\ signing-ring size $n$.}
    \label{fig:ab7c}
\end{figure}

The byte-overhead ratio is 10.59$\times$ the plain-BSM baseline for the
post-quantum pipeline and 10.53$\times$ for the classical one, and both
are exactly flat across all five ring sizes. The 0.6\% difference
between the arms is the striking result: Dilithium signatures are
roughly an order of magnitude larger than P-256 signatures, yet
substituting one for the other moves the communication overhead ratio by
well under one percent.

The explanation is that FHE ciphertext expansion under
Equation~\ref{eq:fhe_rsu_enc} dominates the byte budget so completely
that the signature contribution is nearly invisible within it. Flatness
in $n$ has the same cause --- ciphertext bytes are a per-RSU constant, so
the ratio of encrypted to plain traffic does not depend on how many RSUs
participate. For \textbf{RQ5} this decouples two design decisions that
would otherwise be entangled: the communication cost of the pipeline is
set by the choice to encrypt at all, not by the choice of signature
scheme, so adopting post-quantum signing carries no bandwidth penalty
worth weighing. It also means that if the 10.5$\times$ expansion is ever
found unacceptable for a deployment, the lever to adjust is the FHE
parameterisation of Section~\ref{sec:thresh_fhe}, not the signature
scheme.

\begin{figure}[H]
    \centering
    \includegraphics[width=0.6\linewidth]{images/AB7_results_d.jpeg}
    \caption{AB7 ablation: post-quantum vs.\ classical cryptographic
    overhead --- (d) one-time DKG\,+\,keygen setup latency vs.\
    signing-ring size $n$.}
    \label{fig:ab7d}
\end{figure}

One-time distributed key generation and key setup
(Equations~\ref{eq:dkg_trs} and~\ref{eq:dkg_fhe}) measures 275.1, 156.3,
180.7, 194.3 and 222.3\,ms for the post-quantum arm and 284.3, 194.8,
187.8, 212.8 and 240.5\,ms for the classical one. The two curves track
each other closely for $n \geq 6$ --- within 8\% at every point --- and
both show the same dip-then-rise shape as panel~(a), including the same
elevated $n = 4$ value with markedly higher variance ($\pm 76.4$\,ms and
$\pm 92.8$\,ms respectively, against single-digit standard deviations at
most larger ring sizes). The cross-arm agreement again indicates a
measurement-environment effect rather than a scheme property, consistent
with the warm-up account developed for panel~(a).

Two qualifications belong with this panel. The post-quantum $n = 6$
value rests on a re-run: the original seed-3 measurement recorded
1849.8\,ms, an order of magnitude outside the distribution, and was
replaced by a clean re-run measuring 154.7\,ms. Substituting a re-run for
an outlier is the correct handling when the outlier is attributable to
system noise rather than to the configuration, but it does mean the
$n = 6$ point rests on a smaller effective sample than the others. And
this quantity is a one-time setup cost amortised over the entire
deployment lifetime, so its absolute magnitude --- a few hundred
milliseconds --- is not operationally significant at any measured ring
size; the panel's value is in confirming that DKG does not scale
prohibitively with ring membership rather than in the specific figures.

Panels~(b), (c) and~(d) together answer the overhead half of
\textbf{RQ5} with an unusually clean result: the post-quantum migration
costs a fixed 18.5\,ms per signer per epoch, essentially nothing in
bandwidth, and nothing that scales with ring size or with the BFT
parameter $f$. The costs that do scale belong to FHE, are localised to
one component by panel~(a), and are bounded at 320\,ms per epoch at
$n = 12$.

\subsection{AB9 --- Multi-Controller Architecture Removed}
Figures~\ref{fig:ab9a}--\ref{fig:ab9b} present control-plane decision
error and overall detection quality as controller compromise onset
time $t_{\mathrm{comp}}/T_{\mathrm{sim}}$ varies, comparing the full
multi-controller architecture against AB9 (single fixed controller, no
revocation pathway).

% TODO: the archived AB9 sweep contains the ablated (single-controller)
% arm only -- there is no full-mode comparison arm in the data. Both
% panels are an onset sweep of the ablated configuration. The analysis
% below is written to what was measured; if the plotted figure shows a
% second series, confirm its provenance before final submission.

\begin{figure}[H]
    \centering
    \includegraphics[width=0.6\linewidth]{images/AB9_results_a.jpeg}
    \caption{AB9 ablation: multi-controller architecture removed ---
    (a) CDER vs.\ controller compromise onset
    $t_{\mathrm{comp}}/T_{\mathrm{sim}}$.}
    \label{fig:ab9a}
\end{figure}

Control decision error decays monotonically as compromise onset moves
later in the run: 0.633, 0.625, 0.497, 0.379 and 0.265 at
$t_{\mathrm{comp}}/T_{\mathrm{sim}} = 0, 0.25, 0.5, 0.75$ and 1.0, with
standard deviations rising from 0.034 to 0.068. The interpretation is
that without a revocation pathway, control-decision damage is simply
proportional to the exposure interval $T_{\mathrm{sim}} -
t_{\mathrm{comp}}$: a controller compromised at the start of the run
corrupts decisions for its entire duration, one compromised at the end
corrupts almost none, and nothing in between arrests the process. The
curve is close to linear in onset, which is exactly what unbounded
accumulation predicts and what a functioning revocation pathway would
replace with a plateau at the revocation latency.

The shape of the decay is therefore the argument for the
multi-controller design rather than an incidental observation. Under the
trusted-controller set of Equation~\ref{eq:trusted_ctrl_set}, damage is
bounded by the time required to detect the compromise and reassign, not
by when the compromise happened, so CDER should become approximately
independent of onset once revocation functions. That the ablated
configuration instead shows onset-proportional damage isolates precisely
what the revocation pathway contributes, and it is the substance of the
\textbf{RQ6}~(Section~\ref{sec:rqs}) claim regarding sustained
control-plane attack.

Two features of the measurement bound how strongly this should be read.
The near-equality of the $0.00$ and $0.25$ points (0.633 and 0.625,
against a standard deviation of 0.034) indicates the damage curve
saturates for early onsets --- once the controller has been compromised
for most of the run, additional exposure time adds little, because the
corrupted decisions have already been made. And the standard deviation
grows steadily with onset, from 0.034 to 0.068, because later onsets
leave fewer control decisions in the exposed window and the rate is
computed over a smaller denominator. Neither undermines the trend, but
both mean the endpoints are less precisely determined than the midpoints.

\begin{figure}[H]
    \centering
    \includegraphics[width=0.6\linewidth]{images/AB9_results_b.jpeg}
    \caption{AB9 ablation: multi-controller architecture removed ---
    (b) MCC vs.\ controller compromise onset
    $t_{\mathrm{comp}}/T_{\mathrm{sim}}$.}
    \label{fig:ab9b}
\end{figure}

Detection quality is invariant to compromise onset, and the invariance
is exact rather than approximate. MCC\textsubscript{full} measures
0.8877 at every onset value with an identical standard deviation of
0.0517, because the three seeds return the same three values
(0.828, 0.919, 0.916) at every point in the sweep. The entire variation
in this panel is between seeds; none of it is between onsets.

This exactness is more informative than an approximately flat curve
would be. It confirms that the vehicle- and RSU-level attacks that drive
detection --- variants a1--a4 and a6 --- proceed identically regardless of
controller state, so the detection pipeline is genuinely decoupled from
the control plane. The multi-controller architecture is consequently a
mitigation mechanism and not a detection one: it bounds the damage a
compromised controller inflicts without contributing any additional
detection signal, exactly as its design purpose as a revocation and
reassignment pathway implies. Panels~(a) and~(b) together therefore
answer \textbf{RQ6} in a specific form --- the contribution is real and it
is located in control-decision outcomes, not in detection quality --- and
a reader who evaluated this component on MCC alone would incorrectly
conclude it contributes nothing.

The on-chain nature of this sweep imposes a reliability caveat worth
recording. Three of the fifteen runs failed the Fabric registration gate
outright, with zero RSUs, vehicles and controllers registered, and were
re-run; one of those required two attempts before passing. Because
partial registration corrupts metrics silently rather than failing
loudly, the reported means use the last gate-passing run per
configuration and no gate-failed row contributes to any point in either
panel. The retry rate is itself a finding about on-chain experiment
cost rather than about the architecture under test.

\subsection{AB10 --- Blockchain Trust Layer Removed}
Figures~\ref{fig:ab10a}--\ref{fig:ab10e} present control decision error,
false revocation and demotion rates, on-chain confirmation latency,
and detection quality as attack penetration $\rho_a$ increases,
comparing the full blockchain-enabled trust layer against AB10
(SC-Trust/SC-Revoke disabled, centralized controller-side revocation).

\begin{figure}[H]
    \centering
    \includegraphics[width=0.6\linewidth]{images/AB10_results_a.jpeg}
    \caption{AB10 ablation: blockchain trust layer removed --- (a)
    CDER vs.\ attack penetration $\rho_a$.}
    \label{fig:ab10a}
\end{figure}

Control decision error is indistinguishable between the arms at low
penetration and separates modestly above $\rho_a = 0.6$. The full system
records 0.531, 0.579, 0.611, 0.626, 0.655 and 0.679 across
$\rho_a \in \{0, 0.2, 0.4, 0.6, 0.8, 1.0\}$, against 0.530, 0.579,
0.600, 0.633, 0.667 and 0.700 for the ablated arm. The separation is
consistent in sign from $\rho_a = 0.6$ upward and reaches 0.021 at
$\rho_a = 1.0$, which is directionally what the loss of BFT-conditioned
revocation predicts: without the $2f+1$ endorsement requirement of
Equation~\ref{eq:bft_revoke_rsu}, revocation decisions depend on a
single controller's view and become less reliable as the population of
misbehaving entities grows.

The magnitude of that separation must be reported honestly, because it
is small. A 0.02 difference in CDER sits close to the resolution at
which the ablated arm's values were logged --- several of its rows carry
only one decimal place --- so this panel establishes the direction of the
effect but not its size. The blockchain layer's demonstrated value in
this ablation does not rest here; it rests on panels~(b) and~(c), where
the differences are qualitative rather than marginal. Presenting the
CDER panel as the primary evidence for \textbf{RQ6} would overstate what
it supports.

\begin{figure}[H]
    \centering
    \includegraphics[width=0.6\linewidth]{images/AB10_results_b.jpeg}
    \caption{AB10 ablation: blockchain trust layer removed --- (b)
    FRR\textsubscript{revoke} vs.\ attack penetration $\rho_a$.}
    \label{fig:ab10b}
\end{figure}

\begin{figure}[H]
    \centering
    \includegraphics[width=0.6\linewidth]{images/AB10_results_c.jpeg}
    \caption{AB10 ablation: blockchain trust layer removed --- (c)
    FRR\textsubscript{demote} vs.\ attack penetration $\rho_a$.}
    \label{fig:ab10c}
\end{figure}

The false revocation rate is where the two designs differ
categorically. The full system revokes zero honest vehicles at every
penetration rate tested --- FRR\textsubscript{revoke} is exactly 0.000 in
every run at every $\rho_a$ --- while the ablated arm's rate climbs from
0.080 at $\rho_a = 0$ to 0.200 at $\rho_a = 0.8$. At the highest tested
penetration, one honest vehicle in five is incorrectly revoked by
centralised controller-side revocation.

The mechanism is the $2f+1$ peer-endorsement condition of
Equation~\ref{eq:bft_revoke_rsu}. Under the full design a revocation
commits only when a Byzantine-tolerant quorum of independent RSUs
endorses the evidence, so a single controller --- whether compromised or
merely mistaken --- cannot revoke an honest vehicle unilaterally.
Removing the on-chain layer removes that check, and the rate rises with
$\rho_a$ because the controller's evidence base becomes noisier as the
attacker population grows, producing more spurious revocation triggers
with nothing to filter them. That the ablated arm shows a non-zero
0.080 rate even at $\rho_a = 0$, with no attackers present at all, is
the sharpest form of this point: unilateral revocation misfires on a
purely honest population, and the BFT condition is what prevents it.
This is a direct and unambiguous answer to the safety half of
\textbf{RQ6}~(Section~\ref{sec:rqs}), which asks whether the trust layer
achieves its benefit without incorrectly revoking honest entities.

The full system's false demotion rate on honest RSUs is not zero,
rising from 0.011 at $\rho_a = 0$ through 0.026, 0.088 to a maximum of
0.115 at $\rho_a = 0.6$ before easing to 0.103 at 0.8. This is a real
cost and should not be presented as harmless, but its severity is
categorically lower than a false revocation and the difference is by
design. Demotion is the middle state of the three-state trust lifecycle:
a demoted RSU has its evidence down-weighted but continues to
participate and can recover its standing through subsequent correct
behaviour, whereas revocation is terminal. The framework therefore
trades a recoverable error against an unrecoverable one, and the
measured rates show it making that trade in the intended direction --- a
maximum 11.5\% recoverable demotion rate against a 0\% terminal
revocation rate, versus the ablated arm's 20\% terminal revocation rate.
The ablated arm's flat zero in panel~(c) is not an advantage: it has no
demotion mechanism at all, so the metric is undefined rather than good,
and its honest-entity errors all land in the terminal category.

The non-monotone peak at $\rho_a = 0.6$ is worth noting rather than
smoothing. Demotion requires accumulated misbehaviour evidence against
an RSU under Equation~\ref{eq:rsu_trust}; at very high penetration the
distinction between honest and compromised RSU behaviour compresses, so
fewer RSUs are classified as honest in the denominator of
Equation~\ref{eq:frr_demote} and the rate falls for reasons unrelated to
the mechanism improving. The peak marks where honest RSUs are most
numerous and most likely to be caught by an evidence threshold tuned for
a noisier environment.

\begin{figure}[H]
    \centering
    \includegraphics[width=0.6\linewidth]{images/AB10_results_d.jpeg}
    \caption{AB10 ablation: blockchain trust layer removed --- (d)
    TCL\textsubscript{confirm} vs.\ attack penetration $\rho_a$.}
    \label{fig:ab10d}
\end{figure}

On-chain confirmation latency establishes what the trust layer costs.
The full system's TCL\textsubscript{confirm} measures 3.01, 4.07, 4.67,
3.44, 3.77 and 3.39\,s across the penetration sweep, while the ablated
arm is zero by definition, having no transactions to confirm. A
3--4.7\,s confirmation latency is the price of BFT-conditioned
revocation, and whether it is acceptable depends entirely on what it
gates. It would be prohibitive on the per-beacon detection path, which
is governed by the $T_b = 100$\,ms budget of
Equation~\ref{eq:pbpo}; it is not on the revocation path, which is an
administrative action taken once per entity and whose correctness
matters far more than its speed. The architecture places the on-chain
layer on the latter path specifically, and panel~(d) quantifies why that
placement is necessary.

The latency is non-monotone in $\rho_a$, peaking at 4.67\,s at
$\rho_a = 0.4$ and declining at higher penetration. Confirmation latency
under the tiered commit strategy is dominated by batching --- an event
waits for its batch window to close --- so latency is highest when
revocation events arrive too sparsely to fill batches quickly and falls
once the event rate is high enough that batches fill on volume rather
than on timeout. The peak at intermediate penetration is the signature
of that transition. This account is consistent with the measured shape
but is an inference from the batching design rather than a directly
instrumented decomposition of the confirmation path.

\begin{figure}[H]
    \centering
    \includegraphics[width=0.6\linewidth]{images/AB10_results_e.jpeg}
    \caption{AB10 ablation: blockchain trust layer removed --- (e) MCC
    vs.\ attack penetration $\rho_a$.}
    \label{fig:ab10e}
\end{figure}

Detection quality is effectively identical between the arms --- 0.799
against 0.806 at $\rho_a = 0.2$, 0.827 against 0.833 at 0.4, 0.617
against 0.633 at 0.6 --- with the differences well inside the seed
spread. As with AB9, this null result is the correct outcome and
localises the component's contribution rather than questioning it: the
blockchain layer governs how trust decisions are committed and
endorsed, not how anomalies are detected, so it should leave MCC
untouched. Read across panels, AB10 shows a component that costs
3--4.7\,s of confirmation latency, contributes nothing to detection, and
eliminates a 20\% false revocation rate against honest vehicles. That is
the shape of the answer \textbf{RQ6} requires, and the false-revocation
result is the part that carries it.

Both arms collapse to MCC $= 0$ at $\rho_a = 0$ and $\rho_a = 1.0$ for
the class-degeneracy reason established in AB1: with no attackers or no
honest vehicles, the confusion matrix has an empty class and the
coefficient is undefined. These endpoints should be read as boundary
artefacts rather than as performance measurements. As in AB9, several
runs in this on-chain sweep required retries after Fabric registration
failures, and the reported means exclude every gate-failed row.

One further property of this sweep belongs in the record. Controller
reassignment latency TCL\textsubscript{reassign}
(Equation~\ref{eq:tcl_reassign}) registered its ``never fired'' sentinel
in \emph{both} arms at every penetration rate, which is why it does not
appear as a panel. The cause is a characterised property of
Equation~\ref{eq:ctrl_trust} rather than an implementation fault: the
conflict term is normalised over every RSU assigned to a controller
while only a small subset is positioned to observe any given decision,
so the measured conflict rate falls far below the level the trust update
would need in order to trigger reassignment. This is the same dilution
effect identified for CP-DETECT in Section~\ref{sec:res-exp5}, and it
means the controller-reassignment half of the trust layer is present and
correct in construction but has not been empirically exercised. It is
recorded as outstanding in the Conclusion.

\subsection{AB11 --- LKH Replaced with Unicast Key Distribution}
Figures~\ref{fig:ab11a}--\ref{fig:ab11b} present the empirical rekey
message count and resulting channel occupancy as vehicle zone density
increases, comparing the Logical Key Hierarchy (LKH) against AB11
(direct per-vehicle unicast rekeying).

This ablation is executed as a standalone micro-benchmark against the
production key-management implementation rather than as a full network
simulation, and the reason is a structural property of the scenario
rather than a matter of convenience. With 200 vehicles distributed over
64 RSU zones, the mean zone occupancy is approximately three and the
maximum observed occupancy across all zones is under sixteen --- an order
of magnitude below the $|\mathcal{V}_j|$ range this ablation must probe.
A simulation run therefore cannot produce the measurement no matter how
long it is executed. The benchmark links the production LKH module
directly, so the measured message counts are those the deployed
implementation generates, and correctness is verified by replaying the
receiver side: across 200 runs, every survivor recovered the new group
key and the revoked vehicle did not, with zero failures. Negative
controls --- a corrupted ciphertext and a deliberately broken
forward-secrecy cut --- were detected in 100 of 100 trials each, so the
verification is not vacuous.

\begin{figure}[H]
    \centering
    \includegraphics[width=0.6\linewidth]{images/AB11_results_a.jpeg}
    \caption{AB11 ablation: LKH replaced with unicast key distribution
    --- (a) rekey message count vs.\ vehicles registered per RSU zone
    $|\mathcal{V}_j|$.}
    \label{fig:ab11a}
\end{figure}

The measured LKH rekey cost matches the closed-form prediction of
Equation~\ref{eq:lkh_rekey} exactly rather than approximately. Across
$|\mathcal{V}_j| \in \{40, 80, 120, 160, 200\}$ the implementation emits
6, 7, 7, 8 and 8 messages, which equals $\lceil \log_2 |\mathcal{V}_j|
\rceil$ at every point without deviation. The unicast alternative emits
39, 79, 119, 159 and 199 messages --- exactly $|\mathcal{V}_j| - 1$, the
survivor count after the revoked vehicle is excluded. At
$|\mathcal{V}_j| = 200$ this is 8 messages against 199, a reduction of
24.9$\times$.

Exact agreement with a closed form is a stronger result than close
agreement, and it reflects what the two schemes actually do. LKH rotates
only the keys on the path from the revoked leaf to the root and
broadcasts one message per rotated path node, each carrying the new key
wrapped under both child keys; the number of such nodes is the tree
depth, which is $\lceil \log_2 |\mathcal{V}_j| \rceil$ by construction,
and it is independent of how many vehicles remain. Unicast rekeying must
instead contact every survivor individually, giving $|\mathcal{V}_j| -
1$. The step structure visible in the LKH curve --- flat from 80 to 120,
flat again from 160 to 200 --- is the tree depth changing only at powers
of two, which is further confirmation that the measured quantity is the
structural one rather than an incidental correlate.

The result also reflects a defect that this ablation exposed and that
was corrected before these measurements were taken, which is worth
recording because it changes how the panel should be read. In the
original implementation both branches of the revocation handler pushed
every surviving member into the affected-member set and unicast one
message per entry, so the LKH arm transmitted $|\mathcal{V}_j|$ messages
--- identical to the unicast arm --- while the logarithmic figure existed
only as a returned value and never as network traffic. The
$\log_2|\mathcal{V}_j|$ scaling claim was therefore not satisfied by the
code that implemented it, and the ablation's own acceptance test is what
surfaced this. The corrected implementation no longer rotates survivors'
leaf keys, which is the correct LKH semantics, and the bottom-level
broadcast deliberately omits the path ciphertext --- that omission is the
forward-secrecy cut that excludes the revoked vehicle. Panel~(a)
therefore validates the closed form of Equation~\ref{eq:lkh_rekey} as
implemented, and it is worth noting that a purely analytical treatment
would have missed the defect entirely.

\begin{figure}[H]
    \centering
    \includegraphics[width=0.6\linewidth]{images/AB11_results_b.jpeg}
    \caption{AB11 ablation: LKH replaced with unicast key distribution
    --- (b) rekey burst channel occupancy vs.\ vehicles registered per
    RSU zone $|\mathcal{V}_j|$.}
    \label{fig:ab11b}
\end{figure}

Panel~(b) converts the message counts into channel occupancy against the
$T_b = 100$\,ms beacon budget. LKH airtime grows from 1.58\,ms at
$|\mathcal{V}_j| = 40$ to 2.11\,ms at 200, remaining under 3\,ms
throughout, while unicast airtime grows from 10.27\,ms to 52.40\,ms ---
more than half the entire beacon budget consumed by a single revocation
event at dense zone occupancy.

The operational consequence is a hard constraint rather than an
efficiency preference. A rekey burst occupying 52\,ms of a 100\,ms
budget leaves insufficient channel time for the beacon traffic that
detection depends on, so unicast rekeying at high zone density does not
merely cost more --- it competes directly with the per-beacon processing
path that Equation~\ref{eq:pbpo} governs, and revocation events cluster
precisely when the network is under attack and detection matters most.
LKH's 2.11\,ms is 2\% of the budget and creates no such contention. This
answers the key-management component of \textbf{RQ5} in a form the
message-count panel alone cannot: the logarithmic scaling is not an
abstract complexity advantage but the difference between a rekey scheme
that coexists with real-time beacon processing and one that does not.

One qualification is essential to reading this panel correctly. The
airtime values are an analytical projection computed from the measured
message counts and payload sizes under a stated 802.11p model --- 6\,Mbps
control-channel rate, 72\,B of protocol headers, a 40\,$\mu$s OFDM
preamble and 58\,$\mu$s of DIFS plus mean backoff on an idle channel
--- and are not an ns-3 channel measurement. The message counts they are
derived from are measured; the conversion to airtime is not. The model
assumes an idle channel and therefore gives a lower bound on real
occupancy: under contention both arms would occupy more airtime, and
the unicast arm disproportionately so, since its cost scales with a
message count that is two orders of magnitude larger. The projection is
consequently conservative with respect to the conclusion it supports.

\subsection{Cross-Ablation Synthesis}

The individual ablations attribute performance to components. Read
together, they answer the question of whether those components are
redundant, which no single ablation can address.

AB3 and AB4 jointly establish that the two AI detectors cover
non-overlapping attack classes rather than reinforcing one another. AB3
removes the GAT and detection of coordinated Sybil configurations falls
to chance ($-0.007$ to $-0.115$ MCC\textsubscript{full}) while the
LSTM-AE remains fully active; AB4 removes the LSTM-AE and detection of
stealth temporal drift flattens at 0.025 while the GAT remains fully
active. Neither surviving detector compensates for the other's absence
in either direction. This is the decisive form of the
\textbf{RQ4}~(Section~\ref{sec:rqs}) claim, because a framework whose
two learning components caught overlapping attack classes could drop one
at no cost; these results show that dropping either surrenders an entire
attack class. AB5 completes the argument by removing both at once and
showing the result is not the sum of the two individual losses but a
configuration that fails to localise any genuine detection at all,
consistent with the Detection Infeasible Zone of
Equation~\ref{eq:detection_condition}.

AB1 and AB5 together bound the lightweight mode from both sides and
answer \textbf{RQ3} in a form neither provides alone. AB1 removes the
rule layer while retaining AI and loses roughly half the detection
quality at negligible overhead saving; AB5 removes the AI layer while
retaining rules and never localises a detection. The lightweight mode is
therefore necessary but not sufficient, and the full mode is more
capable but not self-sufficient --- which is precisely the dual-mode
architecture's premise (Section~\ref{sec:dual_mode}) rather than a
compromise forced by resource limits.

AB6 and AB7 jointly characterise the cost of the cryptographic pipeline
and separate two questions that are easily conflated. AB6 varies ring
compromise at fixed ring size and shows what threshold signing buys ---
60\% more poisoned aggregates rejected at $f/n = 0.5$ --- and what it
costs, which below $f/n \approx 0.6$ is nothing, since the full pipeline
runs cheaper than the FHE-only alternative. AB7 varies ring size at
fixed compromise and shows that the post-quantum premium is a fixed
18.5\,ms per signer that does not scale, while the cost that does scale
belongs to FHE and reaches 320\,ms per epoch at $n = 12$. The combined
answer to \textbf{RQ5} is that the security-relevant cryptographic
choices --- threshold signing and post-quantum signatures --- are
essentially free at deployment scale, and the pipeline's real cost
centre is homomorphic encryption, which is a privacy choice rather than
an integrity one. AB11 completes the overhead picture on the
communication side, establishing that key management scales
logarithmically and consumes 2\% of the beacon budget.

AB9 and AB10 jointly characterise the trust layer, and the pairing
matters because each result alone is easy to misread. Both show flat MCC
--- 0.8877 invariant to compromise onset in AB9, statistically
indistinguishable between arms in AB10 --- and a reader assessing either
component on detection quality would conclude it contributes nothing.
The contributions are real and both lie in mitigation: AB9 shows that
without the multi-controller revocation pathway, control-decision damage
accumulates in proportion to exposure time with nothing to arrest it,
and AB10 shows that without BFT-conditioned revocation, one honest
vehicle in five is falsely revoked at high penetration against zero for
the full system. Together they establish the \textbf{RQ6} claim in both
of its halves --- reduced control-decision error and preserved safety for
honest entities --- and they establish that the trust layer must be
evaluated on mitigation metrics rather than on detection metrics.

Two limitations qualify this synthesis and are recorded in full in the
Conclusion rather than dispersed here. The detection ablations rest on
three seeds (five for AB4) with attacker draw varying materially between
seeds, so component orderings are robust while magnitudes are
directional. And the trust layer's two unexercised paths --- CP-DETECT
alerting and controller reassignment --- share a single identified cause
in the conflict dilution of Equation~\ref{eq:ctrl_trust}, so the
control-plane half of the framework is validated in construction but not
yet in measurement.

\chapter{Timeline and Resource Required}


\section{Timeline}
\begin{figure}[h!]
	\centering
	\includegraphics[width=1.1\textwidth, height=0.5\textheight]{images/PROJECT1.pdf}
	\caption{Project Timeline}
	\label{fig:TimeLine}
\end{figure}

\section{Resource Required}

This project mainly requires software tools and computing resources to design, simulate, and evaluate the proposed mitigation framework. No specialized hardware is required since the evaluation is conducted through simulations.

\subsection{Software Resources}
The following software tools are needed for the project:
\begin{itemize}
	\item Network Simulator 3 (NS-3) for simulating the Software Defined Vehicular Network environment and attack scenarios.
	\item Simulation of Urban Mobility (SUMO) for generating realistic vehicular mobility and traffic patterns.
	\item Python programming environment for implementing machine learning models, including the Graph Attention Network and temporal autoencoder.
	\item Machine learning libraries liked PyTorch or TensorFlow for training and evaluating the learning models.
	\item Blockchain frameworks or libraries (e.g., Hyperledger Fabric or Web3 libraries) for implementing the blockchain based trust mechanism.
	\item LaTeX for preparing project documentation and reports.
\end{itemize}

\subsection{Hardware Resources}
The project requires a standard personal computer or laptop.

\subsection{Estimated Cost}
Since all required tools are open-source or freely available for academic use, the estimated financial cost of the project is minimal. The primary resources are time, computational effort, and access to a personal computer. Therefore, no additional budget allocation is required for this project.


\chapter{Conclusion}

This research addresses mobility pattern and trajectory poisoning attacks in software-defined vehicular networks (SDVNs), a security gap in modern intelligent transport systems. SDVNs increasingly rely on a centralized controller to manage a dynamic environment with large numbers of vehicles, which exposes them to attacks that gradually corrupt mobility and trajectory information. Such attacks are dangerous precisely because they slip past the prevailing security model, which verifies who sent a message but not whether its content is true, and can lead the controller to learn incorrect mobility patterns and issue inefficient or unsafe control decisions.

The literature reviewed in this thesis shows a consistent gap: existing SDVN security solutions rarely address both data-plane and control-plane attacks together. Most focus on a single defense strategy, such as rule-based mobility verification, generic intrusion detection, or cryptographic authentication, and none of the reviewed work accounts for an insider attack originating from the controller itself.

To address these shortcomings, this work proposes a mitigation framework built on three technologies:
\begin{enumerate}
	\item Use of blockchain for providing distributed trust, integrity of data, and traceability, which will minimize the need for a centralized controller.
	\item Graph Attention Networks (GAT) for representing spatial relationships between vehicles and roadside units (RSUs).
	\item Temporal Autoencoders learn normal mobility and trajectory patterns over time and detect subtle deviations that indicate poisoning attacks.
\end{enumerate}


The framework has been assessed via NS-3/SUMO simulation experiments
across ten completed ablation studies (AB1--AB7 and AB9--AB11,
Section~\ref{sec:res-ablation}) and a preliminary per-attack-variant
comparison against three external baselines (Experiment~5,
Section~\ref{sec:res-exp5}), simulating attack cases across both data
and control planes. A related engineering finding from this evaluation
round, correcting a defect in the RSU-side detector state retention
path, reduces the honest false positive rate from 0.0666 to 0.0432 at
matched duration and brings it inside the 0.05 target at the 100\,s
evaluation horizon (Section~\ref{sec:res-rsu-fix}); it is applied
throughout the reported results, and the residual growth of this rate
at the 300\,s horizon is recorded among the outstanding items below.
The completed ablation results
consistently confirm the necessity of every framework component. Removing
the rule-based signature layer (AB1) drops MCC\textsubscript{full} from
0.78 to 0.40 at the sweep's design endpoint $\rho_a = 0.40$;
removing HMAC integrity verification (AB2) allows
MitM detection quality to collapse at high penetration; removing the GAT
spatial detector (AB3) and the LSTM temporal autoencoder (AB4) each
independently reduce coordinated-attack and stealth-drift detection
toward zero, confirming that spatial and temporal analysis catch
distinct attack classes that per-beacon rule checks alone cannot; and
removing the full AI layer (AB5) reproduces the predicted Detection
Infeasible Zone at high speed, where lightweight-only detection never
localizes a genuine alert within the observation window. On the
mitigation side, weakening the TRS threshold-signing gate to a
single signer (AB6) roughly halves the poisoned-aggregate rejection
rate under ring compromise, and removing FHE pre-coordination
encryption (AB7) isolates a bounded, ring-size-independent post-quantum
signing premium (approximately 6$\times$ classical latency) that does
not itself grow with the BFT security parameter. Removing the
multi-controller revocation pathway (AB9) leaves control-decision
error uncorrected once a controller is compromised, and removing the
blockchain trust layer (AB10) reintroduces false revocations against
honest vehicles that the BFT-conditioned SC-Revoke design eliminates
entirely. Finally, replacing the Logical Key Hierarchy with unicast
rekeying (AB11) empirically confirms the predicted logarithmic-versus-linear
scaling gap, with LKH requiring 24.9$\times$ fewer rekey messages at dense
RSU zone occupancy. Complementing these component-level results, the
preliminary per-attack-variant comparison shows MPTD-PQS matching or
exceeding all three external baselines across every evaluated attack
type, with the sharpest advantage on the two control-plane variants that
the data-plane-only baselines have no mechanism to detect at all. Taken
together, these completed results substantiate the research project's
central claim: integrating distributed trust, spatial learning, and
temporal anomaly identification measurably improves the security of
learning-based SDVNs.
\newline
\newline
Seven items remain outstanding at the close of this evaluation round.
They are set out here in full rather than dispersed through
Chapter~\ref{chap:results}, because each qualifies the strength of a
specific claim rather than the direction of it, and none of them
reverses a reported result.

\begin{enumerate}

\item \textbf{AB8 --- RSU three-state trust lifecycle.} The ablation
isolating the RSU lifecycle (Section~\ref{sec:ablation}) is not among
the completed studies reported in Chapter~\ref{chap:results}. Its
absence means the lifecycle's contribution to containing insider RSU
threats is argued from design and from the AB10 false-demotion
measurements rather than isolated experimentally, and the
\textbf{RQ6}~(Section~\ref{sec:rqs}) claim regarding trust-management
sub-components is correspondingly supported by AB9 and AB10 alone.

\item \textbf{Residual false-positive growth at the 300\,s horizon.}
The RSU-side detector state correction of
Section~\ref{sec:res-rsu-fix} brings the honest false positive rate
inside the 0.05 target at 100\,s but not at 300\,s, where it measures
0.2077. The growth is time-dependent within a single run rather than
seed-driven, indicating a second accumulation path still present. Until
it is located, the false-positive requirement should be regarded as met
at 100\,s and not met at the full evaluation horizon.

\item \textbf{Control-plane detection not yet exercised in
measurement.} CP-DETECT raised zero controller-compromise alerts across
the archived runs, and controller reassignment latency registered its
``never fired'' sentinel in both arms of AB10. Both share a single
identified cause --- the conflict term of Equation~\ref{eq:ctrl_trust}
is normalised over all assigned RSUs while only a small subset observes
any given decision, diluting the measured conflict far below the
trust-update threshold. The control-plane half of the framework is
therefore validated in construction but not in measurement, which
bounds the a5 and a7 results of Section~\ref{sec:res-exp5} to the claim
that the baselines cannot address control-plane poisoning.

\item \textbf{Seed count in the detection ablations.} AB1, AB2, AB3,
AB5 and the on-chain sweeps use three seeds and AB4 uses five, with the
effective attacker draw varying materially between seeds --- in AB3 the
realised poison rate ranges from 6.3\% to 18.8\% across the three
seeds. Component orderings are robust at this sample size because they
hold independently in every seed; specific magnitudes and curve shapes
are directional and would require a larger sweep with the attacker draw
held fixed across seeds to establish.

\item \textbf{Experiment~5 baseline comparison.} The per-attack-variant
comparison is reported as preliminary. The archived result set covers
variant a1 at a single penetration rate, so the plotted penetration
sweeps and the B1--B3 baseline curves are not reproducible from the
archived data as it stands, and the seed spread at the measured
operating point is wide relative to several of the between-method gaps.
Re-running the full variant grid with the baselines under a single
globally trained model, and archiving the per-variant outputs alongside
the figures, is required before these curves can carry more than
directional weight.

\item \textbf{Panel~(b) of AB11 is an analytical projection.} The
rekey message counts of Section~\ref{sec:res-ablation} are measured
against the production key-management implementation, but their
conversion to channel occupancy uses a stated 802.11p airtime model on
an idle channel rather than an ns-3 channel measurement. The projection
is conservative with respect to its conclusion, since contention would
penalise the unicast arm disproportionately, but an in-simulation
measurement would strengthen it.

\item \textbf{Rural-scenario generalisation.} The rural mobility
scenario and its associated generalisation experiments remain
outstanding. All results in Chapter~\ref{chap:results} are drawn from
the urban scenario, so the framework's behaviour under the lower RSU
density and longer inter-contact intervals characteristic of rural
deployment is not yet established.

\end{enumerate}

These outstanding items are positioned to extend rather than to
establish the conclusions above: each bounds the horizon, scenario or
statistical confidence over which a completed result holds, and none
contradicts one.
\newline
\newline
A critical and previously uncharacterized aspect of this threat class is
the role of vehicle kinematics as an attack amplification vector. As
formally derived in this work, increasing vehicle speed simultaneously
reduces RSU contact duration and shrinks the observable beacon
sampling window. In high-speed highway regimes, the observable beacon
count $N_b$ falls below the minimum detection threshold $N_{\min}$
(Eq.~\ref{eq:detection_condition}), creating a Detection
Infeasible Zone where sequential anomaly detection becomes unreliable,
while the attacker's falsification budget simultaneously peaks. This
mobility-induced dual unfavorability means that existing threshold-based
defenses are not merely insufficient but are systematically miscalibrated
for real SDVN operating conditions, particularly under high-speed
scenarios where Sybil fraction $f_{\mathrm{sybil}}$ also approaches 1
at low traffic density, allowing a compromised RSU to dominate the
controller's entire view of local traffic. The AB5 ablation results
above provide direct empirical confirmation of this predicted Detection
Infeasible Zone.
\newline\newline
To address these challenges, we propose MPTD-PQS, a two-mode
detection and mitigation system for securing SDVNs with six significant
contributions. Firstly, a multi-source, dual-plane threat model incorporates
attacks against malicious vehicles, compromised RSUs, MitM attackers, and
malicious controllers concurrently in both data and control planes. Secondly,
a formal formulation of mobility-induced threat amplification quantifies
closed-form bounds on how the speed of a vehicle, handover frequency, and
congestion jointly reduce the availability of detection information while
expanding the scope of adversary actions. Thirdly, an attack signature
characterization identifies nine formally defined detection signatures
(TP-S1 to MP-S4), which create a composite anomaly score based on
principled reasoning. Fourthly, a dual-mode detection approach combines
simple rule-based RSU-side detection with more complex full-mode detection
using GAT spatial and LSTM temporal autoencoder anomaly scoring, and the
fused detection scores form a decision variable. Fifthly, a post-quantum
TRS and FHE cryptographically mitigated system structurally ensures that a
minority of compromised RSUs cannot reach the controller regardless of AI
detection efficacy. Sixthly, a permissioned blockchain with smart contracts
(SC-Trust and SC-Revoke) encodes a full-state machine with RSU-direct
Byzantine fault-tolerant consensus, thereby breaking the trust bottleneck
that all reviewed defenses require.

% References
\renewcommand{\bibname}{References}
\bibliographystyle{ieeetr}
\addcontentsline{toc}{chapter}{References} % Add to table of contents
\bibliography{bibliography}

\end{document}
